% !TeX program = xelatex
% Compile twice with XeLaTeX. Standard TeX Live Chinese/Fandol or Noto CJK fonts suffice.
\documentclass[11pt,a4paper]{article}
\usepackage[margin=23mm,headheight=16pt]{geometry}
\usepackage{fontspec}
\IfFontExistsTF{Noto Serif CJK SC}{
  \setmainfont{Noto Serif CJK SC}[AutoFakeSlant=.15]
  \setsansfont{Noto Serif CJK SC}
}{\IfFontExistsTF{FandolSong-Regular.otf}{
  \setmainfont{FandolSong-Regular.otf}[BoldFont=FandolSong-Bold.otf,AutoFakeSlant=.15]
  \setsansfont{FandolSong-Regular.otf}[BoldFont=FandolSong-Bold.otf]
}{
  \setmainfont{NotoSerifCJKsc-Regular.otf}[Path=fonts/,BoldFont=NotoSerifCJKsc-Bold.otf,AutoFakeSlant=.15]
  \setsansfont{NotoSerifCJKsc-Regular.otf}[Path=fonts/,BoldFont=NotoSerifCJKsc-Bold.otf]
}}
\setmonofont{Latin Modern Mono}
\XeTeXlinebreaklocale "zh"
\XeTeXlinebreakskip=0pt plus 1pt
\usepackage{amsmath,amssymb,amsthm,bm,mathtools}
\usepackage{booktabs,tabularx,longtable,array,graphicx,xcolor}
\usepackage{xurl,hyperref,fancyhdr,enumitem}
\usepackage{tikz}
\usepackage{pgfplots}
\pgfplotsset{compat=1.18}
\usetikzlibrary{arrows.meta,positioning}
\usepackage{needspace}
\hypersetup{unicode=true,colorlinks=true,linkcolor=blue!40!black,urlcolor=blue!40!black,citecolor=blue!40!black,pdftitle={Anatomical GPS V4.4},pdfsubject={全历史校准资源前沿、真实材料相位与受约束任务合成}}
\setlength{\emergencystretch}{4em}
\setlength{\parskip}{.30em}
\setlength{\parindent}{1.4em}
\linespread{1.06}
\setlist{nosep,leftmargin=2em}
\newcolumntype{Y}{>{\raggedright\arraybackslash}X}
\numberwithin{equation}{section}
\newtheorem{theorem}{定理}[section]
\newtheorem{proposition}[theorem]{命题}
\newtheorem{lemma}[theorem]{引理}
\newtheorem{corollary}[theorem]{推论}
\theoremstyle{definition}
\newtheorem{definition}[theorem]{定义}
\newtheorem{assumption}[theorem]{假设}
\newtheorem{remark}[theorem]{说明}
\renewcommand{\proofname}{证明}
\renewcommand{\abstractname}{摘要}
\renewcommand{\refname}{参考文献}
\renewcommand{\figurename}{图}
\renewcommand{\tablename}{表}
\renewcommand{\contentsname}{目录}
\DeclareMathOperator{\rank}{rank}
\DeclareMathOperator{\range}{range}
\DeclareMathOperator{\diag}{diag}
\newcommand{\R}{\mathbb R}
\newcommand{\norm}[1]{\left\lVert#1\right\rVert}
\pagestyle{fancy}\fancyhf{}
\fancyhead[L]{\small Anatomical GPS V4.4}
\fancyhead[R]{\small 真实相位与可执行校准}
\fancyfoot[C]{\thepage}
\begin{document}
\title{Anatomical GPS V4.4\\[3mm]
\Large 真实材料任务的可执行校准理论\\[2mm]
\large 信息代价、数据后任务合成与真实覆盖保证}
\author{理论研究稿}
\date{2026年10月9日}
\maketitle
\begin{abstract}
本文研究如何用有限实际探查合成满足材料任务与机器人能力的反馈动作。
主结果在单参数仿射响应、确定性总范数噪声、任务多面体和紧凸物理输入集合下，
以参数后验端点与任务顶点刻画数据后的任意分配，
给出覆盖全部观测历史和任意因果自适应探查的精确校准预算。
一个冗余执行器模型具有完全已知的相位中性任务通道，结构暴露为零，
但仍可因输入能力与真实推进要求而需要严格正的校准能量；本文给出这一耦合的闭式前沿与匹配达到策略。
独立增益多通道的扩展进一步给出共享预算与能力的精确界，
并用同一个最低增益模型构造最坏校准历史，连接数据后速度选择与两阶段完成时间，
并给出具有校准幅值、加速度和真实返回约束的显式波形及常数倍时间保证。
一般矩阵模型 $J=H+\Theta N$ 使用数据条件商空间 $D_Q$ 描述分配自由度与已有信息，
结合经典最优恢复工具，得到零残差历史、名义无偏分配下的先验感知联合半定设计，
以及覆盖全部观测历史的可执行修复充分界。
执行层从完整接触运动学推导真实材料相位，联合验收反馈泄漏、参考配准、任务恢复和输入能力，
并给出含失锁、回扫及离散图像足迹的有限覆盖时间界。
一个有限转角解析任务以真实有限探查产生数据，在非零模型余项与图像噪声下同时满足可逆域、采样保持、输入能力和严格正推进，
还给出一次显式中断返回的完整状态界。材料图册、GP forcing 与逆可达性接口在附录定义。
所有精确性、必要性与充分性均注明其模型集合、策略类和观测历史范围。
\end{abstract}
\newpage
\begingroup\footnotesize\setlength{\parskip}{0pt}\tableofcontents\endgroup
\newpage
\section{材料任务、允许实现与保证层级}\label{sec:rev:task}
令 $\mathcal M$ 为有效材料体积图册，任务为 $\mathcal T=(\tau,\mathcal A_\tau,\mathcal Q_\tau)$：$\tau$ 指定血管、器官、材料片段及扫描目的；$\mathcal A_\tau$ 是允许的探头位置、姿态、相位偏差和机器人实现；$\mathcal Q_\tau$ 指定可信目标关联、视野、声窗与特征质量。允许关系必须在优化前定义。若任务要求图像横向坐标严格为零，则把目标改成离中心点不是合法任务保持；若任务只要求同一片段与合格切面、允许横向偏置，则不同偏置可以是等价实现。

对每个实现 $a\in\mathcal A_\tau$，局部可信特征误差 $e\in\R^m$ 描述完成条件；维数 $m$ 必须保持。选取新的 $a$ 可以改变特征参考值，但不能删掉难控制的特征后声称降低了校准维数。下文的保证首先是局部首次单位响应或统一收缩合同；完整长扫描仍需沿路径反复验收。

\subsection{研究对象、三个层级与阅读路径}
同一材料任务可以由不同的接触位置、切面姿态和机器人构型实现。本文关注这些自由度对反馈信息需求的影响：在给定局部响应结构后，控制分配可先消去一部分未知通道；实际探查又使后续分配能够利用已经获知的方向。数据条件暴露 $D_Q$ 描述信息与分配之间的关系；本稿将其连接到具有真实相位约束的校准资源前沿、数据后执行能力与有效材料覆盖。

全文区分三个层级。\emph{局部分配层}固定物理模型索引 $a$，在同一个 $H_a,N_a,\Theta_a$ 内选择右逆 $B$ 或修复 $C$。\emph{候选实现层}主动改变切片、姿态、接触或机器人分支，形成各自的局部模型与数据集。\emph{执行层}在已接受窗口内连续推进，并通过明确的状态运输和有限过渡连接窗口。局部响应最小最大等式与方向数属于第一层；有限候选算法与数据迁移条件属于第二层；恢复管、真实相位及覆盖条件属于第三层。单参数仿射与独立增益多通道前沿覆盖任意自适应校准历史；一般矩阵零历史半定前沿限定为名义无偏线性分配。单参数、独立增益先验与完整矩阵谱范数球分属不同模型类，信息方向计数不在两者之间直接搬用。

第~\ref{v44g:sec:polytope-frontier}--\ref{v44r:sec:redundant}节先给出单参数仿射主结果及冗余执行器的闭式前沿，随后给出多通道、时间及一般矩阵理论。真实相位、配准和覆盖接口见第~\ref{v44:sec:actual-phase}--\ref{v44:sec:hybrid-coverage}节，完整有限几何实例见第~\ref{v44:sec:finite-geometry-example}节。材料图册、内部几何、GP 技能及 IRD 的构造接口在附录。


\subsection{主要结论与保证范围}
\begin{center}\small
\begin{tabularx}{\linewidth}{p{.23\linewidth}YY}
\toprule
结果 & 量词与最优性 & 物理含义\\
\midrule
定理~\ref{v44g:thm:all-history}、推论~\ref{v44g:cor:energy}
& 单参数仿射响应；任务多面体、紧凸能力和真实相位下，全部历史与任意自适应策略的精确预算。
& 先计算后验可执行宽度，再确定最低物理校准成本。\\
定理~\ref{v44r:thm:exact-frontier}
& 冗余执行器闭式前沿；允许数据后任意动作。
& $D_\star=0$ 仍可能需要严格正的探查预算，且预算随推进速度改变。\\
定理~\ref{v44d:thm:adaptive-speed-deadline}、命题~\ref{v44d:prop:sine-two-approximation}
& 独立增益任务的精确两阶段时间；校准加速度与回位下的两倍时间界。
& 同一最坏模型同时限制学习时间和执行速度。\\
定理~\ref{v44:thm:joint-zero-SDP}
& 零残差历史、$HC=I$ 的线性名义无偏架构；物理 Gram 与未来分配的精确联合设计。
& 先验、探查能量、探查相位和最终能力同时进入设计。\\
定理~\ref{v44:thm:all-history-energy-upper}
& 一般矩阵模型中覆盖全部历史的构造性充分保证。
& 给出可执行修复以及同一预算支持的响应误差与分配范数。\\
命题~\ref{v44:prop:joint-certificate}、定理~\ref{v44:thm:hybrid-finish}
& 已验收物理管、参考配准及有限中断预算下的恢复和覆盖充分保证。
& 真实相位、有效图像与反向桥接共同决定何时完成材料任务。\\
命题~\ref{v44:prop:nonzero-calibration-scan}
& 同一有限几何模型族内，对全部允许噪声与非零余项给出的解析构造。
& 有限真实探查后，以严格正速度完成材料段，并验证采样保持和一次重获桥接。\\
\bottomrule
\end{tabularx}
\end{center}

\subsection{固定局部模型与主动实现选择}
\label{rev:sec:scope}

\begin{definition}[局部实现索引与控制分配]
\label{rev:def:realization-levels}
实现索引 $a\in\mathcal A_\tau$ 指定本次分析采用的物理工作状态、任务特征定义与局部可信窗口。其模型为
\begin{equation}
 J_a=H_a+\Theta_aN_a+E_a^{\rm geo}.
 \label{rev:eq:indexed-model}
\end{equation}
在固定 $a$ 的窗口内，$H_a,N_a$、物理度量及允许的不确定集合保持同一含义；$B$ 或 $C$ 是该模型下可重新选择的命令分配。改变分配不改变 $J_a$，因而可以继续使用同一组响应数据。改变切片、接触构型或其它使响应模型改变的物理状态，则记为 $a\to b$，并使用相应的 $H_b,N_b,\Theta_b$ 及可信窗口。
\end{definition}

以下精确下界均在固定 $a$ 下成立，随后省略下标。所谓任意首次共同动作，是校准历史给定后，在这一固定局部模型内选择任意静态规则 $u(v)$。允许先转移到另一个物理状态、继续探查、再完成任务的策略具有不同的动作历史，须在候选模型及其转移条件下另行分析。

候选选择采用逐候选保证：对每个候选建立自己的数据一致集合和物理合同；只有候选及其转移通过验收时才执行。形式上，若对候选 $a$ 的每个数据一致模型都有
\begin{equation}
 \text{该候选的合同成立}
 \ \Longrightarrow\
 \text{该候选的局部任务保证成立},
 \label{rev:eq:per-candidate-quantifier}
\end{equation}
则任何从已验收候选中进行的选择保留所选候选的这一保证。此陈述不赋予候选搜索全局最优性，也不把固定模型的首次响应下界扩展为任意主动转移策略的下界。跨候选数据复用的一个明确充分条件见命题~\ref{rev:prop:data-transfer}。

\subsection{全稿关键符号及单位约定}
\begin{center}\small
\begin{tabularx}{\linewidth}{lY}
\toprule
符号 & 含义与责任\\
\midrule
$\xi=(\varphi,h,d)$、$P,Q$ & 调和材料坐标、身体位置映射与局部姿态标架\\
$s,\nu$、$\sigma$ & 名义进度及其速率；实际成像所对应的材料相位\\
$H,N,\Theta,E_{\rm geo}$ & 已知任务通道、不确定通道、患者响应及结构外误差\\
$\delta,\eta,\varepsilon_{\rm data}$ & 响应范数界、几何余项界及实际探查数据噪声界\\
$B,D_\star,D_Q$ & 任务实现、不可消除暴露与数据条件剩余暴露\\
$X,\widehat J,C,q$ & 实际激励、估计响应、反馈修复与验收误差\\
$e,\zeta,\omega,S$ & 相对已接受参考的任务误差、纠偏速度、反馈带宽及增广包络\\
\bottomrule
\end{tabularx}\end{center}
生成器使用姿态空间角速度 $\omega_{\rm sk}(s)$，反馈带宽统一记为标量 $\omega>0$。图索引用 $\ell$，物理候选用 $a$；候选包含其有效图 $\ell(a)$。探查积分运动与瞬时速度使用相应时间尺度，毫米、弧度及任务容差通过固定物理度量归一化。学习空间中的 GP 输出是虚拟形状驱动，不是患者端口的物理接触力。





% Front section replacement for the self-contained V4.4 manuscript.
\section{与先例的定理级关系及贡献定位}
\label{sec:rev-related}

面向下游控制性能设计辨识试验已经有长期研究基础。
最低代价辨识研究所需试验时间与工作信号扰动\cite{bombois2006}；
确定性集合辨识研究受幅值和变化率约束的最坏信息半径\cite{tanaskovic2014}；
任务导向探索研究决策损失下的信息配置\cite{wagenmaker2021}。
本文研究一个有限校准后即须统一满足任务误差、真实相位与输入能力的合成问题。
在单参数仿射响应类中，后验区间能否执行由任务顶点与参数端点共同决定；
所需校准预算因而由数据后的分配能力反向确定。
冗余执行器的闭式实例进一步显示：即使存在完全已知的相位中性任务通道，
其物理能力不足仍会造成严格正的校准成本。
表~\ref{tab:rev-related}比较相应定理，而不以模块命名判断新颖性。

\begin{center}\small
\renewcommand{\arraystretch}{1.17}
\begin{longtable}{p{.21\linewidth}p{.345\linewidth}p{.345\linewidth}}
\caption{已验证先例与本文结果的精确边界。}\label{tab:rev-related}\\
\toprule
先例 & 已有结论 & 本文采用或增加的结构\\
\midrule
\endfirsthead
\toprule 先例 & 已有结论 & 本文采用或增加的结构\\\midrule
\endhead
\bottomrule\endfoot
De Moor--Golub\cite{demoorgolub1991}，定理6
& 对 $A+BDC$ 的受限低秩扰动，刻画给定秩下降所需的最小扰动范数。
& 精确满射半径是 $A=H,B=I,C=N$ 的任务结构特例。本文用分配商空间显式写成 $1/\|D_\star\|$，作为后续可实现性门槛。\\
El Ghaoui--Lebret\cite{elghaoui1997}，定理3.1
& 范数有界数据扰动下的鲁棒最小二乘具有精确最坏残差表达。
& 首次共同动作证明使用同类秩一最坏扰动。新增任务对象包含可重选分配零空间 $N\ker H$ 和实际信息空间 $Q$。\\
Kang--You\cite{kang2023}，定理1及第6节
& 性质辨识的最低充分输入空间；先验限制下也讨论不同充分空间及交集失去充分性的情形。
& $D_Q$ 针对固定响应模型的数据后任务重新实现，给出有限误差值与精确方向数。不同混合方向都可能充分这一现象本身已有先例。\\
Bombois 等\cite{bombois2006}；Tanaskovi\'c 等\cite{tanaskovic2014}
& 分别研究满足控制性能的低代价辨识，以及确定性有界噪声下的受约束最坏试验设计。
& 定理~\ref{v44g:thm:all-history}及推论~\ref{v44g:cor:energy}直接连接数据后任务可执行性与实际探查预算，具有全部历史的精确达到策略和自适应下界。\\
Barenthin--Hjalmarsson\cite{barenthin2005}，定理2及第3节
& 以统计参数置信椭球与充分鲁棒反馈条件，将下游 $H_\infty$ 要求转换为辨识输入谱与试验功率设计。
& 任务性能改变信息需求这一思想已有明确先例；本文的新增量化结果是在所声明确定性单参数合同下的必要充分预算及匹配自适应下界。\\
Tohidi--Yildiz--Kolmanovsky\cite{tohidi2020}，定理1及第4节
& 正对角未知执行器增益、饱和与自适应分配；无需独立参数估计或持续激励，证明有界性及渐近分配性质。
& 本文研究有限校准后首次即满足统一误差的资源门槛。定理~\ref{v44r:thm:exact-frontier}允许数据后任意分配，精确量化已知冗余通道能力不足时的必要校准成本。\\
Foucart--Liao\cite{foucart2024}，定理1
& 两个中心超椭球交集的信息半径与最优恢复具有精确双乘子表达。
& 引理~\ref{v44:lem:two-ellipsoid}是其有限维特例；本文的矩阵设计进一步约束物理可实现 Gram 矩阵、最终分配及相位。\\
Wagenmaker 等\cite{wagenmaker2021}；Venkata\-subramanian 等\cite{venkata2024}
& 任务相关最小最大探索，以及具有规定概率鲁棒性能保证的顺序学习控制。
& 本文使用冻结局部响应、确定性总噪声合同与实际执行能力；不把随机有限样本风险与本稿的逐历史合同混用。\\
Venkata\-subramanian 等\cite{venkata2025robust}，定理14
& 能量有界非随机扰动下，以预定频率多正弦及充分 SDP 约束设计达到规定参数椭球精度的低能量探查。
& 本文的单参数任务直接约束最终任务多面体、执行能力与真实相位，给出任意因果自适应策略的匹配资源下界；一般多参数矩阵 SDP 的条件范围单独保留。\\
Shakouri 等\cite{shakouri2026}
& 有先验及有界噪声下，通用输入须对全部系统及其允许数据记录保证集合辨识精度。
& 本文保留这一全历史量词，进一步将单参数仿射响应的后验端点、任务顶点与数据后受约束分配连接；没有宣称一般先验集合辨识的新起点。\\
Peregudin--Dinh\cite{peregudin2026scalar,peregudin2026minimal}
& 标量稳定化的精确噪声门槛；一般系统的最短通用试验，以及与因果自适应 oracle 的性能比较。后续反馈不受饱和限制。
& 本文约束数据后的全部任务命令并指定真实相位速率，求得该结构化任务的完整代价--误差--能力边界。全历史与自适应比较本身已有先例。\\
Grauer--Pei\cite{grauer2022}；Kobayashi\cite{kobayashi2026}
& 前者研究参数不确定性下的最小方差分配；后者研究有限步任务兼容校准与终端任务能力。
& 本文量化相位约束造成的信息需求及其可执行代价。对后者仅核实公开摘要与书目信息，未据此排除其全文所有定理的潜在重叠。\\
\end{longtable}
\end{center}

\paragraph{主贡献的准确陈述。}
第一，在单参数仿射响应、任务多面体和紧凸物理输入集合下，
用后验端点与任务顶点刻画数据后任意任务分配，得到全部观测历史与任意因果自适应探查下的精确预算。
冗余执行器的闭式前沿进一步证明 $D_\star=0$ 与严格正的物理校准成本可以同时成立，
并精确给出同一误差目标下相位速度如何改变最低信息代价。
第二，在独立增益多通道中，用同一个最坏模型连接校准与执行，
给出数据后选速的精确两阶段时间，以及满足校准幅值、加速度和回位约束的两倍时间构造。
第三，在一般矩阵模型中，给出零残差历史、名义无偏架构下的先验感知联合设计，
以及对全部历史有效的修复与执行能力充分界；其最优性范围分别陈述。
第四，将真实材料相位、参考配准、有效成像覆盖与反馈分配放入共同执行合同，
并在具有有限转角、真实有限探查、非零余项及采样保持的同一解析任务上完成全部连接。
这些结果的价值在于可计算的任务合成界；商空间、支持函数、半正定设计与公共 Lyapunov 工具分别按其经典来源使用。

\paragraph{技能表达与理论主线。}
弧长参数化及几何 DMP\cite{aldmp2016,gdmp2024}、GP forcing\cite{gpdmp2026}、
流形向量 GP\cite{gpmanifold2021}提供名义技能表达；
超声自适应视觉伺服\cite{krupa2005}提供交互矩阵背景。
附录完整定义材料图册、身体映射、姿态运输、技能生成与逆可达性接口。
这些模块提供候选和前馈，主文中的响应、能力、配准与覆盖证书决定候选能否执行。


% Root-owned general single-parameter theorem, V4.4.
\section{单参数仿射响应：任务多面体与物理约束的精确校准预算}
\label{v44g:sec:polytope-frontier}

以下结果允许多维任务、冗余执行器、椭球能力和线性接触或相位约束。
未知性仍是一维参数；这使观测的区间几何与数据后的受约束动作能够精确连接。

\subsection{响应与实际探查合同}
以下已固定一个物理实现；本节 $a$ 表示其中的标量未知参数，与第~\ref{rev:sec:scope}节的候选实现索引分别使用。
令 $a\in[\ell,h]$，$\ell<h$，且固定局部响应为
\begin{equation}
 J(a)=J_0+aJ_1\in\mathbb R^{m\times n},\qquad
 \mathcal V=\operatorname{conv}\{v_1,\ldots,v_K\}\subset\mathbb R^m.
 \label{v44g:eq:model}
\end{equation}
给定真实相位速率 $\nu$，定义与参数无关的实际输入集合
\begin{equation}
 \mathcal U_\nu=
 \{u:\|W^{1/2}u\|_2\le U,\ A_cu\le b_c,\ B_cu=c_c,\ gu=\nu\},
 \qquad W\succ0.
 \label{v44g:eq:input-set}
\end{equation}
所有矩阵、相位行 $g$ 和物理度量均已固定；$gu$ 是本局部模型的真实相位导数。
空集直接表示无可执行动作。下文也适用于其它已知紧凸输入集合。
在数据后允许任意映射 $u=\Phi(v,\mathcal H)$，要求
\begin{equation}
 \forall a,\ \forall\text{允许历史 }\mathcal H,\ \forall v\in\mathcal V:
 \quad u\in\mathcal U_\nu,\qquad \|J(a)u-v\|_2\le qR,
 \quad R>0, q\ge0.
 \label{v44g:eq:response-contract}
\end{equation}
$R$ 是对整个任务集固定的尺度；不用随顶点变化的归一化来进行凸插值。

校准的全部参数相关观测由
\begin{equation}
 z_i=ap_i+w_i,\qquad
 P^2=\sum_i p_i^2\le E,\qquad
 \|w\|_2\le\varepsilon+\eta P,
 \quad \varepsilon>0,\ \eta\ge0
 \label{v44g:eq:oracle}
\end{equation}
给出。探查和停时可以是全部过去记录的任意确定性因果函数，但须在每个合法模型、噪声历史上有限停止并遵守预算。
式~\eqref{v44g:eq:oracle}必须有已验证的物理实现：例如，沿已知相位中性方向 $b_{\rm cal}$ 施加 $u_i^{\rm cal}=p_i b_{\rm cal}$，
其中 $gb_{\rm cal}=0$、$\|W_{\rm cal}^{1/2}b_{\rm cal}\|=1$，扣除已知响应后得到该标量记录。
声明的实际探查代价至少为 $P^2$，并且对所讨论的预算存在 $P^2=E$、实际代价恰为 $E$ 的探查，
其完整路径保持在同一模型域且满足全部校准阶段约束。
若还有保持输入，其代价须计入；若还有额外参数观测，则须重新建立完整信息合同。
本节不由抽象 $J(a)$ 自动推出该校准端口存在，也不把离散位移二次预算直接当作连续速度能量。
固定时长的能量测量或完整连续记录实现须分别满足后文给出的噪声和运动合同。

\subsection{后验端点与任务顶点}
对先验内的闭区间 $I=[L,H]$，令 $\mathsf F_q(I;\nu)$ 表示以下有限约束可行：
\begin{equation}
 \forall k\in\{1,\ldots,K\}\ \exists u_k\in\mathcal U_\nu:
 \quad
 \|J(L)u_k-v_k\|_2\le qR,\qquad
 \|J(H)u_k-v_k\|_2\le qR.
 \label{v44g:eq:vertex-socp}
\end{equation}
固定 $L,H,\nu$ 后，这是一个有限二阶锥可行性问题。
每个顶点可以使用不同输入；没有预先强制一个线性分配矩阵服务全部顶点。

\begin{lemma}[固定后验区间的精确任务验收]
\label{v44g:lem:vertices}
$\mathsf F_q(I;\nu)$ 当且仅当存在一个对 $I$ 中全部参数统一有效、满足式~\eqref{v44g:eq:response-contract}的任务映射。
\end{lemma}
\begin{proof}
必要性是将任意有效映射限制到每个 $v_k$ 和两个参数端点。
充分性中，写 $v=\sum_k\lambda_kv_k$，$\lambda_k\ge0$、$\sum_k\lambda_k=1$，并选
$u(v)=\sum_k\lambda_ku_k$。输入集合凸，故 $u(v)\in\mathcal U_\nu$。
对 $a=(1-t)L+tH$，$0\le t\le1$，仿射响应和范数凸性给出
\[
 \|J(a)u(v)-v\|_2
 \le\sum_k\lambda_k\big[(1-t)\|J(L)u_k-v_k\|_2
             +t\|J(H)u_k-v_k\|_2\big]\le qR.
\]
可固定多面体的一个三角剖分来选取分段仿射、可测的凸系数；低维任务集在其仿射包中处理。
\end{proof}

\begin{theorem}[全部历史与任意自适应探查的区间宽度前沿]
\label{v44g:thm:all-history}
在上述物理校准和首次统一响应合同下，对有限 $E>0$ 定义
\begin{equation}
 d_E=\min\{h-\ell,\ 2(\varepsilon/\sqrt E+\eta)\}.
 \label{v44g:eq:data-width}
\end{equation}
存在满足式~\eqref{v44g:eq:response-contract}的策略，当且仅当
\begin{equation}
 \boxed{\quad
 \forall L\in[\ell,h-d_E],\quad
 \mathsf F_q([L,L+d_E];\nu).
 \quad}
 \label{v44g:eq:exact-all-intervals}
\end{equation}
$E=0$ 时同一结论使用 $d_0=h-\ell$。
必要性覆盖任意因果自适应探查、数据后非线性分配与有限停时；一个预定探查和后验顶点分配达到充分性。
\end{theorem}
\begin{proof}
先取 $E>0$。充分性可用一个已实现的 $p=\sqrt E$ 探查。
记 $\widehat a=z/\sqrt E$、$t_E=\varepsilon/\sqrt E+\eta$，则真实参数位于非空区间
\[
 I_{\mathcal H}=[\widehat a-t_E,\widehat a+t_E]\cap[\ell,h].
\]
该区间长度不超过 $d_E$，且可以被某个先验内长度恰为 $d_E$ 的区间包含。
式~\eqref{v44g:eq:exact-all-intervals}和引理~\ref{v44g:lem:vertices}因此给出统一有效的动作。
固定区间后的优化可用严格凸的最小范数择优规则或标准可测择优确定一个策略。

为证必要性，假定存在失败区间 $I=[L,L+d_E]$。
对任意声称有效的自适应策略，在合法参数 $a_c=L+d_E/2$ 上运行零噪声记录 $z_i=a_cp_i$。
因此该记录上策略必有限停止并有 $P^2\le E$。
现在分别取实际参数 $a_-=L$、$a_+=L+d_E$，并在完全相同记录上取
$w_i^\pm=\mp(d_E/2)p_i$。
当 $d_E/2\le\eta$ 时噪声界直接成立；否则由
$d_E/2\le\varepsilon/\sqrt E+\eta$ 和 $P\le\sqrt E$ 得
\[
 \|w^\pm\|_2=(d_E/2)P\le\varepsilon+\eta P.
\]
逐前缀同样有效，故因果探查和停时在两个端点世界完全相同。
失败区间至少有一个任务顶点没有对两端都有效的共同动作。
两端世界在该顶点上使用同一个数据后动作，至少一个世界违约。
这与全部模型、全部历史的保证矛盾。
$E=0$ 时全部 $p_i=0$，零观测记录与整个先验相容；引理直接给出结论。
\end{proof}

\subsection{最小物理预算及可计算的认证界}
先验收完全已知参数时的能力条件
\begin{equation}
 \forall a\in[\ell,h],\quad\mathsf F_q([a,a];\nu).
 \label{v44g:eq:known-feasibility}
\end{equation}
若失败，任意校准预算都不能使统一任务可行。
若通过，定义任务、相位与能力共同允许的最大统一区间宽度
\begin{equation}
 d_{\rm adm}(q,\nu)=
 \max\big\{d\in[0,h-\ell]:
     \forall L\in[\ell,h-d],\ \mathsf F_q([L,L+d];\nu)\big\}.
 \label{v44g:eq:admissible-width}
\end{equation}

\begin{corollary}[数据后任务合成的精确最小预算]
\label{v44g:cor:energy}
在式~\eqref{v44g:eq:known-feasibility}成立时，式~\eqref{v44g:eq:admissible-width}中的最大值存在，且
\begin{equation}
 \boxed{
 E_{\rm req}(q,\nu)=
 \begin{cases}
 0,&d_{\rm adm}=h-\ell,\\[2pt]
 \displaystyle\frac{\varepsilon^2}{(d_{\rm adm}/2-\eta)^2},
   &d_{\rm adm}<h-\ell,\ d_{\rm adm}/2>\eta,\\[6pt]
 +\infty,&d_{\rm adm}<h-\ell,\ d_{\rm adm}/2\le\eta.
 \end{cases}}
 \label{v44g:eq:exact-energy}
\end{equation}
有限正门槛在等号处可达到；$+\infty$ 表示没有有限探查能保证首次统一响应。
\end{corollary}
\begin{proof}
长度为 $d'$ 的区间能嵌入任意较大允许长度 $d$ 的一个先验内区间，故可行宽度集合向下闭合。
由式~\eqref{v44g:eq:known-feasibility}它包含零。
设可行宽度 $d_j\to d$，对任意 $L\in[\ell,h-d]$ 取
$L_j=\min\{L,h-d_j\}$，则 $L_j\to L$。
对每个顶点，由紧性从对应 $u_{k,j}\in\mathcal U_\nu$ 选收敛子列。
有限个顶点可依次取同一子列；响应与约束连续，极限满足 $[L,L+d]$ 的端点约束。
因此宽度集合闭合，是 $[0,d_{\rm adm}]$，最大值达到。
定理~\ref{v44g:thm:all-history}现等价于 $d_E\le d_{\rm adm}$；反解这一不等式即得所列各分支。
特别地，在 $d_{\rm adm}<h-\ell$ 的活动分支，若 $d_E\le d_{\rm adm}$，
必有 $2(\varepsilon/\sqrt E+\eta)\le d_{\rm adm}$，故 $d_{\rm adm}>2\eta$。
噪声地板等号不能用有限能量跨过。有限正最小预算的达到结论继承该预算处的实际探查可实现条件；
若实际校准端口另有上限，还须同时验收该条件。
\end{proof}

式~\eqref{v44g:eq:exact-all-intervals}仍含一个滑动区间起点的连续量词，不能直接称为单个有限 SOCP。
其认证可以给出明确误差界。令
\begin{equation}
 \Psi(L,d)=\max_k\min_{u\in\mathcal U_\nu}
  \max\{\|J(L)u-v_k\|_2,\ \|J(L+d)u-v_k\|_2\},
 \quad K_J=U\|J_1W^{-1/2}\|_2.
 \label{v44g:eq:value-function}
\end{equation}
当输入集合非空时，$\Psi$ 对 $L$ 是 $K_J$-Lipschitz：对固定 $u$ 两个响应的变化均不超过
$|\Delta L|\|J_1u\|\le K_J|\Delta L|$，取极小和有限极大保留该界。
因此，若网格覆盖 $[\ell,h-d]$ 且任一点距网格不超过 $\gamma$，则
\begin{equation}
 \max_{L\in\text{网格}}\Psi(L,d)
 \le\max_{L\in[\ell,h-d]}\Psi(L,d)
 \le\max_{L\in\text{网格}}\Psi(L,d)+K_J\gamma.
 \label{v44g:eq:grid-certificate}
\end{equation}
左界超过 $qR$ 就严格不可行；右界不超过 $qR$ 就严格认证全部滑动区间。
每个网格值由有限 SOCP 得到，结合宽度单调性可以夹逼允许预算。
临界等号附近可能需要无限细化；本节没有把有限网格的保守结论冒充精确有限算法。

\paragraph{与无约束暴露的连接。}
这些结论把信息需求放在任务与输入能力之后：先求哪些后验区间能被数据后动作统一执行，再求制造这种区间所需的最小探查代价。
第~\ref{v44r:sec:redundant}节给出此判据的闭式实例。
其中存在完全已知、相位中性的任务通道，故 $D_\star=0$；但该通道的能力不足时，仍有 $E_{\rm req}(q,\nu)>0$。
这里的必要性针对有限校准后首次就满足指定误差的任务合同，不能据此断言所有允许学习过渡误差的自适应控制都需要独立辨识阶段。
若改为多参数非矩形后验、非仿射响应或参数相关输入集合，区间端点约化须重新证明。



% Independent V4.4 redundant-actuator resource theorem.
% Fragment only. Uses the main manuscript's packages and theorem environments.
% Scope: one fixed positive unknown gain, a known phase-neutral actuator,
% deterministic total noise, actual total quadratic calibration budget,
% a task-command interval, and a shared execution input ball.

\section{零结构暴露仍需校准：冗余执行器的精确资源前沿}
\label{v44r:sec:redundant}

下面的模型在相位中性子空间内同时具有一个未知响应通道和一个已知响应通道。
它在没有输入能力限制时可以完全绕过未知增益，故结构最小暴露秩为零；
共享输入能力却可能迫使控制器使用未知通道。
因此校准门槛与最终的任务分配、材料推进速度发生直接耦合，
不能将一个只依赖目标误差的校准门槛与一个单独的执行条件并列替代。

\subsection{模型、实际预算与全部历史量词}

设执行模型为
\begin{equation}
 \dot y=ax+z+bu_s,\qquad \dot\sigma=u_s=\nu,
 \qquad a\in[\ell,h],\quad 0<\ell\le h<\infty,
 \label{v44r:eq:plant}
\end{equation}
其中 $x,z,u_s$ 均为经过同一固定物理白化的实际输入，
$b\in\mathbb R$ 已知，增益 $a$ 在本次校准和执行窗口内固定。
给定 $r>0$、$U>0$、$0\le\nu\le U$ 和 $q\in[0,1)$，要求对所有
$|v|\le r$ 同时满足
\begin{equation}
 x(v)^2+z(v)^2+\nu^2\le U^2,
 \qquad
 \frac{|ax(v)+z(v)-[v-b\nu]|}{R(\nu)}\le q,
 \qquad R(\nu)=r+|b|\nu.
 \label{v44r:eq:execution-contract}
\end{equation}
先将 $\nu$ 视为校准前给定的目标速度。记剩余归一化能力
\begin{equation}
 M(\nu)=\frac{\sqrt{U^2-\nu^2}}{r+|b|\nu}.
 \label{v44r:eq:normalized-capacity}
\end{equation}

校准阶段实际 $u_s=0$。
第 $i$ 次探查可同时采用未知通道系数 $p_i$ 和已知通道系数 $s_i$，其报告为
\begin{equation}
 Y_i=ap_i+s_i+w_i,
 \qquad \mathsf r_i:=Y_i-s_i=ap_i+w_i.
 \label{v44r:eq:calibration-oracle}
\end{equation}
已知通道的响应可精确扣除，但该通道的输入仍计入实际预算：
\begin{equation}
 \|w\|_2\le\varepsilon+\eta\|p\|_2,
 \quad \varepsilon>0,\quad\eta\ge0,
 \qquad
 \sum_i(p_i^2+s_i^2)\le E<\infty.
 \label{v44r:eq:noise-and-total-energy}
\end{equation}
这里是整个实际历史上的确定性噪声范数合同，
并且同一组 $\varepsilon,\eta$ 对全部比较的探查设计和停止历史共同有效。
校准产生的全部未知增益信息均由所声明的观测涵盖，不另行提供无噪声参数或响应观测。

允许任意良定的确定性因果规则选择 $(p_i,s_i)$ 和停止时刻，
最终 $x(v),z(v)$ 可以任意依赖全部报告和当前命令，也可以是非线性规则。
预算、有限终止和执行合同均须对每个先验模型与每个允许噪声历史成立。
固定随机种子后的逐历史保证也属于这个范围；下文不主张随机噪声下的统计最优性。

\subsection{后验区间的精确最小执行范数}

\begin{lemma}[两个相位中性执行器的区间最小范数]
\label{v44r:lem:interval-minimum}
对 $L>0$、区间宽度 $D\ge0$ 和 $q\in[0,1)$，定义
\begin{equation}
 f(L,D,q)=
 \min_{c_x,c_z\in\mathbb R}
 \left\{\sqrt{c_x^2+c_z^2}:
       |ac_x+c_z-1|\le q\quad\forall a\in[L,L+D]\right\}.
 \label{v44r:eq:interval-minimum-definition}
\end{equation}
令 $k=1-q$，并约定 $D=0$ 时 $2q/D=+\infty$。
则一个最小范数分配为
\begin{equation}
 c_x^*(L,D,q)
 =\min\left\{\frac{kL}{1+L^2},\frac{2q}{D}\right\},
 \qquad
 c_z^*(L,D,q)=k-Lc_x^*(L,D,q),
 \label{v44r:eq:interval-minimum-allocation}
\end{equation}
因此
\begin{equation}
 \boxed{
 f(L,D,q)=
 \begin{cases}
 \displaystyle\frac{k}{\sqrt{1+L^2}},
   &\displaystyle D\frac{kL}{1+L^2}\le2q,\\[2mm]
 \displaystyle\sqrt{\left(\frac{2q}{D}\right)^2+
                    \left(k-\frac{2qL}{D}\right)^2},
   &\displaystyle D\frac{kL}{1+L^2}>2q.
 \end{cases}}
 \label{v44r:eq:interval-minimum-closed-form}
\end{equation}
固定 $D,q$ 时 $f$ 随 $L$ 单调不增；固定 $L,q$ 时随 $D$ 单调不减。
并且
\begin{equation}
 \frac{1-q}{\sqrt{1+L^2}}\le f(L,D,q)\le1-q.
 \label{v44r:eq:interval-minimum-bounds}
\end{equation}
特别地，$q=0,D>0$ 时 $f(L,D,0)=1$，而
$f(L,0,0)=1/\sqrt{1+L^2}$。
\end{lemma}

\begin{proof}
由于响应关于 $a$ 仿射，只需约束区间两个端点。
若 $c_x<0$，上端点的下界要求
$c_z\ge k-(L+D)c_x>k$，其范数严格大于可行已知通道解 $(0,k)$。
故可取 $c_x\ge0$。此时端点约束等价于
\[
 k-Lc_x\le c_z\le 1+q-(L+D)c_x,
 \qquad Dc_x\le2q.
\]
仅考虑下端点半空间 $Lc_x+c_z\ge k$ 时，
原点的最小范数投影为
$(kL/(1+L^2),k/(1+L^2))$。
若它满足 $Dc_x\le2q$，它也满足两个端点的全部约束，得到第一支。
否则全部可行非负 $c_x$ 均有
$0\le c_x\le2q/D<kL/(1+L^2)<k/L$。
在该区间上最小可行 $c_z=k-Lc_x$ 为正，且
\[
 c_x^2+(k-Lc_x)^2
 =(1+L^2)c_x^2-2kLc_x+k^2
\]
严格递减，故在 $c_x=2q/D$ 达到最小值。这给出第二支。

增大 $D$ 会缩小可行集合，故 $f$ 随 $D$ 单调不减。
第一支的平方对 $L$ 的导数为
$-2k^2L/(1+L^2)^2\le0$。
第二支中记 $c_x=2q/D$，则其平方对 $L$ 的导数是
$-2c_x(k-Lc_x)\le0$；在该支上 $k-Lc_x>0$。
两支在切换处连续，因此得到关于 $L$ 的全局单调性。
下界来自前述半空间，上界来自 $(0,k)$。
\end{proof}

\subsection{任意因果探查下的精确前沿}

对有限 $E\ge0$，定义
\begin{equation}
 t(E)=
 \begin{cases}
 \varepsilon/\sqrt E+\eta,&E>0,\\
 +\infty,&E=0,
 \end{cases}
 \qquad
 D_E=\min\{h-\ell,2t(E)\},
 \qquad F_E(q)=f(\ell,D_E,q).
 \label{v44r:eq:worst-posterior-width}
\end{equation}

\begin{theorem}[冗余执行器的全部历史精确校准--推进条件]
\label{v44r:thm:exact-frontier}
在式~\eqref{v44r:eq:plant}--\eqref{v44r:eq:noise-and-total-energy} 的模型和量词下，
存在满足所有历史预算及执行合同的策略，当且仅当
\begin{equation}
 \boxed{
 f\bigl(\ell,\min\{h-\ell,2t(E)\},q\bigr)\le M(\nu).}
 \label{v44r:eq:exact-frontier}
\end{equation}
等价地，信息和真实推进能力满足同一个必要与充分不等式
\begin{equation}
 \boxed{\ \nu^2+F_E(q)^2[r+|b|\nu]^2\le U^2.\ }
 \label{v44r:eq:joint-energy-speed-frontier}
\end{equation}
任意因果探查、自适应停止和最终非线性分配均不能改善这个全部历史前沿。
达到策略采用一次预定探查，再针对实际后验区间选择最小范数执行分配。
\end{theorem}

\begin{proof}
先证充分性。$E>0$ 时取一次探查 $p=\sqrt E,s=0$。
令 $\widehat a=\mathsf r/p$，实际数据一致区间为
\begin{equation}
 [L,H]=[\ell,h]\cap[\widehat a-t(E),\widehat a+t(E)].
 \label{v44r:eq:posterior}
\end{equation}
该区间非空，真实 $a$ 属于其中，且 $L\ge\ell$、$H-L\le D_E$。
若 $E=0$，直接使用先验区间。
由引理的单调性，所有历史上都有
\[
 f(L,H-L,q)\le f(\ell,D_E,q)=F_E(q).
\]
按式~\eqref{v44r:eq:interval-minimum-allocation} 对该后验求得 $(c_x^*,c_z^*)$，并取
\begin{equation}
 x(v)=c_x^*(v-b\nu),\qquad
 z(v)=c_z^*(v-b\nu),\qquad u_s=\nu.
 \label{v44r:eq:attaining-policy}
\end{equation}
其误差不超过 $q|v-b\nu|\le qR(\nu)$，且总输入平方不超过
$R(\nu)^2F_E(q)^2+\nu^2\le U^2$。
它因而满足全部历史与全部命令的保证。

再证必要性。考虑任意满足预算和有限终止要求的因果协议。
先令真实增益为
\begin{equation}
 a_0=\ell+D_E/2
 \label{v44r:eq:common-midpoint}
\end{equation}
并在所有报告中取零噪声。
在这条允许历史上，协议产生某个有限的完整 $(p,s)$，且
$P:=\|p\|_2\le\sqrt E$。
将真实增益分别改为
$a_-=\ell$ 与 $a_+=\ell+D_E$，并分别选择
\begin{equation}
 w_-=\frac{D_E}{2}p,\qquad
 w_+=-\frac{D_E}{2}p.
 \label{v44r:eq:indistinguishable-pair-noise}
\end{equation}
两个模型的全部报告都仍然是 $a_0p+s$。
若 $E>0$，由 $D_E/2\le\varepsilon/\sqrt E+\eta$ 和 $P\le\sqrt E$，
有
\[
 \frac{D_E}{2}P\le\varepsilon+\eta P.
\]
同一估计也适用于每个前缀，所以两条噪声历史均合法。
$E=0$ 时 $p=s=0$，两条零噪声历史直接一致。
因果性保证两个模型诱导完全相同的探查、停止与最终共同规则；
这没有将数据后分配预先固定。

现在取 $v=-r\operatorname{sgn}(b)$；$b=0$ 时取 $v=-r$。
于是 $g:=v-b\nu$ 满足 $|g|=R(\nu)>0$。
若同一个动作 $(x,z)$ 同时满足两个模型的执行合同，
归一化 $(c_x,c_z)=(x/g,z/g)$ 必须满足
\[
 |a_-c_x+c_z-1|\le q,
 \qquad |a_+c_x+c_z-1|\le q,
 \qquad\sqrt{c_x^2+c_z^2}\le M(\nu).
\]
由引理，前两个端点约束所需的最小范数恰为 $f(\ell,D_E,q)$，
故必须有式~\eqref{v44r:eq:exact-frontier}。
该论证只用同一命令下的共同动作，因而同时覆盖任意最终非线性规则。
\end{proof}

\begin{corollary}[显式最小能量与边界]
\label{v44r:cor:minimum-energy}
固定 $q,\nu$，令 $k=1-q$，
\begin{equation}
 F_{\rm known}=\frac{k}{\sqrt{1+\ell^2}},
 \qquad F_{\rm prior}=f(\ell,h-\ell,q),
 \qquad M=M(\nu).
 \label{v44r:eq:known-prior-norms}
\end{equation}
以 $\mathcal E_{\rm red}(q,\nu)$ 表示上述全部历史保证所需的最小有限能量；
不存在有限能量策略时令其为 $+\infty$。
\begin{enumerate}
 \item 若 $M\ge F_{\rm prior}$，则 $\mathcal E_{\rm red}(q,\nu)=0$。
 \item 若 $M<F_{\rm known}$，则即使精确知道增益也无法在全部先验模型上执行，
 故 $\mathcal E_{\rm red}(q,\nu)=+\infty$。
 \item 在剩余区域 $F_{\rm known}\le M<F_{\rm prior}$，必有 $h>\ell$。
 若 $q=0$，则 $\mathcal E_{\rm red}(0,\nu)=+\infty$。
 若 $q>0$，定义
 \begin{equation}
 x_-(q,M)=
 \frac{\ell k-\sqrt{(1+\ell^2)M^2-k^2}}{1+\ell^2}>0,
 \qquad t_{\rm req}(q,M)=\frac{q}{x_-(q,M)}.
 \label{v44r:eq:required-uncertainty}
 \end{equation}
 则
 \begin{equation}
 \boxed{
 \mathcal E_{\rm red}(q,\nu)=
 \begin{cases}
 \displaystyle\frac{\varepsilon^2}{[t_{\rm req}(q,M(\nu))-\eta]^2},
       &t_{\rm req}(q,M(\nu))>\eta,\\[2mm]
 +\infty,&t_{\rm req}(q,M(\nu))\le\eta.
 \end{cases}}
 \label{v44r:eq:minimum-energy-closed-form}
 \end{equation}
\end{enumerate}
所有有限门槛均可达到。能力边界 $M=F_{\rm known}$ 可以包含在第三项中，
不要求严格能力余量；但由于 $\varepsilon>0$，
$t_{\rm req}=\eta$ 仍不能由任何有限能量达到。
\end{corollary}

\begin{proof}
前两项直接来自定理和引理的上下界。
在中间区域 $M<F_{\rm prior}\le k$，
方程
\[
 x^2+(k-\ell x)^2=M^2
\]
的较小根正是式~\eqref{v44r:eq:required-uncertainty} 的 $x_->0$，
且 $x_-\le k\ell/(1+\ell^2)$。
由引理的二次式表示，
\[
 f(\ell,D,q)\le M
 \quad\Longleftrightarrow\quad
 D x_-\le2q.
\]
又因 $M<F_{\rm prior}$，先验宽度满足
$(h-\ell)x_->2q$，
所以定理中的 $D_E=\min\{h-\ell,2t(E)\}$ 能通过能力限制，
当且仅当 $t(E)\le q/x_-=t_{\rm req}$。
解这个不等式即得式~\eqref{v44r:eq:minimum-energy-closed-form}。
若 $q=0,h>\ell$，则正的 $\varepsilon$ 使每个有限 $E$ 都有 $D_E>0$，
而 $f(\ell,D_E,0)=1=F_{\rm prior}>M$，所以不存在有限能量策略。
\end{proof}

零能量分支必须单独保留：在 $M=F_{\rm prior}$ 处已可直接用先验，
而从 $M<F_{\rm prior}$ 一侧逼近时，式~\eqref{v44r:eq:minimum-energy-closed-form}
可能趋向严格正值。这是正确定性噪声预算下的门槛现象。
若先验退化为 $h=\ell$，则 $F_{\rm prior}=F_{\rm known}$，
只剩零校准可执行与已知模型也不可执行两种情况。
本节采用 $\varepsilon>0$；无加性噪声 $\varepsilon=0$ 时，
任意小非零探查与完全无探查之间可能出现不同的可达性边界，不能直接套用最小值表述。

\begin{corollary}[能量给定后的最大统一推进速度]
\label{v44r:cor:uniform-speed}
给定有限 $E$，令 $F=F_E(q)>0$。
若 $Fr>U$，则连 $\nu=0$ 都无法保证；若 $Fr=U$，只允许 $\nu=0$。
若 $Fr<U$，最大可保证速度为
\begin{equation}
 \boxed{
 \nu_{\rm red}(E,q)=
 \frac{-F^2r|b|+
       \sqrt{(1+F^2b^2)U^2-F^2r^2}}
      {1+F^2b^2}>0.}
 \label{v44r:eq:uniform-speed-root}
\end{equation}
这个保证也允许控制器在看到数据或当前命令后选择速度，
只要所要求的是每个允许历史、每个命令下的统一下界
$\nu\ge\nu_0\ge0$：保证该下界的条件仍是
式~\eqref{v44r:eq:exact-frontier} 以 $\nu_0$ 代入。
\end{corollary}

\begin{proof}
第一部分由式~\eqref{v44r:eq:joint-energy-speed-frontier} 的二次不等式直接解得。
对最后一项，$M(\nu)$ 在 $[0,U]$ 上单调不增，
且必要性证明中的命令角点与任何非负 $\nu$ 无关。
同一个报告与角点命令使两个模型选择同一速度和同一动作。
若该速度至少为 $\nu_0$，则其归一化动作范数至多为
$M(\nu)\le M(\nu_0)$，所以原来的不可区分对仍给出必要性。
充分性只需使用固定速度 $\nu_0$ 的达到策略。
\end{proof}

\subsection{可算的零秩、正能量分离及物理实现}

将 $a=a_c+\theta$、$a_c=(\ell+h)/2$，则
\begin{equation}
 J=H+\theta N,
 \qquad H=[\,a_c,1,b\,],\qquad N=[\,1,0,0\,],
 \qquad |\theta|\le(h-\ell)/2.
 \label{v44r:eq:HN-realization}
\end{equation}
真实相位行是 $g=[0,0,1]$。
取 $B=(0,1,0)^\mathsf T$，有 $HB=1$、$NB=0$、$gB=0$，
所以即使要求相位中性，结构最小暴露秩仍为零。
这个事实保证存在完全已知的任务通道，并不保证该通道能在指定能力内同时支持给定的任务幅值和推进速度。

例如，取
\begin{equation}
 \ell=1,\quad h=3,\quad r=b=U=1,\quad
 q=\frac1{10},\quad\varepsilon=\frac1{10},\quad\eta=\frac1{50}.
 \label{v44r:eq:numerical-parameters}
\end{equation}
此时
\[
 F_{\rm prior}=\sqrt{\frac{13}{20}},\qquad
 F_{\rm known}=\frac9{10\sqrt2},\qquad
 M(\nu)^2=\frac{1-\nu}{1+\nu}.
\]
因此 $0\le\nu\le7/33$ 时不需校准；
$\nu>119/281$ 时即使完整知道增益，也无法在最小增益模型上执行。
在二者之间校准能够改变是否可执行。
本例 $q=1/10$ 时，中间区域的最小 $t_{\rm req}$ 为 $2/9>\eta$，
所以直到已知模型的能力边界都可用有限能量达到。
例如 $\nu=1/3$ 时
\[
 M=\frac1{\sqrt2},\qquad
 x_-=\frac{9-\sqrt{19}}{20},\qquad
 t_{\rm req}=\frac{9+\sqrt{19}}{31}.
\]
\begin{equation}
 \boxed{
 \mathcal E_{\rm red}\left(\frac1{10},\frac13\right)
 =\frac{1/100}{[(9+\sqrt{19})/31-1/50]^2}
 \simeq0.0592188.}
 \label{v44r:eq:numerical-positive-energy}
\end{equation}
这是所有允许因果探查和数据后分配的匹配下界与达到值，
并非限定某个预选右逆后的证书代价。
若仅将本例的 $\eta$ 改成零，则同一点简化为
\[
 \mathcal E_{\rm red}\left(\frac1{10},\frac13\right)
 =\frac{50-9\sqrt{19}}{200}\simeq0.0538495.
\]
幅值相关误差的影响不限于增大这个能量常数。
仍取 $\nu=1/3$ 和 $\eta=1/50$，但允许目标 $q$ 变化，
则无限能量下仍只能将最坏后验宽度逼近 $2\eta=1/25$。
在相关的第二分支，极限能力条件为
\[
 f(1,1/25,q)^2=(50q)^2+(1-51q)^2\le\frac12.
\]
故不可消除的误差下确界为
\begin{equation}
 q_{\rm red,floor}
 =\frac{102-\sqrt{202}}{10202}
 \simeq0.00860491.
 \label{v44r:eq:example-noise-floor}
\end{equation}
等号对应 $t_{\rm req}=\eta$，不能用有限能量达到；
更小的 $q$ 即使无限增大探查能量也不能保证。
未知参数只有一个，探查也只需一个时间模式；这里的分离来自执行能力与信息的耦合，
不应被解释为新的探查方向秩下界。


\begin{figure}[htbp]
\centering
\begin{tikzpicture}
\begin{axis}[
 width=.88\linewidth,height=73mm,
 xlabel={指定真实相位速率 $\nu$},ylabel={精确最小校准能量 $\mathcal E_{\rm red}$},
 xmin=0,xmax=.46,ymin=-.005,ymax=.27,
 xtick={0,.1,.2,.3,.4},ytick={0,.05,.10,.15,.20,.25},
 grid=major,major grid style={gray!20},
 tick label style={font=\small},label style={font=\small}]
\addplot[draw=none,fill=red!9] coordinates
 {(0.423487544484,-.005) (.46,-.005) (.46,.27) (0.423487544484,.27)} \closedcycle;
\addplot[blue!70!black,very thick] coordinates {(0,0) (0.212121212121,0)};
\addplot[blue!70!black,only marks,mark=*] coordinates {(0.212121212121,0)};
\pgfplotstableread[row sep=\\]{
nu energy\\
0.212121212121 0.0104123281966\\
0.214229591287 0.0108532269859\\
0.216327402135 0.0113024680323\\
0.218414644667 0.0117600118531\\
0.220491318883 0.0122258196606\\
0.222557424782 0.0126998533532\\
0.224612962364 0.0131820755065\\
0.226657931629 0.0136724493653\\
0.228692332578 0.0141709388348\\
0.230716165211 0.0146775084728\\
0.232729429527 0.0151921234816\\
0.234732125526 0.0157147497002\\
0.236724253208 0.0162453535965\\
0.238705812574 0.0167839022601\\
0.240676803623 0.0173303633942\\
0.242637226356 0.0178847053094\\
0.244587080772 0.0184468969155\\
0.246526366872 0.0190169077154\\
0.248455084654 0.0195947077977\\
0.250373234121 0.0201802678302\\
0.25228081527 0.0207735590535\\
0.254177828103 0.0213745532742\\
0.256064272619 0.0219832228588\\
0.257940148819 0.0225995407274\\
0.259805456702 0.0232234803478\\
0.261660196269 0.0238550157292\\
0.263504367519 0.0244941214167\\
0.265337970452 0.0251407724853\\
0.267161005068 0.0257949445345\\
0.268973471368 0.0264566136826\\
0.270775369352 0.0271257565614\\
0.272566699019 0.0278023503107\\
0.274347460369 0.0284863725736\\
0.276117653402 0.0291778014906\\
0.277877278119 0.0298766156956\\
0.27962633452 0.03058279431\\
0.281364822603 0.0312963169387\\
0.28309274237 0.0320171636648\\
0.284810093821 0.0327453150455\\
0.286516876955 0.0334807521072\\
0.288213091772 0.0342234563412\\
0.289898738272 0.0349734096995\\
0.291573816456 0.0357305945901\\
0.293238326324 0.0364949938734\\
0.294892267874 0.0372665908579\\
0.296535641109 0.0380453692958\\
0.298168446026 0.0388313133799\\
0.299790682627 0.039624407739\\
0.301402350911 0.0404246374345\\
0.303003450879 0.0412319879566\\
0.30459398253 0.0420464452209\\
0.306173945864 0.0428679955646\\
0.307743340882 0.0436966257432\\
0.309302167583 0.0445323229272\\
0.310850425968 0.0453750746985\\
0.312388116036 0.0462248690475\\
0.313915237787 0.0470816943696\\
0.315431791222 0.0479455394624\\
0.31693777634 0.0488163935225\\
0.318433193141 0.0496942461426\\
0.319918041626 0.0505790873087\\
0.321392321794 0.0514709073968\\
0.322856033646 0.0523696971708\\
0.324309177181 0.0532754477794\\
0.325751752399 0.0541881507534\\
0.327183759301 0.0551077980034\\
0.328605197886 0.056034381817\\
0.330016068155 0.0569678948565\\
0.331416370107 0.0579083301565\\
0.332806103742 0.0588556811215\\
0.334185269061 0.0598099415236\\
0.335553866063 0.0607711055001\\
0.336911894748 0.0617391675519\\
0.338259355117 0.0627141225407\\
0.339596247169 0.0636959656872\\
0.340922570905 0.0646846925695\\
0.342238326324 0.0656802991202\\
0.343543513426 0.0666827816255\\
0.344838132212 0.0676921367227\\
0.346122182681 0.0687083613986\\
0.347395664833 0.0697314529878\\
0.348658578669 0.0707614091708\\
0.349910924189 0.0717982279726\\
0.351152701391 0.0728419077611\\
0.352383910277 0.0738924472451\\
0.353604550847 0.0749498454734\\
0.354814623099 0.0760141018329\\
0.356014127035 0.0770852160474\\
0.357203062655 0.0781631881763\\
0.358381429958 0.0792480186128\\
0.359549228944 0.0803397080835\\
0.360706459614 0.0814382576461\\
0.361853121967 0.0825436686892\\
0.362989216003 0.0836559429305\\
0.364114741723 0.084775082416\\
0.365229699126 0.0859010895187\\
0.366334088213 0.0870339669381\\
0.367427908983 0.0881737176986\\
0.368511161436 0.0893203451489\\
0.369583845573 0.0904738529613\\
0.370645961393 0.0916342451304\\
0.371697508897 0.0928015259726\\
0.372738488084 0.0939757001252\\
0.373768898954 0.0951567725461\\
0.374788741508 0.0963447485123\\
0.375798015745 0.09753963362\\
0.376796721665 0.0987414337837\\
0.377784859269 0.0999501552358\\
0.378762428556 0.101165804526\\
0.379729429527 0.10238838852\\
0.380685862181 0.1036179144\\
0.381631726518 0.104854389667\\
0.382567022539 0.106097822132\\
0.383491750243 0.107348219926\\
0.38440590963 0.108605591494\\
0.385309500701 0.109869945593\\
0.386202523455 0.111141291298\\
0.387084977893 0.112419637997\\
0.387956864014 0.11370499539\\
0.388818181818 0.114997373495\\
0.389668931306 0.116296782641\\
0.390509112477 0.117603233472\\
0.391338725332 0.118916736945\\
0.39215776987 0.120237304334\\
0.392966246091 0.121564947223\\
0.393764153995 0.122899677514\\
0.394551493584 0.124241507419\\
0.395328264855 0.12559044947\\
0.39609446781 0.126946516508\\
0.396850102448 0.128309721694\\
0.39759516877 0.129680078502\\
0.398329666775 0.131057600721\\
0.399053596463 0.132442302457\\
0.399766957835 0.133834198133\\
0.40046975089 0.135233302487\\
0.401161975628 0.136639630577\\
0.40184363205 0.138053197777\\
0.402514720155 0.13947401978\\
0.403175239944 0.140902112597\\
0.403825191416 0.14233749256\\
0.404464574571 0.143780176321\\
0.40509338941 0.145230180852\\
0.405711635932 0.146687523447\\
0.406319314138 0.148152221724\\
0.406916424027 0.149624293622\\
0.407502965599 0.151103757406\\
0.408078938855 0.152590631664\\
0.408644343794 0.154084935312\\
0.409199180416 0.15558668759\\
0.409743448722 0.157095908069\\
0.410277148711 0.158612616646\\
0.410800280384 0.160136833549\\
0.41131284374 0.161668579339\\
0.411814838779 0.163207874904\\
0.412306265502 0.16475474147\\
0.412787123908 0.166309200595\\
0.413257413998 0.167871274173\\
0.413717135771 0.169440984436\\
0.414166289227 0.171018353951\\
0.414604874366 0.17260340563\\
0.415032891189 0.17419616272\\
0.415450339696 0.175796648814\\
0.415857219886 0.177404887848\\
0.416253531759 0.179020904104\\
0.416639275315 0.180644722208\\
0.417014450555 0.182276367138\\
0.417379057479 0.183915864221\\
0.417733096085 0.185563239133\\
0.418076566375 0.187218517907\\
0.418409468349 0.18888172693\\
0.418731802006 0.190552892944\\
0.419043567346 0.192232043053\\
0.41934476437 0.193919204717\\
0.419635393077 0.195614405764\\
0.419915453467 0.197317674382\\
0.420184945541 0.199029039126\\
0.420443869298 0.200748528922\\
0.420692224738 0.202476173063\\
0.420930011862 0.204212001216\\
0.42115723067 0.205956043422\\
0.42137388116 0.207708330101\\
0.421579963334 0.20946889205\\
0.421775477192 0.211237760447\\
0.421960422733 0.213014966854\\
0.422134799957 0.214800543219\\
0.422298608864 0.21659452188\\
0.422451849455 0.218396935563\\
0.42259452173 0.220207817389\\
0.422726625687 0.222027200875\\
0.422848161329 0.223855119936\\
0.422959128653 0.225691608887\\
0.423059527661 0.227536702447\\
0.423149358352 0.229390435744\\
0.423228620727 0.231252844312\\
0.423297314785 0.233123964099\\
0.423355440526 0.235003831466\\
0.423402997951 0.236892483193\\
0.423439987059 0.238789956482\\
0.423466407851 0.240696288956\\
0.423482260326 0.242611518667\\
0.423487544484 0.244535684096\\
}\redundantfrontierdata
\addplot[blue!70!black,thick] table[x=nu,y=energy]{\redundantfrontierdata};
\addplot[blue!70!black,only marks,mark=o,mark options={fill=white}] coordinates
 {(0.212121212121,0.0104123281966)};
\addplot[blue!70!black,only marks,mark=*] coordinates
 {(0.423487544484,0.244535684096)};
\addplot[gray,dashed] coordinates {(0.423487544484,0) (0.423487544484,.27)};
\node[font=\small,anchor=south] at (axis cs:.10,.013) {无需探查};
\node[font=\small,anchor=south] at (axis cs:.285,.15) {校准与分配共同决定};
\node[font=\small,rotate=90,anchor=center] at (axis cs:.443,.13) {已知参数也能力不足};
\end{axis}
\end{tikzpicture}
\caption{固定同一个误差目标 $q=0.1$，只改变真实相位速率，最低校准预算即发生变化。
参数为式~\eqref{v44r:eq:numerical-parameters}；全程 $D_\star=0$。
$\nu=7/33$ 仍可零校准，右侧的空心点表示正噪声造成的右极限跳跃；
$\nu=119/281$ 的有限能量端点可达到，再提高速度便超出已知模型的执行能力。
曲线直接计算自闭式定理。}
\label{fig:v44:redundant-frontier}
\end{figure}

若需将系数探查落实为固定无漂移模型中的实际往返，
可使用连续速度观测
$\widetilde{\dot y}=ax+z+n(t)$，并要求共同噪声合同
\[
 \|n\|_{L^2(0,T)}\le\varepsilon+\eta\|x\|_{L^2(0,T)},
 \qquad \int_0^T[x(t)^2+z(t)^2]\,dt\le E.
\]
完整连续观测下的必要性证明采用同一比例噪声
$n_\pm=\pm(D_E/2)x$，仍适用于任意良定因果波形和停止。
对有限正门槛 $E_*=\mathcal E_{\rm red}(q,\nu)$，取单位零积分模式
\[
 \phi_T(t)=
 \begin{cases}
 1/\sqrt T,&0\le t<T/2,\\
 -1/\sqrt T,&T/2\le t\le T,
 \end{cases}
 \qquad x(t)=\sqrt{E_*}\phi_T(t),\quad z(t)=u_s(t)=0.
\]
它消耗的实际总二次输入恰为 $E_*$；
投影 $\int\phi_T\widetilde{\dot y}$ 给出同一校准报告和噪声界。
任务特征在终点回到起点，相位始终不变；
若输入确实实现实际积分坐标 $\dot\xi_x=x$，该坐标也在终点回位。
若校准幅值限制为 $|x|\le U_{\rm cal}$、$U_{\rm cal}>0$，
取 $T=E_*/U_{\rm cal}^2$ 即可达到；
峰值任务偏离不超过 $h\sqrt{E_*T}/2$，该完整轨迹须包含在已声明的校准域中。
这里没有附加加速度保证；若有硬加速度限制，可使用命题~\ref{v44d:prop:sine-two-approximation}的正弦回位构造并重新检查物理域。
时间比较要求上述同一 $L^2$ 噪声合同对所比较的时长和波形共同成立，
不能把固定逐帧位置噪声直接当作这个合同。


% V4.4 root derivation checkpoint. Standalone section fragment.
\section{全历史精确校准--推进资源前沿}
\label{sec:v44:scalar-frontier}

本节给出一个允许自适应探查、数据后非线性命令选择且具有输入能力约束的精确结果。一般矩阵情形中的名义无偏约束不用于本节。

\subsection{物理任务与信息模型}
在一个冻结窗口内，考虑两个实际输入和一个横向任务特征：
\begin{equation}
 \dot y=a u_1+u_2,\qquad \dot\sigma=u_2,
 \qquad 0<a_-<a_+,\quad a\in[a_-,a_+].
 \label{eq:v44:scalar-plant}
\end{equation}
这里 $\sigma$ 是实际材料相位。输入已在固定物理椭球度量中白化，能力为 $u_1^2+u_2^2\le U^2$。固定纠偏命令范围 $|v|\le r$，其中 $r>0$；所需真实推进速度为 $0\le\nu<U$。任务要求实际 $u_2=\nu$，并定义归一化最坏响应损失
\begin{equation}
 \mathcal R(\mathsf A;E,\nu)
 =\sup_{a,w}\sup_{|v|\le r}
 \frac{|a u_1(v;z)+\nu-v|}{r+\nu}.
 \label{eq:v44:scalar-risk}
\end{equation}
其中 $\mathsf A$ 同时包含探查协议及数据后的命令规则，后者允许非线性。所有执行命令须满足能力约束。

探查使用固定物理度量下的输入二次代价。具体地，取单位 $L_2$ 范数、支撑互不重叠的时间模式，$p_i$ 是实际输入在相应模式上的系数，因此 $\sum_i p_i^2=\int u_1(t)^2\,dt$。该量是控制二次代价；换算成电机能耗还需要相应物理模型。主定理先把 $E$ 定义为可实现的探查预算：允许输入集合至少包含一个范数为 $\sqrt E$ 的固定探查；其中 $U$ 约束数据后的执行。固定探查时间及瞬时幅值的限制由下文推论显式加入。模式及其测量须满足下面同一个总噪声合同；若窗口长度改变，噪声参数也须重新给定。用记录输入 $p=(p_1,\ldots,p_k)^\top$，扣除已知输入贡献后得到
\begin{equation}
 z_i=a p_i+w_i,\qquad
 \|p\|_2^2\le E,\qquad
 \|w\|_2\le\varepsilon_0+\eta_{\rm cal}\|p\|_2,
 \label{eq:v44:scalar-data}
\end{equation}
其中 $\varepsilon_0>0$、$\eta_{\rm cal}\ge0$ 固定。每个 $p_i$ 和停止时刻可以因果地依赖此前数据，但能量上限须对每个允许历史成立，且协议须有限终止。协议是确定性因果规则；若使用随机种子，以下最坏保证要求对每个种子实现成立，不涉及仅对随机性取期望的另一种博弈。$\eta_{\rm cal}$ 表示不能从本批数据中区分的输入相关测量或跨阶段模型偏差；若某固定余项在校准和执行中完全相同且可被识别，应合并到 $a$，不能再重复计入该偏差。

定义
\begin{align}
 q_{\rm prior}&=\frac{a_+-a_-}{a_++a_-},\\
 t(E)&=\frac{\varepsilon_0}{\sqrt E}+\eta_{\rm cal}\quad(E>0),
 \qquad t(0)=+\infty,\\
 M(\nu)&=\frac{\sqrt{U^2-\nu^2}}{r+\nu}.
\end{align}

\Needspace{11\baselineskip}
\begin{theorem}[全历史与自适应策略的精确资源前沿]
\label{thm:v44:scalar-exact}
在上述信息与物理模型中，所有允许协议的最优最坏损失为
\begin{equation}
 \boxed{\mathcal R_*(E,\nu)
 =\max\left\{(1-a_-M(\nu))_+,
 \min\left(q_{\rm prior},\frac{t(E)}{a_-+t(E)}\right)\right\}.}
 \label{eq:v44:scalar-exact-frontier}
\end{equation}
$E=0$ 时第二项取 $q_{\rm prior}$。该下界覆盖任意因果自适应探查、历史相关停止及数据后非线性命令；一个非自适应探查加区间反馈规则达到该界。
\end{theorem}

\begin{proof}
先证明上界。$E>0$ 时选择任意固定 $p$ 满足 $\|p\|_2^2=E$，令
\[
 \widehat a=\frac{p^\top z}{E},\qquad
 I(z)=[\widehat a-t(E),\widehat a+t(E)]\cap[a_-,a_+]
 =[\ell,h].
\]
由 Cauchy--Schwarz 不等式，$a\in I(z)$，所以该区间对每个允许历史非空。取
\begin{equation}
 c(z)=\min\left\{M(\nu),\frac2{\ell+h}\right\},
 \qquad u(v;z)=\begin{bmatrix}c(z)(v-\nu)\\\nu\end{bmatrix}.
 \label{eq:v44:scalar-policy}
\end{equation}
因为 $|v-\nu|\le r+\nu$，该命令满足输入能力。对任何正区间 $[\ell,h]$，直接最小化两端误差给出
\[
 \min_{|c|\le M}\max_{a\in[\ell,h]}|ac-1|
 =\max\left\{\frac{h-\ell}{h+\ell},\,1-\ell M\right\}.
\]
最优 $c$ 正是式~\eqref{eq:v44:scalar-policy}。区间下端满足 $\ell\ge a_-$，宽度不超过 $\min(a_+-a_-,2t(E))$。宽度比随宽度增大、随下端减小而增大，因此
\[
 \frac{h-\ell}{h+\ell}
 \le\frac{\min(a_+-a_-,2t(E))}{2a_-+\min(a_+-a_-,2t(E))}
 =\min\left\{q_{\rm prior},\frac{t(E)}{a_-+t(E)}\right\}.
\]
式~\eqref{eq:v44:scalar-policy}同时给出更强的逐命令保证
\[
 |a u_1+\nu-v|\le \mathcal R_*(E,\nu)|v-\nu|.
\]
$E=0$ 时直接使用先验区间作同样的选择。

再证明任意协议的下界。令
\[
 d=\min\left\{\frac{a_+-a_-}{2},\ t(E)\right\},\qquad
 a^{(0)}=a_-,\quad a^{(1)}=a_-+2d.
\]
向任意因果协议依次报告 $z_i=(a_-+d)p_i$。这唯一确定该协议的一条共同历史、同一输入序列及同一停止决定；需要其通常的良定性假设。中点 $a_-+d$ 本身属于先验，该历史也是它的零噪声历史，故必然有限终止且满足预算。对于两个端点参数，分别选 $w_i^{(0)}=d p_i$、$w_i^{(1)}=-d p_i$。沿该历史的 $P=\|p\|_2$ 满足 $P\le\sqrt E$，故
\[
 dP\le\left(\frac{\varepsilon_0}{\sqrt E}+\eta_{\rm cal}\right)P
 \le\varepsilon_0+\eta_{\rm cal}P.
\]
因此两个世界都满足同一噪声合同，且不能被任意自适应协议区分。执行时只取 $v=-r$，任何命令都对应某个实数 $c=-u_1/(r+\nu)$；无论该规则在其它 $v$ 上是否线性，其两世界最坏损失至少为
\[
 \min_{c\in\mathbb R}\max\{|a_-c-1|,|(a_-+2d)c-1|\}
 =\frac d{a_-+d}.
\]
这就是式~\eqref{eq:v44:scalar-exact-frontier}的第二项。输入约束还强制 $|c|\le M(\nu)$，故在 $a=a_-$ 时损失至少为 $(1-a_-M(\nu))_+$。两项下界与上界吻合。$E=0$ 时用两个先验端点即可。
\end{proof}

\begin{corollary}[精确探查能量与能力门槛]
\label{cor:v44:scalar-energy}
给定 $0\le q_0<1$，存在满足 $\mathcal R\le q_0$ 的允许协议，当且仅当
\begin{equation}
 M(\nu)\ge\frac{1-q_0}{a_-}
 \label{eq:v44:scalar-cap-threshold}
\end{equation}
且下列信息条件成立。若 $q_0\ge q_{\rm prior}$，无需探查；若 $0<q_0<q_{\rm prior}$，有限能量可行当且仅当
\[
 b(q_0):=\frac{a_-q_0}{1-q_0}-\eta_{\rm cal}>0,
 \qquad
 E\ge E_{\min}(q_0):=\frac{\varepsilon_0^2}{b(q_0)^2}.
\]
当 $q_0=0<q_{\rm prior}$ 时不存在有限能量解。若 $b(q_0)=0$，只能在 $E\to\infty$ 的极限趋近，不能把极限值当作有限可达。
\end{corollary}
\begin{proof}
将式~\eqref{eq:v44:scalar-exact-frontier}分别与 $q_0$ 比较。信息分支等价于 $t(E)\le a_-q_0/(1-q_0)$，直接反解即可。
\end{proof}

若使用能量已归一化的等长离散探查且 $|p_i|\le p_{\max}$，任何协议至少需 $\lceil E_{\min}/p_{\max}^2\rceil$ 个样本；没有额外输入限制时可用固定幅值序列达到该能量门槛。对时长为 $\Delta t$ 的常值实际输入，$p_{\max}=u_{\max}\sqrt{\Delta t}$。连续时间中 $|u_1|\le u_{\max}$ 则直接给 $T_{\rm cal}\ge E_{\min}/u_{\max}^2$；平滑、往返和状态管要求可以进一步增加这个时长。若样本时长或噪声预算随样本数改变，必须先把这些依赖代回式~\eqref{eq:v44:scalar-data}，不能默认重复采样缩小对抗性噪声。

\begin{corollary}[固定时间和探查幅值下的精确前沿]
\label{cor:v44:scalar-duration}
在时长 $T_{\rm cal}>0$ 的固定校准窗口内，令探查期间 $u_2=0$，
$|u_1(t)|\le U_{\rm cal}$，且总二次预算至多 $E$。
假定常值探查及其归一化积分观测满足式~\eqref{eq:v44:scalar-data}，
除此之外没有状态、回位或变化率约束。
则定理~\ref{thm:v44:scalar-exact}中的精确值应使用
\[
 E_{\rm eff}=\min\{E,U_{\rm cal}^2T_{\rm cal}\}.
\]
若校准期间另固定 $u_2=\nu_{\rm cal}$ 且总输入满足同一个球 $\|u\|\le U$，
则以 $U_{\rm cal}=\sqrt{U^2-\nu_{\rm cal}^2}$ 代入。
在同一个噪声合同对所有比较时长都成立时，达到指定可行 $q_0$ 的最短时间为
$E_{\min}(q_0)/U_{\rm cal}^2$；若噪声界随时间改变，则最短时间由
$E_{\min}(q_0;\varepsilon_0(T),\eta_{\rm cal}(T))\le U_{\rm cal}^2T$
的最小可行 $T$ 决定。
\end{corollary}
\begin{proof}
幅值限制使每条自适应历史的实际能量至多 $E_{\rm eff}$，故共同历史下界仍然成立。
常值输入 $u_1=\sqrt{E_{\rm eff}/T_{\rm cal}}$ 满足幅值约束，
在单位模式 $1/\sqrt{T_{\rm cal}}$ 上产生系数 $p=\sqrt{E_{\rm eff}}$，
因此达到同一上界。最短时间是该能量不等式的直接反解。
如果必须回位、保持额外状态管或满足有限加速度，应在允许探查集合中另加这些限制，
此时本推论保留必要下界，达到性由实际构造重新验收。
\end{proof}


\subsection{前沿的定量读取}
取固定归一化参数
\[
 [a_-,a_+]=[0.8,1.2],\quad r=0.04,\quad U=0.30,
 \quad\varepsilon_0=0.002,\quad\eta_{\rm cal}=0.01.
\]
先验误差为 $q_{\rm prior}=0.2$。若要求 $q_0=0.05$，则
\[
 E_{\min}=0.003880677\ldots,
 \qquad \nu\le\nu_{\rm cap}(0.05)=0.168829526\ldots.
\]
在 $\nu=0.12$、$E=0.01$ 时，精确最坏响应为 $0.0361446\ldots$。
在 $\nu=0.17$ 时，执行能力造成的误差下限约为 $0.05835$，因此无论增加多少探查都不能达到 $0.05$。
图~\ref{fig:v44:scalar-frontier}直接绘制定理的精确值；时间和幅值须按推论~\ref{cor:v44:scalar-duration}换成可实现预算。
第~\ref{v44:sec:finite-geometry-example}节进一步给出带回位、有限几何与采样的完整实际探查构造。

\begin{figure}[htbp]
\centering
\begin{tikzpicture}
\begin{semilogxaxis}[
 width=.88\linewidth,height=72mm,
 xlabel={探查二次预算 $E$},ylabel={精确最坏响应 $\mathcal R_*(E,\nu)$},
 xmin=1e-5,xmax=.1,ymin=0,ymax=.40,
 ytick={0,.05,.1,.2,.3,.4},grid=both,
 major grid style={gray!25},minor grid style={gray!12},
 legend style={font=\small,at={(0.97,0.65)},anchor=north east},
 tick label style={font=\small},label style={font=\small}]
\pgfplotstableread[row sep=\\]{
E qA qB qC\\
1e-05 0.2 0.2 0.3724283676\\
1.107756851e-05 0.2 0.2 0.3724283676\\
1.22712524e-05 0.2 0.2 0.3724283676\\
1.359356391e-05 0.2 0.2 0.3724283676\\
1.505836354e-05 0.2 0.2 0.3724283676\\
1.668100537e-05 0.2 0.2 0.3724283676\\
1.847849797e-05 0.2 0.2 0.3724283676\\
2.046968272e-05 0.2 0.2 0.3724283676\\
2.267543126e-05 0.2 0.2 0.3724283676\\
2.511886432e-05 0.2 0.2 0.3724283676\\
2.782559402e-05 0.2 0.2 0.3724283676\\
3.08239924e-05 0.2 0.2 0.3724283676\\
3.414548874e-05 0.2 0.2 0.3724283676\\
3.782489906e-05 0.2 0.2 0.3724283676\\
4.190079106e-05 0.2 0.2 0.3724283676\\
4.641588834e-05 0.2 0.2 0.3724283676\\
5.141751828e-05 0.2 0.2 0.3724283676\\
5.695810811e-05 0.2 0.2 0.3724283676\\
6.309573445e-05 0.2 0.2 0.3724283676\\
6.989473207e-05 0.2 0.2 0.3724283676\\
7.742636827e-05 0.2 0.2 0.3724283676\\
8.576958986e-05 0.2 0.2 0.3724283676\\
9.501185073e-05 0.2 0.2 0.3724283676\\
0.0001052500285 0.2 0.2 0.3724283676\\
0.0001165914401 0.1961606588 0.1961606588 0.3724283676\\
0.0001291549665 0.188628226 0.188628226 0.3724283676\\
0.0001430722989 0.1813395677 0.1813395677 0.3724283676\\
0.0001584893192 0.1742921193 0.1742921193 0.3724283676\\
0.0001755676291 0.1674828728 0.1674828728 0.3724283676\\
0.0001944862439 0.1609084176 0.1609084176 0.3724283676\\
0.000215443469 0.1545649811 0.1545649811 0.3724283676\\
0.0002386589787 0.1484484674 0.1484484674 0.3724283676\\
0.0002643761186 0.142554496 0.142554496 0.3724283676\\
0.0002928644565 0.1368784379 0.1368784379 0.3724283676\\
0.0003244226079 0.1314154507 0.1314154507 0.3724283676\\
0.0003593813664 0.1261605119 0.1261605119 0.3724283676\\
0.0003981071706 0.1211084505 0.1211084505 0.3724283676\\
0.0004410059454 0.1162539761 0.1162539761 0.3724283676\\
0.0004885273572 0.1115917074 0.1115917074 0.3724283676\\
0.0005411695265 0.1071161971 0.1071161971 0.3724283676\\
0.0005994842503 0.1028219566 0.1028219566 0.3724283676\\
0.0006640827851 0.09870347727 0.09870347727 0.3724283676\\
0.0007356422545 0.09475525086 0.09475525086 0.3724283676\\
0.0008149127469 0.09097178785 0.09097178785 0.3724283676\\
0.0009027251779 0.08734763388 0.08734763388 0.3724283676\\
0.001 0.0838773847 0.0838773847 0.3724283676\\
0.001107756851 0.08055569944 0.08055569944 0.3724283676\\
0.00122712524 0.07737731239 0.07737731239 0.3724283676\\
0.001359356391 0.07433704333 0.07433704333 0.3724283676\\
0.001505836354 0.07142980654 0.07142980654 0.3724283676\\
0.001668100537 0.06865061851 0.06865061851 0.3724283676\\
0.001847849797 0.06599460455 0.06599460455 0.3724283676\\
0.002046968272 0.06345700422 0.06345700422 0.3724283676\\
0.002267543126 0.06103317586 0.06103317586 0.3724283676\\
0.002511886432 0.05871860012 0.05871860012 0.3724283676\\
0.002782559402 0.05650888269 0.05834612615 0.3724283676\\
0.00308239924 0.05439975623 0.05834612615 0.3724283676\\
0.003414548874 0.05238708169 0.05834612615 0.3724283676\\
0.003782489906 0.05046684887 0.05834612615 0.3724283676\\
0.004190079106 0.04863517656 0.05834612615 0.3724283676\\
0.004641588834 0.04688831204 0.05834612615 0.3724283676\\
0.005141751828 0.04522263028 0.05834612615 0.3724283676\\
0.005695810811 0.04363463265 0.05834612615 0.3724283676\\
0.006309573445 0.04212094527 0.05834612615 0.3724283676\\
0.006989473207 0.04067831717 0.05834612615 0.3724283676\\
0.007742636827 0.03930361809 0.05834612615 0.3724283676\\
0.008576958986 0.03799383609 0.05834612615 0.3724283676\\
0.009501185073 0.03674607504 0.05834612615 0.3724283676\\
0.01052500285 0.03555755184 0.05834612615 0.3724283676\\
0.01165914401 0.03442559361 0.05834612615 0.3724283676\\
0.01291549665 0.03334763476 0.05834612615 0.3724283676\\
0.01430722989 0.03232121399 0.05834612615 0.3724283676\\
0.01584893192 0.0313439712 0.05834612615 0.3724283676\\
0.01755676291 0.0304136444 0.05834612615 0.3724283676\\
0.01944862439 0.02952806663 0.05834612615 0.3724283676\\
0.0215443469 0.02868516284 0.05834612615 0.3724283676\\
0.02386589787 0.02788294675 0.05834612615 0.3724283676\\
0.02643761186 0.02711951782 0.05834612615 0.3724283676\\
0.02928644565 0.02639305818 0.05834612615 0.3724283676\\
0.03244226079 0.02570182964 0.05834612615 0.3724283676\\
0.03593813664 0.02504417071 0.05834612615 0.3724283676\\
0.03981071706 0.02441849374 0.05834612615 0.3724283676\\
0.04410059454 0.02382328204 0.05834612615 0.3724283676\\
0.04885273572 0.02325708716 0.05834612615 0.3724283676\\
0.05411695265 0.02271852611 0.05834612615 0.3724283676\\
0.05994842503 0.02220627879 0.05834612615 0.3724283676\\
0.06640827851 0.0217190854 0.05834612615 0.3724283676\\
0.07356422545 0.02125574396 0.05834612615 0.3724283676\\
0.08149127469 0.02081510789 0.05834612615 0.3724283676\\
0.09027251779 0.02039608368 0.05834612615 0.3724283676\\
0.1 0.01999762866 0.05834612615 0.3724283676\\
}\frontierdata
\addplot[blue!75!black,thick] table[x=E,y=qA]{\frontierdata};
\addlegendentry{$\nu=0.12$}
\addplot[orange!80!black,thick] table[x=E,y=qB]{\frontierdata};
\addlegendentry{$\nu=0.17$}
\addplot[red!70!black,thick] table[x=E,y=qC]{\frontierdata};
\addlegendentry{$\nu=0.22$}
\addplot[black,dashed,domain=1e-5:.1,samples=2]{.05};
\addlegendentry{目标 $q_0=0.05$}
\addplot[gray,densely dotted] coordinates {(0.0038806772373,0) (0.0038806772373,.4)};
\end{semilogxaxis}
\end{tikzpicture}
\caption{同一任务的精确信息与能力边界。曲线为解析定理值，不使用仿真或实验数据。
竖虚线为 $q_0=0.05$ 的信息预算门槛；只有同时通过最终执行能力的相位速率才能达到该目标。}
\label{fig:v44:scalar-frontier}
\end{figure}

\subsection{从纠偏自由度到真实扫描自由度}
取 $a_0=(a_-+a_+)/2$，$H=[a_0,1]$、$N=[1,0]$。若只要求横向任务响应，选择 $B=(0,1)^\top$ 即得 $HB=1$、$NB=0$，故 $k_*=0$。要求纠偏不改变实际相位后，必须 $[0,1]B=0$，于是 $B=(1/a_0,0)^\top$，$k_{\rm phase}=1$。该变化由物理任务的附加要求引起，不能通过被动坐标变化消除。上面的能量门槛量化了这一个新增方向究竟需要多少信息代价。

\subsection{与二阶恢复及严格正速度的同一参数连接}
执行式~\eqref{eq:v44:scalar-policy}并令 $v=\zeta$，$r=\omega R$。若冻结所得 $c$ 满足 $|ac-1|\le q_0<1$，记 $\Delta=ac-1$，则
\[
 \dot e=(1+\Delta)\zeta-\nu\Delta+w,
 \qquad \dot\sigma=\nu.
\]
故定理~\ref{thm:cal-recovery}直接给出条件
\begin{equation}
 q_0\nu+d_0+\sqrt2\omega n_{\max}
 \le \frac{\omega(1-q_0)}2R.
 \label{eq:v44:scalar-progress-recovery}
\end{equation}
定义 $c_{\rm req}=(1-q_0)/a_-$。由式~\eqref{eq:v44:scalar-cap-threshold}，输入能力所允许的最大推进速度为
\begin{equation}
 \nu_{\rm cap}(q_0)
 =\frac{-c_{\rm req}^2r+
 \sqrt{(1+c_{\rm req}^2)U^2-c_{\rm req}^2r^2}}
 {1+c_{\rm req}^2},
 \label{eq:v44:scalar-cap-speed}
\end{equation}
其中需 $U\ge c_{\rm req}r$，正速度需严格不等式。若 $q_0>0$ 且恢复余量为正，还可取
\[
 \nu\le\nu_{\rm rec}(q_0)
 :=\frac{\omega(1-q_0)R/2-d_0-\sqrt2\omega n_{\max}}{q_0}.
\]
因此，给定 $E$ 后，可在所有满足 $E\ge E_{\min}(q_0)$ 的 $q_0$ 上选择正的
\[
 \nu<\min\{\nu_{\rm cap}(q_0),\nu_{\rm rec}(q_0)\}.
\]
若固定时间与探查幅值另有限制，上述 $E$ 使用推论~\ref{cor:v44:scalar-duration}中的 $E_{\rm eff}$。
式~\eqref{eq:v44:scalar-exact-frontier}是精确静态资源前沿；加入公共 Lyapunov 恢复后得到的速度是一个明确充分保证，不声称为所有动态控制器的最快扫描速度。

\paragraph{非零执行余项。}
若实际横向响应为 $\dot y=a u_1+u_2+e_1(t)u_1+e_2(t)u_2+d_y(t)$，其中 $\|(e_1,e_2)\|_2\le\eta_{\rm exe}$、$|d_y|\le d_y^{\max}$，则对固定 $c$ 可用
\[
 \bar q=q_0+\eta_{\rm exe}c,\qquad
 \alpha_s=q_0+\eta_{\rm exe}\sqrt{1+c^2},\qquad
 d_0=d_y^{\max}
\]
代入通用恢复与推进定理。若另外 $\dot\sigma=u_2+d_\sigma$、$|d_\sigma|\le\bar d_\sigma$，真实推进下界为 $\nu-\bar d_\sigma$。相位漂移和横向余项分别进入各自合同，只计一次。


% Independent V4.4 multichannel resource theorem.
% Fragment only; uses the main manuscript's theorem environments and packages.
% Scope: independent positive diagonal gains, rowwise deterministic total noise,
% task-command box, fixed physical phase speed, and speed-level execution.
% No claim that the number of unknown diagonal gains equals probe-direction rank.

\section{共享探查预算下的多通道精确校准--推进前沿}
\label{v44m:sec:multichannel}

本节给出一个允许任意因果探查分配和停止的多通道资源定理。
其不确定集合是独立正对角增益，不是任意矩阵谱范数球。
精确性同时涉及最坏历史的总探查能量与执行期的共享输入能力；
它不要求控制器在看到数据前固定纠偏分配，也不要求最终规则线性或名义无偏。

\subsection{模型、信息合同与量词}

固定 $m\ge1$ 个任务通道，执行期模型为
\begin{equation}
 \dot y_j=a_j u_j+b_j u_s,\qquad
 \dot\sigma=u_s=\nu,\qquad j=1,\ldots,m,
 \label{v44m:eq:plant}
\end{equation}
其中 $b_j\in\mathbb R$ 已知，$\nu\ge0$ 是校准前固定的真实材料推进速度。
未知增益固定于本次校准和执行窗口，且其先验是笛卡尔积
\begin{equation}
 a_j\in[\ell_j,h_j],\qquad 0<\ell_j\le h_j<\infty.
 \label{v44m:eq:prior}
\end{equation}
允许的任务命令为 $|v_j|\le r_j$，其中 $r_j>0$。
执行输入满足同一物理白化后的能力球
\begin{equation}
 \sum_{j=1}^m u_j^2+\nu^2\le U^2,
 \qquad U>0.
 \label{v44m:eq:execution-cap}
\end{equation}
定义每行的归一化任务幅值
\begin{equation}
 s_j(\nu)=r_j+|b_j|\nu>0.
\end{equation}
给定目标 $q\in[0,1)$，要求校准后对全部命令箱均有
\begin{equation}
 \max_{1\le j\le m}
 \frac{|a_j u_j(v)-[v_j-b_j\nu]|}{s_j(\nu)}\le q.
 \label{v44m:eq:rowbox-objective}
\end{equation}
这是按行归一化的最大误差合同，不是 Euclidean 总输出误差合同。

在校准阶段令 $u_s=0$，或精确扣除已知相位输入的响应。
一次探查可以同时作用于全部通道。写第 $j$ 行的实际探查系数及报告为
\begin{equation}
 z_{ji}=a_j p_{ji}+w_{ji},\qquad
 p_j=(p_{j1},\ldots,p_{jk}),\qquad
 w_j=(w_{j1},\ldots,w_{jk}).
 \label{v44m:eq:calibration-oracle}
\end{equation}
整个实际历史满足行数据合同
\begin{equation}
 \|w_j\|_2\le\varepsilon_j+\eta_j\|p_j\|_2,
 \qquad \varepsilon_j>0,\quad \eta_j\ge0,
 \label{v44m:eq:row-noise}
\end{equation}
以及共享二次探查预算
\begin{equation}
 \sum_{j=1}^m\|p_j\|_2^2\le E.
 \label{v44m:eq:shared-energy}
\end{equation}
这些是整个历史上的确定性范数合同，不是独立随机噪声假设，也不是固定逐样本幅度界。

每个探查向量 $(p_{1i},\ldots,p_{mi})$ 可以由此前全部行的报告决定，
停止时刻也可以自适应。允许任意最终共同规则 $u(v;\text{history})$。
所要求的保证覆盖先验中的每个增益向量及每个允许的噪声历史，
且预算对每个历史都成立，算法必须在每个有效历史上结束校准。
下文的任意算法指确定性因果算法，或固定随机种子后的逐历史保证；
不据此宣称随机算法的期望风险或独立噪声统计复杂度。

\subsection{精确的必要与充分资源条件}

记无数据时各行的相对先验宽度及相应噪声地板为
\begin{equation}
 q_{{\rm prior},j}=\frac{h_j-\ell_j}{h_j+\ell_j},
 \qquad
 q_{{\rm floor},j}=\frac{\eta_j}{\ell_j+\eta_j}.
 \label{v44m:eq:prior-and-floor}
\end{equation}
对 $q\in[0,1)$，定义活跃行集合
\begin{equation}
 \mathcal A(q)=\{j:q<q_{{\rm prior},j}\},
 \qquad d_j(q)=\frac{\ell_jq}{1-q}.
 \label{v44m:eq:active-rows}
\end{equation}
这里的活跃表示先验精度尚不足以满足当前目标，不表示必须分时执行不同探查。
各行的精确能量门槛定义为
\begin{equation}
 e_j(q)=
 \begin{cases}
 0,&q\ge q_{{\rm prior},j},\\[1mm]
 \displaystyle\frac{\varepsilon_j^2}{[d_j(q)-\eta_j]^2},
     &q<q_{{\rm prior},j}\ \text{且}\ d_j(q)>\eta_j,\\[3mm]
 +\infty,&q<q_{{\rm prior},j}\ \text{且}\ d_j(q)\le\eta_j.
 \end{cases}
 \qquad \mathcal E_*(q)=\sum_{j=1}^m e_j(q).
 \label{v44m:eq:row-energy-threshold}
\end{equation}

\begin{theorem}[多通道共享资源的全部历史精确前沿]
\label{v44m:thm:exact-frontier}
在式~\eqref{v44m:eq:plant}--\eqref{v44m:eq:shared-energy} 的模型与量词下，
存在满足全部历史预算、全部任务命令响应以及执行能力的策略，当且仅当
\begin{equation}
 \boxed{
 E\ge\mathcal E_*(q),\qquad
 \nu^2+(1-q)^2
   \sum_{j=1}^m\frac{[r_j+|b_j|\nu]^2}{\ell_j^2}
 \le U^2.}
 \label{v44m:eq:exact-resource-frontier}
\end{equation}
因此，因果地改变各行探查能量及自适应停止不能降低该最坏历史总能量。
达到构造可以预先固定各行能量，在同一次向量探查中取得数据，
然后根据实际后验区间重新选择执行分配。
\end{theorem}

\begin{proof}
先证充分性。若 $\mathcal E_*(q)<\infty$，对每个活跃行选择一项探查
\begin{equation}
 p_j=\sqrt{e_j(q)}
 =\frac{\varepsilon_j}{d_j(q)-\eta_j}>0,
 \label{v44m:eq:attaining-probe}
\end{equation}
所有行可以同时执行；非活跃行不探查。
令 $\widehat a_j=z_j/p_j$，活跃行的数据一致区间为
\begin{equation}
 [L_j,H_j]
 =[\ell_j,h_j]\cap
 [\widehat a_j-d_j(q),\widehat a_j+d_j(q)].
 \label{v44m:eq:posterior-interval}
\end{equation}
真实增益属于该非空区间，因为
$|\widehat a_j-a_j|\le\varepsilon_j/p_j+\eta_j=d_j(q)$。
非活跃行取 $[L_j,H_j]=[\ell_j,h_j]$。
各行均满足
\begin{equation}
 \frac{H_j-L_j}{H_j+L_j}\le q,
 \qquad\text{等价地}\qquad
 \frac{H_j}{L_j}\le\frac{1+q}{1-q}.
 \label{v44m:eq:posterior-ratio}
\end{equation}
活跃行使用 $L_j\ge\ell_j$、$H_j-L_j\le2d_j(q)$ 得到该式；
非活跃行由先验宽度条件得到。

在看到实际数据后选择
\begin{equation}
 \boxed{
 c_j=\frac{1-q}{L_j},\qquad
 u_j(v)=c_j(v_j-b_j\nu),\qquad u_s=\nu.}
 \label{v44m:eq:attaining-controller}
\end{equation}
式~\eqref{v44m:eq:posterior-ratio} 给出
$1-q\le a_jc_j\le1+q$，所以该构造甚至满足更强的逐点合同
\begin{equation}
 |a_ju_j(v)-[v_j-b_j\nu]|
 \le q|v_j-b_j\nu|
 \le q s_j(\nu).
\end{equation}
又有 $c_j\le(1-q)/\ell_j$，故
\begin{equation}
 \sum_j u_j(v)^2+\nu^2
 \le(1-q)^2\sum_j s_j(\nu)^2/\ell_j^2+\nu^2\le U^2.
\end{equation}
能量恰为 $\mathcal E_*(q)$。活跃集为空时，直接跳过探查使用先验区间，
同一控制律仍成立。这完成全部历史的充分性证明。

再证能力必要性。取实际增益向量 $a_j=\ell_j$，并考虑使
$|v_j-b_j\nu|=s_j(\nu)$ 同时成立的命令箱角点。
例如 $b_j\ne0$ 时取 $v_j=-r_j\operatorname{sgn}(b_j)$，
$b_j=0$ 时任选一个端点。
任何满足目标 $q<1$ 的共同动作都必须满足
\begin{equation}
 |u_j|\ge\frac{(1-q)s_j(\nu)}{\ell_j}.
\end{equation}
所有分量在同一个命令角点被同时约束，因而得到式~\eqref{v44m:eq:exact-resource-frontier}
的能力不等式。此论证不要求最终规则线性。

最后证共享能量必要性。反设 $E<\mathcal E_*(q)$。
此时 $\mathcal A(q)$ 非空。因行数有限，可选择 $q'>q$，使
\begin{equation}
 q'<\min_{j\in\mathcal A(q)}q_{{\rm prior},j},
 \qquad \sum_{j\in\mathcal A(q)}e_j(q')>E.
 \label{v44m:eq:strict-target-choice}
\end{equation}
这一选择在无穷门槛处也成立：若 $d_j(q)<\eta_j$，足够近的 $q'$ 仍给出无穷门槛；
若 $d_j(q)=\eta_j$，则 $e_j(q')\to+\infty$；
若 $d_j(q)>\eta_j$，则 $e_j(q')\to e_j(q)$。
因此有限和在 $q'\downarrow q$ 时趋于 $\mathcal E_*(q)$，而活跃集保持不变。

对每个活跃行定义中点
\begin{equation}
 a_j^0=\ell_j+d_j(q').
\end{equation}
非活跃行取任意允许的固定增益。
在这个允许的中点增益向量和零噪声下运行所考虑的任意因果算法。
算法必须结束，形成一个实际历史，报告满足 $z_j=a_j^0p_j$，
各行实际能量记为 $E_j=\|p_j\|_2^2$，且 $\sum_jE_j\le E$。
由式~\eqref{v44m:eq:strict-target-choice}，存在一个活跃行 $j$ 满足
$E_j<e_j(q')$。

只改变这一行的真实增益，构造两个允许模型
\begin{equation}
 a_j^- =\ell_j,\qquad
 a_j^+=\ell_j+2d_j(q'),
 \label{v44m:eq:indistinguishable-pair}
\end{equation}
其它行保持 $a_k^0$。因为 $q'<q_{{\rm prior},j}$，两端点都属于原先验。
在该行分别取噪声
\begin{equation}
 w_j^\pm=(a_j^0-a_j^\pm)p_j,
\end{equation}
其它行噪声为零。这两个模型产生相同的全部报告。
若 $d_j(q')\le\eta_j$，两套噪声对任意有限 $p_j$ 都满足数据合同；
若 $d_j(q')>\eta_j$，则能量不足给出
\begin{equation}
 [d_j(q')-\eta_j]\|p_j\|_2<\varepsilon_j,
\end{equation}
也即 $\|w_j^\pm\|_2<\varepsilon_j+\eta_j\|p_j\|_2$。
该构造还对每个前缀成立，因为前缀范数不超过最终范数。
噪声逐步按当前探查乘以同一个常数产生，因此不需要预知未来探查。
由因果性，两模型诱导完全相同的下一次探查、停止时刻和最终共同动作。

选择第 $j$ 行的任务命令端点，使 $|v_j-b_j\nu|=s_j(\nu)$。
任何共同动作在两模型下的归一化误差至少为
\begin{align}
 \inf_{c\in\mathbb R}
 \max\{|\ell_jc-1|,\ |[\ell_j+2d_j(q')]c-1|\}
 &=\frac{d_j(q')}{\ell_j+d_j(q')}\\
 &=q'>q.
 \label{v44m:eq:pair-relative-error}
\end{align}
这里把实际 $u_j$ 除以该非零命令端点即可得到标量 $c$，
故下界覆盖任意数据后非线性规则；输入能力限制只会进一步缩小可用动作集合。
至少一个允许模型违反目标合同，矛盾。必要性得证。
\end{proof}

\subsection{边界、噪声地板与精确最小最大值}

先验边界按 $q\ge q_{{\rm prior},j}$ 接受，因此一行的门槛在该点降为零。
这不是连续近似或忽略常数：目标一旦允许先验的最坏误差，就真的无需获得该行数据。
当某一行仍活跃时，有限能量可行的严格条件是
\begin{equation}
 q>q_{{\rm floor},j}.
\end{equation}
因为 $\varepsilon_j>0$，在地板等号处有限探查能量不够；
若地板不低于先验误差，则有限能量不能在该行上优于先验。

定义信息预算的精确响应门槛
\begin{equation}
 q_I(E)=\min\{q\in[0,1):\mathcal E_*(q)\le E\}.
 \label{v44m:eq:information-risk}
\end{equation}
该最小值存在：可行集非空，$\mathcal E_*(q)$ 单调不增且右连续，
各活跃行在有限预算下不能从上方达到其不可达地板。
在 $0\le\nu\le U$ 时，再定义
\begin{equation}
 q_C(\nu)=
 \left[1-
 \frac{\sqrt{U^2-\nu^2}}
 {\sqrt{\sum_j s_j(\nu)^2/\ell_j^2}}
 \right]_+.
 \label{v44m:eq:capacity-risk}
\end{equation}
因此，对式~\eqref{v44m:eq:rowbox-objective} 的全部历史最小最大风险，
\begin{equation}
 \boxed{q_*(E,\nu)=\max\{q_I(E),q_C(\nu)\}.}
 \label{v44m:eq:full-minimax-risk}
\end{equation}
当 $\nu=U$ 时右侧为 $1$，以 $u_j=0$ 达到；当 $\nu>U$ 时相位输入本身不可执行。
当 $E=0$ 时，$q_I(0)=\max_jq_{{\rm prior},j}$。此外，
\begin{equation}
 \lim_{E\to\infty}q_I(E)
 =\max_j\min\{q_{{\rm prior},j},q_{{\rm floor},j}\}.
 \label{v44m:eq:infinite-energy-floor}
\end{equation}
该极限不自动在有限能量达到；是否达到仍由式~\eqref{v44m:eq:row-energy-threshold}判断。

严格正的 $\varepsilon_j$ 是本定理关于达到性的重要假设。
若改成 $\varepsilon_j=0$，任意小但非零的激励可能已提供所需信息；
这时零能量往往只是下确界，而真正不探查仍不能区分先验中的多个增益。
因此不能把无噪声极限机械代入式~\eqref{v44m:eq:row-energy-threshold}后保留原来的达到性结论。

\subsection{统一正速度的显式边界}

对固定目标 $q<1$，令
\begin{align}
 A_q&=1+(1-q)^2\sum_j b_j^2/\ell_j^2,\\
 B_q&=(1-q)^2\sum_j r_j|b_j|/\ell_j^2,\\
 C_q&=(1-q)^2\sum_j r_j^2/\ell_j^2-U^2.
 \label{v44m:eq:speed-polynomial}
\end{align}
能力合同等价于 $A_q\nu^2+2B_q\nu+C_q\le0$。
若 $C_q>0$，连零推进也不能满足该响应目标；若 $C_q=0$，仅允许零推进。
若 $C_q<0$，统一正速度的精确上限为
\begin{equation}
 \boxed{
 \nu_{\max}(q)
 =\frac{-B_q+\sqrt{B_q^2-A_qC_q}}{A_q}>0.}
 \label{v44m:eq:uniform-speed-root}
\end{equation}
因此全部历史均可采用同一常速度 $\nu>0$，当且仅当
\begin{equation}
 \mathcal E_*(q)\le E,\qquad
 0<\nu\le\nu_{\max}(q),
\end{equation}
并同时满足其它已声明的速度上限。信息预算决定可达到的目标误差，
执行能力进一步决定在该精度下能维持的真实材料推进速度。

\subsection{总能量合同的连续实现与真实返回}

上述系数模型可由一个明确的连续观测合同实现。
在校准期假定能够观测任务速度信号
\begin{equation}
 \widetilde z_j(t)=a_j u_j(t)+n_j(t),\qquad
 \|n_j\|_{L^2(0,T)}\le
 \varepsilon_j+\eta_j\|u_j\|_{L^2(0,T)}.
 \label{v44m:eq:continuous-noise-oracle}
\end{equation}
这里 $\varepsilon_j$ 对所比较的时间与波形共同有效。
将 $\ell^2$ 换成 $L^2$ 后，定理~\ref{v44m:thm:exact-frontier} 的下界证明完全保留：
公共中点报告为 $\widetilde z_j(t)=a_j^0u_j(t)$，
两端点噪声为 $\pm d_j(q')u_j(t)$，每个时间前缀也满足合同。
因而该下界覆盖在此连续速度观测模型内的任意因果波形选择与停止，
不只覆盖预先指定的回归系数。

现在给出一个既达到能量又真实返回的构造。
假定 $0<\mathcal E_*(q)<\infty$，选定 $T>0$，令
\begin{equation}
 \phi_T(t)=
 \begin{cases}
 T^{-1/2},&0\le t<T/2,\\
 -T^{-1/2},&T/2\le t\le T,
 \end{cases}
 \quad
 u_j(t)=\sqrt{e_j(q)}\,\phi_T(t),\qquad u_s(t)=0.
 \label{v44m:eq:returning-waveform}
\end{equation}
它满足
\begin{equation}
 \int_0^T\phi_T^2dt=1,\qquad
 \int_0^T\phi_Tdt=0,
 \quad\int_0^T\sum_j u_j(t)^2dt=\mathcal E_*(q).
\end{equation}
以 $z_j=\int_0^T\phi_T(t)\widetilde z_j(t)\,dt$ 构造数据，
则 $z_j=a_j\sqrt{e_j(q)}+w_j$，且
$|w_j|\le\varepsilon_j+\eta_j\sqrt{e_j(q)}$。
这正是达到构造~\eqref{v44m:eq:attaining-probe}所需的数据。
返回半程同样参与信息计算，没有把它当作额外无信息回位。

由于校准期的真实系统无未建模漂移、增益固定且 $u_s=0$，
\begin{equation}
 y_j(T)=y_j(0),\qquad \sigma(T)=\sigma(0).
 \label{v44m:eq:exact-return}
\end{equation}
若物理探查速度限为 $\sum_j u_j(t)^2\le U_{\rm cal}^2$，则该构造要求
\begin{equation}
 T\ge\frac{\mathcal E_*(q)}{U_{\rm cal}^2},
 \qquad U_{\rm cal}>0.
 \label{v44m:eq:calibration-time}
\end{equation}
取等号即可在速度级合同内达到。
每行的峰值任务偏离和输入累计位移分别不超过
\begin{equation}
 \frac{h_j\sqrt{e_j(q)T}}2,
 \qquad \frac{\sqrt{e_j(q)T}}2.
 \label{v44m:eq:calibration-wide-domain}
\end{equation}
这些量须处于预先声明的校准宽域；不能由返回终点相同推断中间路径可行。
若还有限制加速度或开关频率，应另行验收新的波形；
本构造采用已声明的速度级模型，不能把瞬时换向称为加速度连续。
当 $\mathcal E_*(q)=0$ 时，跳过整个校准阶段即可。

\begin{corollary}[固定扫描速度的两阶段最短时间]
\label{v44m:cor:two-stage-deadline}
在连续噪声合同~\eqref{v44m:eq:continuous-noise-oracle}、速度级探查界和
能够容纳式~\eqref{v44m:eq:calibration-wide-domain} 的共同模型域内，
考虑以下两阶段协议类：校准期实际相位保持不变；随后使用校准前已固定的常速度
$\nu>0$，完成真实相位长度 $L_\sigma>0$，并对全部任务命令保持目标 $q$。
若式~\eqref{v44m:eq:exact-resource-frontier} 的能力条件成立且
$\mathcal E_*(q)<\infty$，该协议类的最小最坏完成时间恰为
\begin{equation}
 \boxed{
 T_*(q,\nu)
 =\frac{\mathcal E_*(q)}{U_{\rm cal}^2}
    +\frac{L_\sigma}{\nu}.}
 \label{v44m:eq:two-stage-deadline}
\end{equation}
允许设计者预先选择统一常速度时，在没有更小外部速度限的情形下取
$\nu=\nu_{\max}(q)$ 达到该协议类的最短时间。
\end{corollary}
\begin{proof}
如果某协议的最坏总时长小于式~\eqref{v44m:eq:two-stage-deadline}，
因为扫描期速度固定且相位长度固定，其每个历史的校准时间都小于
$\mathcal E_*(q)/U_{\rm cal}^2$。
速度界将给出每个历史都严格小于 $\mathcal E_*(q)$ 的统一能量预算，
与连续版本的资源必要性矛盾。
反过来，使用式~\eqref{v44m:eq:returning-waveform}取
$T=\mathcal E_*(q)/U_{\rm cal}^2$，其校准后状态和相位返回起点，
然后采用式~\eqref{v44m:eq:attaining-controller}，以常速度 $\nu$ 推进，
即达到所给时间、全部历史响应及共享能力合同。
\end{proof}

本推论采用校准前固定速度；定理~\ref{v44d:thm:adaptive-speed-deadline}将其扩展到数据后及当前命令依赖的速度选择，并证明两个阶段下界可在同一个模型上相加。边探查边推进、失锁或额外观测开销需各自的附加合同。
若使用位置帧推导回归噪声，$\varepsilon_j$ 往往随 $T$ 和帧处理方式变化，
此时必须代入真实噪声预算重新分析，不能保留固定 $\varepsilon_j$ 的时间前沿。
实际成像质量、任务状态恢复及有限材料覆盖，仍须按正文对应合同与此端口结论组合。

\subsection{与反馈恢复及方向数的准确关系}

达到策略的纠偏矩阵和前馈可以写成
\begin{equation}
 C=\begin{bmatrix}\operatorname{diag}(c_j)\\0\ \cdots\ 0\end{bmatrix},
 \qquad
 u_{\rm ff}=
 \begin{bmatrix}-\nu\operatorname{diag}(c_j)b\\\nu\end{bmatrix}.
\end{equation}
因此纠偏严格相位中性；令
$\Delta=\operatorname{diag}(a_jc_j-1)$，则 $\|\Delta\|_2\le q$。
执行响应满足
\begin{equation}
 \dot y=(I_m+\Delta)v-\nu\Delta b.
\end{equation}
所以可以把同一个 $q$ 用于二阶恢复证明，并把
$-\nu\Delta b$ 作为推进产生的响应残差验收。
不能只使用 $\Delta$ 的恢复率而把这一前馈项遗漏。

独立未知标量数、活跃行数与独立探查方向数不是同一个量。
由于每行输出单独可测且先验为对角结构，一条各活跃分量非零的物理输入方向，
或式~\eqref{v44m:eq:returning-waveform}的单一共同时间模态，
可以同时提供多个增益的信息。
定理~\ref{v44m:thm:exact-frontier}证明的是共享能量的精确加和门槛，
没有证明必须进行 $m$ 个独立输入方向的探查。
任意 $\Theta$ 谱范数球中的方向数结论不能搬到这个对角先验上。

同样，本模型只有一个已知相位输入列 $b$；
不要求相位中性时，它只提供一个已知任务方向，通常不能无数据完成全部 $m$ 维任务。
因此标量例中的零校准与一方向校准分离，不应被扩大成这里的 $0\to m$ 方向数结论。
本节新增的多通道内容是：面对真实共享输入能力和预算，
任何因果探查都不能绕过各活跃行误差目标所需的总能量，
而数据后分配、真实返回及统一正推进又能同时达到该门槛。

\paragraph{两个防止越界解释的反例。}
若增益不是独立先验而是已知 $a_1=a_2$，对一个通道的探查可以同时约束两个增益，
上述逐行能量加和一般不再必要。
若任务命令从箱形集合改成 Euclidean 单位球，取 $b=0$、$a_j=1$，
能力需求为 $1-q$，而非命令箱给出的 $\sqrt m(1-q)$。
这些改变都形成不同的问题；它们不属于定理的量词。


% Independent adversarial-audit derivation for V4.4.
% Read after v44_multichannel_notes.tex. Uses its model and theorem labels.
% This is a source fragment for integration, not an edit to the main manuscript.
% No general nonlinear-system time-optimality or acceleration claim is intended.

\section{同一个最坏模型、两阶段截止时间与有限加速度回位}
\label{v44d:sec:deadline}

本节仍使用第~\ref{v44m:sec:multichannel}节的独立正对角增益模型、
逐行确定性总噪声合同及整个任务命令箱。
记
\[
 \mathcal A(q)=\{j:q<q_{{\rm prior},j}\},\qquad
 d_j=\frac{\ell_jq}{1-q},\qquad
 \mathcal E_*(q)=\sum_j e_j(q).
\]
校准产生的全部增益信息均由所声明的观测合同涵盖；
不另行提供该合同之外的无噪声响应或参数观测。
执行段第一笔动作发出前及之后每一时刻，都要求对整个命令箱满足
式~\eqref{v44m:eq:rowbox-objective} 的统一响应合同。
这里要求的是每个动作执行前成立的保证，不能先允许一个未满足该合同的辨识动作，
再仅从其后的时刻开始计执行保证。
探查与执行规则均取良定的因果规则，预算和有限终止要求覆盖每个允许历史。

\subsection{各行门槛在同一个最小增益模型上同时必要}

仅分别证明最大校准时间与最大扫描时间，未必允许把二者相加。
下面给出一个更强的共同最坏模型：全部真实增益同时取各自的最小值，
并由一个因果噪声规则迫使每行支付其信息门槛。

\begin{lemma}[同一最小增益模型上的逐行必要能量]
\label{v44d:lem:common-low-model}
固定目标 $q\in[0,1)$。
对任何保证全部模型、全部噪声历史和整个任务命令箱满足目标 $q$ 的协议，
存在以
\begin{equation}
 a_j=\ell_j,\qquad j=1,\ldots,m
 \label{v44d:eq:common-low-model}
\end{equation}
为真实参数的允许因果噪声历史，使其在结束校准时必须同时满足
\begin{equation}
 \boxed{\ \|p_j\|_2^2\ge e_j(q)\quad\text{对每个 }j.\ }
 \label{v44d:eq:rowwise-energy-on-low-model}
\end{equation}
若某个 $e_j(q)=+\infty$，则该保证不存在有限终止的校准协议。

同一结论适用于连续速度观测
\begin{equation}
 \widetilde z_j(t)=a_j u_j(t)+n_j(t),\qquad
 \|n_j\|_{L^2(0,T)}
 \le\varepsilon_j+\eta_j\|u_j\|_{L^2(0,T)},
 \label{v44d:eq:continuous-oracle}
\end{equation}
此时将式~\eqref{v44d:eq:rowwise-energy-on-low-model} 中的
$\|p_j\|_2^2$ 替换为 $\int_0^{T_{\rm cal}}u_j(t)^2\,dt$。
噪声合同须对全部比较的波形、停止时间和观测历史共同有效。
\end{lemma}

\begin{proof}
非活跃行取零噪声，它们的门槛为零。
先考虑活跃且 $e_j(q)<\infty$ 的行。此时
$d_j>\eta_j$，并令
\begin{equation}
 P_j^*=\sqrt{e_j(q)}=\frac{\varepsilon_j}{d_j-\eta_j},
 \qquad d_jP_j^*=\varepsilon_j+\eta_jP_j^*.
 \label{v44d:eq:row-threshold}
\end{equation}
设第 $i$ 次探查后该行的累计系数范数为
$P_{j,i}=(\sum_{k\le i}p_{jk}^2)^{1/2}$。
真实增益始终是 $\ell_j$。因果地选择
\begin{equation}
 w_{ji}=
 \begin{cases}
 d_jp_{ji},&P_{j,i}\le P_j^*,\\
 0,&P_{j,i}>P_j^*.
 \end{cases}
 \label{v44d:eq:threshold-adversary}
\end{equation}
累计能量单调不减，因此第一笔越过门槛的探查以及其后的所有探查都取零噪声。
本次输入已确定后再生成本次报告，故该规则不预知未来探查。
在任意前缀，噪声满足
\[
 \|w_j\|_2\le d_j\min\{P_j,P_j^*\}.
\]
当 $P_j\le P_j^*$ 时，由
$(d_j-\eta_j)P_j\le\varepsilon_j$ 得到所需合同；
当 $P_j>P_j^*$ 时，由
$d_jP_j^*=\varepsilon_j+\eta_jP_j^*
 \le\varepsilon_j+\eta_jP_j$ 得到同一合同。
因此这是一条在每个前缀都合法的真实历史。

假设协议结束时某个活跃行仍有 $P_j<P_j^*$。
这一行从未越过门槛，其全部报告因而为
\begin{equation}
 z_j=(\ell_j+d_j)p_j.
 \label{v44d:eq:unsaturated-midpoint-data}
\end{equation}
由 $q<q_{{\rm prior},j}$，
$h_j-\ell_j-2d_j>0$。
若 $P_j>0$，又有
$\varepsilon_j/P_j+\eta_j-d_j>0$。
因此可以选择 $\kappa_j>0$，使
\begin{equation}
 \kappa_j<h_j-\ell_j-2d_j,\qquad
 \kappa_j<\varepsilon_j/P_j+\eta_j-d_j,
 \label{v44d:eq:strict-alternative-margin}
\end{equation}
其中 $P_j=0$ 时只需第一项。
只把该行真实参数改成
\begin{equation}
 a'_j=\ell_j+2d_j+\kappa_j
 \label{v44d:eq:strict-alternative-gain}
\end{equation}
并取 $w'_j=-(d_j+\kappa_j)p_j$，
就会得到完全相同的全部报告；其它行保持原参数和原报告。
式~\eqref{v44d:eq:strict-alternative-margin} 保证新噪声也满足合同。
两世界诱导相同的跨行探查、停止决定与首次执行规则。

选择任务箱角点，使第 $j$ 行有
$|v_j-b_j\nu|=r_j+|b_j|\nu$。
当 $b_j\ne0$ 时可取 $v_j=-r_j\operatorname{sgn}(b_j)$；
$b_j=0$ 时任选一个端点。
该角点选择与任何非负的 $\nu$ 无关，因此即使推进速度由数据或当前任务命令决定，
两个世界在同一报告与同一命令下也选择相同的 $\nu$。
对任何共同动作，两模型的归一化最坏误差至少为
\begin{equation}
 \inf_{c\in\mathbb R}
 \max\{|\ell_jc-1|,\ |a'_jc-1|\}
 =\frac{2d_j+\kappa_j}{2\ell_j+2d_j+\kappa_j}
 >\frac{d_j}{\ell_j+d_j}=q.
 \label{v44d:eq:strict-pair-lower}
\end{equation}
这与统一响应合同矛盾，故该历史上每个活跃行必须支付 $e_j(q)$。
输入能力限制只会进一步缩小共同动作集合，不削弱此必要性。

若活跃行的 $e_j(q)=+\infty$，则 $d_j\le\eta_j$。
在该行永远选择 $w_{ji}=d_jp_{ji}$，
其噪声对任意有限前缀均合法。
任何有限终止时刻仍有
$\varepsilon_j/P_j+\eta_j-d_j>0$（$P_j=0$ 另按上述方式处理），
所以同一个 $\kappa_j$ 论证再次给出式~\eqref{v44d:eq:strict-pair-lower}。
因此不存在有限终止的成功协议。严格正的 $\varepsilon_j$ 在此处用于保证有限能量时的严格余量。

连续观测时，令
$Q_j(t)=\int_0^t u_j(\tau)^2\,d\tau$，
在有限门槛行选择
\begin{equation}
 n_j(t)=d_ju_j(t)\,
       \mathbf 1_{\{Q_j(t)\le e_j(q)\}},
 \label{v44d:eq:continuous-threshold-adversary}
\end{equation}
在无限门槛行选择 $n_j(t)=d_ju_j(t)$。
这是一条因果、可测的噪声规则。
有限门槛行满足
\[
 \|n_j\|_{L^2(0,t)}
 \le d_j\min\{\sqrt{Q_j(t)},P_j^*\},
\]
故前述合法性证明逐字适用。
若校准停止时某行 $Q_j(T_{\rm cal})<e_j(q)$，
该行的完整连续报告仍是
$\widetilde z_j=(\ell_j+d_j)u_j$，
而替代模型的完整噪声为 $-(d_j+\kappa_j)u_j$。
因此不可区分性适用于全部连续观测轨迹，而非仅适用于某一个投影统计量。
\end{proof}

\subsection{允许数据后选速和时变推进的精确两阶段截止时间}

以下完成时间针对端口模型中的真实材料相位遍历。
若还要求离散图像的连续材料覆盖、模板质量或失锁后的重扫，
须再满足正文的参考注册、帧足迹及混合覆盖合同；
本节不把相位遍历自动等同于这些额外要求。

\begin{theorem}[一般执行选速下的两阶段精确截止时间]
\label{v44d:thm:adaptive-speed-deadline}
在引理~\ref{v44d:lem:common-low-model} 的连续观测合同下，
考虑以下两阶段协议类：
\begin{enumerate}
 \item 校准阶段实际 $u_s=0$，峰值能力为
 $\sum_j u_j(t)^2\le U_{\rm cal}^2$，其中 $U_{\rm cal}>0$。
 波形、各行能量与停止时间均可依赖全部已有数据。
 \item 校准结束后，执行阶段实际 $\dot\sigma=u_s=\nu(t)\ge0$，
 并满足式~\eqref{v44m:eq:execution-cap} 的共享输入能力。
 第一笔执行动作及此后每个时刻都须对整个任务命令箱满足目标 $q$。
 速度可以在校准后根据数据选择，也可以随时间、已有执行历史和当前任务命令变化。
 \item 执行目标是 $\sigma-\sigma(0)=L_\sigma>0$。
 不在校准阶段累计这段执行进度，不另设执行加速度、延迟或观测开销。
 全部增益在两阶段保持相同，噪声只作用于所声明的校准测量。
\end{enumerate}
假定共同有效域能够容纳下述达到构造。
若 $\mathcal E_*(q)<\infty$ 且式~\eqref{v44m:eq:speed-polynomial} 中 $C_q<0$，
则该协议类的最小最坏完成时间为
\begin{equation}
 \boxed{
 T_*(q)=\frac{\mathcal E_*(q)}{U_{\rm cal}^2}
        +\frac{L_\sigma}{\nu_{\max}(q)}.}
 \label{v44d:eq:adaptive-speed-deadline}
\end{equation}
若另有总校准能量上限 $E$，还须且只须 $E\ge\mathcal E_*(q)$。
若 $\mathcal E_*(q)=+\infty$ 或 $C_q\ge0$，则不存在有限完成时间。
\end{theorem}

\begin{proof}
在同一个最小增益真实模型~\eqref{v44d:eq:common-low-model} 上，
引理~\ref{v44d:lem:common-low-model} 的因果噪声历史强迫
\[
 \int_0^{T_{\rm cal}}\sum_j u_j(t)^2\,dt\ge\mathcal E_*(q).
\]
由校准峰值约束，
$T_{\rm cal}\ge\mathcal E_*(q)/U_{\rm cal}^2$。

在此真实模型的整个执行阶段，选择固定的任务箱角点
\begin{equation}
 v_j^\square=
 \begin{cases}
 -r_j\operatorname{sgn}(b_j),&b_j\ne0,\\
 -r_j,&b_j=0.
 \end{cases}
 \label{v44d:eq:persistent-corner-command}
\end{equation}
无论控制器如何选择非负的 $\nu(t)$，
$|v_j^\square-b_j\nu(t)|=r_j+|b_j|\nu(t)$。
响应合同在同一时刻强迫
\[
 |u_j(t)|\ge
 \frac{(1-q)[r_j+|b_j|\nu(t)]}{\ell_j}.
\]
代入共享输入球便得到
\[
 A_q\nu(t)^2+2B_q\nu(t)+C_q\le0,
 \qquad \nu(t)\le\nu_{\max}(q).
\]
因此完成 $L_\sigma$ 的执行时长至少是
$L_\sigma/\nu_{\max}(q)$。
这两个必要时间发生在同一个真实参数、同一个允许校准噪声历史及同一个执行任务序列上，
所以可以相加。这里没有把两个互不相容的分别最坏情形相加。
固定角点~\eqref{v44d:eq:persistent-corner-command} 也说明，
允许 $\nu$ 依赖当前 $v$ 不会绕开该下界。

上界采用第~\ref{v44m:sec:multichannel}节的固定行能量、
式~\eqref{v44m:eq:returning-waveform} 的一次往返速度波形，
并取 $T_{\rm cal}=\mathcal E_*(q)/U_{\rm cal}^2$。
这个波形同时取得全部行的有效数据，并在漂移为零的固定增益模型中恢复实际积分坐标、
任务特征与相位。
其峰值累计输入位移为 $\sqrt{e_j(q)T_{\rm cal}}/2$，
峰值任务偏离为其 $h_j$ 倍，均须由共同校准域容纳。
然后使用式~\eqref{v44m:eq:attaining-controller} 的后验区间规则，
以常速度 $\nu_{\max}(q)$ 执行。
该构造是所考虑的大协议类的一个成员，达到式~\eqref{v44d:eq:adaptive-speed-deadline}。
$\mathcal E_*(q)=0$ 时直接跳过校准。
无限能量门槛由引理排除有限校准；
$C_q\ge0$ 时，同一最小增益模型上的命令角点不允许正速度，因而不能完成正长度。
\end{proof}

该精确时间结论采用速度级执行接口。
将公共 Lyapunov 恢复定理的充分速率上界额外加入后，
所得时间是相应恢复证书的充分保证，不由本定理自动变成任意动态控制器的最优时间。
执行段若也受硬加速度、启动、末端制动或图像有效性约束，
须为这些条件单独构造并验收接续。

\subsection{有限幅值、有限加速度与实际回位下的两倍时间构造}

现在仅对校准阶段增加实际速度与加速度合同
\begin{equation}
 \|u(t)\|_2\le U_{\rm cal},\qquad
 \|\dot u(t)\|_2\le A_{\rm cal},\qquad
 u(0)=u(T)=0,\qquad \int_0^T u(t)\,dt=0,
 \quad A_{\rm cal}>0.
 \label{v44d:eq:rest-to-rest-return-contract}
\end{equation}
导数约束按几乎处处解释，$u$ 绝对连续。
这里 $u=(u_1,\ldots,u_m)$ 是实际积分坐标 $x_j$ 的速度，
即 $\dot x_j=u_j$；所以最后一个等式保证 $x(T)=x(0)$。
在一般机器人 twist、时变坐标重建或身体漂移模型中，
不能仅从一组输入分量的零积分推断完整机器人位姿回位；
使用本构造时须具有上述实际积分坐标，或另外证明同一路径的真实往返。
在当前固定线性端口模型且 $u_s=0$ 时，同时有 $y(T)=y(0)$ 和 $\sigma(T)=\sigma(0)$。

\begin{proposition}[有限加速度回位探查的两倍最坏时间保证]
\label{v44d:prop:sine-two-approximation}
设 $0<\mathcal E_*(q)<\infty$，连续观测合同
式~\eqref{v44d:eq:continuous-oracle} 对全部比较的时长和波形共同成立。
在式~\eqref{v44d:eq:rest-to-rest-return-contract} 下，
令 $T_{\rm cal}^*(q)$ 为全部允许因果协议达到目标响应精度 $q$ 的最小最坏校准时间。
此处成功指校准后的共同规则能够对全部剩余模型和整个命令箱满足该误差合同，
暂不把执行输入球加入这个仅衡量校准精度的优化问题；
若还要求在给定执行输入球内完成任务，须另外满足
式~\eqref{v44m:eq:exact-resource-frontier} 的执行能力条件。
定义
\begin{align}
 T_{\rm lb}
 &=\max\left\{
 \frac{\mathcal E_*(q)}{U_{\rm cal}^2},
 \left(\frac{12\mathcal E_*(q)}{A_{\rm cal}^2}\right)^{1/3}
 \right\},\label{v44d:eq:acceleration-time-lower}\\
 T_{\sin}
 &=\max\left\{
 \frac{2\mathcal E_*(q)}{U_{\rm cal}^2},
 \left(\frac{8\pi^2\mathcal E_*(q)}{A_{\rm cal}^2}\right)^{1/3}
 \right\}.\label{v44d:eq:sine-duration}
\end{align}
若共同校准物理域能够容纳下面的显式正弦轨迹，则
\begin{equation}
 \boxed{
 T_{\rm lb}\le T_{\rm cal}^*(q)\le T_{\sin}
 \le2T_{\rm lb}\le2T_{\rm cal}^*(q).}
 \label{v44d:eq:factor-two-time}
\end{equation}
一个达到所述构造上界的实际速度为
\begin{equation}
 u_j(t)=\sqrt{\frac{2e_j(q)}{T_{\sin}}}
          \sin\left(\frac{2\pi t}{T_{\sin}}\right),
 \qquad 0\le t\le T_{\sin},\qquad u_s=0.
 \label{v44d:eq:sine-returning-probe}
\end{equation}
非活跃行取 $u_j=0$。
\end{proposition}

\begin{proof}
对任何允许协议，引理~\ref{v44d:lem:common-low-model}
在同一个最小增益模型上强迫总校准能量至少为 $\mathcal E_*(q)$。
设该真实历史的校准终止时刻为 $T$。
速度界给出能量至多 $U_{\rm cal}^2T$。
由起末速度为零及加速度界，在任意 $t\in[0,T]$ 有
\[
 \|u(t)\|_2\le A_{\rm cal}\min\{t,T-t\}.
\]
因此
\begin{equation}
 \int_0^T\|u(t)\|_2^2\,dt
 \le A_{\rm cal}^2\int_0^T\min\{t,T-t\}^2\,dt
 =\frac{A_{\rm cal}^2T^3}{12}.
 \label{v44d:eq:acceleration-energy-bound}
\end{equation}
联合两个必要能量界即得 $T\ge T_{\rm lb}$。
这个下界对自适应波形和历史相关停止同样成立；
它没有用未证明的固定波形最优性。

对任意 $T>0$，式~\eqref{v44d:eq:sine-returning-probe}
以 $T$ 代替 $T_{\sin}$ 后满足
\begin{align}
 \int_0^T u_j(t)^2\,dt&=e_j(q),&
 \int_0^T u_j(t)\,dt&=0,\\
 \max_t\|u(t)\|_2&=\sqrt{2\mathcal E_*(q)/T},&
 \max_t\|\dot u(t)\|_2
 &=\frac{2\pi\sqrt{2\mathcal E_*(q)}}{T^{3/2}}.
 \label{v44d:eq:sine-resource-check}
\end{align}
选择 $T=T_{\sin}$ 就满足全部幅值、加速度及起末静止约束。
以单位 $L^2$ 模式
\[
 \phi_T(t)=\sqrt{2/T}\sin(2\pi t/T)
\]
投影完整观测，得到
\[
 z_j=\int_0^T\phi_T(t)\widetilde z_j(t)\,dt
     =a_j\sqrt{e_j(q)}+w_j,\qquad
 |w_j|\le\varepsilon_j+\eta_j\sqrt{e_j(q)}.
\]
所以第~\ref{v44m:sec:multichannel}节的区间达到策略可以直接使用这些实际数据。
两个半程都进入测量和能量计算；零净位移没有被误当成零信息。

沿整个轨迹，实际积分位移和任务偏离满足
\begin{align}
 x_j(t)-x_j(0)
 &=\frac{\sqrt{2e_j(q)T}}{2\pi}
    \left[1-\cos(2\pi t/T)\right],\\
 \max_t|x_j(t)-x_j(0)|
 &=\frac{\sqrt{2e_j(q)T}}{\pi},\\
 \max_t|y_j(t)-y_j(0)|
 &\le\frac{h_j\sqrt{2e_j(q)T}}{\pi}.
 \label{v44d:eq:sine-domain-check}
\end{align}
因此若预先声明的校准观测宽度为
$|y_j-y_j(0)|\le w_{y,j}$，可直接检查
\begin{equation}
 \frac{h_j\sqrt{2e_j(q)T_{\sin}}}{\pi}\le w_{y,j},
 \label{v44d:eq:sine-sound-window}
\end{equation}
实际输入坐标、关节或其它物理边界也按同一完整轨迹验收。
终点确实回位，式~\eqref{v44d:eq:sine-domain-check}
同时验证了中间路径，而不只比较两个端点。

最后，
\[
 \left(\frac{8\pi^2}{12}\right)^{1/3}
 =\left(\frac{2\pi^2}{3}\right)^{1/3}<2.
\]
故 $T_{\sin}\le2T_{\rm lb}$。可行正弦轨迹又给出
$T_{\rm cal}^*(q)\le T_{\sin}$，从而得到
式~\eqref{v44d:eq:factor-two-time}。
\end{proof}

式~\eqref{v44d:eq:sine-sound-window} 是构造可用的明确前提。
若它不成立，本节仍保留必要时间下界，但不能声称该正弦构造通过物理域或具有两倍近似保证。
尤其，该一周期轨迹的累计位移随 $\sqrt T$ 增长；
不能仅靠延长时间就宣称所有声窗限制都会改善。
当 $\mathcal E_*(q)=0$ 时直接跳过校准，不使用含 $T^{-1/2}$ 的波形公式。

\begin{corollary}[有限校准加速度下的两阶段时间夹界]
\label{v44d:cor:smooth-two-stage-bound}
在命题~\ref{v44d:prop:sine-two-approximation} 的物理域条件下，
显式假定 $0<\mathcal E_*(q)<\infty$ 且 $C_q<0$。
若执行阶段仍采用定理~\ref{v44d:thm:adaptive-speed-deadline} 的速度级合同，
令 $\tau_\sigma=L_\sigma/\nu_{\max}(q)$。
该两阶段类的最小最坏总时间满足
\begin{equation}
 T_{\rm lb}+\tau_\sigma
 \le T_{\rm total}^*(q)
 \le T_{\sin}+\tau_\sigma
 \le2\bigl(T_{\rm lb}+\tau_\sigma\bigr)
 \le2T_{\rm total}^*(q).
 \label{v44d:eq:smooth-total-time-bracket}
\end{equation}
下界仍由同一个最小增益模型、门槛噪声和持续任务箱角点共同产生；
上界由实际正弦回位后接常速执行达到。
\end{corollary}

这里的有限加速度约束只施加在校准阶段。
若要求执行阶段同样满足硬加速度约束，
必须另外计入从静止到正推进以及末端制动的真实接续时间和状态域。
上述两倍结论不自动覆盖增加了这些约束后的新优化问题。
连续速度观测的 $L^2$ 合同也不能由位置帧噪声或未经验证的差分、重复采样自动推出；
使用位置图像时，必须把真实测量处理得到的时间与波形依赖噪声界重新代回。


\subsection{三通道共享预算的完整参数实例}
\label{v44m:sec:numeric-example}
以下数值由精确资源公式与相应物理构造直接计算。
取归一化任务与输入参数
\begin{equation}
 q=0.06,\quad U=0.35,\quad U_{\rm cal}=0.20,
 \qquad r=(0.04,0.03,0.05),\quad b=(0.6,-0.4,0.8).
\end{equation}
各通道具有不同先验与噪声：
\begin{center}\small
\begin{tabular}{ccccc}
\toprule
通道 $j$ & $[\ell_j,h_j]$ & $\varepsilon_j$ & $\eta_j$ & 精确所需 $e_j(0.06)$\\
\midrule
1 & $[0.8,1.2]$ & $0.002$ & $0.010$ & $0.0023721442$\\
2 & $[1.0,1.4]$ & $0.003$ & $0.015$ & $0.0037746166$\\
3 & $[1.2,1.8]$ & $0.0025$ & $0.005$ & $0.0012192860$\\
\bottomrule
\end{tabular}
\end{center}
因此 $\mathcal E_*=0.0073660468\ldots$。
式~\eqref{v44m:eq:returning-waveform}取
$T=\mathcal E_*/U_{\rm cal}^2=0.1841511711\ldots$，
三行的实际输入幅值分别为
$(0.11349672,0.14316907,0.08137024)$。
其共同幅值满足 $\|u\|=0.20$，峰值任务偏离分别小于
$(0.012541,0.018456,0.013486)$。
故校准宽域 $|y_j-y_j(0)|\le0.03$ 容纳整个探查；真实任务状态与材料相位在探查结束时返回起点。
这些回位结论采用式~\eqref{v44m:eq:plant}的固定增益与无状态漂移模型，观测噪声仍非零。

在所得到的每一个允许观测历史下，后验分配均满足响应目标与统一能力。
式~\eqref{v44m:eq:uniform-speed-root}给出
\begin{equation}
 \nu_{\max}=0.2097077743\ldots.
\end{equation}
选择 $\nu=0.18$ 时，对整个任务命令盒，实际输入范数不超过
$0.308104<0.35$。
若取二阶反馈 $v=\zeta$、$\omega=4$、$R_s=0.0075$，
则 $\|\zeta\|\le0.03\le\min_jr_j$，任务命令盒确实包含全部恢复状态。
取执行期外部响应漂移界 $d_0=0.0002$、测量误差界 $n_{\max}=0.0001$，
有
\begin{equation}
 q\nu\|b\|+d_0+\sqrt2\omega n_{\max}
 =0.012397642\ldots
 <\frac{\omega(1-q)}2R_s=0.0141.
\end{equation}
校准返回后令 $\zeta=0$，并采用常参考 $y_r=y(0)$，则 $S(0)=0$。
恢复管 $S\le R_s$、任务误差 $\|e\|\le\sqrt2R_s$、
共同输入能力和真实 $\dot\sigma=0.18$ 同时成立。
该全局仿射解析模型中的 $L_\sigma=0.05$ 材料段耗时 $0.27777778\ldots$，
加上上述速度级返回探查，总时长为 $0.46192895\ldots$。
观测中断、成像足迹和有限转角的几何检查由后续非线性实例进一步落实。

若校准还要求起止速度为零及 $\|\dot u\|\le A_{\rm cal}=2$，
可使用后续时间界中的同步正弦构造。
它在 $T_{\sin}=0.5258413559\ldots$ 内完成同一能量的探查并返回，
而所有允许因果探查的时间下界为 $0.2806199721\ldots$；
二者之比小于 $1.874$。
正弦的峰值输入小于 $0.168$、峰值加速度为 $2$，
三行的最大任务偏离均小于 $0.029$，因此同一 $0.03$ 校准宽域仍有效。
这一比较认证带加速度的校准时间；扫描启动和终端若另受加速度硬约束，仍需追加对应的执行过渡。


\section{局部响应模型、物理输入与度量}
\label{sec:cal-local}

本节讨论一个明确且可检验的局部模型类。所有变量已经按物理运动度量与任务容差归一化，因而下文的欧氏范数、正交投影和奇异值具有相同的物理解释。固定允许的机器人实现后，省略其索引，令
\begin{equation}
  \dot y=Ju+w,\qquad
  J=H+\Theta N+E_{\mathrm{geo}},
  \label{eq:cal-model}
\end{equation}
其中 $y\in\mathbb R^m$ 是有效任务特征。$u\in\mathbb R^n$ 是当前受控子空间内的实际刚体运动，经已知约束约简和固定物理白化表示；已知人体保持运动按同一响应模型补偿，其剩余影响进入 $w$。移动 CoP 地址与该刚体输入的精确定义见第~\ref{sec:rev:cop-input}节。$H\in\mathbb R^{m\times n}$ 是已知名义响应，$N\in\mathbb R^{r\times n}$ 是已知的不确定物理通道，$\Theta\in\mathbb R^{m\times r}$ 是患者相关的未知响应。$E_{\mathrm{geo}}$ 汇集几何运输、模型遗漏和局部线性化的误差。扰动 $w$ 可以包括不能解释为命令响应的目标漂移、人体运动估计残差以及参考变化。

\begin{assumption}[局部模型、有效观测与物理执行]
\label{ass:cal-model}
在本次校准和验收的局部窗口内，$H$ 满行秩，$N$ 已知，且
\begin{equation}
  \|\Theta\|_2\le\delta,\quad\delta>0,\qquad
  \|E_{\mathrm{geo}}\|_2\le\eta .
\end{equation}
校准测量来自有效的目标特征；被实际执行的命令、而非理想规划命令，被用于构造数据。模型只在明确的控制与观测可信区间内使用。失锁或关联未确定的图像不产生新的有效任务响应测量。
\end{assumption}

$N$ 必须来自物理响应结构或经过验证的辨识结果。例如，特定接触变形通道可以由明确的接触模型给出。SMPL-X 对应关系和坐标不变性本身不会给出已知 $N$，也不会保证小的 $\eta$。若所有命令方向都任意未知，应取 $N=I_n$；此时不能预期以下结构产生省校准收益。首先在 $E_{\mathrm{geo}}=0$、无噪声且局部 $\Theta$ 固定的模型内证明精确结论，再给出带误差的验收条件。

\subsection{固定物理度量与被动坐标变换}

输入、任务、不确定通道分别赋正定度量 $W,M,S$。对响应矩阵定义
\begin{equation}
 \widetilde H=M^{1/2}HW^{-1/2},\quad
 \widetilde N=S^{1/2}NW^{-1/2},\quad
 \widetilde\Theta=M^{1/2}\Theta S^{-1/2}.
\end{equation}
全文伪逆、正交投影和奇异值都在这个白化模型中理解。为简洁，全文省略波浪号。度量需要固定到真实的毫米、弧度、任务容差与通道尺度，不能为了获得收益而重设。

被动换坐标 $u'=Tu$, $e'=Pe$, $z'=Vz$ 给出
\begin{equation}
 H'=PHT^{-1},\quad N'=VNT^{-1},\quad
 \Theta'=P\Theta V^{-1},\quad B'=TBP^{-1}.
\end{equation}
度量同步作协变变换，例如 $W'=T^{-\top}WT^{-1}$。白化前后算子由正交基变换关联，故诱导奇异值与任务下界不变。本文比较在同一物理度量下的不同动作实现；被动重参数保持相应信息阈值。

若接触的线性约束是 $E_cu=0$，先取 $\ker E_c$ 的正交基 $Z_c$，将输入写成 $u=Z_c\bar u$，对 $HZ_c,NZ_c$ 推导。实际能力多面体、关节约束和碰撞仍需后续验收。

\subsection{保持任务的控制分配}

名义任务保持的线性实现 $B\in\mathbb R^{n\times m}$ 满足
\begin{equation}
  HB=I_m,\qquad u=Bv ,
  \label{eq:cal-lift}
\end{equation}
其中 $v\in\mathbb R^m$ 是同一个材料任务命令。实际响应为
\begin{equation}
  JB=I_m+\Theta NB+E_{\mathrm{geo}}B .
  \label{eq:cal-exposure}
\end{equation}
下文精确谱、首次下界与必要充分方向数采用 $E_{\mathrm{geo}}=0$、无噪声的冻结局部模型；有余项时使用后文 $q$ 的验收证书。

因此，需要校准的是 $\Theta$ 在 $\operatorname{range}(NB)$ 上的作用，称 $NB$ 为该实现的不确定暴露。这里不要求未知矩阵 $\Theta$ 低秩；低维性来自可选择的实际机器人实现。

若精确接触条件要求 $Eu=0$，令 $Z_E$ 是 $\ker E$ 的正交基，在允许的命令空间内使用
\begin{equation}
  H_E=HZ_E,\qquad N_E=NZ_E .
  \label{eq:cal-contact-reduction}
\end{equation}
只有 $H_E$ 满行秩时，才能继续实现全部任务方向。以下代数结果可以对 $(H_E,N_E)$ 使用。实际探查所用输入来自运动测量；若只能获得速度接口指令，应将指令与实际运动的差异纳入独立执行误差界。速度、关节、碰撞和接触的不等式能力限制通过后续能力集合验收。若总输入为 $u_{\rm ff}+C\zeta$，精确接触条件同时要求 $Eu_{\rm ff}=0$ 和 $EC=0$；作用于实际输入的等式还要求执行误差属于 $\ker E$。





\section{控制分配与不可消除暴露}
\label{sec:cal-realization}

\subsection{同时最小化所有奇异值的实现}

取 $Z_H$ 为 $\ker H$ 的正交基。由于 $H$ 满行秩，$HH^\dagger=I_m$，而任意右逆都有形式
\begin{equation}
  B=H^\dagger+Z_HK .
  \label{eq:cal-all-lifts}
\end{equation}
记
\begin{equation}
  A=NH^\dagger,\qquad C_0=NZ_H,\qquad
  \mathcal C=\operatorname{range}(C_0)=N\ker H .
\end{equation}
令 $P_{\mathcal C}=C_0C_0^\dagger$ 为到 $\mathcal C$ 的正交投影，并定义
\begin{equation}
  B_*=
  H^\dagger-Z_HC_0^\dagger A,\qquad
  D_*=
  NB_*=(I-P_{\mathcal C})A .
  \label{eq:cal-optimal-lift}
\end{equation}

\begin{theorem}[最小不确定暴露]
\label{thm:cal-singular-opt}
在无额外物理不等式约束的右逆集合 $HB=I_m$ 中，$B_*$ 同时最小化 $NB$ 的所有奇异值：
\begin{equation}
  \sigma_j(NB)\ge\sigma_j(D_*),
  \qquad j=1,\ldots,\min\{r,m\}.
  \label{eq:cal-singular-order}
\end{equation}
特别地，它同时最小化谱范数、任意奇异值尾和暴露的秩。
\end{theorem}

\begin{proof}
对任意 $B=H^\dagger+Z_HK$，
\begin{align}
  NB
  &= A+C_0K\\
  &= (I-P_{\mathcal C})A+
     \big(P_{\mathcal C}A+C_0K\big)\\
  &=D_*+E_B .
\end{align}
$D_*$ 的列属于 $\mathcal C^\perp$，$E_B$ 的列属于 $\mathcal C$，故 $D_*^\top E_B=0$。于是
\begin{equation}
  (NB)^\top NB
  =D_*^\top D_*+E_B^\top E_B
  \succeq D_*^\top D_* .
\end{equation}
由对称矩阵的特征值单调性得到~\eqref{eq:cal-singular-order}。选择 $K=-C_0^\dagger A$ 时 $E_B=0$，从而各个下界均达到。
\end{proof}

这个结论区分了两种改变：被动坐标变化不会创造新的暴露收益；主动改变物理命令实现，可以用 $N\ker H$ 中的额外自由度抵消不确定暴露，而同时保留 $HB=I_m$。

\subsection{最低精确校准维数}

\begin{theorem}[任务保持实现的结构秩]
\label{thm:cal-rank}
最低不确定暴露秩为
\begin{equation}
  \boxed{
  k_*=
  \min_{HB=I_m}\operatorname{rank}(NB)
  =m-\dim(H\ker N)
  =m+\operatorname{rank}N-
  \operatorname{rank}
  \begin{bmatrix}H\\N\end{bmatrix}.}
  \label{eq:cal-kstar}
\end{equation}
\end{theorem}

\begin{proof}
因为 $HB=I_m$，$B$ 是单射，而且 $H$ 在 $\operatorname{range}(B)$ 上是到 $\mathbb R^m$ 的同构。因此
\begin{align}
  \operatorname{rank}(NB)
  &=m-\dim\big(\operatorname{range}(B)\cap\ker N\big)\\
  &\ge m-\dim(H\ker N).
  \label{eq:cal-rank-lower}
\end{align}
另一方面，$D_*v=0$ 当且仅当 $NH^\dagger v\in N\ker H$。这等价于存在 $z\in\ker H$ 使
\begin{equation}
  N(H^\dagger v-z)=0,\qquad
  H(H^\dagger v-z)=v ,
\end{equation}
即 $v\in H\ker N$。因而
\begin{equation}
  \ker D_*=H\ker N,\qquad
  \operatorname{rank}D_*=m-\dim(H\ker N),
\end{equation}
下界由 $B_*$ 达到。最后，对从命令空间到联合输出 $(Hu,Nu)$ 的线性映射，先取 $Nu$ 的像，再加入 $Nu=0$ 时的 $Hu$，得到
\begin{equation}
  \operatorname{rank}
  \begin{bmatrix}H\\N\end{bmatrix}
  =\operatorname{rank}N+\dim(H\ker N).
\end{equation}
\end{proof}

当 $N=I_n$ 时，$H\ker N=\{0\}$，故 $k_*=m$。因此，结构收益依赖实际不确定通道与任务通道之间的关系，而不能由改变参数化自动获得。$B_*$ 也可能增加所需物理运动幅度；受约束的实际问题必须保留这一代价。

\subsection{完整范数球的精确鲁棒满射阈值}
\label{rev:sec:surjectivity}

最小暴露还确定了结构化响应失去瞬时任务满射性的精确距离。受限奇异值分解早已给出 $A+BDC$ 的最小致秩亏扰动框架\cite{demoorgolub1991}；本节将其结构特例写成任务商空间 $D_\star$ 的显式表达，并给出直接证明。定义
\begin{equation}
 \rho_{\rm sing}(H,N)
 =
 \inf\left\{
 \|\Theta\|_2:
 \operatorname{rank}(H+\Theta N)<m
 \right\},
 \qquad \inf\varnothing=+\infty .
 \label{rev:eq:sing-radius}
\end{equation}
这里暂不施加物理输入幅度限制，且 $E_{\rm geo}=0$。“满射”指线性瞬时任务映射的像为 $\mathbb R^m$。

\begin{theorem}[最小致秩亏扰动与闭范数球阈值]
\label{rev:thm:robust-surjectivity}
设 $H\in\mathbb R^{m\times n}$ 满行秩，$N\in\mathbb R^{r\times n}$，并采用固定物理度量白化后的谱范数。令 $s_*=\|D_*\|_2$，其中 $D_*$ 由式~\eqref{eq:cal-optimal-lift} 定义。则
\begin{equation}
 \boxed{
 \rho_{\rm sing}(H,N)=
 \begin{cases}
  1/s_*,&s_*>0,\\
  +\infty,&s_*=0.
 \end{cases}}
 \label{rev:eq:exact-sing-radius}
\end{equation}
当 $s_*>0$ 时，该下确界由一个秩一扰动达到。因而对任意有限 $\delta\ge0$，
\begin{equation}
 \boxed{
 \left[
 \forall\,\Theta\in\mathbb R^{m\times r},\
 \|\Theta\|_2\le\delta:
 \operatorname{rank}(H+\Theta N)=m
 \right]
 \iff
 \delta\|D_*\|_2<1.}
 \label{rev:eq:robust-rank-iff}
\end{equation}
\end{theorem}

\begin{proof}
对任意 $\Theta$，
\begin{equation}
 (H+\Theta N)B_*=I_m+\Theta D_* .
 \label{rev:eq:robust-right-product}
\end{equation}
若 $\|\Theta\|_2s_*<1$，则右侧可逆，故 $H+\Theta N$ 满行秩。这证明 $s_*>0$ 时 $\rho_{\rm sing}\ge1/s_*$。若 $s_*=0$，同一式的右侧恒为 $I_m$，不存在任何有限范数的致秩亏扰动，故 $\rho_{\rm sing}=+\infty$。

下面设 $s_*>0$，取单位最大奇异向量 $v_*\in\mathbb R^m$、$w_*\in\mathbb R^r$，满足
\begin{equation}
 D_*v_*=s_*w_*,
 \qquad
 D_*^\top w_*=s_*v_* .
 \label{rev:eq:worst-singular-pair}
\end{equation}
由 $D_*$ 的值域属于 $\mathcal C^\perp$ 可知 $w_*\perp\mathcal C=N\ker H$。于是 $w_*^\top N$ 在 $\ker H$ 上为零，并且
\begin{align}
 w_*^\top N
 &=w_*^\top NH^\dagger H\\
 &=w_*^\top D_*H
 =s_*v_*^\top H .
 \label{rev:eq:rank-loss-factor}
\end{align}
构造
\begin{equation}
 \Theta_{\rm sing}
 =-\frac1{s_*}v_*w_*^\top .
 \label{rev:eq:singular-witness}
\end{equation}
其谱范数为 $1/s_*$，且
\begin{equation}
 H+\Theta_{\rm sing}N
 =(I_m-v_*v_*^\top)H .
 \label{rev:eq:rank-deficient-response}
\end{equation}
由于 $H$ 满射，右侧像空间恰为 $v_*^\perp$，秩为 $m-1$。因此 $\rho_{\rm sing}\le1/s_*$，且最小值被达到。若 $\delta s_*\ge1$，该扰动属于允许的闭范数球，鲁棒满射性即失败；若 $\delta s_*<1$，前述可逆性论证适用于球内全部模型。这证明了等价式。
\end{proof}

\begin{remark}[阈值的集合与物理范围]
\label{rev:rem:rank-scope}
式~\eqref{rev:eq:robust-rank-iff} 的必要性使用了完整闭谱范数球；边界 $\delta s_*=1$ 已包含致秩亏模型，因此不能改成非严格不等式。若实际物理不确定集合是该球的真子集，$\delta s_*<1$ 仍是充分条件，但上述秩一扰动可能不属于物理集合，必要性需要重新分析。非零 $E_{\rm geo}$ 使用后文的响应误差合同，不能直接沿用本定理的必要性。

本定理刻画瞬时任务满射性。输入幅度、接触约束及有限时间可达性仍须分别验收；精确接触条件应先按式~\eqref{eq:cal-contact-reduction} 降到允许命令子空间，再对降维后的矩阵应用本定理。主动改变物理状态 $a$ 后，应对新模型重新计算阈值。
\end{remark}





\section{探查信息与首次任务响应}
\label{sec:cal-minimax}

这一节不预先锁定 $u=Bv$。校准以后，控制器可以根据数据重新选择右逆，甚至使用任意非线性的静态规则 $u(v)$。因此，需要剔除任务零空间提供的所有绕开机会，不能把固定 $B$ 的谱尾直接当作全体实现的任务下界。

\subsection{校准信息与正确的剩余空间}

一条无噪声局部校准记录给出实际响应增量 $\Delta y_i$（或对冻结 Jacobian 积分的局部响应）。输入 $u_i$ 对应相同时间尺度的实际积分运动，因此
\begin{equation}
  r_i=\Delta y_i-Hu_i=\Theta x_i,\qquad x_i=Nu_i .
  \label{eq:cal-probe-oracle}
\end{equation}
记 $Q=\operatorname{span}\{x_1,\ldots,x_s\}$。这里计数的是独立激励方向 $\dim Q$，不是图像帧数、探查持续时间或动作片段数。长片段可以提供多个独立激励。

考虑所有校准残差均为零的历史。数据一致模型满足
\begin{equation}
  \Theta Q=0,\qquad \|\Theta\|_2\le\delta .
  \label{eq:cal-zero-history}
\end{equation}
定义
\begin{equation}
  D_Q=P_{\mathcal C+Q}^{\perp}NH^\dagger,
  \qquad \mathcal C=N\ker H .
  \label{eq:cal-DQ}
\end{equation}
正确剩余矩阵是 $D_Q$，一般不是 $(I-P_Q)D_*$。实际控制器还能够通过 $\ker H$ 中的命令，把未校准暴露移动到已校准空间。

\begin{lemma}[给定名义任务输出的最小剩余暴露]
\label{lem:cal-DQ}
对任意 $z\in\mathbb R^m$，
\begin{equation}
  \min_{Hu=z}\|(I-P_Q)Nu\|_2
  =\|D_Qz\|_2 .
  \label{eq:cal-exposure-z}
\end{equation}
\end{lemma}

\begin{proof}
所有满足 $Hu=z$ 的命令为 $u=H^\dagger z+Z_H\ell$。记 $P_Q^\perp=I-P_Q$，则
\begin{align}
  \min_{Hu=z}\|P_Q^\perp Nu\|_2
  &=
  \min_\ell
  \|P_Q^\perp Az+P_Q^\perp C_0\ell\|_2\\
  &=
  \operatorname{dist}
  \big(P_Q^\perp Az,P_Q^\perp\mathcal C\big)\\
  &=\|P_{\mathcal C+Q}^\perp Az\|_2 .
\end{align}
最后一步因为 $Q^\perp$ 内再剔除 $P_Q^\perp\mathcal C$，等价于在原空间同时剔除 $\mathcal C+Q$。
\end{proof}

\subsection{非线性共同动作也不能绕开的下界}
固定动作的最坏范数扰动采用鲁棒最小二乘中的标准机制\cite{elghaoui1997}；零历史不可区分性也是最优恢复与信息复杂度中的经典证明方法\cite{plaskota2023}。这里需额外消去可自由选择的任务零空间，并将下界覆盖到数据后的任意共同动作。


定义零校准历史之后首次重新实现任务的最小最大误差
\begin{equation}
  \mathcal L(Q)=
  \inf_{u(\cdot)}
  \sup_{\|v\|_2=1}
  \sup_{\substack{\Theta Q=0\\\|\Theta\|_2\le\delta}}
  \|(H+\Theta N)u(v)-v\|_2 .
  \label{eq:cal-first-loss}
\end{equation}
这里允许任意静态非线性命令规则 $u(v)$，且暂时放松所有物理输入限制。因而获得的下界对有约束系统仍然有效；匹配的上界只有在相应命令确实可执行时才能使用。

\begin{theorem}[首次任务响应的精确条件最小最大误差]
\label{thm:cal-first-loss}
在~\eqref{eq:cal-zero-history} 的数据一致模型中，
\begin{equation}
  \boxed{\mathcal L(Q)=
  \min\{1,\delta\|D_Q\|_2\}.}
  \label{eq:cal-first-exact}
\end{equation}
\end{theorem}

\begin{proof}
对固定物理命令 $u$，令 $z=Hu$。由于 $\Theta Q=0$，$\Theta Nu=\Theta P_Q^\perp Nu$。范数球内的矩阵可以把 $P_Q^\perp Nu$ 映射到任意半径为 $\delta\|P_Q^\perp Nu\|_2$ 的任务输出向量。因此
\begin{equation}
  \sup_{\Theta Q=0,\;\|\Theta\|_2\le\delta}
  \|(H+\Theta N)u-v\|_2
  =
  \|z-v\|_2+\delta\|P_Q^\perp Nu\|_2 .
  \label{eq:cal-task-ball}
\end{equation}
利用引理~\ref{lem:cal-DQ}，逐个任务命令的最优值为
\begin{equation}
  \inf_z\big(\|z-v\|_2+\delta\|D_Qz\|_2\big).
  \label{eq:cal-task-prox}
\end{equation}
令 $v_*$ 为 $D_Q$ 的最大右奇异向量，$w_*$ 为相应左奇异向量，$s=\|D_Q\|_2$。当 $s>0$ 时，对任意 $z$，
\begin{align}
  \|z-v_*\|_2+\delta\|D_Qz\|_2
  &\ge
  |v_*^\top z-1|
  +\delta|w_*^\top D_Qz|\\
  &=
  |t-1|+\delta s|t|,
  \qquad t=v_*^\top z\\
  &\ge\min\{1,\delta s\}.
\end{align}
若 $\delta s\le1$，选择满足 $HB_Q=I_m$ 且 $P_Q^\perp NB_Q=D_Q$ 的最小暴露右逆，例如
\begin{equation}
  B_Q=
  H^\dagger-
  Z_H(P_Q^\perp NZ_H)^\dagger
  P_Q^\perp NH^\dagger ,
  \label{eq:cal-BQ}
\end{equation}
并令 $u(v)=B_Qv$，其最大误差不超过 $\delta s$。若 $\delta s\ge1$，选择零增量命令 $u(v)=0$，其误差为 $1$。当 $s=0$，$B_Q$ 实现零误差。故上下界匹配。
\end{proof}

可以直接构造一对限制任何物理动作的模型。选 $w_*\in(\mathcal C+Q)^\perp$ 满足 $D_Q^\top w_*=sv_*$。因为 $w_*^\top N$ 在 $\ker H$ 上为零，
\begin{equation}
  w_*^\top N
  =w_*^\top NH^\dagger H
  =sv_*^\top H .
  \label{eq:cal-adversary-factor}
\end{equation}
令
\begin{equation}
  \Theta_\pm=\pm\delta v_*w_*^\top .
\end{equation}
两模型都满足校准数据，但
\begin{equation}
  J_\pm=
  (I_m\pm\delta s v_*v_*^\top)H .
  \label{eq:cal-adversary-J}
\end{equation}
对任意物理命令 $u$，任务 $v_*$ 的误差投影都满足
\begin{equation}
  \max_{\pm}
  \left|
  (1\pm\delta s)v_*^\top Hu-1
  \right|
  \ge\min\{1,\delta s\}.
  \label{eq:cal-anyphysical-lower}
\end{equation}
这个构造同时约束全部任务零空间动作，消除了预先固定 $B$ 所产生的漏洞。

\begin{corollary}[零残差数据后的鲁棒满射性]
\label{rev:cor:conditional-surjectivity}
在数据一致集合 $\{\Theta:\Theta Q=0,\ \|\Theta\|_2\le\delta\}$ 内，
\begin{equation}
 \left[
 \forall\,\Theta Q=0,\ \|\Theta\|_2\le\delta:
 \operatorname{rank}(H+\Theta N)=m
 \right]
 \iff
 \delta\|D_Q\|_2<1 .
 \label{rev:eq:conditional-rank}
\end{equation}
若 $D_Q\ne0$，满足 $\Theta Q=0$ 的最小致秩亏扰动范数为 $1/\|D_Q\|_2$；若 $D_Q=0$，该距离为 $+\infty$。
\end{corollary}
\begin{proof}
式~\eqref{eq:cal-BQ} 给出
$(H+\Theta N)B_Q=I_m+\Theta D_Q$，从而得到充分性与距离下界。若 $s_Q=\|D_Q\|_2>0$，沿用式~\eqref{eq:cal-adversary-factor} 的最大奇异向量，并取
$\Theta_{\rm sing}=-v_*w_*^\top/s_Q$。由于
$w_*\in(\mathcal C+Q)^\perp$，有 $\Theta_{\rm sing}Q=0$；对应响应为 $(I_m-v_*v_*^\top)H$，其秩为 $m-1$。$D_Q=0$ 时 $B_Q$ 是全部数据一致模型的共同右逆。
\end{proof}

\subsection{最优探查方向与谱尾}

\begin{theorem}[最优独立校准方向的条件谱尾]
\label{thm:cal-width}
在无物理探查限制的零校准历史模型内，
\begin{equation}
  \min_{\dim Q\le k}\|D_Q\|_2
  =\sigma_{k+1}(D_*),
  \label{eq:cal-quotient-width}
\end{equation}
其中超过矩阵秩的奇异值取为零。因此
\begin{equation}
  \boxed{
  \min_{\dim Q\le k}\mathcal L(Q)
  =\min\{1,\delta\sigma_{k+1}(D_*)\}.}
  \label{eq:cal-minimax-width}
\end{equation}
\end{theorem}

\begin{proof}
因为 $D_*$ 的列在 $\mathcal C^\perp$ 中，
\begin{equation}
  D_Q=P_{\mathcal C+Q}^\perp D_* .
\end{equation}
投影掉 $\mathcal C$ 后，$Q$ 只能提供至多 $k$ 维的额外空间 $P_{\mathcal C}^\perp Q$。因此，左投影消除 $D_*$ 的至多 $k$ 个方向后，谱残差至少为 $\sigma_{k+1}(D_*)$。选择 $Q$ 为 $D_*$ 的前 $k$ 个左奇异向量张成的空间可以达到该下界。这些向量位于 $\operatorname{range}(N)$ 中，故存在物理命令满足 $Nu_i=q_i$。定理~\ref{thm:cal-first-loss} 完成最后一步。
\end{proof}

若验收要求首次任务误差不超过 $\epsilon$，其中 $0\le\epsilon<1$，零历史模型中最低方向数为
\begin{equation}
  k_\epsilon=
  \#\{j:\delta\sigma_j(D_*)>\epsilon\}.
  \label{eq:cal-epsilon-dimension}
\end{equation}
式~\eqref{eq:cal-epsilon-dimension} 对应 $\mathcal L\le\epsilon$，因此边界等号被接受。对 $0<\epsilon<1$，若合同改为 $\mathcal L<\epsilon$，计数条件相应改为 $\delta\sigma_j(D_*)\ge\epsilon$。

\begin{remark}[自适应探查、非零历史与闭环边界]
\label{rem:cal-width-scope}
对任何确定性自适应校准算法，或随机种子固定后的逐历史确定性保证，考虑其全部残差为零的历史。随机算法的期望风险或高概率复杂度不由此自动推出。最终 $Q$ 的维数不超过独立激励次数，而上述对手在该算法已经执行的全部探查上均返回零。因此~\eqref{eq:cal-minimax-width} 给出原始全模型类的必要最坏情况首次任务误差下界。

条件等式及其匹配上界仅针对零历史。非零观测会改变已知响应，并可能使数据一致不确定集合变为非对称或椭球集合，不能继续把原始 $D_*$ 的谱尾当作所有历史的精确上界。此外，首次任务响应下界不等于任意最终闭环恢复性能的下界；容许等待、继续适应或接受暂态误差的控制器解决的是另一个任务。
\end{remark}





\section{精确单位任务响应的校准复杂度}
\label{sec:cal-exact-complexity}

前一节已经给出零历史的精确近似响应结果。下面证明在完整无噪声模型类内，一个限定明确的任务保证具有匹配的校准复杂度。

\begin{theorem}[精确任务重新实现的必要与充分方向数]
\label{thm:cal-exact-task}
假定 $\delta\|D_*\|_2<1$。考虑校准后对一个非空任务命令邻域实现精确单位响应
\begin{equation}
  Ju(v)=v
  \label{eq:cal-unit-task}
\end{equation}
的规则，并要求该保证对所有 $\|\Theta\|_2\le\delta$ 成立。无约束模型类内，最坏情况下必须且只需 $k_*=\operatorname{rank}(D_*)$ 个独立校准方向。

物理限制不会削弱必要下界；充分性还要求相应探查和修复后的任务命令位于实际可执行可信区间内。
\end{theorem}

\begin{proof}
必要性可由定理~\ref{thm:cal-width} 在零误差下直接得到，也可用维数证明。取有效真实模型 $\Theta=0$。全部校准残差为零，任何满足 $\Delta Q=0$ 的小范数矩阵均与数据一致。若对全部这些模型都满足~\eqref{eq:cal-unit-task}，必须有
\begin{equation}
  Hu(v)=v,\qquad Nu(v)\in Q .
\end{equation}
因此能实现的任务命令属于 $H(N^{-1}Q)$，而
\begin{equation}
  \dim H(N^{-1}Q)
  \le\dim(H\ker N)+\dim Q .
\end{equation}
覆盖一个 $m$ 维邻域要求 $\dim Q\ge m-\dim(H\ker N)=k_*$。该证明允许任意静态非线性 $u(v)$，并适用于自适应算法的零历史。

充分性选用 $B_*$，令 $Q=\operatorname{range}(D_*)$，并校准 $\Theta$ 在 $Q$ 上的作用。由数据可完整知道
\begin{equation}
  G=JB_*=I_m+\Theta D_* .
\end{equation}
因为 $\|\Theta D_*\|_2\le\delta\|D_*\|_2<1$，$G$ 可逆。修复实现
\begin{equation}
  C_{\mathrm{exact}}=B_*G^{-1}
  \label{eq:cal-exact-repair}
\end{equation}
满足 $JC_{\mathrm{exact}}=I_m$。只需 $k_*$ 个独立校准方向即可识别 $G$，无需识别整个 $\Theta$。
\end{proof}

修复后的实际实现运动代价不能省略。例如
\begin{equation}
  \|C_{\mathrm{exact}}\|_2
  \le
  \frac{\|B_*\|_2}{1-\delta\|D_*\|_2}.
  \label{eq:cal-exact-command-bound}
\end{equation}
当分母小或 $B_*$ 运动幅度大时，校准方向减少仍可能以不可执行的命令为代价。

\begin{corollary}[鲁棒可实现性与精确信息预算]
\label{rev:cor:rank-information-link}
考虑 $E_{\rm geo}=0$、$\delta>0$、无噪声和无物理输入限制的固定局部模型。要求校准后对一个非空 $m$ 维任务命令邻域实现精确响应，并覆盖完整闭范数球内的全部模型。记满足这一要求的最坏情况最低独立校准方向数为 $K_{\rm exact}$，不可实现时记为 $+\infty$。则
\begin{equation}
 \boxed{
 K_{\rm exact}=
 \begin{cases}
 k_*=m-\dim(H\ker N),&\delta\|D_*\|_2<1,\\
 +\infty,&\delta\|D_*\|_2\ge1.
 \end{cases}}
 \label{rev:eq:rank-information-frontier}
\end{equation}
特别地，$D_*=0$ 时对任意有限 $\delta$ 均有 $K_{\rm exact}=0$。
\end{corollary}
\begin{proof}
第一种情形由定理~\ref{thm:cal-exact-task} 得到。第二种情形由定理~\ref{rev:thm:robust-surjectivity} 得到允许集合内的秩亏模型；其像空间不包含任何 $m$ 维开集，因而即使响应矩阵完全已知，也无法满足所要求的精确任务响应。$D_*=0$ 时 $k_*=0$ 且 $JB_*=I_m$。
\end{proof}
这一结论同时区分了信息不足与瞬时任务方向本身缺失。若允许先主动改变物理状态，必须在新状态下重新定义模型和信息预算。

\subsection{名义任务中性探查的范围}

\begin{proposition}[名义任务中性探查的最坏历史限制]
\label{prop:cal-neutral}
若所有探查均满足 $Hu_i=0$，则在零校准历史上 $Q\subseteq\mathcal C=N\ker H$，因而
\begin{equation}
  D_Q=D_*,
  \qquad
  \mathcal L(Q)=\min\{1,\delta\|D_*\|_2\}.
  \label{eq:cal-neutral-no-width}
\end{equation}
无论增加多少这种探查，都无法缩减该历史上的不可避免首次任务暴露。
\end{proposition}

\begin{proof}
$Hu_i=0$ 意味着 $u_i\in\ker H$，所以 $Nu_i\in\mathcal C$。于是 $\mathcal C+Q=\mathcal C$，直接得到结论。
\end{proof}

这使允许的有限任务偏离获得明确职责：它提供激励 $\mathcal C^\perp$ 中任务耦合通道的权限。该结论是最坏历史限制，而不是断言名义任务中性探查在每个实际患者上都无用。非零观测可能揭示能够用于新实现的任务零空间响应；仍须重新计算数据一致模型和实际能力。

\subsection{所有观测历史的一个充分上界}

取 $Q$ 为 $D_*$ 的前 $k$ 个左奇异方向，在任意无噪声历史上定义
\begin{equation}
  \widehat\Theta=\Theta P_Q,\qquad
  \widehat G=I_m+\widehat\Theta D_*,
  \qquad C_k=B_*\widehat G^{-1}.
\end{equation}
若 $\delta\|P_QD_*\|_2<1$，则
\begin{align}
  JC_k-I_m
  &=\Theta(I-P_Q)D_*\widehat G^{-1},\\
  \|JC_k-I_m\|_2
  &\le
  \delta\sigma_{k+1}(D_*)\,
  \|\widehat G^{-1}\|_2\\
  &\le
  \frac{\delta\sigma_{k+1}(D_*)}
  {1-\delta\|P_QD_*\|_2}.
  \label{eq:cal-nonzero-upper}
\end{align}
这是所有历史上的充分上界。它与零历史的必要下界通常存在间隙；不能将二者写成一个未经证明的精确样本复杂度。





% V4.4 independent mathematical fragment.
% Uses the main manuscript's amsmath/amssymb/theorem environments.
% Not a claim of priority: the two-ellipsoid radius lemma is classical;
% see Foucart--Liao (2024), already present as foucart2024 in the main draft.
% Author of this fragment: independent_math_review agent.

\section{物理探查预算与数据后任务分配的条件前沿}
\label{v44:sec:energy-frontier}

本节区分三个结果。第一，在固定未来任务分配后，给出完整激励其未知暴露所需的精确二次探查预算。第二，在零残差历史上，把未来分配、物理探查协方差、纠偏能力与材料相位约束联合写成一个精确半定可行性问题。第三，给出覆盖全部观测历史的增益修复上界。前两项的精确性不能替代第三项的全部历史量词；第二项特别限定为满足名义无偏条件 $HC=I_m$ 的线性分配架构。允许带偏差或任意非线性共同动作的问题需要另行分析。

\subsection{物理尺度、数据预算与代价的确切含义}
本节允许显式保留固定物理权重 $W$；若输入已按主文的物理度量白化，取 $W=I$。$W$ 不随设计改变，$M=NW^{-1/2}$ 仅为本节的加权通道记号。


固定一个局部模型、任务特征与允许输入子空间。除非另行说明，本节先令结构外余项为零：
\begin{equation}
 J=H+\Theta N,\qquad \|\Theta\|_2\le\delta,
 \qquad R_{\rm data}=\Theta X+V,\quad X=NU,
 \quad\|V\|_2\le\varepsilon,
 \label{v44:eq:energy-data}
\end{equation}
其中 $\delta>0$、$\varepsilon>0$，$U=[u_1,\ldots,u_p]$ 由实际执行的积分运动组成。输入物理度量 $W\succ0$ 固定，令
\begin{equation}
 P=W^{1/2}U,\qquad M=NW^{-1/2},\qquad
 K=MM^\top=NW^{-1}N^\top,
 \qquad \overline C=W^{1/2}C.
 \label{v44:eq:probe-whitening}
\end{equation}
于是 $X=MP$，未来分配的未知暴露为
\begin{equation}
 A=NC=M\overline C.
 \label{v44:eq:future-exposure}
\end{equation}
本节的二次探查预算是
\begin{equation}
 \mathcal E(U)=\|P\|_F^2
 =\sum_{i=1}^p u_i^\top W u_i.
 \label{v44:eq:displacement-energy}
\end{equation}
它是归一化平方积分运动的预算，不自动等于机械功，也不自动等于连续速度的平方积分。若每条探查持续相同时间 $\tau$，并以常速度 $u_i/\tau$ 执行，则其连续二次速度代价恰为 $\mathcal E(U)/\tau$。不等时长情形为 $\sum_i u_i^\top W u_i/\tau_i$，应使用该真实代价重新优化。

式~\eqref{v44:eq:energy-data} 的 $\varepsilon$ 是整个实际数据矩阵的确定性谱范数界。所有关于改变探查数量或幅值的结论，均以该界在所比较设计和历史上共同有效为前提。若只有每列噪声界 $\|v_i\|\le\varepsilon_{\rm col}$，不能把同一个 $\varepsilon_{\rm col}$ 直接用作任意列数的数据矩阵谱范数界；一般只能推出 $\|V\|_2\le\varepsilon_{\rm col}\sqrt p$。末节给出这种区别导致的不可消除误差下界。

物理输入椭球可包括每次积分运动的界
\begin{equation}
 \|W^{1/2}u_i\|_2\le a_{\rm p},
 \label{v44:eq:individual-probe-amplitude}
\end{equation}
以及未来纠偏的能力集合。例如，在已扣除保持与推进所需能力后，若纠偏半径为 $r_\zeta$、可用输入半径为 $u_{\rm res}$，则
\begin{equation}
 \|W^{1/2}C\|_2\le c_{\rm cap}:=u_{\rm res}/r_\zeta
 \label{v44:eq:ellipsoid-correction-cap}
\end{equation}
保证全部 $\|\zeta\|_2\le r_\zeta$ 下的纠偏输入可执行。零前馈时这是椭球能力的精确条件；使用预留半径处理非零前馈时，它是明示的充分内近似。

若 $F$ 是已知材料相位或其它需限制的运动映射，还可以要求
\begin{equation}
 FC=0 \quad\text{或}\quad \|FC\|_2\le\gamma_{\rm corr},
 \qquad \|FU\|_F^2\le b_{\rm probe}^2.
 \label{v44:eq:phase-budgets}
\end{equation}
最后一个量是探查相位的二次预算，不是净相位位移，也不蕴含每个探查中间状态的点态相位界。实际回位、声窗与共同冻结窗口的存在性仍须单独验收。

\subsection{固定分配下的精确完整激励预算}

\begin{theorem}[固定未来分配的最小完整激励能量]
\label{v44:thm:fixed-C-energy}
固定 $C$，令 $A=NC$，并给定 $\kappa>0$。在式~\eqref{v44:eq:energy-data} 的数据模型中，考虑同时满足
\begin{equation}
 \operatorname{range}(A)\subseteq\operatorname{range}(X),
 \qquad \varepsilon\|X^\dagger A\|_2\le\kappa
 \label{v44:eq:fullcoverage-target}
\end{equation}
的全部有限实际探查矩阵。其最小二次物理预算为
\begin{equation}
 \boxed{
 \mathcal E_{\rm full}(C,\kappa)
 =\frac{\varepsilon^2}{\kappa^2}
       \operatorname{tr}(A^\top K^\dagger A)
 =\frac{\varepsilon^2}{\kappa^2}\|M^\dagger NC\|_F^2.}
 \label{v44:eq:fixed-C-exact-energy}
\end{equation}
若 $A\ne0$，最优探查只需 $d=\operatorname{rank}(A)$ 个独立方向。若 $A=0$，空探查达到零预算。

另施加每次积分运动界~\eqref{v44:eq:individual-probe-amplitude}，并允许重复方向及任意有限列数，则所需最少列数恰为
\begin{equation}
 p_{\min}
 =\max\left\{d,
 \left\lceil\frac{\mathcal E_{\rm full}(C,\kappa)}{a_{\rm p}^2}\right\rceil
 \right\}
 \quad(A\ne0).
 \label{v44:eq:fixed-C-minimum-columns}
\end{equation}
该列数结论继续使用同一个总谱噪声界 $\varepsilon$。
\end{theorem}

\begin{proof}
记 $G=XX^\top$。条件~\eqref{v44:eq:fullcoverage-target} 等价于
\begin{equation}
 AA^\top\preceq\frac{\kappa^2}{\varepsilon^2}G.
 \label{v44:eq:Gram-domination}
\end{equation}
这由 $\|X^\dagger A\|_2^2=\lambda_{\max}(A^\top G^\dagger A)$ 及相应的值域条件得到；也可直接由广义 Schur 补证明。

固定 $X$ 后，满足 $MP=X$ 的最小 Frobenius 范数解是 $P=M^\dagger X$。任意其它解为 $M^\dagger X+P_0$，其中 $MP_0=0$，两项逐列正交。因此
\begin{equation}
 \mathcal E(U)\ge\|M^\dagger X\|_F^2
 =\operatorname{tr}(K^\dagger G)
 \ge\frac{\varepsilon^2}{\kappa^2}
       \operatorname{tr}(K^\dagger AA^\top).
 \label{v44:eq:fixed-energy-lower}
\end{equation}
最后一步使用式~\eqref{v44:eq:Gram-domination} 和 $K^\dagger\succeq0$。

对 $A$ 取紧致 SVD $A=L_d\Sigma_d V_d^\top$，并选择
\begin{equation}
 X_* =\frac{\varepsilon}{\kappa}L_d\Sigma_d,
 \qquad P_*=M^\dagger X_*,
 \qquad U_*=W^{-1/2}P_*.
 \label{v44:eq:energy-optimal-probes}
\end{equation}
因为 $\operatorname{range}(A)\subseteq\operatorname{range}(M)$，有 $MP_*=X_*$。又有
\begin{equation}
 X_*^\dagger A=\frac{\kappa}{\varepsilon}V_d^\top,
 \qquad X_*X_*^\top=\frac{\varepsilon^2}{\kappa^2}AA^\top.
\end{equation}
故下界达到，并且 $\operatorname{rank}(X_*)=d$。

若每列满足幅值界，则任何可行矩阵都有
$p\ge d$ 和 $p a_{\rm p}^2\ge\mathcal E(U)\ge\mathcal E_{\rm full}$。
反过来，对任何满足这两项的整数 $p$，在 $P_*$ 后补零列，并作右正交变换，使所有列的平方范数都等于 $\mathcal E_{\rm full}/p$。这种正交变换存在：对 $P_*^\top P_*$ 使用实对称矩阵的常对角构造，或逐次用二维平面旋转平衡一个高于均值和一个低于均值的对角元。右正交变换不改变 $P_*P_*^\top$、$X_*X_*^\top$ 及上述证书。因而每列幅值均不超过 $a_{\rm p}$。
\end{proof}

\begin{remark}[这是证书的精确预算，不是任意策略的无条件信息下界]
\label{v44:rem:fullcoverage-scope}
定理~\ref{v44:thm:fixed-C-energy} 精确优化的是既消除未激励项、又满足指定最小二乘放大界的设计。先验球已足够小、允许残余未激励方向，或者采用不同的数据一致集合控制器时，可能需要更少能量。因此它不能直接命名为所有策略的最小校准能量。下面的零历史集合半径将显式纳入先验球，而最后的全部历史上界将处理数据后增益修复。
\end{remark}

\begin{corollary}[与纠偏能力和相位相容的联合完整激励预算]
\label{v44:cor:joint-fullcoverage}
令 $\mathcal K$ 是一个固定闭凸分配集合，包含物理能力、精确接触以及式~\eqref{v44:eq:phase-budgets} 中针对 $C$ 的约束。考虑零残差历史后的名义无偏分配，即 $HC=I_m$。若 $\mathcal K$ 有界，则
\begin{equation}
 \boxed{
 \min_{\substack{HC=I_m,\ C\in\mathcal K\\ U\ {
m 满足}\ \eqref{v44:eq:fullcoverage-target}}}
 \mathcal E(U)
 =\frac{\varepsilon^2}{\kappa^2}
 \min_{HC=I_m,\ C\in\mathcal K}\|M^\dagger NC\|_F^2.}
 \label{v44:eq:joint-energy-C}
\end{equation}
若分配集合非空，右侧是有达到解的凸二次优化，并可由定理~\ref{v44:thm:fixed-C-energy} 构造达到左侧的实际输入；分配集合为空时两侧按 $+\infty$ 解释。
\end{corollary}

这里最小化的是未来实际采用的 $NC$，而非预先锁定的一个 $NB$。能力或相位约束可以使原本零暴露的右逆不可执行，从而重新产生非零信息预算。此耦合是物理任务分配问题的一部分，不能仅由已知回归目标下的试验设计公式替代。

\paragraph{一个可算的能力分离例。}
取 $m=1,n=2$、$H=[1\ \ 1]$、$N=[1\ \ 0]$、$W=I_2$，并施加 $\|C\|_2\le b$，其中 $1/\sqrt2\le b\le1$。任何名义无偏分配为 $C=(c,1-c)^\top$，所以允许的最小未知暴露为
\begin{equation}
 a_b=\min|c|=\frac{1-\sqrt{2b^2-1}}2.
 \label{v44:eq:capacity-coupling-example}
\end{equation}
因此联合完整激励预算为 $(\varepsilon a_b/\kappa)^2$。当 $b=1$，可用 $C=e_2$，预算为零；当 $1/\sqrt2\le b<1$，同一零暴露分配不再满足能力合同，且 $a_b>0$。若还要求 $FC=0$、$F=[0\ \ 1]$，则必须 $C=e_1$，需要 $b\ge1$。这个例子说明校准、输入能力和相位相容性是不同但相互制约的条件。

\subsection{从二次探查预算到有限探查时间}

下面只讨论执行探查本身的运动时间。假定允许恒定方向的运动，每个探查从可重复的初始状态开始，速度满足 $\|W^{1/2}\dot u\|\le v_{\rm p}$；这里 $u$ 表示该探查积累的位移。必要回位也按同一速度执行，暂不计固定观测开销和加速度过渡。若这些运动不能留在共同模型窗口内，则本段时间结论不能使用。

\begin{proposition}[固定分配完整激励的运动时间]
\label{v44:prop:probe-time}
记
\begin{equation}
 E_*:=\mathcal E_{\rm full}(C,\kappa),\qquad
 L_*:=\frac{\varepsilon}{\kappa}\|M^\dagger NC\|_*.
\end{equation}
其中 $\|\cdot\|_*$ 是核范数。无每次位移幅值限制时，最少单程探查运动时间恰为 $L_*/v_{\rm p}$。有式~\eqref{v44:eq:individual-probe-amplitude} 时，最优单程时间 $T_*$ 满足
\begin{equation}
 \frac1{v_{\rm p}}\max\{L_*,E_*/a_{\rm p}\}
 \le T_*
 \le\frac1{v_{\rm p}}(L_*+E_*/a_{\rm p})
 \le\frac2{v_{\rm p}}\max\{L_*,E_*/a_{\rm p}\}.
 \label{v44:eq:time-factor-two}
\end{equation}
若显式采用每次原路回位、且不把回程记录用于本证书数据的构造，则该构造的往返运动时间是其单程运动时间的两倍。因而无幅值限制时有构造上界 $2L_*/v_{\rm p}$，有幅值限制时有构造上界 $2(L_*+E_*/a_{\rm p})/v_{\rm p}$。若允许回程也提供同一冻结模型的有效数据，这些仍是忽略回程信息的一种构造上界，不能把上述单程最优值或下界一概乘以二作为全部策略的结论。已知的观测开销和可行加速度过渡时间另行加上。
\end{proposition}

\begin{proof}
任何输入矩阵都有 $\sum_i\|p_i\|\ge\|P\|_*$。又由 $P=M^\dagger X+P_0$ 的正交分解及式~\eqref{v44:eq:Gram-domination}，有
\begin{equation}
 \|P\|_*\ge\|M^\dagger X\|_*
 \ge\frac{\varepsilon}{\kappa}\|M^\dagger A\|_*=L_*.
\end{equation}
第二个不等式来自对相应半正定 Gram 矩阵的平方根迹单调性。对最优矩阵 $P_*$ 作右正交变换，使其非零列是物理输入协方差的特征向量乘以奇异值，列长度之和恰为 $L_*$，从而达到无幅值限制的下界。

若所有列长不超过 $a_{\rm p}$，则
$E_*\le\sum_i\|p_i\|^2\le a_{\rm p}\sum_i\|p_i\|$，得到另一项下界。设最优物理协方差的各特征方向长度为 $b_j$。把该方向拆成 $k_j=\lceil b_j^2/a_{\rm p}^2\rceil$ 个等长探查，每个长度为 $b_j/\sqrt{k_j}$，保持协方差不变且满足幅值界。其总长度满足
\begin{equation}
 \sum_j b_j\sqrt{k_j}
 \le\sum_j(b_j+b_j^2/a_{\rm p})
 =L_*+E_*/a_{\rm p}.
\end{equation}
用最大允许速度分别执行即可。
\end{proof}

若 $\varepsilon$ 中含有随总探查时间增长的漂移，或随列数增长的观测误差，上述比较应使用真实的 $\varepsilon(U,T,p)$，而不能固定旧预算后宣称缩短或延长时间仍然满足同一合同。

\subsection{纳入先验球的零历史响应半径}

在全零残差历史上，允许的未知响应为
\begin{equation}
 \mathfrak T_X=
 \{\Theta:\|\Theta\|_2\le\delta,\ \|\Theta X\|_2\le\varepsilon\}.
 \label{v44:eq:zero-history-consistency}
\end{equation}
对满足 $HC=I_m$ 的分配，其精确共同响应半径为
\begin{equation}
 \mathfrak R_X(C)=\sup_{\Theta\in\mathfrak T_X}\|\Theta NC\|_2.
 \label{v44:eq:exact-zero-radius}
\end{equation}

\begin{lemma}[双居中椭球的经典半径工具]
\label{v44:lem:two-ellipsoid}
令 $A=NC$、$G=XX^\top$。则
\begin{align}
 \mathfrak R_X(C)^2
 &=\max_{\substack{z^\top z\le\delta^2\\ z^\top Gz\le\varepsilon^2}}
          z^\top AA^\top z\nonumber\\
 &=\min_{\lambda,\mu\ge0}
 \{\lambda\delta^2+\mu\varepsilon^2:
       AA^\top\preceq\lambda I_r+\mu G\}.
 \label{v44:eq:two-ellipsoid-dual}
\end{align}
这是 Foucart--Liao 最优恢复半径结果的有限维特例\cite{foucart2024}。以下保留一个自包含证明，以说明本稿谱范数矩阵模型如何化为该经典半径。
\end{lemma}

\begin{proof}
若 $\Theta$ 可行，对任何任务输出单位向量 $b$，$z=\Theta^\top b$ 都满足前两项二次约束。因此 $\|\Theta A\|=\sup_{\|b\|=1}\|A^\top\Theta^\top b\|$ 不超过右侧向量极大值。反过来，对任意可行 $z$ 和任意单位 $b$，秩一矩阵 $\Theta=bz^\top$ 满足原始两项谱范数约束，并有 $\|\Theta A\|=\|A^\top z\|$。这证明第一行。

把 $zz^\top$ 提升为 $Y\succeq0$，考虑
\begin{equation}
 \max\ \operatorname{tr}(AA^\top Y)
 \quad\text{s.t.}\quad
 \operatorname{tr}Y\le\delta^2,
 \operatorname{tr}(GY)\le\varepsilon^2,
 \quad Y\succeq0.
\end{equation}
可行集紧，故线性目标在极点达到。任何秩至少为二的 $Y$ 都不是极点：在其值域内，对称扰动空间维数至少为三，因而存在非零对称扰动同时保持上述两个迹；足够小的正、负扰动均保持半正定。故存在秩至多为一的最优解，提升精确。由于 $\delta,\varepsilon>0$，充分小的正定 $Y$ 满足严格可行性，半定规划强对偶给出第二行且对偶达到。
\end{proof}

\begin{theorem}[零历史下的物理探查--分配联合精确前沿]
\label{v44:thm:joint-zero-SDP}
给定目标 $q\in[0,1)$、二次能量预算 $E\ge0$、探查相位预算 $b_{\rm probe}\ge0$，并令 $\mathcal K$ 为针对 $C$ 的固定凸能力与相位合同。允许任意有限列数的探查。以下两项等价：
\begin{enumerate}
 \item 存在实际输入矩阵 $U$ 及分配 $C$，满足
 \[
 HC=I_m,\quad C\in\mathcal K,\quad
 \mathcal E(U)\le E,\quad\|FU\|_F^2\le b_{\rm probe}^2,
 \quad\mathfrak R_{NU}(C)\le q.
 \]
 \item 存在 $C$、$Z\succeq0$ 及 $\lambda,\mu\ge0$，满足
 \begin{equation}
 \boxed{
 \begin{gathered}
 HC=I_m,\qquad C\in\mathcal K,\\
 \begin{bmatrix}
 \lambda I_r+MZM^\top&NC\\
 C^\top N^\top&I_m
 \end{bmatrix}\succeq0,\\
 \delta^2\lambda+\varepsilon^2\mu\le q^2,\qquad
 \operatorname{tr}Z\le\mu E,\\
 \operatorname{tr}(\Gamma Z\Gamma^\top)
       \le\mu b_{\rm probe}^2,
 \qquad\Gamma=FW^{-1/2}.
 \end{gathered}}
 \label{v44:eq:joint-zero-SDP}
 \end{equation}
\end{enumerate}
若 $\mathcal K$ 由半定可表示的椭球、行范数和线性约束构成，则固定预算 $E$ 后，这是一个半定可行性问题。最小预算可以通过对 $E$ 的可行性二分求取；不将 $E$ 与 $\mu$ 同时作为线性变量。
\end{theorem}

\begin{proof}
对可行的 $U,C$，令 $S=PP^\top$。引理~\ref{v44:lem:two-ellipsoid} 给出达到半径界的 $\lambda,\mu$。取 $Z=\mu S$，各个预算约束直接成立，块半正定约束由 Schur 补等价于
$AA^\top\preceq\lambda I_r+\mu MSM^\top$。

反过来，若式~\eqref{v44:eq:joint-zero-SDP} 可行且 $\mu>0$，令 $S=Z/\mu$，取任意有限因子 $P$ 使 $PP^\top=S$，并令 $U=W^{-1/2}P$。输入能量和相位预算成立，双椭球引理给出 $\mathfrak R_{NU}(C)\le q$。若 $\mu=0$，迹约束和 $Z\succeq0$ 强制 $Z=0$；此时 $\delta^2\|NC\|_2^2\le q^2$，空探查已经可行。
\end{proof}

\begin{remark}[列数、幅值与自适应量词]
\label{v44:rem:joint-SDP-scope}
若预先允许 $p\ge n$ 列，并要求式~\eqref{v44:eq:individual-probe-amplitude}，可以在式~\eqref{v44:eq:joint-zero-SDP} 中另加 $\operatorname{tr}Z\le\mu p a_{\rm p}^2$；常对角构造给出满足每列幅值界的因子，并且保留两项总预算。若 $p<n$，还需要 $\operatorname{rank}(Z)\le p$，一般不再是凸约束。因此本定理没有同时解决最少方向数的非凸问题。

该前沿也给出一个自适应算法的必要界，但范围必须准确：对任何最终在零历史采用 $HC=I_m$ 的线性分配、且满足相同能力合同的因果探查算法，沿全部报告为零的历史记录最终 $U,C$。每个满足~\eqref{v44:eq:zero-history-consistency} 的 $\Theta$ 都可以配以 $V=-\Theta NU$，使算法经历同一历史。所以若上述 SDP 在预算 $E$ 下不可行，该架构不能以该预算保证全部模型、全部历史上的响应误差不超过 $q$。这不是对允许 $HC\ne I_m$ 的带偏差修复，或任意非线性共同动作的普遍必要界；有限能力下，这些更大策略类可能有更好的权衡。

反方向只给零历史上的可达到性。非零历史会改变已知响应与一致集合的几何，不能从该 SDP 的零历史可行性直接推出所有历史保证。
\end{remark}

\paragraph{固定 $C$、不限制探查相位时的闭式降维。}
上述 SDP 还给一个便于比较先验与噪声的单变量公式。若 $q\ge\delta\|A\|_2$，空探查可行，最小预算为零。若 $0<q<\delta\|A\|_2$，固定 $C$ 的零历史最小预算为
\begin{equation}
 \mathcal E_{\rm prior}(C,q)
 =\varepsilon^2
 \min_{0\le\lambda<q^2/\delta^2}
 \frac{
 \operatorname{tr}\!\left[
 K^{\dagger/2}(AA^\top-\lambda I_r)K^{\dagger/2}
 \right]_+
 }{q^2-\delta^2\lambda}.
 \label{v44:eq:prior-energy-profile}
\end{equation}
方括号下标 $+$ 表示对称矩阵的正谱部分。为验证该式，先固定双椭球乘子 $\lambda$，取允许的最大 $\mu=(q^2-\delta^2\lambda)/\varepsilon^2$；再在 $\operatorname{range}(N)$ 内最小化探查 Gram 的物理加权迹。变量代换后，该最小值是所给对称矩阵正谱部分的迹。$\lambda=0$ 恢复完整激励预算；$\lambda>0$ 允许先验承担一部分方向，不要求未来暴露被全部探查覆盖。该式也是经典半径引理在本结构模型中的代价剖面，而非额外的全部历史最小最大定理。

\subsection{全部观测历史上的可执行增益修复上界}

\begin{theorem}[完整激励后的全部历史恢复与能力上界]
\label{v44:thm:all-history-energy-upper}
在式~\eqref{v44:eq:energy-data} 的模型中，预先选择 $B$ 满足 $HB=I_m$，记
\begin{equation}
 A_B=NB,\quad F_B=M^\dagger A_B,\quad
 \rho_B=\delta\|A_B\|_2<1,\quad
 b_B=\|W^{1/2}B\|_2.
\end{equation}
给定 $0<\beta<1-\rho_B$。用定理~\ref{v44:thm:fixed-C-energy} 构造探查，使
\begin{equation}
 \operatorname{range}(A_B)\subseteq\operatorname{range}(X),
 \qquad\varepsilon\|X^\dagger A_B\|_2\le\beta,
 \qquad\mathcal E(U)=\frac{\varepsilon^2}{\beta^2}\|F_B\|_F^2.
 \label{v44:eq:all-history-probe-budget}
\end{equation}
对任意允许的真实 $\Theta$ 与噪声矩阵 $V$，令
\begin{equation}
 \widehat\Theta=R_{\rm data}X^\dagger,\qquad
 \widehat G=I_m+\widehat\Theta A_B,
 \qquad C=B\widehat G^{-1}.
 \label{v44:eq:all-history-repair}
\end{equation}
则 $\widehat G$ 可逆、$\widehat JC=I_m$，并且
\begin{equation}
 \boxed{
 \|JC-I_m\|_2
 \le\frac{\beta}{1-\rho_B-\beta},\qquad
 \|W^{1/2}C\|_2
 \le\frac{b_B}{1-\rho_B-\beta}.}
 \label{v44:eq:all-history-q-cap}
\end{equation}
任何 $FB=0$ 的线性相位或接触约束也满足 $FC=0$。若另有针对探查的相位或位姿预算，仍须独立验收所构造的 $U$；$FB=0$ 不推出 $FU=0$。
\end{theorem}

\begin{proof}
完整激励给出
\begin{equation}
 \widehat G=I_m+\Theta A_B+VX^\dagger A_B,
 \qquad \|\widehat G-I_m\|_2\le\rho_B+\beta<1.
\end{equation}
所以 $\|\widehat G^{-1}\|_2\le(1-\rho_B-\beta)^{-1}$。又有
\begin{equation}
 JC-I_m=-VX^\dagger A_B\widehat G^{-1}.
\end{equation}
取范数即得响应界和能力界。最后 $FC=FB\widehat G^{-1}=0$。
\end{proof}

因此，若要求误差不超过 $q_{\rm goal}\in(0,1)$ 且修复能力满足 $\|W^{1/2}C\|\le c_{\rm cap}$，可以选择
\begin{equation}
 0<\beta\le
 \min\left\{
 \frac{q_{\rm goal}(1-\rho_B)}{1+q_{\rm goal}},
 1-\rho_B-\frac{b_B}{c_{\rm cap}}
 \right\}.
 \label{v44:eq:all-history-feasible-beta}
\end{equation}
右侧严格为正时，式~\eqref{v44:eq:all-history-probe-budget} 给出有限、全部历史有效的探查预算。可以在满足这些严格余量的候选 $B$ 中最小化此充分预算，但该候选优化一般不是上一节的凸问题。

忽略能力约束时，取 $\beta=q_{\rm goal}(1-\rho_B)/(1+q_{\rm goal})$，该构造相对于同一 $B$ 的零历史完整激励证书多付因子
\begin{equation}
 \left(\frac{1+q_{\rm goal}}{1-\rho_B}\right)^2.
\end{equation}
这个比较不声称相对于所有历史、所有分配和所有控制器的最优能量有同一近似比。特别是先验本身足够小或允许部分激励时，零历史最优预算可能更低；上节的精确前沿与本节的充分上界应分别陈述。

\subsection{同一固定结构外余项可以抵消哪些误差}

现在恢复结构外误差，并明确区分是否为同一个矩阵。若校准和随后执行共享同一个固定 $E_0$，令
\begin{equation}
 J=H+\Theta N+E_0,\qquad
 \|E_0W^{-1/2}\|_2\le\eta,
 \qquad R_{\rm data}=\Theta X+E_0U+V.
 \label{v44:eq:shared-geometry-error}
\end{equation}
仍取 $\widehat\Theta=R_{\rm data}X^\dagger$、$\widehat J=H+\widehat\Theta N$。对任意数据后选取的 $C$，有精确恒等式
\begin{equation}
 \boxed{
 (J-\widehat J)C
 =\Theta(I-P_X)NC-VX^\dagger NC
       +E_0(C-UX^\dagger NC).}
 \label{v44:eq:shared-remainder-identity}
\end{equation}
因此可使用比逐项重复计入余项更紧的凸证书
\begin{equation}
 \boxed{
 q_{\rm shared}(C,U)
 =\delta\|(I-P_X)NC\|_2
 +\varepsilon\|X^\dagger NC\|_2
 +\eta\|W^{1/2}(C-UX^\dagger NC)\|_2.}
 \label{v44:eq:shared-remainder-certificate}
\end{equation}
固定已执行数据 $U,X$ 后，该式仍然是线性矩阵映射的谱范数之和。若允许名义响应有偏差，总响应证书为
$\|\widehat JC-I_m\|_2+q_{\rm shared}(C,U)$，而非只保留后者。

\begin{proposition}[最小能量探查的共享余项恢复界]
\label{v44:prop:shared-remainder-all-history}
沿用定理~\ref{v44:thm:all-history-energy-upper} 的 $B$。对 $NB=L_d\Sigma_dV_d^\top$ 的紧致 SVD，给定幅值比例 $t>0$，取最小能量探查 $X=tL_d\Sigma_d$、$U=W^{-1/2}M^\dagger X$，并定义 $\beta=\varepsilon/t$；$NB=0$ 时取空探查。令
\begin{equation}
 \Pi=M^\dagger M,\qquad
 \overline B=W^{1/2}B,\qquad
 f_B=\|\Pi\overline B\|_2=\|M^\dagger NB\|_2,
 \qquad r_B=\|(I-\Pi)\overline B\|_2.
\end{equation}
若 $\rho_B+\eta f_B+\beta<1$，则对式~\eqref{v44:eq:shared-geometry-error} 中全部允许模型和全部噪声，式~\eqref{v44:eq:all-history-repair} 的修复满足
\begin{equation}
 \boxed{
 \|JC-I_m\|_2
 \le\frac{\beta+\eta r_B}{1-\rho_B-\eta f_B-\beta},
 \qquad
 \|W^{1/2}C\|_2
 \le\frac{b_B}{1-\rho_B-\eta f_B-\beta}.}
 \label{v44:eq:shared-remainder-gain}
\end{equation}
\end{proposition}

\begin{proof}
最小能量构造满足
\begin{equation}
 W^{1/2}UX^\dagger NB=M^\dagger NB=\Pi\overline B.
\end{equation}
记 $\overline E_0=E_0W^{-1/2}$，则
\begin{align}
 \widehat G
 &=I_m+\Theta NB+\overline E_0\Pi\overline B+VX^\dagger NB,\\
 JB-\widehat G
 &=\overline E_0(I-\Pi)\overline B-VX^\dagger NB.
\end{align}
分别界定 $\widehat G-I_m$ 与 $JB-\widehat G$，再右乘 $\widehat G^{-1}$ 即得。
\end{proof}

该式准确说明了提高幅值的作用。$\beta=\varepsilon/t$ 随幅值减小；共享余项作用在未由 $N$ 通道探查覆盖的物理分量上，留下 $\eta r_B$。若 $r_B=0$，同一固定余项对未来分配的作用已包含在有效探查中，不能笼统宣称任意非零 $\eta$ 都造成不可校准的绝对误差地板。

若校准余项与执行余项不同，分别为 $E_{\rm cal}$ 和 $E_{\rm exec}$，且只有
\[
 \|E_{\rm cal}W^{-1/2}\|\le\eta_{\rm cal},
 \qquad \|E_{\rm exec}W^{-1/2}\|\le\eta_{\rm exec},
\]
则上面的抵消一般不成立，只能用
\begin{equation}
 \|E_{\rm exec}B-E_{\rm cal}UX^\dagger NB\|_2
 \le\eta_{\rm exec}b_B+\eta_{\rm cal}f_B.
 \label{v44:eq:cross-stage-remainder-floor}
\end{equation}
这类跨阶段偏差不随探查幅值增加自动衰减。逐样本余项变化、图像噪声进入在线矩阵以及随接触状态改变的几何误差，除非另有共同矩阵模型，不能套用共享余项证书。

\subsection{逐样本噪声与重复探查的反例边界}

以标量模型 $z_i=\theta x_i+v_i$ 为例，设 $|\theta|\le\delta$、$|v_i|\le\varepsilon_{\rm col}$、$|x_i|\le a_{\rm p}$。在全零观测历史上，任意数量的探查都无法区分
\begin{equation}
 |\theta|\le\min\{\delta,\varepsilon_{\rm col}/a_{\rm p}\}.
\end{equation}
确实可取 $v_i=-\theta x_i$，逐列噪声均合法。因此对固定未来暴露 $A\ne0$，任何该历史上的共同估计都至少有
\begin{equation}
 |A|\min\{\delta,\varepsilon_{\rm col}/a_{\rm p}\}
 \label{v44:eq:per-column-noise-floor}
\end{equation}
的最坏误差。重复相同幅值的探查可以使总二次位移预算无限增大，却不会降低这个确定性地板。它不反驳定理~\ref{v44:thm:fixed-C-energy}，因为后者使用的是在全部列上固定的总谱噪声合同。

更一般地，若仅有每列界 $\|v_i\|\le\varepsilon_{\rm col}$ 和物理探查椭球 $\|W^{1/2}u_i\|\le a_{\rm p}$，则所有有限探查在零历史上至少保留
\begin{equation}
 \left\{z:\|z\|\le\delta,
      z^\top Kz\le(\varepsilon_{\rm col}/a_{\rm p})^2\right\}
\end{equation}
给出的秩一未知响应。由此得到任何数量探查都不能低于的固定分配响应半径
\begin{equation}
 \max_{\substack{\|z\|\le\delta\\
        z^\top Kz\le(\varepsilon_{\rm col}/a_{\rm p})^2}}
       \|C^\top N^\top z\|_2.
 \label{v44:eq:per-column-ellipsoid-floor}
\end{equation}
该结论不使用噪声独立性。只有另外给出随机噪声集中、独立重复或其它可验证的误差衰减规律，才能从重复次数推出更小的数据误差预算。

\paragraph{贡献与范围的可审查陈述。}
本节的标准工具是最小范数输入提升、Gram 支配、双椭球最优恢复与 Schur 补。新增的任务对象是把实际输入度量、未来可重选的 $NC$、纠偏能力和相位合同放入同一个资源可行域，并严格区分零历史精确前沿与全部历史修复保证。共享固定余项的抵消恒等式则解释了哪些结构外误差真正形成未辨识地板。上述结论不宣称解决任意非线性控制、未知相位映射、一般跨窗口变化或所有噪声模型下的最小总扫描时间。


\section{有噪声的实际探查与凸响应修复}
\label{sec:cal-noisy}

\subsection{使用实际执行量构造估计}

先设实际积分运动被准确记录，并将有效探查堆叠为 $U$。输入记录误差由第~\ref{rev:sec:input-measurement}节统一加入数据预算。令
\begin{equation}
  X=NU,\qquad
  R_{\mathrm{data}}=\Delta Y-HU=\Theta X+W,
  \qquad\|W\|_2\le\varepsilon_{\rm data} .
  \label{eq:cal-data}
\end{equation}
$\varepsilon_{\rm data}$ 属于具体探查输入 $U$ 和局部时间窗口；改变幅度、探查库或窗口后应重算该界。$W$ 包含测量噪声、重复探查间的未抵消漂移以及 $E_{\mathrm{geo}}U$。如果两次探查共享相同初始状态或有可信的基线消除，可降低漂移项；不能默认它自动消失。由此定义
\begin{equation}
  \widehat\Theta=R_{\mathrm{data}}X^\dagger,
  \qquad
  \widehat J=H+\widehat\Theta N,
  \qquad P_X=XX^\dagger .
  \label{eq:cal-estimate}
\end{equation}
因为
\begin{equation}
  \widehat\Theta
  =\Theta P_X+WX^\dagger,
\end{equation}
对任意修复实现 $C$，
\begin{align}
  (J-\widehat J)C
  &=
  \Theta(I-P_X)NC
  -WX^\dagger NC+E_{\mathrm{geo}}C,\\
  \|(J-\widehat J)C\|_2
  &\le
  \underbrace{
  \delta\|(I-P_X)NC\|_2}_{\text{未激励方向}}
  +
  \underbrace{
  \varepsilon_{\rm data}\|X^\dagger NC\|_2}_{\text{估计误差放大}}
  +
  \underbrace{
  \eta\|C\|_2}_{\text{结构外误差}} .
  \label{eq:cal-q}
\end{align}
记右侧为 $q(C,X)$。无探查时取空数据矩阵、$P_X=0$ 和 $\widehat\Theta=0$；此时 $q(C,\varnothing)=\delta\|NC\|+\eta\|C\|$。它同时检查是否激励了正确方向，以及激励是否具有足够的数值条件。单纯提高 $\operatorname{rank}(X)$ 不保证少量低幅度探查在噪声中有效。

\subsection{输入测量误差如何进入数据合同}
\label{rev:sec:input-measurement}

实际执行运动与其记录值分别记为 $U$ 和
$\widetilde U=U+\Delta U$。在同一冻结窗口中，设
\begin{equation}
 \Delta Y=(H+\Theta N+E_0)U+V,\qquad
 \|V\|_2\le\varepsilon_Y,\quad
 \|\Delta U\|_2\le\varepsilon_U,\quad
 \|E_0\|_2\le\eta .
 \label{rev:eq:measured-input-model}
\end{equation}
$V$ 包括输出测量误差和已经有界的未抵消漂移。由记录量构造
\begin{equation}
 \widetilde X=N\widetilde U,\qquad
 \widetilde R=\Delta Y-H\widetilde U .
 \label{rev:eq:measured-data}
\end{equation}
则精确代数关系为
\begin{align}
 \widetilde R
 &=\Theta\widetilde X+\widetilde W,\\
 \widetilde W
 &=V+E_0U-(H+\Theta N)\Delta U .
 \label{rev:eq:input-error-residual}
\end{align}
利用 $\|U\|_2\le\|\widetilde U\|_2+\varepsilon_U$，可选
\begin{equation}
 \boxed{
 \varepsilon_{\rm data}
 =
 \varepsilon_Y+\eta\|\widetilde U\|_2
 +\big(\|H\|_2+\delta\|N\|_2+\eta\big)\varepsilon_U .
 }
 \label{rev:eq:input-error-budget}
\end{equation}
以 $(X,R_{\rm data})=(\widetilde X,\widetilde R)$ 代入估计和 $q$ 的公式即可，无须假定输入测量没有误差，也无须对输入误差与输出噪声作独立性假设。

若每条探查对应不同余项 $E_i$，仅有逐项界 $\|E_i\|_2\le\eta$，则应使用
\begin{equation}
 \|[E_1u_1,\ldots,E_pu_p]\|_2
 \le\eta\|U\|_F .
 \label{rev:eq:varying-remainder-bound}
\end{equation}
此时若另有 $\|\Delta U\|_F\le\varepsilon_{U,F}$，一个有效替代为
\begin{equation}
 \varepsilon_{\rm data}
 =
 \varepsilon_Y+\eta\|\widetilde U\|_F
 +\eta\varepsilon_{U,F}
 +(\|H\|_2+\delta\|N\|_2)\varepsilon_U .
 \label{rev:eq:varying-remainder-budget}
\end{equation}
响应参数在探查间的变化同样需要给出数据残差矩阵的界；仅声明每次误差较小，不能省略堆叠数据后的算子范数验收。

\subsection{固定探查数据后的可计算修复}

设真实输入的一个物理内近似由 $a_i^\top u_{\rm act}\le b_i^{\rm act}$ 给出。若执行误差属于独立有效集合 $\mathcal D_u$，令 $b_i=b_i^{\rm act}-h_{\mathcal D_u}(a_i)$；其中 $h_{\mathcal D_u}$ 是支持函数。以下 $\mathcal U$ 始终表示已经紧化的命令内近似：
\begin{equation}
  \mathcal U=
  \{u:a_i^\top u\le b_i,\ i=1,\ldots,p\}.
\end{equation}
局部误差可信区间为 $\|e\|_2\le r_{\mathrm{task}}$，二阶纠偏的预定可信速度半径为 $r_\zeta$（下文由增广状态界验收），已知前馈命令为 $u_{\mathrm{ff}}$，并先检查 $Eu_{\mathrm{ff}}=0$。固定实际探查数据后可解
\begin{equation}
\begin{aligned}
  \min_C\quad&
  \|C-C_{\mathrm{transported}}\|_{W_C}^2+
  \lambda\|C\|_F^2\\
  \text{s.t.}\quad&
  \widehat JC=I_m,\\
  &q(C,X)\le q_{\mathrm{goal}}<1,\\
  &EC=0,\\
  &a_i^\top u_{\mathrm{ff}}
   +r_\zeta\|a_i^\top C\|_2
   \le b_i,\quad i=1,\ldots,p.
\end{aligned}
\label{eq:cal-convex-repair}
\end{equation}
其中 $W_C$ 为固定正定代价度量，$\lambda\ge0$，$EC=0$ 与预先检查的 $Eu_{\rm ff}=0$ 共同保持名义总输入的线性接触约束。该问题是凸的：等式为仿射约束，$q$ 是线性矩阵映射的谱范数之和，能力条件是行范数约束。最后一组不等式保证所有可信纠偏速度状态下
\begin{equation}
  u=u_{\mathrm{ff}}+C\zeta\in\mathcal U\quad\text{当 }\|\zeta\|_2\le r_\zeta .
\end{equation}

若模型与物理检查使可行集为空，应增加有依据的探查、改变允许实现或报告该保证不可实现，而不是放宽未验证的误差半径。dIRD 的梯度可以用于提出候选实现，但其密度或可达性代理值不代替 $\mathcal U$ 的真实验收。




\paragraph{允许名义带偏差的修复。}
输入能力紧张时，可以不要求 $\widehat JC=I_m$，而直接使用
\begin{equation}
 q_{\rm gen}(C,X)=\|\widehat JC-I_m\|_2+q(C,X),
 \qquad \|JC-I_m\|_2\le q_{\rm gen}(C,X).
 \label{v44:eq:general-biased-certificate}
\end{equation}
固定数据后，$q_{\rm gen}$ 仍为凸函数。以 $q_{\rm gen}\le q_{\rm goal}$ 替换
名义无偏等式与原响应约束，得到更大的凸验收集合。
后续恢复定理只需实际 $\Delta=JC-I_m$ 的范数界，因此同样适用。
第~\ref{sec:v44:scalar-frontier}节的截断区间逆允许这种带偏差实现；
第~\ref{v44:sec:energy-frontier}节的零历史精确 SDP 则保留其 $HC=I_m$ 架构假设。
同一固定余项在校准与执行之间共享时，还可采用式~\eqref{v44:eq:shared-remainder-certificate}。



\section{经校准的二阶图像纠偏、能力验收与滚动未来管}
\label{sec:cal-recovery}

\subsection{参考生成与实际反馈是两个不同对象}

第~\ref{sec:rev:skill}节的材料几何生成器描述材料技能的名义形状、姿态偏好和推进节奏；它的内部状态不是实际超声特征的观测值。令 $y_r(t)$ 为已接受实现对应的任务特征参考，$y(t)$ 为当前可信图像及相关传感器提供的 $m$ 维特征，定义 $e=y-y_r$。技能输出可以包含更多自由度，但不能把不可观测的技能变量加入 $e$ 后声称实现了反馈。

在共同的物理白化、平行运输任务标架中，实际图像误差测量为 $\widehat e=e+n$。推荐控制器只增加一个任务纠偏速度状态 $\zeta\in\R^m$：
\begin{equation}
 \boxed{u=u_{\mathrm{ff}}+C\zeta,\qquad
 \dot\zeta=-\omega^2(e+n)-2\omega\zeta,\quad \omega>0.}
 \label{eq:cal-servo-second-order}
\end{equation}
它替换原来的静态误差到速度映射。$\zeta$ 是虚拟纠偏速度，$\omega$ 是固定设计带宽；这不是对患者施加的弹簧力。GP 学习参考生成器的形状驱动项；实际图像反馈的 $\omega$ 不由普通示教轨迹辨识，避免将两套弹簧重复作用于同一误差。

无偏修复使用式~\eqref{eq:cal-q}，带偏差修复使用式~\eqref{v44:eq:general-biased-certificate}。统一定义
\begin{equation}
 \Delta=JC-I_m,\qquad
 \|\Delta(t)\|_2\le q<1,
 \qquad
 w=Ju_{\mathrm{ff}}-\dot y_r+w_{\mathrm{ext}}.
 \label{eq:cal-feedback-residual}
\end{equation}
在有效局部窗口内，实际误差动力学是
\begin{equation}
 \boxed{\dot e=(I_m+\Delta)\zeta+w,\qquad
 \dot\zeta=-\omega^2(e+n)-2\omega\zeta.}
 \label{eq:cal-error-dynamics}
\end{equation}
若采用 $u_{\mathrm{ff}}=C\dot y_r$，仍有 $w=\Delta\dot y_r+w_{\mathrm{ext}}$；不能默认名义前馈完全抵消未知响应。身体运动补偿误差、目标漂移、图册运输误差和实际执行偏差应合并进入经过验证的 $w$ 或 $n$。

在 $\Delta=0,w=n=0$ 时，$\dot e=\zeta$，所以
\begin{equation}
 \ddot e+2\omega\dot e+\omega^2e=0,
 \quad
 e(t)=\bigl[e(0)+(\zeta(0)+\omega e(0))t\bigr]e^{-\omega t}.
 \label{eq:cal-nominal-second-order}
\end{equation}
名义临界阻尼意味着这种局部模型的双实根，不意味着任意初始速度、执行延迟或耦合误差下所有分量都永不超调。无需对噪声图像误差做数值微分。

\subsection{从实际探查证书到二阶指数恢复}

\begin{theorem}[时变耦合与测量噪声下的二阶恢复]
\label{thm:cal-recovery}
假定在同一有效任务标架内 $\omega$ 固定、$\Delta=JC-I_m$ 且 $\|\Delta(t)\|_2\le q<1$，且
$\|w(t)\|_2\le w_{\max}$、$\|n(t)\|_2\le n_{\max}$。定义
\begin{equation}
 v=\zeta/\omega,\quad r=e+v,\quad
 \chi=\begin{bmatrix}r\\v\end{bmatrix},\quad
 V_s=\tfrac12\|\chi\|_2^2,\quad S=\sqrt{2V_s},
 \quad \lambda_s=\frac{\omega(1-q)}2.
 \label{eq:cal-servo-energy}
\end{equation}
令 $b_{\rm drv}=w_{\max}+\sqrt2\,\omega n_{\max}$，则
\begin{equation}
 \begin{aligned}
 S(t)&\le e^{-\lambda_st}S(0)
       +\frac{b_{\rm drv}}{\lambda_s}(1-e^{-\lambda_st}),\\
 \|e(t)\|_2&\le\sqrt2S(t),\qquad
 \|\zeta(t)\|_2\le\omega S(t).
 \end{aligned}
 \label{eq:cal-continuous-tube}
\end{equation}
此结论不要求 $\Delta$ 对称，也不要求其随时间恒定。
\end{theorem}

\begin{proof}
由式~\eqref{eq:cal-error-dynamics}直接变换得
\begin{equation}
 \dot r=-\omega r+\omega\Delta v+w-\omega n,
 \qquad
 \dot v=-\omega r-\omega v-\omega n.
 \label{eq:cal-servo-rv}
\end{equation}
于是
\begin{align}
 \dot V_s
 &=-\omega(\|r\|_2^2+\|v\|_2^2)
   +\omega r^\top(\Delta-I_m)v
   +r^\top w-\omega(r+v)^\top n\\
 &\le-\omega(\|r\|_2^2+\|v\|_2^2)
   +\omega(1+q)\|r\|_2\|v\|_2
   +S\bigl(\|w\|_2+\sqrt2\omega\|n\|_2\bigr)\\
 &\le-\omega(1-q)V_s+Sb_{\rm drv}.
 \label{eq:cal-servo-lyapunov}
\end{align}
最后一步使用 $2\|r\|_2\|v\|_2\le\|r\|_2^2+\|v\|_2^2$。在 $S>0$ 时除以 $S$，零点使用上 Dini 导数，得到
\begin{equation}
 D^+S\le-\lambda_sS+b_{\rm drv}.
\end{equation}
比较积分即得~\eqref{eq:cal-continuous-tube}。因为 $e=r-v$、$\zeta=\omega v$，得到误差与速度界。
\end{proof}

任务的最低恢复率 $\lambda_{\mathrm{task}}$ 与容许稳态误差 $\epsilon_{\mathrm{task}}$ 可以直接验收为
\begin{equation}
 \lambda_s\ge\lambda_{\mathrm{task}},\qquad
 \frac{\sqrt2b_{\rm drv}}{\lambda_s}\le\epsilon_{\mathrm{task}}.
 \label{eq:cal-task-rate}
\end{equation}
因此 $q$ 不只是响应模型的预测误差；它决定恢复速度、残差管与所需执行能力。提高 $\omega$ 会加快名义恢复，也会增加噪声驱动和速度需求。GP 后验方差描述技能或解剖先验的不确定性；$q(C,X)$ 描述探查后未知反馈响应的可验收误差。这两者不能互相替代。

\subsection{增广状态域与输入能力}

取可信状态域 $S\le R_s$。若
\begin{equation}
 \sqrt2R_s\le r_{\mathrm{task}},\qquad
 b_{\rm drv}\le\lambda_sR_s,
 \label{eq:cal-servo-invariant}
\end{equation}
且域内观测、模型和接触均有效，则式~\eqref{eq:cal-continuous-tube}保证该状态域前向不变。定义 $r_\zeta=\omega R_s$，式~\eqref{eq:cal-convex-repair}的输入验收必须使用
\begin{equation}
 \boxed{a_i^\top u_{\mathrm{ff}}
       +r_\zeta\|a_i^\top C\|_2\le b_i,
       \quad i=1,\ldots,p.}
 \label{eq:cal-servo-capability}
\end{equation}
因为 $\|\zeta\|_2\le r_\zeta$，这保证所有该速度域中的命令都在已紧化命令内近似 $\mathcal U$ 中。固定 $r_\zeta$ 与探查数据时，能力条件仍为凸行范数约束；若将带宽、状态半径、实现与 $C$ 同时作为变量，则不能沿用原凸性结论。接触条件 $EC=0$ 与 $Eu_{\rm ff}=0$ 共同保证名义总输入的线性相容性，但不自动保证非线性声窗或患者端口无源。

若实际输入为 $u_{\rm act}=u_{\rm cmd}+d_u$，将 $Jd_u$ 加入 $w$ 的独立有效界。第~\ref{sec:cal-noisy}节的 $b_i$ 已按 $\mathcal D_u$ 紧化一次；使用原始物理界 $b_i^{\rm act}$ 时，等价地验收 $a_i^\top u_{\rm ff}+r_\zeta\|a_i^\top C\|+h_{\mathcal D_u}(a_i)\le b_i^{\rm act}$。实际接触等式还需 $Ed_u=0$ 或明示接触容差。即使 $\zeta$ 连续，突变的 $C$ 或 $u_{\mathrm{ff}}$ 仍可使物理命令跳变，因此实现切换要连续接续并重新验收。时变正交标架应对 $e,\zeta$ 使用同一个度量兼容导数，或将二者平行运输到共同标架；其反对称连接项不改变二次能量。非正交坐标和时变度量不能忽略连接与度量导数。

\subsection{采样证书与围绕未来名义计划的集合管}

连续时间定理不覆盖任意数字实现。以显式 Euler 为例，在冻结一步模型中令 $a=h\omega$，有
\begin{equation}
 \chi_{j+1}=T_{\Delta_j}\chi_j+\eta_j,\qquad
 T_\Delta=
 \begin{bmatrix}(1-a)I_m&a\Delta\\-aI_m&(1-a)I_m\end{bmatrix}.
 \label{eq:cal-discrete-general}
\end{equation}
其中 $\eta_j$ 包括 $h[w_j-\omega n_j;-\omega n_j]$ 与有界采样、保持和模型变化余项。对任意采用的实际一步矩阵 $T_A$，可先检查
\begin{equation}
 T_A^\top P T_A\preceq\rho^2P,
 \quad P\succ0,\quad\rho<1
 \label{eq:cal-discrete-matrix-certificate}
\end{equation}
是否在全部允许模型上成立；只检查一个名义矩阵不足以认证未知模型。对于式~\eqref{eq:cal-discrete-general}及 $0<a<1$，两行分别取范数给出无需枚举 $\Delta$ 的充分界：
\begin{equation}
 \rho_h^2=\lambda_{\max}(G_h),\qquad
 G_h=\begin{bmatrix}
 (1-a)^2+a^2&a(1-a)(1+q)\\
 a(1-a)(1+q)&(1-a)^2+a^2q^2
 \end{bmatrix}.
 \label{eq:cal-discrete-rho}
\end{equation}
确实，$\|\chi_{j+1}-\eta_j\|_2^2$ 不超过
$[\|r_j\|_2,\|v_j\|_2]G_h[\|r_j\|_2,\|v_j\|_2]^\top$。并且
\begin{align}
 (I_2-G_h)_{11}&=2a(1-a)>0,\\
 \det(I_2-G_h)&=a^2(1-a)(1-q)
                   [3+q-a(1-q)]>0.
\end{align}
所以 $q<1,0<h\omega<1$ 时 $\rho_h<1$。这是冻结模型下的二阶证书，不能继续使用一阶公式 $1-h\kappa(1-q)$。实际采样内变化与图像延迟还必须进入 $\eta_j$ 的有效界。

为区分恢复包络与预测不确定性，令 $\bar\chi_{j+1}=T_0\bar\chi_j$ 为当前接受的名义未来，预测集合为
$\mathcal T_j^\chi=\bar\chi_j\oplus\mathbb B_2(d_j)$。由于
$T_{\Delta_j}-T_0$ 仅含右上块 $a\Delta_j$，若 $\|\eta_j\|_2\le\beta_j$，则
\begin{equation}
 \boxed{d_{j+1}=\rho_{h,j}d_j
             +h\omega q_j\|\bar v_j\|_2+\beta_j.}
 \label{eq:cal-tube-recursion}
\end{equation}
若参考生成器或已知外生驱动也进入名义预测，只需同步从 $\eta_j$ 中扣除同一已知项；不得重复计入。记 $b_j=h\omega q_j\|\bar v_j\|_2+\beta_j$，展开递推得
\begin{equation}
 d_{t+k}\le
 \Bigl(\prod_{j=t}^{t+k-1}\rho_{h,j}\Bigr)d_t+
 \sum_{\ell=t}^{t+k-1}
 \Bigl(\prod_{j=\ell+1}^{t+k-1}\rho_{h,j}\Bigr)b_\ell.
 \label{eq:cal-future-tube}
\end{equation}
对任务误差 $e=[I_m,-I_m]\chi$，相应管半径至多为 $\sqrt2d_j$；对纠偏速度，半径至多为 $\omega d_j$。能力验收可以使用整个预测状态域的 $R_{s,j}=\|\bar\chi_j\|_2+d_j$，或使用更紧的 $r_{\zeta,j}=\omega(\|\bar v_j\|_2+d_j)$。这给出沿名义计划滚动展开的未来管，而非仅围绕零误差的一条标量界。

近处有当前观测支持，远处累积响应变化、人体运输和目标漂移，因此经常呈现近窄远宽的形状；收缩强或噪声小时也可变窄，不能要求管宽永远单调增加。新的图像与探查更新初始集合、$X,\widehat J,q_j$ 及剩余任务参考，应从当前时刻重新递推\emph{整条}未来管。若要宣称新管是旧预测管条件化后的有效收紧，必须使用有效集合相交及一致模型；单纯替换 GP 均值不具有这种集合保证。

材料与姿态误差管经当前有效的身体映射运输到 CoP、姿态和机器人候选。局部映射若具有经过验证的 Lipschitz 常数 $L_P$，则材料位置半径 $d$ 给出物理位置半径至多 $L_Pd$，另加注册误差。GP 样本可生成候选未来形状与统计置信带，但未经覆盖率校准及模型验证的 GP 方差不能作为上述认证物理管半径。

\subsection{沿途有效区间与失锁模式}

一处探查不能保证整段长血管。若在共同运输标架中独立验证了
\begin{equation}
 \|\Theta(\xi)-\Theta(\xi_0)\|_2
      \le Ld_g(\xi,\xi_0),
 \label{eq:cal-spatial-variation}
\end{equation}
则当前 $N(\xi),C(\xi)$ 的证书加入
\begin{equation}
 q_{\mathrm{variation}}=
    Ld_g(\xi,\xi_0)\|N(\xi)C(\xi)\|_2.
 \label{eq:cal-variation-q}
\end{equation}
由此沿未来路径增大 $q_j$ 并重算整条管。新探查更新锚点可以收紧该项；调和坐标自身并不产生此变化界。没有有效变化界时，仅在局部窗口内提供保证。

目标丢失、耦合不足或关联不可信时，关闭图像误差驱动，使用
\begin{equation}
 \dot\zeta=-2\omega\zeta,\qquad\dot s=\nu_{\rm stop}(t),
 \qquad \nu_{\rm stop}(t)=0\quad(t\ge t_{\rm stop}).
 \label{eq:cal-loss-mode}
\end{equation}
缺失检测不能填写成零误差；无观测时不使用定理~\ref{thm:cal-recovery}宣称任务恢复，图像特征出了视野甚至可能不再定义。
制动期间的分配与前馈采用事先验收、无需新增图像的计划，例如在其有效范围内冻结最后一次分配并继续已知身体保持运输；不使用缺失图像更新 $C$。停止距离和整个制动状态管属于该计划的验收内容。

\paragraph{失锁时的有限制动与暂停。}
观测失锁后采用 $\dot\zeta=-2\omega\zeta$。若 $\nu$ 已为零，则名义技能完全暂停；若 $\nu>0$ 且物理加速度有限，则执行事先验收的 $C^1$ 减速计划 $\nu_{\rm stop}(t)\to0$，并保留其停止距离 $\int\nu_{\rm stop}\,dt$ 及整个制动状态管。不能同时要求有限加速度和从正速度瞬时冻结进度。制动后进入 $\dot s=0$ 的暂停模式，名义技能的材料参考状态固定，实时身体运输及接触控制继续；真实相位仍由式~\eqref{v44:eq:actual-phase-port}与中断预算检查。

若单独的潜在目标模型在制动域内有效，且其误差满足
\begin{equation}
 \|\dot e_{\rm lat}\|\le M_C\|\zeta\|+M_s\nu_{\rm stop}+d_{\rm loss},
\end{equation}
则失锁 $\tau$ 时间后的潜在误差包络可取
\begin{equation}
 d_e(t+\tau)\le d_e(t)
 +\frac{M_C\|\zeta(t)\|}{2\omega}(1-e^{-2\omega\tau})
 +M_s\int_t^{t+\tau}\nu_{\rm stop}(r)\,dr+d_{\rm loss}\tau.
 \label{eq:rev:loss-braking-tube}
\end{equation}
此式由速度制动解与三角不等式积分得到。它只传播已声明的潜在目标集合，不构成缺失图像误差的反馈收敛结论。没有该模型或制动管超出其有效域时，停止认证未来任务管，并按可用接触与执行合同处理。

重获观测后首先验收目标关联及窗口，再运输状态、更新参考与重新修复。血管分叉需要保留互斥关联假设或报告歧义；单一高斯均值不能保证落在真实分支上。本文的响应校准定理不保证永不失锁，也不保证关联识别永远正确。





% V4.4 integration fragment.  Independent derivation; no edits to the main manuscript.
% Requires the packages, theorem environments and macros of V4.3_revised.
% Suggested use: replace the nominal-completion part of old Sec. 9 by the first
% sections, and put the finite-geometry construction after the original example.

\section{真实材料相位、反馈泄漏与联合推进验收}
\label{v44:sec:actual-phase}

名义生成器的进度仍记为 $s$，$\dot s=\nu\ge0$。
另以 $\sigma$ 记录实际成像或接触所对应的任务材料相位。
二者只有在另行证明的执行关系下才相同。
例如，一条已关联血管的切片交点相位可以同时依赖接触位置与切片姿态；
单独的 CoP 地址不必就是该交点的相位。

\subsection{从完整 CoP 运动学得到相位端口}
在一个已声明的有效逆分支和同一任务分支内，给出 $C^1$ 相位函数
\begin{equation}
 \sigma=\pi_a(\xi_C,R,c,f,t).
 \label{v44:eq:physical-phase-map}
\end{equation}
这个函数由任务材料对应及成像几何定义；若内部结构不确定，应给出其集合误差，
不能把名义路径参数直接当作观测到的相位。
设 $\ell_R$ 是相位的世界角速度微分，即
\begin{equation}
 D_R\pi_a[\widehat\omega^{W}R]=\ell_R\omega^{W}.
\end{equation}
对当前约简、白化的实际机器人输入取已知重建
$[v_C^{\rm rb};\omega^{W}]=T_a u_{\rm act}$。
已知保持运动也可以另列；下式把它统一包含在 $u_{\rm act}$ 中。
由 $\dot\xi_C=J_P^{-1}(v_C^{\rm rb}+R\dot r_C-b_p)$ 得到精确关系
\begin{align}
 \dot\sigma&=g_a u_{\rm act}+d_{\pi,a},
 \label{v44:eq:actual-phase-port}\\
 g_a&=\begin{bmatrix}D_\xi\pi_a J_P^{-1}&\ell_R\end{bmatrix}T_a,\\
 d_{\pi,a}
 &=D_\xi\pi_a J_P^{-1}(R\dot r_C-b_p)
   +D_c\pi_a\dot c+D_f\pi_a\dot f+\partial_t\pi_a.
 \label{v44:eq:actual-phase-drift}
\end{align}
体表模型使用对应的切向逆或已声明的接触投影。
当相位仅为 $\pi_a(\xi_C)$ 时，$\ell_R,D_c\pi_a,D_f\pi_a,\partial_t\pi_a$
相应消失；姿态影响不能在一般成像相位中被默认删除。
式~\eqref{v44:eq:actual-phase-port}同时显示：$\nu=0$ 不推出 $\dot\sigma=0$，
因为纠偏、压力中心迁移、身体补偿残差和执行误差仍可改变真实相位。

以下令 $u_{\rm act}=u_h+\nu b_s+C\zeta+d_u$。
在一个事先声明的物理状态域上，独立给出
\begin{equation}
 \|\zeta\|\le r_\zeta=\omega R_s,\qquad
 d_u\in\mathcal D_u,\qquad
 d_\pi\in[\underline d_\pi,\overline d_\pi].
 \label{v44:eq:phase-uncertainty-contract}
\end{equation}
其中 $\mathcal D_u$ 为非空紧凸集，因此下面使用的执行误差支持函数有限且能取得极值。
身体与 CoP 项中已知的部分可以保留在区间中心内；未补偿部分及相位几何误差进入区间宽度。
这些界是在整个声明域上对模型族取界，不是由希望得到的轨迹不变性反推。

\begin{lemma}[固定状态下的精确相位最坏值]
\label{v44:lem:phase-support}
对固定的 $g,u_h,b_s,C,\nu$，在式~\eqref{v44:eq:phase-uncertainty-contract}
给出的笛卡尔积合同上，真实相位速度的最小值为
\begin{equation}
 \boxed{\underline v_\sigma(C,\nu)
  =g(u_h+\nu b_s)-r_\zeta\|gC\|_2
   +\underline d_\pi-h_{\mathcal D_u}(-g^\top).}
 \label{v44:eq:exact-phase-lower}
\end{equation}
其中 $h_{\mathcal D_u}(z)=\sup_{d\in\mathcal D_u}z^\top d$。
因此 $\underline v_\sigma\ge v_*>0$ 当且仅当该乘积合同内的所有状态均满足
$\dot\sigma\ge v_*$。
\end{lemma}
\begin{proof}
将式~\eqref{v44:eq:actual-phase-port}代入控制分解。
在球 $\|\zeta\|\le r_\zeta$ 上，线性函数 $gC\zeta$ 的最小值为
$-r_\zeta\|gC\|$；若 $gC\ne0$，它由
$\zeta=-r_\zeta(gC)^\top/\|gC\|$ 取得。
执行误差项的最小值是 $-h_{\mathcal D_u}(-g^\top)$，相位漂移项的最小值是
$\underline d_\pi$。各不确定量在本合同中独立取值，故各最小值可以同时实现。
若真实系统中这些量相关，该公式仍是有效的外包充分证书，其必要性只针对声明的乘积集合。
\end{proof}

对变化的状态，式~\eqref{v44:eq:exact-phase-lower}须在整个候选状态管上成立。
例如已知 $g$ 属于有限多面体而其它依赖被一致包络时，可逐顶点检查相应约束。
对固定 $C$，$g\mapsto g(u_h+\nu b_s)-r_\zeta\|gC\|$ 是凹函数，
故其在一个多面体上的最小值可以在顶点取得。
一般集合也可以直接保留鲁棒约束，不以单个名义 $g$ 代替整个物理管。

\subsection{名义参考与真实材料任务模板的注册}
令实际材料任务在相位 $\sigma$、时刻 $t$ 要求的特征模板为
$y_{\rm tar}(\sigma,t)$。
控制器中的误差仍是 $e=y-y_r(s,t)$；真实任务误差另记为
\begin{equation}
 e_\tau=y-y_{\rm tar}(\sigma,t).
 \label{v44:eq:actual-template-error}
\end{equation}
二者使用同一个任务标架和固定物理度量。
正的 $\dot\sigma$ 与小的 $e$ 本身不足以证明 $e_\tau$ 小。

\begin{lemma}[真实任务模板的参考注册界]
\label{v44:lem:reference-registration}
若在整个声明窗口与物理状态域上有独立有效的参考注册界
\begin{equation}
 \|y_r(s,t)-y_{\rm tar}(\sigma,t)\|\le\varepsilon_{\rm reg},
 \label{v44:eq:reference-registration-budget}
\end{equation}
则 $S\le R_s$ 给出
\begin{equation}
 \|e_\tau\|\le\sqrt2R_s+\varepsilon_{\rm reg}.
 \label{v44:eq:actual-template-guarantee}
\end{equation}
一个可计算的充分构造是：在涉及的相位域中，$y_{\rm tar}$ 对相位一致 Lipschitz，常数为
$L_{\rm tar}$，并独立给出
\begin{equation}
 \|y_r(s,t)-y_{\rm tar}(s,t)\|\le\varepsilon_{\rm tpl},\qquad
 |s-\sigma|\le r_{\rm reg}.
\end{equation}
此时可取 $\varepsilon_{\rm reg}=\varepsilon_{\rm tpl}+L_{\rm tar}r_{\rm reg}$。
\end{lemma}
\begin{proof}
由 $e_\tau=e+[y_r(s,t)-y_{\rm tar}(\sigma,t)]$ 与 $\|e\|\le\sqrt2S$ 得到第一式。
在参考差中加减 $y_{\rm tar}(s,t)$，再用模板偏差与 Lipschitz 界即得所述构造。
\end{proof}

相位注册半径也不必由实际轨迹的正确性反推。
令 $r=s-\sigma$，其精确动力学为
\begin{equation}
 \dot r=\nu-g(u_h+\nu b_s+C\zeta+d_u)-d_\pi.
 \label{v44:eq:phase-registration-dynamics}
\end{equation}
与式~\eqref{v44:eq:exact-phase-lower}相应的相位速度上界为
\begin{equation}
 \overline v_\sigma
 =g(u_h+\nu b_s)+r_\zeta\|gC\|+\overline d_\pi
       +h_{\mathcal D_u}(g^\top).
\end{equation}
用整个声明域上的有效上下界包络，便有可直接积分的证书
\begin{equation}
 |r(t)|\le r_0+
 \int_{t_0}^{t}
 \max\{|\nu(\tau)-\underline v_\sigma(\tau)|,
        |\nu(\tau)-\overline v_\sigma(\tau)|\}\,d\tau
 =:r_{\rm reg}(t).
 \label{v44:eq:phase-registration-integral}
\end{equation}
这里 $r_0$ 是独立初值上界；例如有相位测量时可取
$r_0=|s(t_0)-\widehat\sigma(t_0)|+\epsilon_{\sigma,0}$。
相位同步 $\sigma=s$ 是使该半径为零的特例。
暂停与桥接期间仍用实际端口动力学运输 $r$，下一窗口重新检查注册预算；
不能只恢复小的图像误差而忽略已累积的相位错配。
式~\eqref{v44:eq:phase-registration-integral}与其它状态边界一起使用首次退出论证，
因此其上界是接受域上的计算结果，而不是预先假定真实相位跟随名义相位。

\subsection{真实速度与数据修复使用同一个分配变量}
固定候选、实际数据、物理度量、$\omega,R_s,\bar q<1$，并固定前馈方向 $b_s$。
令 $\lambda=\omega(1-\bar q)/2$，且在整个声明域上有独立有效的前馈误差界
\begin{equation}
 \|Jb_s-\partial_s y_r\|\le\alpha_s,
 \qquad
 \|Ju_h-\partial_t y_r+w_{\rm ext}+Jd_u\|\le d_0.
\end{equation}
若 $b_s$ 由数据构造，应先固定该方向，再计算其数据条件误差界；
不把 $\nu q(C,X)$ 这样的乘积悄悄当成仿射项。
可以求解
\begin{equation}
\begin{aligned}
 \max_{C,\nu,v_*}\quad &v_*\\
 \text{s.t.}\quad
 &\widehat JC=I_m,\qquad E_cC=0,\\
 &q(C,X)\le\bar q,\qquad 0\le\nu\le\nu_{\rm pref},\\
 &\nu\alpha_s+d_0+\sqrt2\omega n_{\max}\le\lambda R_s,\\
 &a_i^\top(u_h+\nu b_s)+r_\zeta\|a_i^\top C\|_2\le b_i,
       \quad i=1,\ldots,p,\\
 &v_*+r_\zeta\|gC\|_2
   \le g(u_h+\nu b_s)+\underline d_\pi-h_{\mathcal D_u}(-g^\top).
 \label{v44:eq:joint-real-progress-program}
\end{aligned}
\end{equation}
这里 $b_i$ 仍是按同一个执行误差集合紧化一次的命令界，
$E_cu_h=E_cb_s=0$ 已预先验收；实际等式另要求 $E_cd_u=0$ 或明确接触容差。
整个管上有多个相位行或多个模型顶点时，添加对应的同类约束。
固定上述数据和方向后，该问题是凸的：每个非线性左侧都是线性映射的范数之和，
等式为仿射约束，目标为线性函数。
加速度和完整物理状态管仍须使用接受路径上的导数或实际一步执行模型验收；
此处不宣称这些外层路径变量也构成凸问题。
真实任务质量另外验收
$\sqrt2R_s+\varepsilon_{\rm reg}\le\epsilon_{\rm task}$，
其中注册预算由引理~\ref{v44:lem:reference-registration}及实际相位运输得到。
固定参考、状态半径或误差反馈合同不能自动替代这个注册验收。

\begin{proposition}[恢复、能力与真实正推进的共同证书]
\label{v44:prop:joint-certificate}
若式~\eqref{v44:eq:joint-real-progress-program}在一个声明窗口的整个物理状态域上成立，
且 $v_*>0$、$S(t_0)\le R_s$，并有独立有效的注册预算
$\sqrt2R_s+\varepsilon_{\rm reg}\le\epsilon_{\rm task}$，则在该域内
\begin{equation}
 \begin{aligned}
 &S(t)\le R_s,\qquad u_{\rm act}(t)\in\mathcal U_{\rm act},
 \qquad \sigma(t)-\sigma(t_0)\ge v_*(t-t_0),\\
 &\|y(t)-y_{\rm tar}(\sigma(t),t)\|
      \le\sqrt2R_s+\varepsilon_{\rm reg}\le\epsilon_{\rm task}.
 \end{aligned}
 \label{v44:eq:joint-guarantee}
\end{equation}
若域的其它边界也通过独立的前向管证书，则这些结论在整个窗口有效。
\end{proposition}
\begin{proof}
数据修复与恢复不等式给出
$D^+S\le-\lambda S+\nu\alpha_s+d_0+\sqrt2\omega n_{\max}$，
因此边界 $S=R_s$ 上的上导数非正。
行范数约束和一次执行误差紧化给出实际输入能力。
引理~\ref{v44:lem:phase-support}给出 $\dot\sigma\ge v_*$，积分得到相位推进不等式。
引理~\ref{v44:lem:reference-registration}再把相对于名义参考的误差界转成真实材料模板误差界。
把这些边界与独立认证的其余物理边界一起应用首次退出论证，即得到整段结论。
\end{proof}

\paragraph{什么构成不循环的完整物理管证书。}
一种可检查的方式是在候选开始前声明闭域
$\mathcal D(t)=\{z:B_j(z,t)\le0,\ j=1,\ldots,k\}$，
在该域及所需边界邻域上给出完整状态动力学 $\dot z=F_a(z,t,\upsilon)$ 的不确定集合。
令时空图为 $\operatorname{graph}\mathcal D=\{(t,z):z\in\mathcal D(t)\}$，
完整的 Nagumo 切锥条件是：在窗口终点以前，对每个允许的 $\upsilon$ 检查
\begin{equation}
 (1,F_a(z,t,\upsilon))\in
 T_{\operatorname{graph}\mathcal D}(t,z),\qquad t<T.
 \label{v44:eq:full-state-graph-tangency}
\end{equation}
只有当时空约束满足 MFCQ 或其它保证该交集的切锥等于线性化锥的正规条件时，
才能用各活动边界上的不等式实现这一检查：
\begin{equation}
 \sup_{\upsilon}\bigl[\partial_t B_j+D_zB_jF_a(z,t,\upsilon)\bigr]\le0.
 \label{v44:eq:full-state-boundary-certificate}
\end{equation}
单个边界函数分别光滑且梯度非零，并不足以保证交集满足这个替换条件。
在相应闭集、正则性和动力学局部唯一性条件下，
也可以使用能直接验证包含关系的
区间积分、微分包含可达集或离散一步管。
该步骤检查的是事先给定的域，不是先假定真实轨迹不会离开。
后面的有限几何例用解析积分与严格余量完成这个步骤。

\section{相位中性纠偏的校准代价}
\label{v44:sec:neutral-calibration}

若任务要求纠偏不改变实际相位，施加 $gC=0$。
令 $E_0$ 为既有接触等式，$G$ 为一个或多个相位约束行，取固定物理白化下的正交基
\begin{equation}
 Z_0:\ \operatorname{range}Z_0=\ker E_0,\qquad
 Z:\ \operatorname{range}Z=\ker E_0\cap\ker G.
 \label{v44:eq:phase-reduction}
\end{equation}
纠偏输入写为 $u_{\rm corr}=Z\bar C\zeta$，因此局部响应模型变为
\begin{equation}
 J_Z=HZ+\Theta NZ+E_{\rm geo}Z.
 \label{v44:eq:reduced-phase-model}
\end{equation}
前馈不在 $\ker G$ 内；它另外满足已规定的正相位作用，例如 $gb_s=1$。
必须分别检查 $HZ$ 满行秩以及存在合法推进方向，二者不能互相替代。

\begin{theorem}[新增相位中性约束的结构代价与精确复杂度]
\label{v44:thm:phase-rank-cost}
在无噪声、无结构外余项、固定完整响应范数球的局部模型内，
允许约简输入空间中的独立可观测探查方向，并假定 $HZ_0$ 与 $HZ$ 均满行秩。
在相同任务维数 $m$ 下，相位中性实现的最小不确定暴露秩为
\begin{equation}
 k_{\rm phase}=m-\dim H(\ker N\cap\ker E_0\cap\ker G),
 \label{v44:eq:phase-rank}
\end{equation}
并使用约简对 $(HZ,NZ)$ 的 $D_{\star,Z},D_{Q,Z}$ 计算相应首次误差与鲁棒满射阈值。
令 $k_0=m-\dim H(\ker N\cap\ker E_0)$，则
\begin{equation}
 0\le k_{\rm phase}-k_0\le\operatorname{rank}G.
 \label{v44:eq:phase-rank-increase}
\end{equation}
尤其，一个标量相位中性要求可以严格增加一个结构暴露方向；
也可能使 $HZ$ 失去满行秩，此时全部任务方向的中性纠偏不可行，不能继续使用该公式宣称可实现。
若约简完整球还满足统一满射条件 $\delta\|D_{\star,Z}\|<1$，
上述 $k_{\rm phase}$ 才等于相位中性全历史精确任务重新实现的最低独立校准方向数。
若该球包含相位中性响应的秩亏成员，则全历史精确任务复杂度为无穷；
有限的结构秩增量不能解释为增加有限探查就必然可恢复。
\end{theorem}
\begin{proof}
$Z$ 是允许纠偏空间的等距坐标，因此原局部信息定理可直接用于 $(HZ,NZ)$。
由 $Z\ker(NZ)=\ker N\cap\ker E_0\cap\ker G$ 得到式~\eqref{v44:eq:phase-rank}。
置 $V=\ker N\cap\ker E_0$、$V_1=V\cap\ker G$。
因为 $V_1\subset V$，有 $HV_1\subset HV$，给出结构秩差值非负。
再取 $V=V_1\oplus V_2$，其中
$\dim V_2=\operatorname{rank}(G|_V)\le\operatorname{rank}G$。
于是 $HV=HV_1+HV_2$，故
$\dim HV-\dim HV_1\le\dim V_2$，得到上界。
结构秩等式与增量不等式不要求不确定球内所有成员可控。
球内统一满射成立时，再应用该约简模型的精确任务方向定理获得复杂度等式。
若球内有一个秩亏成员，无论提供多少数据也不能为它实现全部任务方向，因此全历史复杂度为无穷。
满射条件属于同一个约简模型，不能沿用未约简模型的阈值。
\end{proof}

\paragraph{最小的严格分离例。}
取 $m=1$、$H=[1\ \ 1]$、$N=[1\ \ 0]$、$g=[0\ \ 1]$、$E_0=0$，
即 $\dot y=(1+\theta)u_1+u_2$、$\dot\sigma=u_2$，$|\theta|\le\delta<1$。
不限制纠偏相位时，$B=(0,1)^\top$ 满足 $HB=1,NB=0$，所以 $k_0=0$。
要求 $gC=0$ 后，只有 $u_1$ 可用于任务纠偏，故 $k_{\rm phase}=1$。
新增方向不是坐标变化的产物，而是禁止用实际任务进度来替代未知纠偏作用的代价。
若把该例改成 $\delta=1$，未约简响应因第二列恒为 $1$ 仍对整个球满行秩；
相位中性响应却在 $\theta=-1$ 时变为零。
此时结构秩仍是 $0\to1$，全历史精确复杂度则是 $0\to\infty$。
实际数据能量与这一代价的精确关系，应与同一分配变量及同一真实相位合同一起计算。

\section{采样、失锁和反向桥接下的材料覆盖记账}
\label{v44:sec:hybrid-coverage}

完成判据应同时包含真实相位与可用成像记录。
仅在失锁期间前进到段末端，并不说明中间材料已经得到有效扫描。
以下在一个固定任务分支的共同相位坐标中表述；被动换图保留同一个 $\sigma$，
主动换参考不得通过重置相位标签制造进度。

把执行时间分为有效扫描区间 $\mathsf G$ 和其它区间 $\mathsf B$。
后者包括探查、等待、失锁制动、重新关联及物理桥接。
在有效扫描区间使用命题~\ref{v44:prop:joint-certificate}，给出统一
$\dot\sigma\ge v_g>0$。
对其它区间给出独立可验收的总时间与总反向变化预算
\begin{equation}
 |\mathsf B|\le T_B,\qquad
 \int_{\mathsf B}(-\dot\sigma)_+\,dt\le B_\sigma<\infty.
 \label{v44:eq:bad-mode-budgets}
\end{equation}
正向失锁位移不直接计为已扫描覆盖。
每次恢复后，第一段有效成像须与此前连续已覆盖前缀重叠；
若实际相位越过了尚未有效扫描的区间，先执行已验收的返回或重扫桥接。
这种桥接的时间和反向相位长度均计入式~\eqref{v44:eq:bad-mode-budgets}。
如果只能验证各次完整中断的净向后相位损失，可以用其和代替更保守的总负变差；
不能只把桥接时间加到名义完成时间而遗漏额外需要重扫的材料长度。

\subsection{离散图像的覆盖语义}
若图像仅在时刻 $t_j$ 采集，给每帧一个与任务质量对应的实际材料覆盖半宽 $a_j>0$。
若相位估计满足 $|\widehat\sigma_j-\sigma(t_j)|\le\epsilon_{\sigma,j}$，则一个保证被该帧覆盖的区间是
\begin{equation}
 I_j=[\widehat\sigma_j-a_j+\epsilon_{\sigma,j},\,
       \widehat\sigma_j+a_j-\epsilon_{\sigma,j}].
 \label{v44:eq:certified-frame-interval}
\end{equation}
实际保存这些区间的并及其从任务起点开始的最大连续前缀。
连续接受的扫描帧须通过 $I_j\cap I_{j+1}\ne\varnothing$ 的检查；
中断后首帧须与已保存前缀相交。
没有正的覆盖宽度或其它有效插值合同，有限个采样点本身不构成连续材料覆盖证明。
也可直接用实际速度上界、帧间隔及真实声束厚度认证相邻实际覆盖重叠。

\begin{theorem}[真实有限段覆盖的混合执行充分条件]
\label{v44:thm:hybrid-finish}
任务为覆盖相位区间 $[0,L]$。
初始有效覆盖包含 $0$，初始真实相位为 $\sigma_0\le0$；
混合执行良定，切换没有有限时间聚集。
在所有完成的有效扫描区间上，独立域、响应与参考注册证书共同保证
真实材料模板误差 $\|y-y_{\rm tar}(\sigma,t)\|\le\epsilon_{\rm task}$、实际能力及
$\dot\sigma\ge v_g>0$，而不仅是相对于名义参考的恢复。
假定式~\eqref{v44:eq:bad-mode-budgets}成立，并且中断后的覆盖重叠规则被实际执行。
对连续有效成像，任务覆盖完成时间满足
\begin{equation}
 T_{\rm cov}-t_0\le T_B+\frac{L-\sigma_0+B_\sigma}{v_g}.
 \label{v44:eq:hybrid-continuous-time}
\end{equation}
对离散采样，把只有起点而尚无合格终点图像的未完成片段计入 $\mathsf B$；
其余完成区间的最大时长不超过 $h_{\max}$。
若式~\eqref{v44:eq:certified-frame-interval}相邻重叠，且
$a_j\ge2\epsilon_{\sigma,j}$，则有
\begin{equation}
 T_{\rm cov}-t_0\le T_B+\frac{L-\sigma_0+B_\sigma}{v_g}+h_{\max}.
 \label{v44:eq:hybrid-sampled-time}
\end{equation}
物理完成还要求末端制动时，另加已验收的 $T_{\rm end}$；
制动空间及其相位停止距离须留在扩展的有效物理域中。
\end{theorem}
\begin{proof}
在任意时刻对真实相位积分，得到
\begin{equation}
 \sigma(t)\ge\sigma_0+v_g|\mathsf G\cap[t_0,t]|-B_\sigma.
 \label{v44:eq:actual-phase-accounting}
\end{equation}
连续成像时，每个有效扫描区间上的 $\sigma$ 连续且严格递增，故成像覆盖该区间两端间的所有相位。
重启重叠规则归纳保证已覆盖材料始终包含从 $0$ 开始的连续前缀。
当有效扫描时间达到 $(L-\sigma_0+B_\sigma)/v_g$，
当前有效相位已不小于 $L$，因而覆盖完成。
全部其它时间至多 $T_B$，得到式~\eqref{v44:eq:hybrid-continuous-time}。

离散采样时，在累计有效扫描时间第一次越过该阈值的合格终点帧检查覆盖。
等待该完成区间的额外时长不超过 $h_{\max}$。
该帧的保证覆盖右端满足
$\widehat\sigma_j+a_j-\epsilon_{\sigma,j}
\ge\sigma(t_j)+a_j-2\epsilon_{\sigma,j}\ge\sigma(t_j)$。
因此当前前缀的右端至少为 $L$；重叠检查保证中间没有遗漏区间。
最后加上其它模式时间及可选末端制动时间。
\end{proof}

上述定理不假定失锁必然在有限时间消失。
如果没有累计中断时间或反向相位损失的有限预算，就不声称有限完成时间。
有限候选策略还须实际执行其已认证的进度承诺，不能反复重规划并把推进推迟到未来。

\subsection{暂停与制动的真实相位损失}
失锁时设 $\dot\zeta=-2\omega\zeta$，并执行已认证的 $\nu_{\rm stop}$ 与保持计划。
若整个制动域内 $\|gC\|\le M_{gC}$，且扣除可精确计算的名义项后相位残差幅度至多 $d_\sigma$，
则一个时长为 $\tau$ 的中断的反向变化满足
\begin{align}
 B_{\rm loss}(\tau)
 \le{}&\frac{M_{gC}\|\zeta(t_0)\|}{2\omega}
                   (1-e^{-2\omega\tau})+d_\sigma\tau\\
 &+\int_{t_0}^{t_0+\tau}
       [-g(u_h+\nu_{\rm stop}b_s)-d_{\pi,\rm known}]_+\,dt.
 \label{v44:eq:loss-phase-budget}
\end{align}
这是对 $(-\dot\sigma)_+$ 逐项取上界再积分的结果。
即使 $gC=0$ 且 $\nu=0$，仍保留 $d_\sigma\tau$；
无限等待不能被计作零相位代价。
若没有控制反馈恢复保证，制动域的物理有效性使用独立潜在目标与接触模型验收。
返回桥接另按其实际反向路径长度积分，不用两个可行端点代替中间路径验证。



\section{有限探查、真实相位中性与有限转角扫描的定量例}
\label{v44:sec:finite-geometry-example}

本例使新增相位约束确实改变校准需求，并从有限实际输入产生数据，
随后完成一个非零噪声、非零模型余项下的真实材料段。
采用 $L_0=10\,\mathrm{mm}$ 的固定长度尺度，
$q_p=(d,\alpha,h)$ 中 $d,h$ 已归一化，$\alpha$ 为弧度。
局部平面接触、TCP 与固定 CoP 重合，实际速度输入为
$u=(\dot d,\dot\alpha,\dot h)^\top$。
本例使用实际速度级执行合同，$d_u=0$；一般执行与身体误差已经由
式~\eqref{v44:eq:actual-phase-drift}--\eqref{v44:eq:exact-phase-lower}处理。

\subsection{一个具有非零未知接触响应的精确有限几何族}
取
\begin{equation}
 y_0=(1,3)^\top,\qquad \beta=(0,0.5)^\top,
 \qquad \|\theta\|\le\delta=0.04,
 \label{v44:eq:geometry-parameters}
\end{equation}
并定义
\begin{equation}
 \boxed{y=R_2(-\alpha)
 \{y_0+\beta h-e_1d+\theta d+\mathsf Dq_p+g(t)\}.}
 \label{v44:eq:finite-exact-geometry}
\end{equation}
其中
\begin{equation}
 \|\mathsf D\|\le\eta_c=0.003,\qquad g(0)=0,
 \qquad\|\dot g\|\le w_0=0.001\,\mathrm{s}^{-1}.
 \label{v44:eq:geometry-remainder}
\end{equation}
$\theta d$ 是随面内接触平移产生的未知目标滑移，
不是材料相位变化；$\mathsf D$ 保留小的其它接触响应。
例如 $\theta=(0.02,0.01)^\top$ 与 $\mathsf D=0.003[\,I_2\ \ 0\,]$
给出该模型族的非零成员。
该模型保持交点的材料面外坐标 $h$，故真实相位可以定义为 $\sigma=h/\ell_h$，
其中扫描长度 $\ell_h=0.05$。

令 $\widehat y=y+n$，$\|n\|\le n_0=0.0005$。
直接微分式~\eqref{v44:eq:finite-exact-geometry}得到
\begin{align}
 J&=H+R_2(-\alpha)\theta N+E,
 \qquad N=[1\ \ 0\ \ 0],\\
 H&=\begin{bmatrix}A(\widehat y,\alpha)&R_2(-\alpha)\beta\end{bmatrix},
 \qquad
 A(\widehat y,\alpha)=
 \begin{bmatrix}-\cos\alpha&\widehat z\\
                 \sin\alpha&-\widehat x\end{bmatrix},\\
 E&=R_2(-\alpha)\mathsf D+
 \begin{bmatrix}0&-n_z&0\\0&n_x&0\end{bmatrix},
 \qquad \|E\|\le\eta_c+n_0=0.0035=:\eta,
 \label{v44:eq:finite-model-factorization}\\
 \|w_{\rm ext}\|&\le w_0.
\end{align}
转角列使用实际形变后图像坐标，因此没有重复漏掉旋转导数。
噪声在分配矩阵与误差反馈中的出现可以相关，后面不使用独立性假设。

\subsection{相位中性要求使方向数严格从零变为一}
在 $\alpha=0$、$y=y_0$、$E=0$ 的结构化模型内，
\begin{equation}
 H_0=\begin{bmatrix}-1&3&0\\0&-1&0.5\end{bmatrix},
 \qquad N=[1\ \ 0\ \ 0].
\end{equation}
若允许用相位方向参与纠偏，取
\begin{equation}
 B_0=\begin{bmatrix}0&0\\1/3&0\\2/3&2\end{bmatrix}.
\end{equation}
则 $H_0B_0=I_2$、$NB_0=0$，故无需校准。
但是 $gB_0\ne0$，其中真实相位行 $g=[0\ \ 0\ \ 1]/\ell_h$。
施加 $gC=0$ 后，只能用前两个输入纠偏，
\begin{equation}
 H_Z=\begin{bmatrix}-1&3\\0&-1\end{bmatrix},\qquad
 N_Z=[1\ \ 0],\qquad
 D_{\star,Z}=N_ZH_Z^{-1}=[-1\ \ -3].
 \label{v44:eq:geometry-rank-cost}
\end{equation}
所以 $k_0=0$、$k_{\rm phase}=1$，而
\begin{equation}
 \delta\|D_{\star,Z}\|=0.04\sqrt{10}
 =0.1264911\ldots<1.
\end{equation}
例如精确结构化模型中的零数据首次误差目标 $0.1$ 不能由相位中性共同动作达到，
一个独立方向是必要的；该结论不是由模型中存在不可控成员造成的。

\subsection{有限能量的实际探查与误差预算}
在正式扫描前保持 $\alpha=h=0$，执行给定往返位移
\begin{equation}
 d(t)=a\sin^2(\pi t/T_p),\qquad
 a=0.2,\quad T_p=1.2\,\mathrm{s},\qquad 0\le t\le T_p.
 \label{v44:eq:actual-probe-path}
\end{equation}
该动作从静止开始、在 $t=T_p/2$ 达到 $a$、并返回 $d=0$ 与零速度。
实际相位始终为零。
其峰值速度、加速度和总归一化输入能量分别为
\begin{align}
 \max|\dot d|&=a\pi/T_p=0.523599\ldots\,\mathrm{s}^{-1},\\
 \max|\ddot d|&=2a\pi^2/T_p^2=2.741557\ldots\,\mathrm{s}^{-2},\\
 E_p:=\int_0^{T_p}\|u(t)\|^2dt
 &=\frac{a^2\pi^2}{2T_p}=0.164494\ldots\,\mathrm{s}^{-1}.
 \label{v44:eq:actual-probe-energy}
\end{align}
即峰值平移速度小于 $5.236\,\mathrm{mm/s}$。
只用前半程的实际位移与观测构造数据；返回段用来恢复起点，不把往返净位移零误当成已获得的信息。
由于 $\alpha=h=0$，有限几何在这段上精确给出
\begin{equation}
 R_p=\widehat y(T_p/2)-\widehat y(0)+e_1a
       =\theta a+W_p,\qquad X=a.
\end{equation}
因此
\begin{align}
 \|W_p\|&\le\eta_ca+w_0T_p/2+2n_0=0.0022,\\
 \widehat\theta&=R_p/a,
 \qquad\|\widehat\theta-\theta\|\le\varepsilon_\theta=0.011,
 \qquad \|\widehat\theta\|\le0.051.
 \label{v44:eq:actual-probe-estimate}
\end{align}
这里采用实际输入，没有无穷小探查或无噪声辨识。
在正式扫描时，将数据响应按已知 $R_2(-\alpha)$ 运输，
即 $\widehat\Theta_\alpha=R_2(-\alpha)\widehat\theta$；
这是同一几何模型中的明确运输关系。

探查中 $0\le d\le0.2$，故 $x$ 保持在
$[0.7902,1.0012]$ 内，而 $z\in[2.9902,3.0098]$。
因此整个探查位于所声明的声窗 $|x_{\rm phys}|\le12\,\mathrm{mm}$、
$25\,\mathrm{mm}\le z_{\rm phys}\le35\,\mathrm{mm}$ 内。
准备探查使用这个独立的宽观测域；正式扫描的 $0.1\,\mathrm{mm}$ 误差合同从返回后开始。
返回时 $d=\alpha=h=0$，真实初始图像误差为 $g(T_p)$，范数至多 $0.0012$。

\subsection{实际数据产生的相位中性修复}
令
\begin{equation}
 \widehat A=A(\widehat y,\alpha)
        +R_2(-\alpha)\widehat\theta[1\ \ 0],\qquad
 C=\begin{bmatrix}\widehat A^{-1}\\0\quad0\end{bmatrix}.
 \label{v44:eq:repaired-neutral-allocation}
\end{equation}
于是 $\widehat JC=I_2$ 且 $gC=0$。
取正式扫描域
\begin{equation}
 \|y-y_0\|\le0.01,\quad |\alpha|\le0.09,\quad |d|\le0.3,
 \quad 0\le h\le0.05.
 \label{v44:eq:new-scan-domain}
\end{equation}
为简化统一界，可以把实际测量箱扩大为
$|\widehat x-1|,|\widehat z-3|\le0.022$。
在此箱内
\begin{equation}
 \widehat x\cos\alpha-\widehat z\sin\alpha\ge0.7020591,
\end{equation}
且未修复矩阵满足
\begin{equation}
 \|A^{-1}\|\le4.77,\qquad
 \|(A^{-1})_{1,:}\|\le4.56,\qquad
 \|(A^{-1})_{2,:}\|\le1.43.
\end{equation}
Neumann 逆界与 Sherman--Morrison 公式分别给出
\begin{align}
 \|C\|&\le\frac{4.77}{1-4.77(0.051)}<6.31=:M_C,\\
 \|C_{1,:}\|&\le\frac{4.56}{1-4.56(0.051)}<5.943=:M_1,\\
 \|C_{2,:}\|&\le1.43\left(1+0.051\frac{4.56}{1-4.56(0.051)}\right)
                      <1.864=:M_2.
 \label{v44:eq:new-inverse-bounds}
\end{align}
所有分母有严格正余量，因此修复在整个声明域有定义。
实际数据证书为
\begin{equation}
 q\le\varepsilon_\theta M_1+\eta M_C
   =0.087458<\bar q:=0.1.
 \label{v44:eq:new-q-certificate}
\end{equation}
它正是满秩单通道数据中的
$\varepsilon_{\rm data}\|X^\dagger NC\|+\eta\|C\|$，
而不是把参数估计误差作为一个无来源的小常数。

\begin{remark}[固定共享余项给出的更紧响应界]
\label{v44:rem:shared-remainder-tightening}
本例的 $\mathsf D$ 在探查和执行中固定，故其首列被同一实际探查数据吸收。
具体地，由有限探查等式可写成
\begin{equation}
 \widehat\theta=\theta+\mathsf D_{:,1}+\epsilon_p,
 \qquad
 \|\epsilon_p\|\le\frac{w_0T_p/2+2n_0}{a}=0.008.
\end{equation}
令 $m_n=(-n_z,n_x)^\top$，则在 $C_{3,:}=0$ 时
\begin{equation}
 (J-\widehat J)C
 =-R_2(-\alpha)\epsilon_p C_{1,:}
   +\bigl(R_2(-\alpha)\mathsf D_{:,2}+m_n\bigr)C_{2,:}.
\end{equation}
因此共享余项的一致性可进一步给出
\begin{equation}
 q\le0.008M_1+0.0035M_2=0.054068.
 \label{v44:eq:shared-remainder-tighter-q}
\end{equation}
这展示了在同一个数据一致模型中处理固定余项的收益：
已被探查识别的首列不再作为独立执行误差重复放大。
后续数值与完整证明仍沿用式~\eqref{v44:eq:new-q-certificate}的保守值 $0.087458$，
因此无需改变既有速度、状态域和桥接参数。
\end{remark}

\subsection{正真实速度区间及完整首次退出证明}
取 $y_r=y_0$、$\omega=4\,\mathrm{s}^{-1}$，并设
\begin{equation}
 R_s=\frac{0.01}{\sqrt2},\qquad
 r_\zeta=\omega R_s=0.028284271\ldots,\qquad
 \lambda=\frac{\omega(1-\bar q)}2=1.8\,\mathrm{s}^{-1}.
\end{equation}
本例的真实材料任务模板同样为 $y_{\rm tar}(\sigma,t)\equiv y_0$，
因此参考注册误差 $\varepsilon_{\rm reg}=0$。
下面的图像误差合同也就是实际扫描相位处的任务模板合同。
令
\begin{equation}
 b_s=\ell_h\begin{bmatrix}
          -\widehat A^{-1}R_2(-\alpha)\beta\\1
        \end{bmatrix},\qquad
 u=\nu b_s+C\zeta.
 \label{v44:eq:new-calibrated-feedforward}
\end{equation}
因此 $\widehat Jb_s=0$，且不依赖未知响应与图像噪声地有
\begin{equation}
 \dot h=\ell_h\nu,\qquad \dot\sigma=\nu.
 \label{v44:eq:new-exact-actual-progress}
\end{equation}
未知响应仍作用在前馈的面内分量上，但现在其作用由同一实际数据控制。
具体地，令
\begin{align}
 M_b&=\ell_h\sqrt{1+M_C^2\|\beta\|^2}
          =0.165484327\ldots,\\
 \alpha_s&=\varepsilon_\theta\ell_hM_1\|\beta\|+\eta M_b
          =0.0022135202\ldots .
 \label{v44:eq:new-feedforward-data-bound}
\end{align}
则 $\|Jb_s\|\le\alpha_s$，且恢复驱动满足
\begin{equation}
 b_{\rm drv}(\nu)\le\nu\alpha_s+w_0+\sqrt2\omega n_0.
\end{equation}
输入幅度合同取
\begin{equation}
 |\dot d_{\rm phys}|\le6\,\mathrm{mm/s},\qquad
 |\dot\alpha|\le10^\circ/\mathrm{s},\qquad
 |\dot h_{\rm phys}|\le1.5\,\mathrm{mm/s}.
\end{equation}
对常数
\begin{equation}
 \boxed{\nu\in[1.4,2.4]\,\mathrm{s}^{-1}}
 \label{v44:eq:new-speed-interval}
\end{equation}
有以下向外取整的统一界：
\begin{align}
 b_{\rm drv}(\nu)&<0.009141<\lambda R_s=0.01272792\ldots,\\
 |\dot d_{\rm phys}|&\le L_0M_1(\nu\ell_h\|\beta\|+r_\zeta)<5.247\,\mathrm{mm/s},\\
 |\dot\alpha|&\le M_2(\nu\ell_h\|\beta\|+r_\zeta)<9.429^\circ/\mathrm{s},\\
 \dot h_{\rm phys}&\le1.2\,\mathrm{mm/s}.
 \label{v44:eq:new-speed-margins}
\end{align}
在 $T=1/\nu$ 内，积分同一分量界还给出
\begin{align}
 |d(t)|&\le M_1\left(\ell_h\|\beta\|+\frac{r_\zeta}{\nu}\right)
                       <0.268642<0.3,\\
 |\alpha(t)|&\le M_2\left(\ell_h\|\beta\|+\frac{r_\zeta}{\nu}\right)
                       <0.084259<0.09.
 \label{v44:eq:new-finite-pose-margins}
\end{align}

\begin{proposition}[有限能量探查后的真实短段扫描]
\label{v44:prop:nonzero-calibration-scan}
在式~\eqref{v44:eq:finite-exact-geometry}的整个模型族和上述速度级合同内，
先执行式~\eqref{v44:eq:actual-probe-path}，使用其实际数据构造
式~\eqref{v44:eq:repaired-neutral-allocation}，并令扫描开始时 $\zeta=0$。
对每个式~\eqref{v44:eq:new-speed-interval}中的常数 $\nu$，
二阶任务反馈与式~\eqref{v44:eq:new-calibrated-feedforward}在 $T=1/\nu$ 内
完成 $0.5\,\mathrm{mm}$ 的真实材料段，保持图像误差不超过 $0.1\,\mathrm{mm}$，
并满足完整位姿范围及实际输入幅度。
其中未知响应、非零余项与图像噪声均保留，实际探查总能量为
式~\eqref{v44:eq:actual-probe-energy}的有限正数。
\end{proposition}
\begin{proof}
探查几何、输入能量、观测域及数据误差已由显式轨迹直接验收。
正式扫描初始时 $S(0)=\|g(T_p)\|\le0.0012<R_s$。
在域~\eqref{v44:eq:new-scan-domain}内，逆界与实际数据证书有效，故
\begin{equation}
 S(t)\le e^{-\lambda t}(0.0012)
       +\frac{\nu\alpha_s+w_0+\sqrt2\omega n_0}{\lambda}
            (1-e^{-\lambda t})<R_s.
 \label{v44:eq:new-finite-state-envelope}
\end{equation}
所以 $\|y-y_0\|\le\sqrt2S<0.01$。
假定 $T$ 之前出现首次退出；在该时刻之前，
式~\eqref{v44:eq:new-finite-pose-margins}排除 $d$ 与 $\alpha$ 边界退出，
式~\eqref{v44:eq:new-exact-actual-progress}给出
$h(t)=\ell_h\nu t\in[0,\ell_h)$，恢复界排除图像误差边界退出。
因此不能存在这样的首次退出，解可延续至 $T$，且 $h(T)=\ell_h$。
输入幅度在整个域内由式~\eqref{v44:eq:new-speed-margins}保证。
这个证明同时构造数据、修复、物理域与真实推进，没有把整条轨迹留在域内作为未展开前提。
\end{proof}

以 $\nu=2\,\mathrm{s}^{-1}$ 为例，正式扫描耗时 $0.5\,\mathrm{s}$，
实际推进速度为 $1\,\mathrm{mm/s}$，全段图像误差上界小于
$0.045391\,\mathrm{mm}$；峰值面内速度小于 $4.653\,\mathrm{mm/s}$，
峰值摇摆角速度小于 $8.361^\circ/\mathrm{s}$。
探查与扫描合计 $1.7\,\mathrm{s}$，启动、末端及可能中断的物理接续另按对应合同记账。

\subsection{一个实际采样实现，而非把连续观测当作数字实现}
在正式扫描中，每隔 $h_m=0.005\,\mathrm{s}$ 读取一次图像，
单帧读数误差满足 $\|n_{\rm meas}\|\le0.00025$，帧间保持该读数。
实时姿态 $\alpha$ 仍由执行测量给出。
内部反馈状态在帧间按
$\dot\zeta=-\omega^2(\widehat y_j-y_0)-2\omega\zeta$
的精确解更新；因此它不等于未经验收的显式 Euler。
分配矩阵使用保持的图像读数与当前姿态。

在所声明域与有效噪声包络内，由已证明的闭环方程直接得到
\begin{equation}
 \|\dot y\|\le(1+\bar q)r_\zeta+2.4\alpha_s+w_0
   <0.037426=:M_y.
\end{equation}
因此帧间等效误差满足
\begin{equation}
 \|\widehat y_j-y(t)\|
 \le0.00025+M_yh_m<0.000438<n_0.
 \label{v44:eq:sampling-noise-closure}
\end{equation}
这不是从希望成立的噪声界反推自身：在物理域或噪声域首次退出前，
上述速度界成立，而式~\eqref{v44:eq:sampling-noise-closure}严格排除噪声域先退出；
命题~\ref{v44:prop:nonzero-calibration-scan}的其它边界论证同时排除物理域先退出。
于是这个真实采样及保持实现也满足相同的连续恢复包络。
它仍是速度级执行合同；分配在读图时刻的跳变若受到实际加速度限制，
须改用已验收的平滑重建或实际一步执行模型。

若本解析成像模型进一步给每帧一个已知材料轴向厚度，
例如覆盖半宽 $a_h=0.001$（物理 $0.01\,\mathrm{mm}$），
且 $h$ 的执行测量精确，则相邻帧的真实材料间距至多
\begin{equation}
 \ell_h\nu h_m\le0.0006<2a_h.
\end{equation}
因此连续帧的材料覆盖确实相交。
扫描起点取得首帧；到达 $h=\ell_h$ 后令 $\nu=0$，并在不超过 $h_m$ 的时间内取得末帧。
首帧与末帧的保证足迹分别须包含 $0$ 与 $\ell_h$，足迹也可直接截断到任务段。
等待末帧时反馈仍满足恢复合同，额外面内位移与转角分别不超过
$h_mr_\zeta M_1<0.000841$、$h_mr_\zeta M_2<0.000264$，
仍小于式~\eqref{v44:eq:new-finite-pose-margins}保留的余量。
因此首尾和全部中间足迹组成整段保证覆盖；对一般 $\nu$，覆盖时间至多比遍历时间多 $h_m$。
在 $\nu=2$ 时 $T=0.5\,\mathrm{s}$ 恰为采样时刻，可以直接取得末帧。
这个厚度是本例明确声明的成像几何部分；
如果只采用零厚度理想切片，则本例认证真实段遍历与采样点任务误差，
不把有限帧数自动解释为连续体积覆盖。

\subsection{一个显式失锁减速与反向返回桥接}
取 $\nu=2$，考虑在离段末还保留至少 $0.0025$ 的材料 $h$ 余量时，
一次失锁在某个有效采样时刻之后立即发生。
把该最后有效帧的实际相位记为 $h_{\rm loss}$；本具体例的失锁信号无检出延迟。
一般情况下，最近有效帧以后及失锁检出延迟中的位移应另加到时间、停止距离和回扫预算，
不能把未观测的当前相位当作最后覆盖端点。
冻结 $d,\alpha$ 的实际位置，关闭图像纠偏命令，并令内部 $\zeta$ 按
$\dot\zeta=-2\omega\zeta$ 制动。
本例使用速度级接口，所以这个命令切换不额外声称加速度连续。
在 $T_s=0.05\,\mathrm{s}$ 内用
\begin{equation}
 \nu_{\rm stop}(r)=2(1-3r^2+2r^3),\qquad r\in[0,1],
\end{equation}
停止相位前馈。其未经有效成像覆盖的前进距离为
\begin{equation}
 \Delta h=\ell_h\int_0^{T_s}\nu_{\rm stop}\,dt
         =0.0025.
\end{equation}
随后在 $T_b=0.1\,\mathrm{s}$ 内保持 $d,\alpha$，执行显式返回
\begin{equation}
 h(r)=h_{\rm loss}+\Delta h(1-3r^2+2r^3),\qquad r\in[0,1].
\end{equation}
该桥接从零 $h$ 速度开始并以零速度到达最后有效相位 $h_{\rm loss}$，
峰值返回速度为 $1.5\Delta h/T_b=0.0375$，即 $0.375\,\mathrm{mm/s}$。
它的持续时间为 $0.15\,\mathrm{s}$，反向相位总变差为
$\Delta h/\ell_h=0.05$，两者分别记账。
本具体例把图像中断定义为这 $0.15\,\mathrm{s}$ 的外生事件：
在返回末端重新取得目标关联正确且质量合格的一帧，随后恢复正常采样。
该时刻距失锁时刻恰为 $30$ 个采样间隔。
几何回位本身不推出观测必然恢复；若还需要等待，应另加已知等待时间和对应的漂移、状态预算。

给定整个模型族，冻结 $d,\alpha$ 时有限几何直接给出
\begin{equation}
 \|y(t)-y(t_{\rm loss})\|
 \le(\|\beta\|+\eta_c)\Delta h+w_0(T_s+T_b)
 =0.0014075.
\end{equation}
在第一次失锁前，由式~\eqref{v44:eq:new-finite-state-envelope}在 $t\le0.5$ 的界可取
$S_{\rm loss}\le0.003209576$。
当 $e$ 固定而 $v=\zeta/\omega$ 缩放为 $kv$、$0\le k\le1$ 时，
状态 $[e+v;v]$ 的对应线性变换范数至多 $\sqrt2$：
该变换是恒等映射与 $[r;v]\mapsto[r-v;0]$ 的凸组合，二者范数分别为 $1$ 与 $\sqrt2$。
因此整个减速与桥接段都满足
\begin{equation}
 S(t)\le\sqrt2(0.003209576)+0.0014075
     <0.005947<R_s.
 \label{v44:eq:explicit-bridge-state-bound}
\end{equation}
返回末端因 $h$ 也恢复原值，几何位移项相消，只剩目标漂移，故
\begin{equation}
S^+\le\sqrt2(0.003209576)+0.00015<0.004690<R_s.
\end{equation}
含观测噪声的可计算初值检查也通过，因为
$\|\widehat\chi^+\|+n_0\le S^++2n_0<0.005690<R_s$。
$d,\alpha$ 没有额外累计移动，$h$ 未超过预留有效范围，且全部输入幅度通过合同。
重新取得有效目标关联后，由当前图像重新构造 $C$，旧数据仍按同一模型有效，
新的初值通过恢复合同；返回到最后有效相位也保证采样覆盖不留失锁缺口。
此后恢复扫描不需要重新假定初始误差为零。

这个具体桥接给出扫描部分的实际总时长 $0.65\,\mathrm{s}$；
采用只知道负变差总量的保守式~\eqref{v44:eq:hybrid-continuous-time}则给出
$0.15+(1+0.05)/2=0.675\,\mathrm{s}$，二者一致。
如果失锁持续更久、发生多次或实际执行另有约束，应使用新的累计时间、负相位变化和物理管预算，
不能复用这一次短中断的数值结论。

\paragraph{结果的贡献范围。}
相位支持函数、凸范数约束与子空间约简本身采用标准工具。
本节提供的具体连接是：真实相位中性要求改变可用纠偏空间，
从而严格增加必要响应信息；同一组有限实际数据又同时进入恢复误差、前馈误差与真实推进能力。
有限几何例把方向数变化、有限探查能量、非零噪声与余项、可逆域、采样保持及一次显式反向桥接
落实在同一任务和同一物理尺度上。


\section{名义速率调度与真实覆盖接口}
\label{sec:rev:positive-progress}
以下调度只在已验收的短窗口内选择名义推进，不改变固定局部模型的首次响应下界。本节 $b_i$ 继续采用第~\ref{sec:cal-noisy}节已经紧化的命令界。令
\begin{equation}
 u_{\rm ff}=u_h+\nu b_s,\qquad
 Eu_h=0,\quad Eb_s=0,
 \label{eq:rev:progress-feedforward}
\end{equation}
其中 $u_h$ 是已知保持或运输动作，$b_s$ 是每单位名义进度的刚体输入方向。其实际材料相位作用由式~\eqref{v44:eq:actual-phase-port}验收。允许参考写为 $y_r(s,t)$。对接受候选，分别给出独立有效的界
\begin{equation}
 \|Jb_s-\partial_s y_r\|\le\alpha_s,
 \qquad
 \|Ju_h-\partial_t y_r+w_{\rm ext}+Jd_u\|\le d_0.
 \label{eq:rev:progress-residual}
\end{equation}
这些界应在整个拟接受状态管上成立，不能从尚未证明的管不变性反推。若 $Hb_s=\partial_s y_r$，可用 $\alpha_s=\delta\|Nb_s\|+\eta\|b_s\|$；有名义不匹配时再加 $\|Hb_s-\partial_s y_r\|$。若由数据构造前馈，可使用相应数据条件方向界。

固定 $\omega>0$、$\bar q<1$、$R>0$，令
\begin{equation}
 \lambda=\frac{\omega(1-\bar q)}2,\qquad
 \sqrt2R\le\min(r_{\rm task},\epsilon_{\rm task}),\qquad
 r_\zeta=\omega R.
 \label{eq:rev:progress-radius}
\end{equation}
以下 $\bar q$ 始终上界实际 $\|JC-I_m\|$：无偏修复以 $q(C,X)\le\bar q$ 验收，带偏差修复以 $q_{\rm gen}(C,X)\le\bar q$ 验收；共享固定余项可替换为相应共享证书。若此界成立且 $S(t_0)\le R$，一个充分的推进条件为
\begin{align}
 \nu\alpha_s+d_0+\sqrt2\omega n_{\max}&\le\lambda R,
 \label{eq:rev:progress-invariant}\\
 a_i^\top(u_h+\nu b_s)+\omega R\|a_i^\top C\|_2
&\le b_i.
 \label{eq:rev:progress-capability}
\end{align}
因为恢复定理给出 $D^+S\le-\lambda S+\nu\alpha_s+d_0+\sqrt2\omega n_{\max}$，式~\eqref{eq:rev:progress-invariant}使 $S\le R$ 前向不变，式~\eqref{eq:rev:progress-capability}同时保证域内实际输入能力。仅满足稳态半径条件不足以跳过初始状态验收。由测量 $\widehat e=e+n$ 可计算 $\widehat\chi=[\widehat e+\zeta/\omega;\zeta/\omega]$；$\|\widehat\chi\|+n_{\max}\le R$ 是一个直接可用的初值充分条件。

\paragraph{固定候选的显式速度区间。}
定义恢复余量及各输入面的余量
\begin{align}
 \mu_r&=\lambda R-d_0-\sqrt2\omega n_{\max},\\
 \mu_i&=b_i-a_i^\top u_h-\omega R\|a_i^\top C\|_2
.
 \label{eq:rev:progress-margins}
\end{align}
当 $\mu_r\ge0$、$\mu_i\ge0$ 时，在不考虑加速度约束前，一个充分区间是
\begin{equation}
 0\le\nu\le
 \min\left\{\nu_{\rm pref},
 \frac{\mu_r}{\alpha_s},
 \min_{i:a_i^\top b_s>0}\frac{\mu_i}{a_i^\top b_s}\right\}.
 \label{eq:rev:progress-upper}
\end{equation}
零分母且分子非负的项不限制速度；负余量表示该候选未通过合同，不能裁成零后视为可行。若接触前馈位于已约简输入子空间外，也必须先修正该候选。

\paragraph{速度变化与命令加速度。}
式~\eqref{eq:rev:progress-upper}只给瞬时上界。若接受的物理参考为 $p_C=\pi(s,t)$，其精确运动学为
\begin{equation}
 \dot p_C=\pi_s\nu+\pi_t,\qquad
 \ddot p_C=\pi_s\dot\nu+\pi_{ss}\nu^2+2\pi_{st}\nu+\pi_{tt}.
 \label{eq:rev:physical-progress-acceleration}
\end{equation}
姿态和机器人关节采用各自实际微分运动学检查。对 $u_{\rm cmd}=u_h+\nu b_s+C\zeta$，在共同输入坐标中还有
\begin{align}
 \dot u_{\rm cmd}
 ={}&\dot u_h+\dot\nu b_s+\nu D_t b_s+(D_t C)\zeta+C\dot\zeta,\\
 \|\dot\zeta\|&\le\omega^2(\sqrt2R+n_{\max}).
 \label{eq:rev:command-acceleration}
\end{align}
$D_t$ 是沿实际接受轨迹的完整总导数，包含材料位置、图像状态、身体与接触状态、实现分配及观测处理等所有依赖。例如 $C=C(s,t,\widehat y)$ 时有 $D_tC=\partial_tC+\nu\partial_sC+D_{\widehat y}C\,\dot{\widehat y}$；仅在预先指定 $C=C(s,t)$ 的特例中才能简化成 $\partial_tC+\nu\partial_sC$。幅度有界的测量噪声不保证 $\widehat y$ 可微，因此在线状态相关分配需要可微观测处理及其导数合同，或采用能直接验收的实际一步输入模型；缺少该条件时只保留响应与速度合同，不额外认证加速度。

物理坐标重建随时间变化时，还要包含该重建矩阵的导数。由有限窗口上相应完整导数界可验收 $|\dot\nu|$、纠偏加速度与总命令变化率。实际加速度若是硬约束，使用实际执行误差的导数界进一步紧化；只检查材料生成器的加速度不够。上述导数要求来自执行能力，二阶公共 Lyapunov 恢复证明本身仍只要求 $\Delta$ 的有效范数界。

\paragraph{推进优先的有限计划选择。}
在共同短时长 $T_p$ 上生成有限组 $C^1$ 速度计划及其实现候选。每个计划从当前 $(s,\nu)$ 和当前技能、反馈状态开始，沿整段检查~\eqref{eq:rev:progress-invariant}、\eqref{eq:rev:progress-capability}、速度与加速度、有效图、关联、接触、完整物理状态管和终端接续条件。记通过者为有限集 $\mathcal P_{\rm acc}$，计划 $j$ 的名义推进量为
\begin{equation}
 \Delta s_j=\int_0^{T_p}\nu_j(t)\,dt.
\end{equation}
若目标仅为名义生成器进度，可首先取 $\Delta s_j$ 最大；材料覆盖任务则应首先最大化本窗口认证的连续覆盖前缀增量 $\Delta\ell_j^{\rm cov}$，再在相同保证进度的候选中最小化切换、探查、技能与实现代价。该增量由真实相位下界、参考配准及有效帧足迹验收得到；坏模式和回扫使用第~\ref{v44:sec:hybrid-coverage}节记账。这给出明确的推进优先级，且最大值在有限非空集合上存在。它只相对于本次验收的计划库作选择，不声称连续全局最优。有限完成的计数以实际完成的窗口为单位：执行所选完整窗口；若中途重新规划，新计划必须保持原指定窗口末端的累计进度下界，不能反复把承诺的推进推迟到新的预测末端。对于覆盖分支，只有当全部接受计划的认证新增覆盖 $\Delta\ell_j^{\rm cov}$ 均为零时，才按没有保证覆盖进度处理；$\nu=0$ 本身不排除其它已认证端口项带来的真实相位运动。

\begin{proposition}[严格余量给出正名义速率及名义段完成]
\label{prop:rev:finite-progress}
考虑名义参数长度为 $L_s>0$ 的一段已声明任务参考。假定该段存在一条平滑接受的实现族，或已给出有限次、各自耗时有界的可行桥接；以下界均在整段声明物理管上独立成立：
\begin{enumerate}
 \item $q\le\bar q<1$，$S(t_0)\le R$，且恢复余量满足 $\mu_r\ge\underline\mu_r>0$；$\alpha_s\le\bar\alpha<\infty$。
 \item 各非等式输入面的余量满足 $\mu_i\ge\underline\mu_i>0$，且 $(a_i^\top b_s)_+\le\bar\beta_i<\infty$；$\nu_{\rm pref}\ge\underline{\nu}_{\mathrm{pref}}>0$。精确接触等式已按~\eqref{eq:rev:progress-feedforward}满足。
 \item 身体映射、关联、成像质量及完整物理状态管保持有效；速度与加速度检查在零推进附近具有统一严格余量，相关完整总导数存在预先给出的统一界。在线状态相关分配按前述观测处理合同计算这些界。
\end{enumerate}
则存在常数 $\nu_0>0$，使巡航计划 $\nu\equiv\nu_0$ 满足上述全部合同。若初始和末端需要速率过渡，且已按相同合同构造了总耗时不超过 $T_{\rm tr}$ 的有限平滑过渡，则名义参数段的巡航及过渡完成时间满足
\begin{equation}
 T_{\rm finish}-t_0\le T_{\rm tr}+\frac{L_s}{\nu_0}.
 \label{eq:rev:finite-progress-time}
\end{equation}
若采用前述有限计划库，并且每个非末端窗口都包含 $\Delta s_j\ge\Delta s_{\min}>0$ 的接受计划、窗口与桥接总时长不超过 $T_{\max}$，所选窗口的进度下界被实际执行或在中途重规划时保留，末端另有耗时不超过 $T_{\rm end}$ 的精确截断与制动计划，则推进优先策略满足
\begin{equation}
 T_{\rm finish}-t_0\le
 \left\lceil\frac{L_s}{\Delta s_{\min}}\right\rceil T_{\max}+T_{\rm end}.
 \label{eq:rev:finite-window-progress}
\end{equation}
\end{proposition}
\begin{proof}
前两项给出正的统一速度上界
\[
 \nu_c=\min\left\{\underline\nu_{\rm pref},
 \frac{\underline\mu_r}{\bar\alpha},
 \min_{i:\bar\beta_i>0}\frac{\underline\mu_i}{\bar\beta_i}\right\}>0,
\]
其中零分母项省略。第三项的严格余量与一致导数界，使速度和加速度关于 $\nu,\dot\nu$ 的新增项在它们足够小时保持可行。具体地，式~\eqref{eq:rev:command-acceleration}可界为 $A_0+B_0|\dot\nu|+B_1\nu+B_2\nu^2$；若 $A_0$ 比相应上限小 $\mu_a>0$，先取 $0<\nu_0\le\nu_c$ 使 $B_1\nu_0+B_2\nu_0^2\le\mu_a/2$，再令 $B_0|\dot\nu|\le\mu_a/2$。物理参考的其余加速度项同理处理。因此存在可行的正巡航速度及足够平缓的局部速率变化。

第一类结论对巡航积分 $\dot s=\nu_0$，再加已经验收的有限过渡时间即得。有限过渡的路径和端点必须实际可行；局部小速率结论不能代替任意实现切换的可达性。第二类结论中，推进优先策略的每个完整非末端窗口都至少增加 $\Delta s_{\min}$，故完整窗口数有上述上界，末端计划再给出 $T_{\rm end}$。两类执行期间的 $S\le R$ 由~\eqref{eq:rev:progress-invariant}和恢复比较不等式保证。
\end{proof}

本命题对 $\dot s=\nu$ 积分，因而证明名义参考段完成。要推出真实材料任务完成，应同时使用命题~\ref{v44:prop:joint-certificate}中的真实相位及参考配准条件，以及定理~\ref{v44:thm:hybrid-finish}中的有效覆盖、中断时间和反向相位预算。只有已另行证明 $s=\sigma$ 且有效成像连续覆盖时，两个完成计数才可以直接等同。对于启动与制动，一个简单的 $C^1$ 速率过渡为 $\nu(t)=\nu_0(3r^2-2r^3)$、$r=t/T_r\in[0,1]$，其 $\max|\dot\nu|=3\nu_0/(2T_r)$；应按剩余材料距离与上述约束验收该段，而不是在非零速度时瞬时置零。





\section{有限探查库相对于固定凸合同的条件完备性}
\label{sec:cal-finite}

固定一个实现 $a$ 及其局部响应窗口。令 $\mathcal U=\{u:a_i^\top u\le b_i\}$ 为同一选定且已紧化的命令输入内近似；固定 $u_{\rm ff}$、$r_\zeta$ 和精确接触矩阵 $E$。若总输入要求 $Eu=0$，先验收 $Eu_{\rm ff}=0$。定义与式~\eqref{eq:cal-convex-repair} 完全相同的能力合同
\begin{equation}
 \mathcal K_{\rm cap}=
 \left\{
 C:
 EC=0,\ 
 a_i^\top u_{\rm ff}
 +r_\zeta\|a_i^\top C\|_2\le b_i
 \quad\forall i
 \right\}.
 \label{rev:eq:same-convex-contract}
\end{equation}
该集合保证 $u_{\rm ff}+C\zeta\in\mathcal U$ 对全部
$\|\zeta\|_2\le r_\zeta$ 成立。

对该固定模型给定有限探查库
$\{u_1,\ldots,u_M\}$，并令
\begin{equation}
 \mathcal L_N=
 \operatorname{span}\{Nu_1,\ldots,Nu_M\},
 \qquad
 r_{\rm lib}=\dim\mathcal L_N .
 \label{rev:eq:library-space}
\end{equation}

\begin{assumption}[共同窗口中的可执行库基]
\label{rev:ass:executable-library}
存在由库中动作组成的有限实际探查序列，其有效激励张成 $\mathcal L_N$；该序列及必要的回位、接触保持与观测过程均能在所声明的响应模型窗口内执行。全部有效数据对应同一冻结 $\Theta,H,N$。校准序列的可执行性独立于后续纠偏合同 $\mathcal K_{\rm cap}$ 验收。
\end{assumption}

\begin{theorem}[指定库与指定凸合同的条件完备性]
\label{rev:thm:library-completeness}
在假设~\ref{rev:ass:executable-library} 下，设
$E_{\rm geo}=0$ 且数据无噪声，证书中取 $\varepsilon_{\rm data}=\eta=0$，并取 $0\le q_{\rm goal}<1$。对实际真实响应 $J$，若存在
\begin{equation}
 C_{\rm lib}\in\mathcal K_{\rm cap},\qquad
 JC_{\rm lib}=I_m,\qquad
 \operatorname{range}(NC_{\rm lib})
 \subseteq\mathcal L_N,
 \label{rev:eq:library-witness}
\end{equation}
则校准一个完整库基后，式~\eqref{eq:cal-convex-repair} 的可行集包含
$C_{\rm lib}$，且其 $q$ 值为零。所需不同的独立激励方向至多为 $r_{\rm lib}$。
\end{theorem}
\begin{proof}
完整库基给出 $\operatorname{range}(X)=\mathcal L_N$。无噪声下
$\widehat\Theta=\Theta P_X$，并且
$(I-P_X)NC_{\rm lib}=0$，故
\begin{equation}
 \widehat J C_{\rm lib}
 =
 H C_{\rm lib}
 +\Theta P_XNC_{\rm lib}
 =
 JC_{\rm lib}
 =I_m .
 \label{rev:eq:library-exact-repair}
\end{equation}
又因 $\varepsilon_{\rm data}=\eta=0$，有
$q(C_{\rm lib},X)=0$。其余能力约束由
$C_{\rm lib}\in\mathcal K_{\rm cap}$ 保证。因此该矩阵是同一凸问题的一个可行点。
\end{proof}

这一完备性相对于指定的探查库、同一局部模型和同一凸内近似成立。真实物理可行域中的一个矩阵可能不属于选定的 $\mathcal K_{\rm cap}$，因而不能据其物理可行性直接推出本定理的前提。独立方向数也不计入必要的回位、切换动作或观测时长。

有噪声时，仅对满足 $\operatorname{range}(NC)\subseteq\mathcal L_N$ 的修复实现，满库秩才消除未激励项；仍需验收
\begin{equation}
 \varepsilon_{\rm data}\|X^\dagger NC\|_2+\eta\|C\|_2
 \le q_{\rm goal}.
 \label{rev:eq:noisy-library-residual}
\end{equation}
有限方向维数不保证有限次重复后必然成功：确定性系统误差可能不随重复衰减，低幅激励也可能放大测量误差。只有另行给定重复采样误差规律或具有严格余量的数据条件，才能推出相应有限样本保证。

实际算法可以用当前候选的 $(I-P_X)NC$ 的主要左奇异方向提出下一条探查，再由能力检查和有限库决定是否执行；这一贪心顺序不具有本文所未证明的最少探查最优性。若所选库基不可共同执行，或全部候选均无法满足合同，算法返回当前保证不可验收。





\section{合法混合探查与后续实现可以相互适配}
以下是代数机制例，尚不是人体实验。取
\begin{equation}
 H=\begin{bmatrix}1&0&0\\0&1&0\end{bmatrix},\qquad
 N=\begin{bmatrix}1&0&1\\0&1&0\end{bmatrix}.
\end{equation}
任务零空间通道为 $\mathcal C=\operatorname{span}(1,0)$，最小暴露实现及其矩阵为
\begin{equation}
 B_\star=[e_1-e_3,e_2],\qquad
 D_\star=NB_\star=\begin{bmatrix}0&0\\0&1\end{bmatrix}.
\end{equation}
假设最便宜或唯一合法的探查为 $u_p=e_2+e_3$，则实际激励是 $q=Nu_p=(1,1)^\top$。若固定 $B_\star$，投影后暴露仍有
\begin{equation}
 \norm{(I-P_Q)D_\star}=1/\sqrt2,
 \qquad Q=\operatorname{span}(q).
\end{equation}
但 $\mathcal C+Q=\R^2$，正确剩余 $D_Q=0$。选择
\begin{equation}
 B_Q=[e_1-e_3,e_2+e_3],\quad
 HB_Q=I_2,\quad NB_Q=\begin{bmatrix}0&1\\0&1\end{bmatrix}
\end{equation}
即可把未来暴露放到已校准的混合方向。测到的向量 $g=\Theta q$ 给出
\begin{equation}
 JB_Q=I_2+[0,g],\qquad
 C_{\rm repair}=B_Q(I_2+[0,g])^{-1}.
\end{equation}
当该增益可逆且实际能力验收通过时，一次合法混合方向探查足够修复全部任务响应；若 $\delta\sqrt2<1$，可逆性对整个范数球成立。代价是 $\norm{B_Q}=\sqrt3$ 而 $\norm{B_\star}=\sqrt2$。

这里不要求 IRD 产生恰好等于理想 SVD 的动作。实际探查可以受约束；机器人也可以改变后续实现，使未来所需反馈落入真实已经校准的方向。这一双向匹配由 $D_Q$ 定理描述，而不是额外增加一个学习模块。




\clearpage
\appendix
\part*{附录：几何、技能及执行接口的完整定义}
\addcontentsline{toc}{part}{附录：几何、技能及执行接口的完整定义}
以下保留超声几何、补充解析例、材料图册、身体反演、技能生成与 IRD 连接的完整推导。
正文的局部理论在其明示输入与观测合同下成立；这些附录给出构造该合同的具体接口。
\section{超声交互矩阵与有限运动几何}\label{sec:rev:ultrasound}
\subsection{面内动作和面外动作需要不同的未知量}
在有效、关联正确的局部窗口，血管切向为 $t$，探头原点为 $p$，切片法向为 $n$，面内单位基为 $E=[e_x,e_z]$。局部交点由
\begin{equation}
 \gamma(s)=\gamma_0+ts,\qquad
 n^\top(\gamma(s)-p)=0,\qquad
 y=E^\top(\gamma(s)-p)
\end{equation}
给出。固定姿态对时间微分，若 $n^\top t\ne0$，则
\begin{align}
 \dot s&=\frac{n^\top v}{n^\top t},\\
 \dot y&=-E^\top v+\beta n^\top v,
 &\beta&=\frac{E^\top t}{n^\top t}.
\end{align}
二维单切片看不到完整面外血管切向，因此 $\beta$ 需要先验或校准。若 $n^\top t$ 太小，交点模型病态，应拒绝该窗口而非宣布小样本可修复。绕切片法向的转角符号取 $\dot y=(z,-x)^\top\dot\alpha$，该动作不改变切片本身。

局部平面皮肤下，选择面内切向平移、绕切片法向摇摆、面外皮肤切向平移三种控制，得到
\begin{equation}
 J(y)=
 \begin{bmatrix}-1&z&\beta_x\\0&-x&\beta_z\end{bmatrix}
 =H+\theta N,\quad
 H=\begin{bmatrix}-1&z&\bar\beta_x\\0&-x&\bar\beta_z\end{bmatrix},\quad
 N=\begin{bmatrix}0&0&1\end{bmatrix}.
\end{equation}
这里 $\theta=\beta-\bar\beta\in\R^{2\times1}$，不是人为添加的低秩未知模型：面外斜率只在面外位移时进入，具有实际几何来源。若两分量均未知，任意有限 $\beta$ 都对应真实单位切向
\begin{equation}
 t(\beta)=\frac{E\beta+n}{\sqrt{1+\norm{\beta}^2}},
\end{equation}
因此局部斜率球可作为明示的物理模型类。若另有解剖约束使未知量只有一个标量或更小集合，谱范数球只是保守包络：修复上界仍有效，精确最小最大下界不能直接宣称在缩小后的物理类中仍紧。未知的接触变形、传感标定误差、关联错误不能由这个因式分解自动消除；前三者的可界定部分进入剩余项，关联错误则使模型无效。

中心实现 $x=0$、$\bar\beta_z\ne0$ 时，$H$ 满行秩，但 $H\ker N$ 只能产生一维面内任务作用。因此最小暴露秩为 $k_*=1$；若要求覆盖完整球的精确任务响应，还须满足 $\delta\|D_*\|<1$。近似首次响应所需方向数按式~\eqref{eq:cal-epsilon-dimension} 的容差判定。若任务允许 $x=\ell\ne0$，则
\begin{equation}
 B_\ell=
 \begin{bmatrix}
 -1&-z/\ell\\0&-1/\ell\\0&0
 \end{bmatrix},\qquad
 JB_\ell=I_2,\qquad NB_\ell=0.
\end{equation}
两维任务仍被控制，而未知 $\beta$ 不被调用。这不是被动换坐标，而是换了实际机器人实现与真实杠杆。另一方面，$1/|\ell|$ 放大了运动需求；太小的偏移可能违反角速度、接触和关节能力约束。

\subsection{有限动作必须用精确接触几何验收}
固定切片平面、静止目标交点和局部平面皮肤；CoP 沿固定皮肤切向平移 $d$，探头在该平面中转动 $\alpha$。则
\begin{equation}
 y'=R_2(-\alpha)\begin{bmatrix}x-d\\z\end{bmatrix},\qquad
 R_2(\alpha)=\begin{bmatrix}\cos\alpha&-\sin\alpha\\\sin\alpha&\cos\alpha\end{bmatrix}.
\end{equation}
目标 $y'=(\ell,z_\star)$ 等价于
\begin{align}
 z&=\ell\sin\alpha+z_\star\cos\alpha,\\
 d&=x-\ell\cos\alpha+z_\star\sin\alpha.
\end{align}
第一式在零角附近的导数为 $\ell$；$\ell\ne0$ 时由隐函数定理得到局部小角解。无角度限制下仍有正深度偏差硬界
\begin{equation}
 z-z_\star\le\sqrt{\ell^2+z_\star^2}-z_\star.
\end{equation}
若有 $|\alpha|\le\alpha_{\max}$，需要取该区间上 $\ell\sin\alpha+z_\star\cos\alpha$ 的实际极值。不能只用一阶 Jacobian 宣称任意有限动作可达。

全路径接触约束下的已知交互矩阵为
\begin{equation}
 H_0(\alpha)=\begin{bmatrix}-\cos\alpha&z'\\\sin\alpha&-x'\end{bmatrix},\qquad
 \det H_0=x'\cos\alpha-z'\sin\alpha=x-d.
\end{equation}
持续可控需要物理横向杠杆 $|x-d|$ 全程有下界；图像横向偏移本身不是完整条件。有限控制可以直接求解上述精确几何，避免把已知的大转角非线性误差塞成不必要的校准。

\noindent\textbf{纯几何核验。}取 $x=\ell=5\,\mathrm{mm}$、$z_\star=30\,\mathrm{mm}$，正深度硬界为 $0.41381265\,\mathrm{mm}$。$\pm0.5\,\mathrm{mm}$ 的对称修复不成立。对 $\pm0.1\,\mathrm{mm}$，表中显示舍入解，终点误差以未舍入解计算：
\begin{center}
\begin{tabular}{rrrr}
\toprule
初始深度差 & 转角 & CoP 平移 & 终点目标误差\\
\midrule
$-0.1\,\mathrm{mm}$ & $-1.08441^\circ$ & $-0.566866\,\mathrm{mm}$ & $<10^{-12}\,\mathrm{mm}$\\
$+0.1\,\mathrm{mm}$ & $+1.22452^\circ$ & $+0.642248\,\mathrm{mm}$ & $<10^{-12}\,\mathrm{mm}$\\
\bottomrule
\end{tabular}
\end{center}
这些是精确刚体模型的数值核验，不是真实探头实验。声学耦合、压力、组织变形、IK 和碰撞仍需独立验收。

\subsection{同一任务、首次合同与总成本的边界}
面外移动会改变交点的材料相位，故共同任务必须允许同一血管片段内的有限相位偏差。若要求严格固定材料截面，则不能将面外中心纠偏与固定平面偏心纠偏当成同一任务。

中心时的一方向下界只对应冻结局部模型与首次合同。任意多步策略可以先平移形成杠杆再摇摆，因此不能声称所有从中心出发的扫描策略都必须校准 $\beta$。这条重构路径本身也是可选实现，其耗时必须计入。

长血管的推进通常仍含面外前馈 $u_{\rm ff}$，产生未知漂移 $\Theta Nu_{\rm ff}$；它进入恢复定理中的 $w$，例如需要纳入 $\delta\norm{Nu_{\rm ff}}$ 的界。$NB_\ell=0$ 免去的是这两个局部特征纠偏的斜率校准，不是免去全扫描的所有未知量。若任务合同还要求精确控制推进相位或额外切面特征，应把这些要求加入任务模型，重新计算 $H,N,D_\star$，不能搬用二维结论。

按横向速度上限 $15\,\mathrm{mm/s}$，从中心转到 $5\,\mathrm{mm}$ 偏移至少要 $0.333\,\mathrm{s}$。免去一个校准方向不保证总时间降低；只有
\begin{equation}
 T_{\rm switch}+T_{\rm probe}+T_{\rm repair}+T_{\rm execute}
\end{equation}
在同一任务质量和真实预算下更小，才构成工程收益。



% 可插入原超声解析节之后；仅使用主稿已经定义的宏与定理环境。


\section{同时满足校准、能力和推进条件的定量解析实例}
\label{rev:quantitative-example}

本节给出一个具有严格余量的有限扫描段。任务要求跟随同一关联正确的血管材料片段，并允许中心与偏心成像两种实现；两个图像特征均保留。允许声窗取 $|x_{\rm phys}|\le12\,\mathrm{mm}$、$25\,\mathrm{mm}\le z_{\rm phys}\le35\,\mathrm{mm}$。实现选择完成后，再启动本节的短窗口；选择实现及启动、结束桥接的时间另行计入总代价。

\subsection{固定物理尺度下的精确几何与非零模型余项}
取长度尺度 $L=10\,\mathrm{mm}$。本节的 $y=(x,z)^\top$、面内 CoP 位移 $d$ 和面外扫描坐标 $h$ 均已除以 $L$，$\alpha$ 保持弧度。速度输入为
\begin{equation}
 u=(\dot d,\dot\alpha,\dot h)^\top,
 \qquad q_{\rm p}=(d,\alpha,h)^\top.
 \label{rev:example-units}
\end{equation}
三个分量均具有 $\mathrm{s}^{-1}$ 的量纲；输入第二分量的物理权重由这一固定长度尺度确定，后续比较不再重选尺度。局部皮肤为平面，CoP 保持法向高度，摇摆围绕当前 CoP 和切片法向进行。取固定刚体接触附着 $r_C\equiv0$，即 TCP 与 CoP 重合，因此 $\dot p_C=v_C^{\rm rb}$，$u$ 是接触约束约简后的实际刚体输入。本例的速度执行接口为 $u_{\rm act}=u_{\rm cmd}$，即 $d_u=0$。

设世界切片基准中的血管斜率为
\begin{equation}
 \beta=\bar\beta+\theta,\qquad
 \bar\beta=(0,0.5)^\top,\qquad \norm{\theta}\le\delta=0.1.
 \label{rev:example-slope}
\end{equation}
因此未加形变的斜率满足 $\norm{\beta}\le0.6$。
为给出具有明确物理含义的非零余项，令局部血管在与切片相交的材料坐标 $h$ 处，基准切片坐标中的附加位移为 $\mathsf Dq_{\rm p}+g(t)$，其中
\begin{equation}
 \norm{\mathsf D}\le\eta_c=0.003,
 \qquad g(0)=0,
 \qquad\norm{\dot g(t)}\le w_0=0.001\,\mathrm{s}^{-1}.
 \label{rev:example-contact-error}
\end{equation}
这是一个明示的局部接触形变族，例如
$\mathsf D=0.003[\,I_2\ \ 0\,]$ 即为非零成员。其前两列描述随 CoP 平移与摇摆产生的位移，第三列允许附加材料斜率。因此实际斜率为 $\beta+\mathsf D_{:,3}$，满足
$|n^\top t|\ge1/\sqrt{1+0.603^2}>0.856$，切片交点一致横截。该形变保持材料的面外坐标 $h$，所以实际扫描的材料段仍可由 $h$ 直接记账。

偏心实现的初始交点为 $y_0=(1,3)^\top$。其有限几何严格写为
\begin{equation}
 y=R_2(-\alpha)
 \left[y_0+\beta h-\begin{bmatrix}d\\0\end{bmatrix}
       +\mathsf Dq_{\rm p}+g(t)\right].
 \label{rev:example-exact-geometry}
\end{equation}
直接求导得到
\begin{align}
 \dot y&=J_{\rm true}u+w_{\rm ext},
 &\norm{w_{\rm ext}}&\le w_0,\nonumber\\
 J_{\rm true}&=
 \left[\,A(y,\alpha)\ \ R_2(-\alpha)\beta\,\right]
             +R_2(-\alpha)\mathsf D,
 &A(y,\alpha)&=
 \begin{bmatrix}-\cos\alpha&z\\ \sin\alpha&-x\end{bmatrix}.
 \label{rev:example-exact-jacobian}
\end{align}
这里的角速度列 $(z,-x)^\top$ 使用式~\eqref{rev:example-exact-geometry} 的真实 $y$，已经包含形变后的完整旋转导数。因而 $R_2(-\alpha)\mathsf D$ 是额外接触响应，已知有限转角几何未被重复计入余项。

令图像估计为 $\widehat y=y+n$，$\norm{n}\le n_0=0.002$，对应 $0.02\,\mathrm{mm}$。在线实现由 $\widehat y$ 计算，采用
\begin{equation}
 H=\left[\,A(\widehat y,\alpha)\ \ \bar\beta_\alpha\,\right],
 \quad\bar\beta_\alpha=R_2(-\alpha)\bar\beta,
 \quad N=\begin{bmatrix}0&0&1\end{bmatrix}.
 \label{rev:example-known-model}
\end{equation}
于是仍有严格分解
\begin{align}
 J_{\rm true}&=H+\Theta_\alpha N+E,
 &\Theta_\alpha&=R_2(-\alpha)\theta,\nonumber\\
 E&=R_2(-\alpha)\mathsf D+
 \begin{bmatrix}0&-n_z&0\\0&n_x&0\end{bmatrix},
 &\norm{E}&\le\eta_c+n_0=:\eta=0.005.
 \label{rev:example-error-factorization}
\end{align}
因此图像误差在控制矩阵计算中的作用也已纳入 $\eta$。它与反馈中的 $n$ 可以相关；后续界不使用独立性或 $\dot n$ 的假设。所有扰动取可测信号，$g$ 绝对连续，动力学按 Carathéodory 解理解；在下述可逆域上，关于状态的局部 Lipschitz 性与这些有界输入保证局部解存在。

\subsection{结构化校准方向的严格比较}
先固定 $\alpha=0$，在 $E=0$ 的结构化主模型内比较两个允许实现。中心参考 $y_0^c=(0,3)^\top$ 给出
\begin{equation}
 H_c=\begin{bmatrix}-1&3&0\\0&0&0.5\end{bmatrix},
 \quad H_c\ker N=\operatorname{span}(e_1),
 \quad D_{\star,c}=\begin{bmatrix}0&2\end{bmatrix}.
 \label{rev:example-centred-rank}
\end{equation}
所以 $k_{\star,c}=1$，且 $\delta\norm{D_{\star,c}}=0.2<1$。全部允许斜率均保持中心瞬时任务满射；这一区别不依赖允许模型中存在不可控成员。零数据的最坏首次单位响应误差为 $0.2$，若要求其不超过 $0.05$，至少需要一个独立通道方向。

偏心参考 $y_0=(1,3)^\top$ 的已知前两列可逆，故
\begin{equation}
 C=\begin{bmatrix}A(\widehat y,\alpha)^{-1}\\0\quad0\end{bmatrix},
 \qquad HC=I_2,
 \qquad NC=0.
 \label{rev:example-allocation}
\end{equation}
在无噪声结构化主模型中，$k_{\star,o}=0$。下面保留 $\eta>0$ 与图像噪声，通过鲁棒误差管验收同一个偏心实现。精确方向数用于结构化主模型，有限精度执行则由 $q$ 证书给出。若实际策略先从中心形成偏心杠杆，应将其作为一次主动实现改变计入；中心的首次响应下界不限制这条多步重构策略。

\subsection{可逆域、反馈包络和正推进上限}
取图像参考 $y_r=y_0$，任务容差 $\epsilon=0.02$，即 $0.2\,\mathrm{mm}$。考虑闭域
\begin{equation}
 \begin{gathered}
 \mathcal D=\left\{(q_{\rm p},y):
 \norm{y-y_r}\le\epsilon,
 \ |\alpha|\le\alpha_{\max},
 \ |d|\le d_{\max},
 \ 0\le h\le\ell_h
 \right\},\\
 \alpha_{\max}=0.09,
 \qquad d_{\max}=0.3,
 \qquad\ell_h=0.05.
 \end{gathered}
 \label{rev:example-domain}
\end{equation}
由 $\cos\alpha\ge1-0.09^2/2$ 与 $|\sin\alpha|\le0.09$，真实物理杠杆在该域满足
\begin{equation}
 x\cos\alpha-z\sin\alpha
 \ge0.98(1-0.09^2/2)-3.02(0.09)
 =0.704231.
 \label{rev:example-physical-lever}
\end{equation}
计算控制所用的测量域为
$|\widehat x-1|,|\widehat z-3|\le\epsilon+n_0=0.022$，故
\begin{equation}
 \widehat x\cos\alpha-\widehat z\sin\alpha
 \ge0.978(1-0.09^2/2)-3.022(0.09)
 =0.7020591.
 \label{rev:example-measured-lever}
\end{equation}
使用
\begin{equation}
 A(\widehat y,\alpha)^{-1}
 =\frac1{\widehat x\cos\alpha-\widehat z\sin\alpha}
 \begin{bmatrix}-\widehat x&-\widehat z\\-\sin\alpha&-\cos\alpha\end{bmatrix}
 \label{rev:example-inverse}
\end{equation}
得到向上取整的统一界
\begin{equation}
 \norm{C}\le M_C=4.77,
 \qquad\norm{C_{1,:}}\le M_1=4.56,
 \qquad\norm{C_{2,:}}\le M_2=1.43.
 \label{rev:example-inverse-bounds}
\end{equation}
例如 $M_C$ 由
$\sqrt{1+1.022^2+3.022^2}/0.7020591$ 的解析上界得到。
由于 $NC=0$，不作斜率探查即可获得
\begin{equation}
 JC=I_2+EC,
 \qquad q\le\eta M_C=0.02385<q_g:=0.05.
 \label{rev:example-zero-probe-certificate}
\end{equation}

扫描段长度为 $L\ell_h=0.5\,\mathrm{mm}$，归一化进度满足
$s\in[0,1]$、$\dot s=\nu>0$。令
\begin{equation}
 b_s=\ell_h\begin{bmatrix}-A(\widehat y,\alpha)^{-1}\bar\beta_\alpha\\1\end{bmatrix},
 \quad u_{\rm ff}=\nu b_s,
 \quad Hb_s=0=y_r'(s),
 \quad Nb_s=\ell_h.
 \label{rev:example-feedforward}
\end{equation}
于是 $\dot h=\ell_h\nu$ 精确成立，而未校准斜率仍作为推进扰动进入
\begin{equation}
 \begin{aligned}
 \norm{w}&\le\nu a_s+w_0,
 \qquad a_s:=\delta\ell_h+\eta M_b,\\
 M_b&:=\ell_h\sqrt{1+M_C^2\norm{\bar\beta}^2}
 <0.129309,\qquad a_s<0.00564655.
 \end{aligned}
 \label{rev:example-progress-error}
\end{equation}
选二阶反馈增益 $\omega=4\,\mathrm{s}^{-1}$，令
\begin{equation}
 R_s=\frac{0.02}{\sqrt2},
 \quad r_\zeta=\omega R_s=\frac{0.08}{\sqrt2}\,\mathrm{s}^{-1},
 \quad\lambda=\frac{\omega(1-q_g)}2=1.9\,\mathrm{s}^{-1}.
 \label{rev:example-feedback-constants}
\end{equation}
这里 $r_\zeta\approx0.0565686\,\mathrm{s}^{-1}$。反馈包络的不变性充分条件为
\begin{equation}
 \nu a_s+w_0+\sqrt2\omega n_0\le\lambda R_s.
 \label{rev:example-recovery-inequality}
\end{equation}
采用式~\eqref{rev:example-progress-error} 中 $M_b$ 的精确表达式，其推进上限为
\begin{equation}
 \nu\le
 \frac{\lambda R_s-w_0-\sqrt2\omega n_0}{a_s}
 =2.5779236\ldots\,\mathrm{s}^{-1}.
 \label{rev:example-recovery-speed}
\end{equation}

\subsection{输入能力和有限窗口共同给出非空速度区间}
采用速度级执行合同
\begin{equation}
 |\dot d_{\rm phys}|\le6\,\mathrm{mm/s},
 \quad|\dot\alpha|\le10^\circ/\mathrm{s},
 \quad|\dot h_{\rm phys}|\le1.5\,\mathrm{mm/s}.
 \label{rev:example-physical-input-set}
\end{equation}
由 $u=u_{\rm ff}+C\zeta$ 和 $\norm{\zeta}\le r_\zeta$，三个分量满足
\begin{equation}
 |u_i|\le M_i(\nu\ell_h\norm{\bar\beta}+r_\zeta),\quad i=1,2,
 \qquad u_3=\ell_h\nu.
 \label{rev:example-component-inputs}
\end{equation}
对应推进上限依次为
$3.0004161\ldots$、$2.6192981\ldots$、$3\,\mathrm{s}^{-1}$。

固定 $\nu$ 时完成时间为 $T=1/\nu$。从 $d(0)=\alpha(0)=0$ 出发，积分同一分量界得到
\begin{align}
 |\alpha(t)|&\le\ell_h\norm{\bar\beta}M_2+\frac{r_\zeta M_2}{\nu},\nonumber\\
 |d(t)|&\le\ell_h\norm{\bar\beta}M_1+\frac{r_\zeta M_1}{\nu},
 \qquad 0\le t\le T.
 \label{rev:example-finite-pose-bounds}
\end{align}
因此本例的有限段证书还包含正的推进下限：
\begin{equation}
 \nu>
 \max\left\{
 \frac{r_\zeta M_2}{\alpha_{\max}-\ell_h\norm{\bar\beta}M_2},
 \frac{r_\zeta M_1}{d_{\max}-\ell_h\norm{\bar\beta}M_1}
 \right\}
 =1.4911154\ldots\,\mathrm{s}^{-1}.
 \label{rev:example-minimum-speed}
\end{equation}
这一正下限来自对有限时间内全部可能纠偏动作的积分上界。它是当前充分证书的性质，不代表机器人低于该速度必然不可行。结合恢复与输入上限，可选取具有严格余量的闭区间
\begin{equation}
 \boxed{\nu\in[1.6,2.5]\,\mathrm{s}^{-1}.}
 \label{rev:example-nonempty-speed-interval}
\end{equation}
其实际推进速度为 $0.8$--$1.25\,\mathrm{mm/s}$，完成时间为 $0.4$--$0.625\,\mathrm{s}$。

\begin{proposition}[含非零余项与图像噪声的有限扫描段]
\label{rev:example-finite-scan-proposition}
在式~\eqref{rev:example-exact-geometry}--\eqref{rev:example-error-factorization} 的模型类、速度级执行合同及参数范围内，取 $y(0)=y_r$、$\zeta(0)=0$、$d(0)=\alpha(0)=h(0)=0$。窗口初始实际速度由状态接续步骤设置为 $u(0)=\nu b_s(\widehat y(0),0)$。对任意式~\eqref{rev:example-nonempty-speed-interval} 中的常数 $\nu$，在线分配~\eqref{rev:example-allocation} 与二阶反馈在 $T=1/\nu$ 内完成 $0.5\,\mathrm{mm}$ 的材料扫描段，全程满足 $\norm{y-y_r}\le0.02$、式~\eqref{rev:example-domain} 的位姿范围与式~\eqref{rev:example-physical-input-set} 的输入界，无需斜率校准数据。
\end{proposition}
\begin{proof}
令 $v=\zeta/\omega$、$S=\norm{[\,(e+v)^\top,v^\top\,]^\top}$，其中 $e=y-y_r$。在 $\mathcal D$ 内，式~\eqref{rev:example-measured-lever} 保证 $C$ 始终有定义；式~\eqref{rev:example-zero-probe-certificate} 与二阶恢复定理给出
\begin{equation}
 D^+S\le-\lambda S+\nu a_s+w_0+\sqrt2\omega n_0.
 \label{rev:example-scalar-envelope}
\end{equation}
初始 $S(0)=0$。区间内最不利恢复右端在 $\nu=2.5$ 时小于
$0.026431<\lambda R_s=0.0268700\ldots$，故 $S(t)<R_s$，并有 $\norm{e(t)}\le\sqrt2S(t)<0.02$。

设存在 $T$ 之前首次离开 $\mathcal D$ 的时刻。在此前各个界成立，式~\eqref{rev:example-finite-pose-bounds} 在区间内最不利的 $\nu=1.6$ 给出
$|\alpha|<0.086309<0.09$、$|d|<0.275221<0.3$。输入界在区间内最不利的 $\nu=2.5$ 给出
$|\dot d_{\rm phys}|<5.430\,\mathrm{mm/s}$、$|\dot\alpha|<9.756^\circ/\mathrm{s}$、$\dot h_{\rm phys}\le1.25\,\mathrm{mm/s}$。又因 $NC=0$，始终有 $h(t)=\ell_h\nu t$。这些严格余量排除首次退出，解可延续至 $T$，且 $h(T)=\ell_h$。
\end{proof}

\Needspace{13\baselineskip}
例如选取 $\nu=2\,\mathrm{s}^{-1}$，解析证书如下。表中有限时间误差由式~\eqref{rev:example-scalar-envelope} 的指数解计算，各数值向外取整。
\begin{center}
\small
\begin{tabular}{lrr}
\toprule
量 & 解析上界 & 允许值\\
\midrule
响应误差 $q$ & $0.02385$ & $0.05$\\
全段图像误差 & $0.108\,\mathrm{mm}$ & $0.2\,\mathrm{mm}$\\
面内 CoP 速度 & $4.860\,\mathrm{mm/s}$ & $6\,\mathrm{mm/s}$\\
摇摆角速度 & $8.732^\circ/\mathrm{s}$ & $10^\circ/\mathrm{s}$\\
面外推进速度 & $1.000\,\mathrm{mm/s}$ & $1.5\,\mathrm{mm/s}$\\
累计摇摆角 & $0.076197\,\mathrm{rad}$ & $0.09\,\mathrm{rad}$\\
累计面内位移 & $2.430\,\mathrm{mm}$ & $3\,\mathrm{mm}$\\
\bottomrule
\end{tabular}
\end{center}

本例在一个有限模型窗口内，同时给出了结构化方向数 $1\to0$、非零 $\eta$、图像噪声、实际输入幅度、真实杠杆下界和正推进区间。初始两种实现的理论比较各自使用对应的局部模型；偏心执行没有复用中心探查数据。沿更长血管拼接多个窗口时，仍按前述窗口更新、状态接续与观测关联条件验收。

本节采用速度级合同，给出速度幅度与位姿范围。若执行器另有加速度限制，需对可微观测 $\widehat y$、姿态及进度的导数给界，或使用实际离散步进模型另行验收。仅有幅度界的 $n$ 不推出 $\dot C$ 的界，也不自动证明启动和停止桥接可行；这些桥接按主文的速度变化与状态接续条件独立处理。




\section{材料图册、身体映射与姿态纤维}\label{sec:rev:geometry}
\subsection{保留圆柱设计，并区分图册和内部解剖}
在每个有效环壳材料子域 $\Omega_\ell$ 上，保留
\begin{equation}
 \xi=(\varphi,h,d),\qquad
 \Delta_{g_0}h=0,\quad \Delta_{g_0}d=0,\quad
 \Delta_{g_0}\widetilde\varphi=0.
\end{equation}
轴向与径向场取各自端面或边界上的 Dirichlet 条件；角度场在切开的域上取两岸相差 $2\pi$ 的周期 lift，跨缝使用转移图。圆柱核心退化域需要独立图，腿和躯干可以是不同圆柱图；不强行制造一个覆盖整个人体的单图。调和性不保证三维坐标无折叠，必须验收 $\det D_\xi F_{\ell,0}\ne0$、方向一致性及自交。

SMPL-X 表面采用标准蒙皮形式\cite{pavlakos2019}
\begin{align}
 \bar v_i(c)&=\bar v_i+B_{S,i}(\beta)+B_{E,i}(\psi)+B_{P,i}(\vartheta),\\
 v_i(c)&=\sum_k w_{ik}\bigl(R_k(c)\bar v_i(c)+t_k(c)\bigr).
\end{align}
这里 $c$ 含体型、姿态、表情和世界位姿，骨架变换含静止姿态校正。体积内部需要另外给出权重 $w_k(p_0)$ 和修正场 $b_c(p_0)$；表面 SMPL-X 不会自动提供它们，也不会给出患者内部血管的准确位置。体积延拓与接触形变为
\begin{align}
 L_c(p_0)&=\sum_k w_k(p_0)\bigl[R_{k,c}(p_0+b_c(p_0))+t_{k,c}\bigr],\\
 P_{\ell,c,f}(\xi)&=C_{c,f}\bigl(L_c(F_{\ell,0}(\xi))\bigr).
\end{align}
已知身体运输只保持材料地址。内部滑移、压迫、注册误差与目标形态差异仍进入独立未知量或误差界。

\subsection{调和圆柱的边界条件、缝合与跨身体运输}
\label{sec:v43:harmonic-detail}
圆柱坐标是对材料拓扑的建模选择，而不是把人体强行拟合成几何圆柱。对腿、躯干分别取参考材料环壳 $\Omega_\ell$，设其轴向边界为 $\Gamma_h^-,\Gamma_h^+$，外表面与内参考壳为 $\Gamma_d^0,\Gamma_d^1$。一种明确的场构造是
\begin{align}
 \Delta_{g_0}h&=0,& h|_{\Gamma_h^-}&=0,&h|_{\Gamma_h^+}&=1,\\
 \Delta_{g_0}d&=0,& d|_{\Gamma_d^0}&=0,&d|_{\Gamma_d^1}&=1.
\end{align}
未指定边界使用与该域一致的无通量条件。角度不能作为一个全局单值实函数绕圆柱递增；在切口两岸指定 $\widetilde\varphi^+-\widetilde\varphi^-=2\pi$，并匹配通量，再将 $\varphi=\widetilde\varphi\bmod 2\pi$ 视为 $S^1$ 值场。有效局部图为 $U_\ell\subset S^1\times(0,1)\times(0,1)$。这保留轴向、环向及内部层次的材料语义，同时避免跨角度接缝的数值大跳变。

本节图索引记为 $\ell$，与物理候选索引 $a$ 分开。这里所谓拓扑保持，指有效图册及其转移映射不改变所覆盖材料域的对应关系。它不等于在任意组织变形、压迫或人体个体差异下都有准确的内部同胚。坐标映射 $F_{\ell,0}:U_\ell\to\Omega_\ell$ 必须局部可逆；环壳内核、腋下、腹股沟及部位连接处可以需要额外图。腿图和躯干图的重叠处若有真实对应，定义 $\xi_k=\psi_{k\ell}(\xi_\ell)$；互不重叠的扫描部位保留部位与图标签 $\ell$，不虚构一条图间转换。图间同一几何量的表示满足
\begin{equation}
 F_{k,0}\circ\psi_{k\ell}=F_{\ell,0},\qquad
 g_k=(D\psi_{k\ell})^{-\top}g_\ell(D\psi_{k\ell})^{-1}.
\end{equation}
本文使用的材料技能尺度是事先固定的参考度量 $g_0$ 或固定任务容差度量；新身体的物理速度与能力由映射的真实 Jacobian 重建。不能把每个身体中的世界欧氏距离都暗中当作同一材料度量，然后声称严格迁移不变。

\subsection{完整 SMPL-X 蒙皮公式及内部延拓的责任边界}
\label{sec:v43:smpl-detail}
为使体表对应可实现，进一步明确前述简写。以 $N_v$ 个顶点的静止网格 $\bar V$、形状参数 $\beta$、表情参数 $\psi$ 和关节旋转参数 $\vartheta$ 定义
\begin{equation}
 T(\beta,\psi,\vartheta)=\bar V+B_S(\beta)+B_E(\psi)+B_P(\vartheta),
 \quad B_S=\sum_l\beta_l S_l,\quad B_E=\sum_j\psi_j E_j.
\end{equation}
姿态修正可由旋转偏离静止姿态的元素构造，例如
\begin{equation}
 B_P(\vartheta)=\sum_{k,j}P_{kj}\,[\operatorname{vec}(R_k(\vartheta)-I)]_j.
\end{equation}
关节位置由模型提供的关节回归器及形状得到 $J(\beta)$。沿运动学树的齐次变换为
\begin{equation}
 G_k(\vartheta,J)=\prod_{j\in\operatorname{anc}(k)}
 \begin{bmatrix}R_j(\vartheta)&J_j-J_{\operatorname{par}(j)}\\0&1\end{bmatrix},
 \qquad \widehat G_k=G_k(\vartheta,J)G_k(0,J)^{-1}.
\end{equation}
根节点按其世界位姿处理。线性混合蒙皮输出
\begin{equation}
 \begin{bmatrix}v_i\\1\end{bmatrix}
 =\sum_k w_{ik}\widehat G_k
 \begin{bmatrix}T_i(\beta,\psi,\vartheta)\\1\end{bmatrix},
 \qquad \sum_k w_{ik}=1.
\end{equation}
这个表达显示了 $R_k,t_k$ 中必须包含静止骨架校正的原因。跨身体的同顶点、同面片及重心坐标提供体表材料对应；它们本身不观测内部结构，也不规定肌肉、血管和脏器随皮肤一起运动。

体积延拓可以在四面体域中插值权重，也可以解弹性或调和延拓问题，但必须作为独立实现记录。接触映射 $C_{c,f}$ 若只由近似组织模型给出，误差进入 $E_{\rm geo}$ 及其可信界。绝不能从表面 SMPL-X 的拟合误差直接推导内部血管误差的大小。使用实际 CoP、人体运动补偿和独立观测校准，是这一边界下保持理论有效的条件。

\subsection{完整身体映射的有效域与有界局部反演}
\label{sec:rev:full-map}
参考材料图的可逆性与变形后身体映射的可逆性分别验收。固定图索引 $\ell$，将身体状态记为 $c$，接触状态记为 $f$，定义
\begin{equation}
 P_{\ell,c,f}=C_{c,f}\circ L_c\circ F_{\ell,0}.
 \label{eq:rev:full-map}
\end{equation}
因子 Jacobian 的范数使用相容的已声明材料与物理度量。下述条件是一个\emph{局部几何模型合同}，其准确性不能由表面拟合误差代替。

\begin{assumption}[完整映射的有效窗口]
\label{ass:rev:map-window}
在声明的材料邻域及时间窗口上，式~\eqref{eq:rev:full-map}的三个因子均具有所需的 $C^2$ 正则性，且它们的空间 Jacobian 满足
\begin{equation}
 \sigma_{\min}(DF_{\ell,0})\ge \underline\sigma_F>0,\quad
 \sigma_{\min}(DL_c)\ge \underline\sigma_L>0,\quad
 \sigma_{\min}(DC_{c,f})\ge \underline\sigma_C>0.
 \label{eq:rev:factor-nondegeneracy}
\end{equation}
每个被使用的逆分支具有唯一对应，且当前材料位置与其预测集合均保留在该分支内。内部真实解剖与所采用映射之间的误差另有明示界；式~\eqref{eq:rev:factor-nondegeneracy}只认证所采用几何映射的非退化性。
\end{assumption}

\begin{proposition}[完整反演与一个可检查的单射充分条件]
\label{prop:rev:map-inverse}
在假设~\ref{ass:rev:map-window}下，$P_{\ell,c,f}$ 在每一点局部可逆，且
\begin{equation}
 \sigma_{\min}(J_P)\ge
 \underline\sigma_P:=\underline\sigma_C\underline\sigma_L\underline\sigma_F,
 \qquad \|J_P^{-1}\|_2\le\underline\sigma_P^{-1}.
 \label{eq:rev:map-inverse-bound}
\end{equation}
若还需在一个凸的局部实坐标域 $U$ 上认证唯一逆分支，下列有限坐标差使用该图内事先选定的固定坐标内积；它与紧致窗口中的 $g_0$ 一致等价。固定 $\xi_0\in U$，令 $A_0=D_\xi P(\xi_0)$，并检查
\begin{equation}
 \sup_{\xi\in U}\|D_\xi P(\xi)-A_0\|_2
 \le\kappa\sigma_{\min}(A_0),\qquad 0\le\kappa<1.
 \label{eq:rev:injective-variation}
\end{equation}
这时，对任意 $\xi_1,\xi_2\in U$，
\begin{equation}
 \|P(\xi_1)-P(\xi_2)\|_2
 \ge(1-\kappa)\sigma_{\min}(A_0)\|\xi_1-\xi_2\|_2.
 \label{eq:rev:bilipschitz-lower}
\end{equation}
因此该分支单射，且其逆具有式~\eqref{eq:rev:bilipschitz-lower}对应的 Lipschitz 上界。
\end{proposition}
\begin{proof}
链式法则给出 $J_P=(DC_{c,f})(DL_c)(DF_{\ell,0})$。最小奇异值的乘积下界与逆函数定理给出~\eqref{eq:rev:map-inverse-bound}。对于第二个结论，在连接 $\xi_1,\xi_2$ 的线段上积分 Jacobian，得到
\[
 P(\xi_1)-P(\xi_2)=A_0(\xi_1-\xi_2)
 +\int_0^1[D_\xi P(\xi_2+t(\xi_1-\xi_2))-A_0](\xi_1-\xi_2)\,dt.
\]
三角不等式与~\eqref{eq:rev:injective-variation}给出结论。逐点非奇异本身不排除较大域上的多个逆分支，因此局部奇异值界与域上的单射验收承担不同职责。
\end{proof}

对体表 $d=0$，使用带边界图 $U^+=V\times[0,d_{\max})$。若完整体积拉回需要三维逆，则明确给出 $P$ 在该边界图附近的 $C^2$ 延拓，并在延拓上应用命题~\ref{prop:rev:map-inverse}；实际接触点仍限制在所用体表。若仅建模体表，采用 $P_{\mathcal B}:V\subset\R^2\to\R^3$ 的嵌入及切向条件 $\sigma_{\min}(D P_{\mathcal B})>0$，只对切向速度求解，法向接触自由度单独建模。带噪声或具有有限接触斑的 CoP 若不恰好位于曲面上，须先用一个已声明、在管状邻域内唯一的接触地址投影，并把投影误差纳入映射误差界。不能对任意三维速度使用一个不存在的二维矩阵逆。

线性混合蒙皮并不自动满足上述条件。例如 $\tfrac12I+\tfrac12R_z(\pi)=\operatorname{diag}(0,0,1)$ 是奇异的；因此只检查 $F_{\ell,0}$ 不足以认证 $P$。角度 lift 的跨缝跳变和通量条件还保留一个加性常数，本文以固定材料标志的角度值或一个固定均值条件选定该自由度。

\subsection{CoP 地址速度、刚体输入与执行误差}
\label{sec:rev:cop-input}
设探头 TCP 位姿为 $(p_T,R)$，压力中心在探头坐标中的位置为 $r_C(t)$，世界坐标为
\begin{equation}
 p_C=p_T+Rr_C(t),\qquad \omega^W=(\dot R R^\top)^\vee.
 \label{eq:rev:cop-position}
\end{equation}
\emph{刚体在当前 CoP 位置处的速度}与\emph{压力中心轨迹的速度}分别是
\begin{align}
 v_C^{\rm rb}&=\dot p_T+\omega^W\times Rr_C,\\
 \dot p_C&=v_C^{\rm rb}+R\dot r_C.
 \label{eq:rev:cop-velocities}
\end{align}
压力分布可以在探头静止时发生变化，故 $\dot r_C$ 不由机器人刚体 twist 单独指定。

令当前接触地址满足 $p_C=P_{\ell,c,f}(\xi_C)$，并定义 $b_p=D_cP\dot c+D_fP\dot f$。在有效逆分支内，完整地址运动为
\begin{equation}
 \boxed{\dot\xi_C=J_P^{-1}(v_C^{\rm rb}+R\dot r_C-b_p).}
 \label{eq:rev:cop-pullback}
\end{equation}
若需要用材料坐标表示机器人刚体作用，使用
\begin{equation}
 \dot\xi_C^{\rm rb}:=J_P^{-1}(v_C^{\rm rb}-b_p)
 =\dot\xi_C-J_P^{-1}R\dot r_C.
 \label{eq:rev:rigid-material-input}
\end{equation}
示教地址及姿态标架的运输使用完整的 $\dot\xi_C$；后续响应矩阵的输入则由实际刚体 twist 经已知约束约简和固定物理白化得到。特别地，$u_{\rm act}$ 不能直接用未经修正的压力中心地址导数代替。

一个明确的约定是，在当前候选的受控子空间上定义已知线性重建 $[v_C^{\rm rb};\omega^W]=T_a u_{\rm act}$；若采用已知保持运动，则另加该运动，并从任务漂移中按同一模型扣除。由此推导
\begin{equation}
 \dot y=J_a u_{\rm act}+w_a,\qquad J_a=H_a+\Theta_aN_a+E_{{\rm geo},a}.
 \label{eq:rev:actual-rigid-model}
\end{equation}
改变输入坐标时同步变换 $H_a,N_a$ 及物理度量。式~\eqref{eq:rev:cop-pullback}用于记录移动接触地址，不能额外向同一个刚体响应重复加入 $R\dot r_C$。

更具体地，若特征表达为 $y=\mathcal Y(\xi_C,R,r_C,c,f)$，其显式接触中心迁移项为
\begin{equation}
 w_C=\bigl(D_{\xi_C}\mathcal Y\,J_P^{-1}R+D_{r_C}\mathcal Y\bigr)\dot r_C.
 \label{eq:rev:cop-feature-drift}
\end{equation}
已知项按该等式补偿，剩余项进入 $w_a$ 的明示界。在纯刚体成像几何中，若 $r_C$ 只是 $p_T=P(\xi_C)-Rr_C$ 的表达变量，这两项可以精确相消；接触形变改变真实成像关系时，其剩余影响由相应接触或几何误差模型承担。这避免把压力中心迁移当作一个虚构的可控平移。

命令与实际执行分开记为
\begin{equation}
 u_{\rm cmd}=u_{\rm ff}+C\zeta,\qquad
 u_{\rm act}=u_{\rm cmd}+d_u,\qquad d_u\in\mathcal D_u,
 \label{eq:rev:execution-contract}
\end{equation}
其中 $\mathcal D_u$ 是在整个声明状态域内独立有效的执行误差集合。由
\begin{equation}
 \|J_a\|_2\le M_J:=\|H_a\|_2+\delta\|N_a\|_2+\eta
\end{equation}
得到 $\|J_ad_u\|_2\le M_J\varepsilon_u$，$\varepsilon_u=\sup_{d\in\mathcal D_u}\|d\|_2$。这项加入恢复扰动界。若真实输入满足 $a_i^\top u_{\rm act}\le b_i^{\rm act}$，命令必须满足
\begin{equation}
 a_i^\top u_{\rm ff}+r_\zeta\|a_i^\top C\|_2
 +h_{\mathcal D_u}(a_i)\le b_i^{\rm act},
 \quad h_{\mathcal D_u}(a_i)=\sup_{d\in\mathcal D_u}a_i^\top d.
 \label{eq:rev:tightened-input}
\end{equation}
精确线性接触约束还要求 $EC=0$、$Eu_{\rm ff}=0$；若约束作用于实际输入，则另需 $Ed_u=0$，或改用事先声明的接触容差集合。仅有执行误差幅度界不推出加速度界；需要认证实际加速度时，另行给出 $\dot d_u$ 或实际一步执行模型的有效界。本文的二阶恢复定理以这些输入合同为条件，不据此认证任意未知延迟内环。

\subsection{姿态纤维与完整接触地址运输}
姿态纤维取 $R=Q_{\ell(a),c,f}(\xi)\rho$，$\rho\in SO(3)$。定义
\begin{equation}
 A_i=(Q^\top\partial_iQ)^\vee,\qquad
 a_c=(Q^\top(D_cQ\dot c+D_fQ\dot f))^\vee .
\end{equation}
若 $\omega_R=(R^\top\dot R)^\vee$，则纤维的相对角速度为
\begin{equation}
 \omega_\rho=\omega_R-\rho^\top(A_\xi\dot\xi_C+a_c).
\end{equation}
这使探查的平移、姿态和材料位置能够共同记录。更换局部标架 $Q'=QS$ 时，$\rho'=S^\top\rho$，实际姿态仍为同一个 $R$。

\subsection{坐标设计服务的理论对象是一个任务耦合商空间}
在材料位置、探头姿态和实现组成的基点 $a$ 上，取允许命令空间 $\mathcal V_a$、任务反馈空间 $\mathcal Y_a$ 与不确定通道空间 $\mathcal Z_a$，映射为
\begin{equation}
 H_a:\mathcal V_a\longrightarrow\mathcal Y_a,\qquad
 N_a:\mathcal V_a\longrightarrow\mathcal Z_a.
\end{equation}
任务零空间可以抵消的通道是 $\mathcal C_a=N_a\ker H_a$。将它商掉，定义与右逆选择无关的诱导映射
\begin{equation}
 \overline D_a:\mathcal Y_a\longrightarrow\mathcal Z_a/\mathcal C_a,
 \qquad \overline D_av=[N_aB_av],\quad H_aB_a=I.
\end{equation}
任何两个右逆之差都在 $\ker H_a$ 中，因此上述等价类相同。在物理度量下选取正交代表，就得到前文 $D_\star$。完成探查后的信息子空间 $Q_a$ 再被商掉，对应 $\mathcal Z_a/(\mathcal C_a+Q_a)$ 和前文 $D_Q$。

这才是调和图册与纤维的理论用途：在各人体、各图和姿态标架中记录同一个任务与通道等价类；允许不同机器人用不同动作重新实现它。各映射及其像保持常秩的区域可以构成商向量丛；秩变化处应分层处理并重新验收，不能声称存在覆盖所有姿态的光滑固定秩丛。超声的中心与偏心切换恰是这样的秩变化例。





\section{内部解剖、任务几何与示教偏好}
\label{sec:v43:anatomy-rules}
\subsection{血管和脏器共享接口，但不共享一条中心线}
腿部与躯干分别使用各自有效调和圆柱图；内部任务仍以部位、结构和进度为索引。血管先验可以在每个离散分支 $b$ 内写为
\begin{equation}
 \xi_b(s)=\bar\xi_b(s;c)+g_b(s),\quad
 g_b\sim\mathcal{GP}(0,K_{{\rm anat},b}),\qquad
 \Pr(b\mid\mathcal D)=\pi_b,
\end{equation}
再由 $P_{\ell(a),c,f}(\xi_b)$ 映射到物理空间。图接缝附近使用周期角核或连续 lift；各分支保留关联，不能将两个血管分支的高斯均值混成一条实际不存在的血管。若先验来自解剖知识而非已采集数据，也可以直接给概率分布或确定性范围，不要求先训练 GP。

脏器任务可由材料域中的隐式边界 $\phi_o(\xi)=0$、目标区域 $\{\phi_o\le0\}$ 和一族指定扫掠切面表示。扫描参数 $s$ 索引一个当前扫掠段，切面编号保留为离散任务标签。脏器不是一条长血管；使用相同图册、纤维和任务实现接口，并不要求把脏器改写成中心线。呼吸滑移或接触压迫导致的内部位移须保留独立先验和在线估计。

\subsection{投影规则保留完整内部几何，GP 只学习相对动作}
对一个内部目标 $p_{\rm int}(s)$，沿已定义的局部路径或法向场找到体表作用点 $p_{\rm surf}(s)$。投影必须在所用片区内唯一、连续且满足接触可行性；若不存在这样的射线交点，切换片区或报告该规则不可用。世界空间中的血管切向为
\begin{equation}
 t(s)=\frac{\partial_s P_{\ell(a),c,f}(\xi_b(s))}
 {\|\partial_s P_{\ell(a),c,f}(\xi_b(s))\|},
 \quad R_{\rm geom}(s)=\mathcal R\bigl(t(s),n_{\rm surf}(s),\mathcal T\bigr),
 \quad \rho_{\rm geom}=Q^\top R_{\rm geom}.
\end{equation}
$\mathcal R$ 是任务指定的成像平面与接触姿态规则，例如横断面、长轴面或允许的混合切面；该规则不必对所有 $t,n_{\rm surf}$ 都有精确同时满足的解，接近退化时需重新验收。不能先把内部路径的深度分量删除，再声称扭曲血管的完整姿态信息已经保留。内部切向、切面目的、体表法向和接触允许角共同决定几何默认动作。

示教提供的是相对于该默认动作的可测量偏好：体表材料位置偏置、相对旋转、在有力测量时的压力规律，以及推进节奏。定义
\begin{equation}
 p_{\rm CoP}^{\rm pref}=P_{\ell(a),c,f}(\xi_{\rm surf}^{\rm pref}),\qquad
 R^{\rm pref}=Q\rho_{\rm geom}\Upsilon,\quad\Upsilon\in SO(3).
\end{equation}
当 $Q'=QS$、$\rho_{\rm geom}'=S^\top\rho_{\rm geom}$ 时，$\Upsilon$ 不变。这个相对量很适合学习，它把几何导致的转动与专家额外动作分开。它不保证某个内部中心线先验正确，也不把几何转动自动认定为专家手感。

GP forcing 学习技能形状，解剖 GP 学习结构位置；两者使用同一种 GP 推断方法但数据语义和后验分开。动捕没有测力时，不报告已学到压力或物理阻抗。法向力参考若由同步示教得到，也要先投影到任务规定范围；机器人接触稳定性由已有接触控制与实际验收负责。

\subsection{技能容差不替代任务完成条件}
记 GP 名义技能为 $Y_{\rm pref}(s)$，可信参考度量中的偏离代价为
\begin{equation}
 J_{\rm style}(a)=\sum_j
 \big\|\operatorname{Log}_{Y_{\rm pref}(s_j)}Y(a,s_j)\big\|^2_{\Lambda_j}.
\end{equation}
$\Lambda_j$ 可来自预先定义的技能容差；若来自示教变化，须区分有效动作多样性、均值估计不确定性与测量噪声。后验方差大不表示任务允许随意偏离。真正必须保持的姿态、声窗、接触和相位条件仍写入 $\mathcal A_\tau,\mathcal Q_\tau$。IRD 可以在该集合中改变实现，不能通过丢弃一个任务特征来降低校准维数。


\paragraph{几何规则的窗口语义。}
以上内部几何与投影规则在窗口起点使用技能上下文 $\kappa$ 的快照，产生 $\gamma_\kappa$ 与几何姿态截面 $\rho_{{\rm geom},\kappa}(s,\xi)$。窗口内固定它们的材料表达，实时身体和接触变化只通过 $P,Q$ 运输该参考。若新观测需要重新投影基准、改动材料姿态截面或改变目标关联，则建立新窗口并验收状态及速度桥接。由此，后续时间重参数结论针对同一固定材料空间向量场。




\section{材料几何上的 GP forcing 技能}\label{sec:rev:skill}
\label{sec:v43-geometric-skill}

\subsection{取舍：在材料空间学习，相同技能允许不同物理实现}
本版固定一条技能路线：GP 学习相对于解剖几何基准的形状驱动项，固定增益的几何二阶系统生成名义技能。机器人世界笛卡尔位姿仅用于动作执行和对照，不作为迁移技能的主要学习坐标。原因不是笛卡尔 GP 不能拟合轨迹，而是人体的旋转、体型、姿态和体表曲率会被混入其训练目标；材料拉回则明确扣除这些已知变化。AL-DMP 已给出空间与速度分离的经典基础\cite{aldmp2016}；几何 DMP 已给出流形技能生成的基础\cite{gdmp2024}；GP 学习 DMP forcing 也已有超声接近动作先例\cite{gpdmp2026}。这些基础在本版中的用途是稳定表达和复现技能，不单独计为原创理论。

内部血管或脏器先验仍由解剖模型及在线估计提供，不能从专家探头动作反推为准确内部结构。给定任务及该先验，在体表材料子域 $\mathcal B\subset\mathcal M$ 上生成一条几何基准曲线 $\gamma(s)$，并由内部目标方向、合格切面与接触几何定义名义姿态 $\rho_{\rm geom}(s)$。这一步必须保留内部几何，单独将深度置零的体表投影一般不能保留切面朝向。专家示教学习的是相对于上述基准的偏置和姿态残余，而不是再次学习整个人体的世界位置。

令 $g_0$ 为预先固定的材料度量，$g_{\mathcal B}$ 为其在体表子域上的诱导度量，使用局部法邻域中的表达
\begin{equation}
 \xi_{\rm pref}(s)=\operatorname{Exp}_{\gamma(s)} b(s),\qquad
 \rho_{\rm pref}(s)=\rho_{\rm geom}(s,\xi_{\rm pref}(s))\Upsilon(s),
 \label{eq:v43-skill-state}
\end{equation}
其中 $b(s)\in T_{\gamma(s)}\mathcal B$ 是体表切向偏置，$\Upsilon(s)\in SO(3)$ 是几何基准之上的姿态残余。$\rho_{\rm geom}(s,\xi)$ 是定义在允许位置邻域上的几何姿态截面，而非只有中心线位置才能评价的姿态序列。$\operatorname{Exp}$ 只在验收过的单射半径内使用；大范围运动用重叠图和状态运输。腿部与躯干可以使用不同圆柱图，但可以共享此接口。共同接口不意味着腿部血管技能无需新数据就能外推到腹部器官。

曲线的进度取材料几何基准的归一化弧长，而非含纠偏往返的实际 TCP 弧长：
\begin{equation}
 \ell(u)=\int_0^u\|\gamma_0'(r)\|_{g_{\mathcal B}}\,dr,\qquad
 L_\gamma=\ell(1),\qquad s=\frac{\ell(u)}{L_\gamma},\qquad
 \dot s=\nu(t)\ge0.
 \label{eq:v43-material-progress}
\end{equation}
任务片段为 $s\in[0,1]$；先验的内部材料弧长也可以作为进度，但所有示教和生成必须使用相同定义。分叉是离散任务分支，不应被单个高斯均值连成穿过两支之间的虚假中心线。

\subsection{规范技能上下文、公共锚点与后验冻结}
\label{sec:rev:canonical-skill}
使用 $\kappa$ 表示\emph{技能上下文}，使用 $c(t),f(t)$ 表示实时身体与接触状态。一个学习模型首先限定同类任务、相同材料部位、相同相位定义及规范基准 $\gamma$；不同身体的示教通过声明的材料对应运输到这一共同参考。此时全部平移标签确实位于同一个 $T_{\gamma(0)}\mathcal B$。不具备材料对应的不同部位分别建模。

以下 $\mathsf P$ 均为材料体表度量的 Levi--Civita 平行运输。若模型允许上下文改变规范基准为 $\gamma_\kappa$，则先指定公共材料锚点 $\xi_\ast$ 及从 $\gamma_\kappa(0)$ 到 $\xi_\ast$ 的运输路径。令
\begin{equation}
 A_\kappa(s)=
 \mathsf P_{\xi_\ast\leftarrow\gamma_\kappa(0)}
 \mathsf P_{\gamma_\kappa(0)\leftarrow\gamma_\kappa(s)},\qquad
 b^\ast=A_\kappa b,\quad F_b^\ast=A_\kappa F_b.
 \label{eq:rev:common-anchor}
\end{equation}
共同锚点上的 GP 以 $(s,\kappa)$ 为输入，以 $(F_b^\ast,F_\rho)$ 为输出，生成时用 $A_\kappa(s)^{-1}$ 运输回相应切空间。对于固定 $\kappa$，第一段运输不随 $s$ 变化，沿 $\gamma_\kappa$ 的第二段运输满足
\begin{equation}
 \frac{d}{ds}(A_\kappa b)=A_\kappa D_sb.
 \label{eq:rev:anchor-derivative}
\end{equation}
因此原有 forcing 标签与协变动力学保持一致。运输路径是模型的一部分；不同路径的 holonomy 不被忽略。

一个执行窗口内固定 $\kappa$、基准曲线、目标、增益、用于生成的 GP 数据及后验。此时 forcing 是同一个 $m_F(s;\kappa)$，固定初态的空间动力学解在任意单调时间重参数下保持相同；$\nu=0$ 暂停全部空间状态。实时 $c(t),f(t)$ 只通过已声明的 $P,Q$ 运输这个材料参考。若新的内部解剖观测需要改变基准或技能上下文，将其作为新窗口，先进行状态及速度桥接，再重新验收。在线任意更改 $m_F(s;\kappa(t))$ 会改变空间向量场，不能沿用固定空间曲线的不变性结论。

\subsection{真实 connection 与姿态残余的 gauge 消去}
材料图册的被动换图为 $\xi'=\phi(\xi)$。沿 $\gamma$ 的切向量 $b$ 与其协变导数变换为
\begin{equation}
 b'=D\phi\,b,\qquad
 (D_sb)^i=\frac{db^i}{ds}
 +\Gamma^i{}_{jk}(\gamma)\gamma'^j b^k,
 \label{eq:v43-coordinate-connection}
\end{equation}
其中 $\Gamma$ 是 $g_{\mathcal B}$ 的 Levi--Civita connection；若使用三维材料表达，应将环境 connection 投影到体表切空间。对 $z_b=D_sb$ 再求 $D_sz_b$ 时也必须包含相同的 connection。调和坐标不是全局欧式坐标，不能将不同图中的分量普通二阶导数直接当成同一个加速度。本文的几何方程使用真实协变导数。几何 DMP 提供流形技能的相关基础\cite{gdmp2024}；此处直接采用所声明度量的协变导数。

姿态纤维的标架改变为 $Q'=QS$、$\rho'=S^\top\rho$。连接一形式 $\widehat A=Q^\top dQ$ 满足
\begin{equation}
 \widehat A'=S^\top\widehat A S+S^\top dS.
 \label{eq:v43-gauge-connection}
\end{equation}
几何姿态同样运输为 $\rho_{\rm geom}'=S^\top\rho_{\rm geom}$，于是
\begin{equation}
 \Upsilon'=\rho_{\rm geom}'{}^\top\rho'
 =\rho_{\rm geom}^\top\rho=\Upsilon.
 \label{eq:v43-invariant-residual}
\end{equation}
这是选择相对姿态残余的重要好处：局部体表标架的任意旋转不会被学习成专家手感。它仍不是把整个物理姿态丛变成一个全局乘积；只是在存在有效几何参考截面的任务窗口内，对相对技能使用这一平凡化。

对实际身体映射后的基准姿态
\begin{equation}
 R_0(t)=Q_{\ell(a),c,f}(\xi(t))\rho_{\rm geom}(s(t),\xi(t)),\qquad
 R(t)=R_0(t)\Upsilon(t),
\end{equation}
设 $\omega_R=(R^\top\dot R)^\vee$ 和 $\omega_0=(R_0^\top\dot R_0)^\vee$，则示教的残余角速度为
\begin{equation}
 \omega_\Upsilon=\omega_R-\Upsilon^\top\omega_0.
 \label{eq:v43-human-angle-compensation}
\end{equation}
$\xi(t)$ 在示教时取实际 CoP 材料位置，在生成时取名义偏置位置。$\omega_0$ 包括沿材料路径、位置偏置变化以及实际人体运动引起的标架变化。平移示教使用式~\eqref{eq:rev:cop-pullback}记录完整材料接触地址，并按式~\eqref{eq:rev:rigid-material-input}区分刚体作用和压力中心迁移；同步人体姿态与探头位姿用于已知运输补偿。没有这个补偿，GP 会把被试移动误记为医生的主动纠偏。

\subsection{固定增益的空间二阶系统}
将偏置目标 $b_g(s)$ 定义为沿 $\gamma$ 平行运输的固定向量，即 $D_sb_g=0$；常用选择是 $b_g=0$。姿态残余目标 $\Upsilon_g$ 固定，常用选择为单位阵。令 $z_b=D_sb$，$\Upsilon'=\Upsilon\widehat\omega_{\rm sk}$，使用
\begin{align}
 D_sb&=z_b,\\
 D_sz_b&=-k_b(b-b_g)-d_bz_b+F_b(s),\\
 \frac{d\Upsilon}{ds}&=\Upsilon\widehat\omega_{\rm sk},\\
 \frac{d\omega_{\rm sk}}{ds}
 &=k_\rho\operatorname{Log}(\Upsilon^\top\Upsilon_g)^\vee
   -d_\rho\omega_{\rm sk}+F_\rho(s).
 \label{eq:v43-space-generator}
\end{align}
所有 $k_b,d_b,k_\rho,d_\rho$ 是固定正数，不由 GP 学习。在局部线性模型下可取 $d_b=2\sqrt{k_b}$、$d_\rho=2\sqrt{k_\rho}$，但这不保证任意大角度姿态、饱和与物理执行下无超调。$SO(3)$ 使用与 $\|\omega_{\rm sk}\|^2$ 对应的双不变度量；沿同一曲线对其速度求协变导数时，body 表达中的 connection 项为 $\tfrac12[\omega_{\rm sk},\omega_{\rm sk}]=0$，故此处的 $d\omega_{\rm sk}/ds$ 是精确的该速度协变加速度表达，而非忽略 connection 的一般近似。姿态 Log 使用 $\|\operatorname{Log}(\Upsilon^\top\Upsilon_g)^\vee\|<\pi-\epsilon$ 的有效窗口。

增益各向同性使被动正交标架变化直接保持方程。若需要矩阵增益，则 $K_b,D_b$ 必须是协变的正定张量，并随平行运输更新；不允许在每个图中重新指定相同数值矩阵而宣称同一物理技能。

时间执行直接乘进度速度：
\begin{align}
 \dot s&=\nu,\qquad D_tb=\nu z_b,\\
 D_tz_b&=\nu[-k_b(b-b_g)-d_bz_b+F_b(s)],\\
 \dot\Upsilon&=\nu\Upsilon\widehat\omega_{\rm sk},\\
 \dot\omega_{\rm sk}&=\nu[k_\rho\operatorname{Log}(\Upsilon^\top\Upsilon_g)^\vee
                    -d_\rho\omega_{\rm sk}+F_\rho(s)].
 \label{eq:v43-time-generator}
\end{align}
对所有状态同时缩放是关键。$\nu=0$ 时材料参考与其内部形状状态全部暂停，避免只冻结 canonical phase、弹簧仍向终点拉动。实际人体变化仍通过 $P_{\ell(a),c,f}$ 和 $Q_{\ell(a),c,f}$ 运输该材料参考，所以暂停材料推进不等于将世界位姿钉死。接触导纳与有效观测下的局部纠偏仍继续执行。

空间曲线对任意单调时间重参数保持相同；其物理速度与加速度并不保持相同。例如
\begin{equation}
 D_t^2 b=\dot\nu z_b+\nu^2D_sz_b.
 \label{eq:v43-time-acceleration}
\end{equation}
因而 $\nu$ 和 $\dot\nu$ 仍需按实际加速度、接触与 IRD 能力限制平滑调度。任务反向重扫应切换到独立反向技能或重新初始化；本版不对 $\nu<0$ 宣称耗散稳定。

\subsection{forcing 示教标签：共同切空间与导数噪声}
平移向量在不同 $s$ 属于不同切空间。统一使用式~\eqref{eq:rev:common-anchor} 的公共锚点与指定运输路径，定义
\begin{equation}
 \bar b(s)=A_\kappa(s)b(s),\qquad
 \bar z_b=\frac{d\bar b}{ds},\qquad
 F_b(s)=A_\kappa(s)^{-1}\bar F_b(s).
 \label{eq:v43-common-tangent}
\end{equation}
并令 $\bar b_g=A_\kappa b_g$，该向量在固定技能上下文中不随 $s$ 改变。若所有数据共享同一 $\gamma$，取 $\xi_*=\gamma(0)$，则 $A_\kappa(s)=\mathsf P_{0\leftarrow s}$，退化为沿基准的共同起点表达。选取 $T_{\xi_*}\mathcal B$ 的固定正交基后，位置 forcing 标签是
\begin{equation}
 \boxed{\bar F_{b,d}(s)=\bar b_d''(s)
      +d_b\bar b_d'(s)+k_b[\bar b_d(s)-\bar b_g].}
 \label{eq:v43-position-label}
\end{equation}
这里的导数是完成材料拉回、平行运输及相位对齐之后的空间导数。$\dot s$ 太小的重复样本先合并，不在停顿期间用 $\dot b/\dot s$ 制造无限标签。

姿态标签由
\begin{equation}
 \boxed{F_{\rho,d}=\omega_d'
       -k_\rho\operatorname{Log}(\Upsilon_d^\top\Upsilon_g)^\vee
       +d_\rho\omega_d}
 \label{eq:v43-orientation-label}
\end{equation}
得到。若平滑示教使用 $r_d(s)=\operatorname{Log}(\Upsilon_g^\top\Upsilon_d(s))^\vee$，则
\begin{align}
 \omega_d&=J_r(r_d)r_d',\\
 \omega_d'&=J_r(r_d)r_d''+
              DJ_r(r_d)[r_d']r_d',\\
 J_r(r)&=I-\frac{1-\cos\theta}{\theta^2}\widehat r
              +\frac{\theta-\sin\theta}{\theta^3}\widehat r^2,
       \qquad\theta=\|r\|.
 \label{eq:v43-right-jacobian}
\end{align}
零角附近用解析极限。对原始欧拉角直接差分，不等于这个角速度与协变加速度。

GP 采用共同切空间中的输出 $(\bar F_b,F_\rho)$。旋转输出位于固定的 Lie algebra $\mathfrak{so}(3)\simeq\R^3$，再用左平移 $F_\rho\mapsto\Upsilon\widehat F_\rho\in T_\Upsilon SO(3)$。这是明确的 tangent trivialization，不是把处在不同姿态的切向量未经说明相加。位置运输与旋转 left trivialization 可组成一个低维矩阵值 GP；共用一种学习方法不要求将解剖位置先验和专家 forcing 混成同一个后验。

二阶差分容易放大动捕噪声。对 $\bar b$ 先拟合可二次均方微分的 GP 平滑器，记
\begin{equation}
 \mathcal L_b=\partial_s^2+d_b\partial_s+k_b.
\end{equation}
若平滑后的条件均值和协方差为 $m_b,K_b^{\rm post}$，则线性导数推断给出
\begin{align}
 m_{F_b}(s)&=\mathcal L_bm_b(s)-k_b\bar b_g,\\
 \operatorname{Cov}(F_b(s),F_b(s'))
 &=\mathcal L_{b,s}\mathcal L_{b,s'}K_b^{\rm post}(s,s').
 \label{eq:v43-derivative-noise}
\end{align}
这一标签噪声在邻近相位之间相关，训练 forcing GP 时应保留该协方差或明确给出其近似。Mat\'ern $5/2$ 核满足本段所需的二阶均方导数。姿态式~\eqref{eq:v43-right-jacobian} 是非线性变换，不能把其后验称为精确高斯；可用局部 Jacobian 传播或后验样本生成标签及噪声估计。

示教平滑、标签变换、再训练 forcing GP 是两阶段 plug-in／经验贝叶斯构造。下文 GP 后验条件于所采用的标签和噪声模型；除非另行建立联合概率模型，不宣称它是原始示教的精确联合后验。这不影响一致运输后的标架协变结论。

本版不使用 $(g-y_0)$ 的除法归一化。零起终点偏置和闭环姿态残余都可能有 $g=y_0$；forcing 标签仍由式~\eqref{eq:v43-position-label}--\eqref{eq:v43-orientation-label}直接定义。若希望无限进度下 forcing 渐消，令
\begin{equation}
 \mathcal K_F(s,s')=\chi(s)\chi(s')k_f((s,\kappa),(s',\kappa'))\,\Sigma_F,
 \qquad \chi(s)=e^{-\alpha s},\quad\alpha>0,
 \label{eq:v43-gated-kernel}
\end{equation}
并直接对完整 forcing 标签建模，不把接近零的标签再除以 $\chi$。$\Sigma_F\succeq0$ 可保留偏置与旋转驱动间的相关性。有限扫描末端 $s=1$ 的 $\chi$ 一般不为零，也不需要为零。

\subsection{可计算 GP 与 atlas/gauge 协变}
堆叠标签 $\mathbf f$，令 $K$ 为上述矩阵核的训练 Gram 矩阵，$\Sigma_\epsilon$ 为示教标签噪声协方差。预测为
\begin{align}
 m_F(x_*)&=m_0(x_*)+K_{*D}(K+\Sigma_\epsilon)^{-1}
                         [\mathbf f-\mathbf m_0],\\
 K_F^{\rm post}(x_*,x_*')&=K_{**'}-
     K_{*D}(K+\Sigma_\epsilon)^{-1}K_{D*'}.
 \label{eq:v43-forcing-gp}
\end{align}
生成器使用 $m_F$，不需要从预测分布随机采样实际探头动作。选择新的技能上下文 $\kappa$ 时，训练数据不足覆盖的部分会保留不确定性；少量专家数据只能支持明确的低维技能类，不能保证任意器官泛化。

若共同参考正交基更换，输出向量为 $\mathbf f'=O\mathbf f$，均值和协方差同步变为 $m_F'=Om_F$、$K_F'=OK_FO^\top$。在原始相位相关局部标架中，一般运输矩阵为 $T(s)$，对应矩阵核必须满足
\begin{equation}
 \mathcal K_F'(s,s')=T(s)\mathcal K_F(s,s')T(s')^\top.
 \label{eq:v43-gauge-kernel}
\end{equation}
标签噪声按协方差变换，增益按式~\eqref{eq:rev:gain-transform}作为线性算子变换，故式~\eqref{eq:v43-forcing-gp}的预测与变换交换，生成的物理偏置保持相同。流形向量值 GP 文献已有相关的 gauge-independent 建模基础\cite{gpmanifold2021}；本版使用任务窗口内的显式运输即可，不额外引入大网络。运输沿指定曲线定义；不同绕行路径可能存在 holonomy，因此不声称平行运输与路径无关。

\paragraph{协方差、增益与度量的变换类型。}
在切向分量变换 $b'=T(s)b$ 下，协方差核与标签噪声的对应块满足
\begin{equation}
 K_F'(s,s')=T(s)K_F(s,s')T(s')^\top.
 \label{eq:rev:covariance-transform}
\end{equation}
作为线性算子的增益则满足
\begin{equation}
 K_b'=TK_bT^{-1},\qquad D_b'=TD_bT^{-1},\qquad
 g'=T^{-\top}gT^{-1}.
 \label{eq:rev:gain-transform}
\end{equation}
矩阵增益应相对于 $g$ 自伴正定，并沿基准平行运输。$T$ 随 $s$ 变化时，其导数由协变导数中的 connection 处理。正交标架下两类公式具有相同的形式，但一般坐标下不能混用。标量增益对应 $k_bI,d_bI$，自动满足算子变换规则。

\begin{lemma}[门控 GP 后验均值的渐消条件]
\label{lem:rev:gated-mean}
固定技能上下文和有限训练集，设 $K+\Sigma_\epsilon$ 可逆，核为
\begin{equation}
 K_F((s,\kappa),(s',\kappa'))
 =\chi(s)\chi(s')K_0((s,\kappa),(s',\kappa')),
 \qquad\chi(s)=e^{-\alpha s},\quad\alpha>0.
\end{equation}
若每个训练输入与预测点之间的 $K_0$ 在 $s\ge0$ 上一致有界，且先验均值满足 $\|m_0(s;\kappa)\|\le M_0\chi(s)$，则存在有限常数 $M_F$ 使
\begin{equation}
 \|m_F(s;\kappa)\|\le M_Fe^{-\alpha s}.
 \label{eq:rev:forcing-decay}
\end{equation}
因此生成器使用的后验均值趋零且平方可积。
\end{lemma}
\begin{proof}
记 $\beta=(K+\Sigma_\epsilon)^{-1}(\mathbf f-\mathbf m_0)$。后验均值修正为有限和 $\chi(s)\sum_i\chi(s_i)K_0((s,\kappa),(s_i,\kappa_i))\beta_i$。对每项取界并加入先验均值界，即得~\eqref{eq:rev:forcing-decay}。
\end{proof}
该引理针对实际用于生成的均值；它不把所有 GP 样本或全部导数自动认证为满足同一个确定性界。若只声明有限扫描 $s\in[0,1]$，端点仍按有限段姿态、偏置与速度容差单独验收。平滑示教、非线性姿态标签与 forcing GP 的两阶段概率语义保持原定义。

\subsection{名义生成器的能量性质与终点边界}
令 forcing 是预测均值，取能量
\begin{equation}
 \mathcal E=\tfrac12\|z_b\|^2+
 \tfrac{k_b}{2}\|b-b_g\|^2+
 \tfrac12\|\omega_{\rm sk}\|^2+
 \tfrac{k_\rho}{2}d_{SO(3)}(\Upsilon,\Upsilon_g)^2.
 \label{eq:v43-skill-energy}
\end{equation}
向量范数均为固定材料度量与旋转度量下的范数，距离仅在上述 Log 有效域内求导。$D_sb_g=0$ 和 connection 的度量相容性给出
\begin{equation}
 \mathcal E'=-d_b\|z_b\|^2-d_\rho\|\omega_{\rm sk}\|^2
               +\langle z_b,F_b\rangle+\omega_{\rm sk}^\top F_\rho.
 \label{eq:v43-energy-identity}
\end{equation}
设 $d_{\min}=\min(d_b,d_\rho)$，Young 不等式得到
\begin{equation}
 \mathcal E'\le -\tfrac{d_{\min}}2
          (\|z_b\|^2+\|\omega_{\rm sk}\|^2)
             +\frac{\|F_b\|^2+\|F_\rho\|^2}{2d_{\min}}.
 \label{eq:v43-energy-inequality}
\end{equation}

\begin{proposition}[名义稳定与有限段能量界]
\label{prop:v43-nominal-skill}
固定基准、固定目标、固定增益下，若轨迹保留在一个不接触姿态 cut locus 的紧致有效邻域内，则零 forcing 的平衡 $(b_g,0,\Upsilon_g,0)$ 在该局部域中渐近稳定。若 forcing 有界且趋零、平方可积，并且解、空间速度及其一阶导数均匀连续（例如相应一、二阶空间导数有界），则延拓到 $s\to\infty$ 的解趋向该平衡。在任意有限进度段 $[0,S]$，只需 forcing 有界，即有
\begin{equation}
 \mathcal E(S)\le\mathcal E(0)+
  \frac1{2d_{\min}}\int_0^S
     (\|F_b\|^2+\|F_\rho\|^2)\,ds.
 \label{eq:v43-finite-energy}
\end{equation}
\end{proposition}
\begin{proof}
式~\eqref{eq:v43-energy-identity}在 forcing 为零时只允许零速度留在 $\mathcal E'=0$ 的最大不变集；动力学随后要求偏置目标和姿态目标同时满足。有效域中唯一的势能临界点是上述平衡，故局部 LaSalle 论证成立。式~\eqref{eq:v43-energy-inequality}直接积分得到有限段能量界。对平方可积 forcing，该式给出速度平方可积；速度均匀连续使其趋零，速度导数均匀连续及速度极限存在进一步使速度导数趋零。代回动力学并使用 forcing 趋零，得到 $b-b_g\to0$ 和 $\operatorname{Log}(\Upsilon^\top\Upsilon_g)\to0$。所有论证均以有效域保持为前提；可用初始能量与 forcing 能量小于域边界上的最低势能来验收一个充分条件。
\end{proof}

时间能量变化为 $\dot{\mathcal E}=\nu\mathcal E'$。因此暂停不增加名义生成器能量。但若累计进度 $\int_0^\infty\nu\,dt$ 有限，尤其到 $s=1$ 就停，不能用 $s\to\infty$ 的结果宣称已经到达吸引子。有限扫描末端必须另外规定允许的偏置、姿态和速度容差；可采用经检查的末端条件化或短段接续。端点位姿或速度不兼容时，不能仅靠 forcing 渐消自动修复。基准目标更新、人体重注册和分叉切换同样应做状态运输与平滑接续，不能搬用固定基准的能量证明。

\subsection{与 IRD、反馈修复及失锁规则的完整接口}
技能生成器只输出 $\xi_{\rm pref},\rho_{\rm pref}$ 及其导数；它不取消允许集合 $\mathcal A_\tau$。物理候选仍为
\begin{equation}
 p_{\rm CoP}=P_{\ell(a),c,f}(\xi),\qquad
 R=Q_{\ell(a),c,f}(\xi)\rho,\qquad a\in\mathcal A_\tau.
\end{equation}
IRD 在这些允许实现中提出候选。可用材料偏置与姿态残余距离给出技能偏好，例如
\begin{equation}
 \mathcal L_{\rm skill}(a,s)=
 \lambda_b\|b(a)-b_{\rm pref}(s)\|^2+
 \lambda_\rho d_{SO(3)}(\Upsilon(a),\Upsilon_{\rm pref}(s))^2,
 \label{eq:v43-skill-preference}
\end{equation}
但同一个有效任务、允许偏离、真实观测质量、接触、输入能力和碰撞条件仍是硬验收。候选内逐一重算 $H,N,D_\star,D_Q$，探查数据仍由实际执行产生，反馈修复仍由 $q<1$ 的证书验收。因此更换技能编码不改变不可消除耦合商空间、最小校准方向及任务恢复的核心主张。

名义 forcing 的能量证明不等于患者端口无源、物理力稳定，也不证明 GP 方差构成可靠误差管。外层图像纠偏、接触导纳、QP 限幅和机器人延迟仍遵守各自证书。二阶参考生成与二阶图像纠偏作用在不同状态：前者表达希望的材料技能，后者修复实际观测误差；不把两套弹簧无说明地串在同一误差通道。图像失锁时，按式~\eqref{eq:cal-loss-mode}执行已验收减速计划，制动完成后令 $\nu=0$ 冻结全部材料技能参考；图像纠偏状态制动，接触控制保持，未来误差管按缺少新观测的界扩张。恢复关联和声窗后才重新验收推进。

示教仅提供低维动作残余及可选推进偏好。合成内部血管可以训练解剖趋势、目标几何与仿真先验，但不能替代专家真实扫描动作作为“手感”证据。少量专家示教是否足够，应以留出人体的技能保持误差、暂停接续、有效声窗、接触变化和重新校准预算验证，而非规定一个未经实验支持的数据量。最强取舍对照是同一材料坐标、同一 IRD、同一观测、相同带宽及算力下的 GP 材料参考加二阶跟踪；若 forcing 编码没有额外收益，应承认二者的功能等价，不将符号丰富度当成新能力。





\section{可微逆可达性与有限候选生成}
\subsection{候选生成与硬验收分开}
令 $a$ 包括材料允许偏离、探头姿态纤维、rail 配置和机器人实现分支。若 IRD 使用 TCP 标架 $T$，由 $p_C=p_T+Rr_C$ 构造
\begin{equation}
 p_T=p_C-Rr_C,\qquad
 T_{\rm probe}=\begin{bmatrix}R&p_T\\0&1\end{bmatrix}.
\end{equation}
若库使用法兰或其它固定工具标架，先应用其已知标定外参。移动 CoP 的位置不能直接替代库中 TCP 的平移。变换到 IRD 的完整目标为
\begin{equation}
 T_{\rm rel}(a)=T_{\rm base}(a)^{-1}T_{\rm probe}(a).
\end{equation}
对可微 IRD 得分 $\mathcal I(T_{\rm rel},a)$，用链式法则
\begin{equation}
 D_a\mathcal I=D_T\mathcal I\,D_aT_{\rm rel}+\partial_a\mathcal I
\end{equation}
搜索能保留任务且具有足够纠偏能力的候选。校准尾谱可作目标项：
\begin{equation}
 \min_{a\in\mathcal A_\tau}
 c_{\rm switch}(a)+c_{\rm probe}(a,k)
 +\lambda\,\sigma_{k+1}(D_\star(a))
 -\mu\,\mathcal I(T_{\rm rel}(a),a).
\end{equation}
奇异值重合处用次梯度、光滑近似或直接枚举小候选库。IRD 是提议器；接受条件来自完整 IK 分支、接触、多面体输入能力、碰撞和有效观测，而非一个软得分。选姿态时 $H,N$ 也可能改变，必须逐候选重算，不搬用旧的 $D_\star$。

\subsection{有限候选库中的可审查算法}
\begin{enumerate}
\item 由材料任务及纤维记录运输任务目的，建立当前允许实现集合；检查动捕、图册与可信目标关联。
\item 由实际几何计算 $H_a,N_a$ 与剩余误差界，先纳入接触约束，再白化度量。
\item 对每个真实可行候选计算 $D_\star(a)$，在从先验起始、要求覆盖全部历史的最坏情况首次响应合同下，用谱尾计算必要激励方向下界；明确的零残差条件历史使用相应 $D_Q$。对任意非零已观测历史，须从其数据一致集合重新推导，不能沿用原始谱作为当前剩余成本的下界。
\item 生成可执行的短探查库。每条候选存储实际可激励通道、时间、幅度、接触变化和观测有效区间；不是只存一列理想 SVD 向量。
\item 在库内枚举或分支定界探查组合；使用预期噪声界与激励条件排除无法达成目标性能的组合。这里的最优性仅相对于声明的有限库。
\item 执行选定短探查，记录实际刚体输入 $u_{\rm act}$ 以及 CoP 地址和运输项；数据矩阵使用统一白化后的实际输入。失锁段不能作为可信特征回归数据；模型有效性被破坏时重新获取目标。
\item 用实际数据求解修复实现 $C$ 的凸验收问题。证书不通过时，补探查、换实现或停止推进；不能把不可行结果剪掉后声称完备。
\item 执行闭环短段，按沿途变化界更新未来管，并在证书到期前重算或探查。
\end{enumerate}
对每个候选得到可行总代价上界 $U_a$。仅在上述相同历史与保证量词下，无约束谱才给出对应方向下界；若每有效方向确有统一代价下界 $c_{\min}>0$，可转为 $L_a=c_{\rm switch}^{\rm lb}+c_{\min}k_{\rm necessary}$。没有实际时间或能量下界时，不能把方向数直接当作时间下界。一般非零历史若没有独立推导的有效下界，只报告可行上界 $U_a$。相同合同下选出的 $U$ 与有效放宽 $L$ 才构成可报告差距；它不是整个连续非线性扫描问题的全球最优证明。





\section{GP 几何技能、IRD 与最小反馈校准的完整组合}
\label{sec:v43:integration}
\subsection{推进优先的候选选择与局部核心}
GP forcing 与几何生成器给出名义偏好，IRD 提议允许的物理实现。第~\ref{sec:rev:positive-progress}节的有限计划库对每个候选保存 $a_j,C_j(t),\nu_j(t)$ 及其来源明确的数据。先检查完整窗口的模型、恢复、参考配准、真实相位、成像足迹、接触、输入和终端接续条件，得到 $\mathcal P_{\rm acc}$。对覆盖任务，令 $\Delta\ell_j^{\rm cov}$ 为该计划在当前已覆盖前缀上认证的新增长度，定义
\begin{align}
 \Delta\ell_{\max}^{\rm cov}&=\max_{j\in\mathcal P_{\rm acc}}\Delta\ell_j^{\rm cov},\\
 j^*&\in\operatorname*{arg\,min}_{j\in\mathcal P_{\rm acc}:
            \,\Delta\ell_j^{\rm cov}=\Delta\ell_{\max}^{\rm cov}} c(j),\\
 c(j)&=c_{\rm switch}(j)+c_{\rm probe}(j)
       +\lambda_{\rm style}J_{\rm style}(j)
       -\lambda_I\mathcal I(j)
       +\lambda_u\int_0^{T_p}\|C_j(t)\|_F^2\,dt.
 \label{eq:v43:joint-choice}
\end{align}
有限非空接受集合保证最大推进量和后续最小代价都能取得。每个固定候选、固定数据与固定速度包络的反馈修复仍使用式~\eqref{eq:cal-convex-repair}的凸子问题；候选提议、窗口计划和整体选择本身不宣称凸性。若 $\Delta\ell_{\max}^{\rm cov}=0$，可以执行恢复、探查或等待，但不把名义推进计为新增覆盖；真实完成时间使用定理~\ref{v44:thm:hybrid-finish}。

\begin{proposition}[技能生成表达不改变局部校准核心]
\label{prop:v43:core-unchanged}
固定相同材料任务维数、允许物理实现、度量及局部响应类。将固定轨迹、GP 参考或 GP-forcing 几何生成器用作名义偏好，不改变 $D_\star,D_Q$ 的定义与其首次任务下界。对一个被接受的实现，只要由实际数据获得 $\widehat JC=I$、$q<1$，且生成的前馈残差及控制执行误差满足已声明扰动界，二阶恢复合同仍成立。
\end{proposition}
\begin{proof}
首次下界仅由局部 $H,N$、任务零空间及实际信息子空间 $Q$ 定义；名义偏好不是这些对象的额外观测。生成器可以影响候选 $a$，故每个被选择候选须重算 $H_a,N_a,D_Q(a)$。接受以后，将 $Ju_{\rm ff}-\dot y_r$ 与执行残差放入 $w$，即可使用二阶恢复定理。该推导没有将名义生成器稳定性转化为未经校准的真实反馈稳定性。
\end{proof}

这给出系统连接的具体含义：调和图册与纤维保留可迁移任务及相对技能；GP 学偏好；几何二阶生成器实现形状和状态接续；IRD 利用允许偏离改变物理实现；商空间谱说明该实现仍必须知道哪些响应；真实探查与凸修复最后给出反馈性能验收。它们的职责相接，但任何一个模块都不能替代另一个模块的证据。

\subsection{最小运行状态与闭环图}
记名义技能的产品流形状态为 $Y=(b,\Upsilon)$，速度状态为 $Z=(z_b,\omega_{\rm sk})$；$\Upsilon\in SO(3)$ 的运动始终按 $\dot\Upsilon=\Upsilon\widehat{\nu\omega_{\rm sk}}$ 计算，$\dot Y=\nu Z$ 仅作对应产品切丛上的运动学简写。

在一个可信短窗口，实际保存的状态为：部位与目标分支标签、名义进度 $s$、实际相位及其估计误差、已覆盖材料前缀与累计中断预算、名义技能 $(Y,Z)$、二阶图像纠偏状态 $\zeta$、实现 $a$、实际探查数据 $(X,R_{\rm data})$、修复 $C$、响应可信界及未来管。GP 数据与后验在窗口边界更新，新的技能上下文经状态桥接后进入下一窗口。

\begin{figure}[htbp]
\centering
\begin{tikzpicture}[node distance=9mm and 10mm,>=Latex,
 box/.style={draw=blue!45!black,rounded corners,fill=blue!4,align=center,text width=56mm,minimum height=12mm,font=\small},
 arrow/.style={->,thick,draw=blue!50!black}]
 \node[box] (skill) {材料任务、内部先验\\GP forcing 与几何技能状态};
 \node[box,right=of skill] (ird) {允许实现与 dIRD 提议\\接触、IK、碰撞验收};
 \node[box,below=of skill] (data) {实际动作与有效图像\\身体运动补偿、探查数据};
 \node[box,below=of ird] (repair) {任务商空间与响应验收\\$D_Q$, $\|JC-I\|\le q_{\rm act}<1$};
 \node[box,below=of data] (servo) {名义前馈与二阶纠偏\\$u=u_{\rm ff}+C\zeta$};
 \node[box,below=of repair] (tube) {增广未来管与模式判断\\推进、补探查、制动重捕获};
 \draw[arrow] (skill)--(ird);
 \draw[arrow] (ird)--(repair);
 \draw[arrow] (data)--(repair);
 \draw[arrow] (repair)--(tube);
 \draw[arrow] (repair)--(servo);
 \draw[arrow] (servo)--(data);
 \draw[arrow] (tube)--(servo);
 \draw[arrow] (tube.east)--++(6mm,0)|-(ird.east);
\end{tikzpicture}
\caption{材料技能、候选模型、实际数据与执行证书的连接。图像失锁时停止证据更新，执行已验收制动后暂停名义进度，接触环与纠偏制动继续。}
\end{figure}

\subsection{数据运输的充分条件}
\label{rev:sec:data-transfer}

固定模型内只改变 $B$ 或 $C$ 时，保留相同的
$(X,R_{\rm data},H,N)$ 即可。被动换图时，设白化后的任务、命令和通道变换分别为正交矩阵
$T,V,S$，则
\begin{equation}
 H_b=TH_aV^\top,\qquad
 N_b=SN_aV^\top,\qquad
 \Theta_b=T\Theta_aS^\top .
 \label{rev:eq:passive-model-transform}
\end{equation}
这些变换保持物理范数、响应不确定集合和实际观测内容；此时
\begin{equation}
 Q_b=SQ_a,\qquad
 D_{Q_b,b}=S D_{Q_a,a}T^\top .
 \label{rev:eq:passive-quotient-transform}
\end{equation}
因此奇异值及阈值保持不变。一般可逆坐标变化应同时变换物理度量，再在相应白化空间中使用上述形式。

主动改变物理状态后，式~\eqref{rev:eq:passive-model-transform} 不自动成立。下面给出一个无需假定全局共享患者参数的局部充分条件。

\begin{proposition}[有误差界的响应数据运输]
\label{rev:prop:data-transfer}
设候选 $a,b$ 具有相同的任务维数和通道维数，且一个独立成立的局部响应对应关系给出已知正交矩阵
$T\in O(m)$、$S\in O(r)$ 以及
\begin{equation}
 \Theta_b=T\Theta_aS^\top+F_{ba},
 \qquad
 \|F_{ba}\|_2\le\varepsilon_{ba}.
 \label{rev:eq:response-transfer-assumption}
\end{equation}
若真实来源为 $a$ 的数据满足
$R_a=\Theta_aX_a+W_a$、$\|W_a\|_2\le\varepsilon_a$，定义
\begin{equation}
 X_{a\to b}=SX_a,\qquad
 R_{a\to b}=TR_a .
 \label{rev:eq:transported-calibration-data}
\end{equation}
则它们是关于 $\Theta_b$ 的有效线性观测，并满足
\begin{equation}
 R_{a\to b}=\Theta_bX_{a\to b}+W_{a\to b},
 \qquad
 \|W_{a\to b}\|_2
 \le\varepsilon_a+\varepsilon_{ba}\|X_a\|_2 .
 \label{rev:eq:transported-data-budget}
\end{equation}
\end{proposition}
\begin{proof}
代入式~\eqref{rev:eq:response-transfer-assumption}，得到
\begin{equation}
 W_{a\to b}=TW_a-F_{ba}SX_a.
\end{equation}
由 $T,S$ 正交及谱范数次乘性得到所述界。
\end{proof}

运输数据始终保留实际来源和对应关系的有效窗口。目标候选的估计采用自己的 $H_b,N_b$，并把
式~\eqref{rev:eq:transported-data-budget} 作为数据误差界；例如
$\delta_b=\delta_a+\varepsilon_{ba}$ 是一个保守响应范数界。该命题以式~\eqref{rev:eq:response-transfer-assumption} 已经成立为前提，材料坐标对应和 GP 后验本身不能替代这一响应关系。若不能给出该关系或其有效误差界，候选分别维护数据并重新探查。

\subsection{物理速度重建与当前状态接续}
归一化材料进度不等于相同物理扫描速度。对新身体中接受的非退化作用点曲线 $\xi(s)$，将 $v_{\rm pref}(s)$ 定义为相对身体的名义接触地址推进速度，不含保持运动与压力中心迁移补偿，则
\begin{equation}
 \nu_{\rm pref}=\frac{v_{\rm pref}(s)}{\|D_\xi P\,\xi'(s)\|},\qquad
 \dot p_C=D_\xi P\,\xi'(s)\nu+b_p,\qquad v_C^{\rm rb}=\dot p_C-R\dot r_C.
\end{equation}
实际 $\nu$ 再受式~\eqref{eq:rev:progress-invariant}及机器人能力限制；曲线导数接近零时转入停留或姿态模式，不作该除法。压力、姿态和速度偏好是不同的示教量，不能由速度归一化推导出接触阻抗。

接触地址的加速度还包括身体映射的曲率及时间变化项：
\begin{equation}
 \ddot p_C=J_P\ddot\xi+\dot J_P\dot\xi+\dot b_p.
\end{equation}
刚体作用点的速度变化还需从 $\ddot p_C$ 中扣除 $\dot R\dot r_C+R\ddot r_C$。$\dot J_P$ 含沿材料运动的 Hessian 作用和人体、接触状态的变化；$\dot b_p$ 也需要相应估计或可信界。不能仅验收材料生成器中的 $\dot\nu z_b+\nu^2D_sz_b$ 就宣称物理加速度已保证。

切换图、分支或几何参考时，在局部可逆表示 $Y=\Phi_{\rm new}(t,s,x)$ 中，可以按旧参考状态 $Y_o,\dot Y_o$ 初始化
\begin{equation}
 x_n=\Phi_{\rm new}^{-1}(Y_o),\qquad
 D_x\Phi_{\rm new}\,\dot x_n
 =\dot Y_o-\partial_s\Phi_{\rm new}\,\nu-\partial_t\Phi_{\rm new}.
\end{equation}
等式中的向量先运输到同一切空间。姿态参考 $R_o$ 的具体匹配为
\begin{equation}
 \Upsilon_n=R_{0,n}^\top R_o,\qquad
 \omega_{\Upsilon,n}=\omega_{R_o}-\Upsilon_n^\top\omega_{0,n}.
\end{equation}
在 $\nu>0$ 时，时间速度还可转换为空间生成速度；在暂停时不除以 $\nu$，先保持当前材料状态，按有限短段桥接后恢复推进。这个初始化只在映射有效、速度匹配可行且接续后的能力及初始增广状态域验收通过时提供连续性。无法匹配时采用显式有限桥接或停止，不宣称任意目标切换自动达到 $C^1$ 或 $C^2$。实际 $C,u_{\rm ff}$ 的接续也要避免命令突变，并单独验收。

\subsection{可执行算法：提议、探查、修复与接续}
\begin{enumerate}
 \item 将当前材料任务、目标分支与 GP 技能后验运输到有效图。检查周期角 lift、旋转对数域及动捕可信区间。
 \item 在材料进度上积分名义技能；从内部几何与残余姿态构造作用点和探头姿态，计算物理速度及 $\dot\nu Z$ 加速度项。该结果是偏好，不锁定唯一世界 TCP 轨迹。
 \item dIRD 结合技能代价与校准谱提出少量允许实现；逐个重算 $H_a,N_a$，纳入接触约束并在固定物理度量下白化。
 \item 使用已有实际数据试解修复。若可行，检查二阶状态包络下的输入能力、前馈残差与短窗口观测条件。
 \item 若不可行，在有限安全探查库中选择能补足 $(I-P_X)N_aC$ 主要方向的短动作；切换实现与探查的时间和扰动一并计入预算。
 \item 仅从关联正确、质量合格的观测和实际材料动作更新 $X,R_{\rm data}$。名义命令与实际执行不同的部分要记录，而不能把名义命令当真值。
 \item 求解响应修复，按增广二阶证书重算整段未来管；通过才推进。失败则补探查、换允许实现或停止。
 \item 失锁时平滑令 $\nu\to0$，完整暂停名义技能，并制动 $\zeta$；接触环继续。重捕获后先验证目标关联和旧证书适用性，再按当前状态接续。
\end{enumerate}
候选库穷尽后的失败是“在声明的库和模型内不可验收”，不是原任务必然不可完成。计算预算耗尽、秩变化、图册失效或模型界失效也应作为明确返回状态。

\subsection{未来位置锥不要求概率，但需要正确的证据语义}
由二阶增广状态管得到任务误差集合；若任务到位姿的局部映射 $\Phi_a$ 在可信域满足 Lipschitz 界 $L_\Phi$，则由误差引起的物理位置偏差可包络为 $L_\Phi\sqrt2 R$。完整未来集合还要加上内部解剖位置集合、名义技能候选变化及人体运输误差：
\begin{equation}
 \mathcal P_{t:j}=\bigcup_{b,a}\Phi_{a,b}\bigl(\mathcal A_{{\rm anat},b,j},\mathcal Y_{{\rm skill},j}\bigr)
 \oplus\mathcal B\bigl(L_\Phi\sqrt2 R_j+\epsilon_{{\rm map},j}\bigr).
\end{equation}
这里的 $\Phi_a$ 必须是经过验收的实际局部实现映射；$m<6$ 的超声特征误差并不约束全部位姿零空间，$HB=I$ 也不自动给出可积逆映射。不存在该映射时只认证任务特征管，物理未来管须由完整位姿状态的前向传播另行得到，不能直接使用 $L_\Phi\sqrt2 R$。
并集保留分叉；近未来可能因为刚校准和观测有效而较窄，远未来因响应变化界、关联和几何误差而变宽。新数据会更新锚点、候选和所有 $R_j$。若集合跨过图边界或 Log 域，应拆成多图候选，不能用一个无限扩大的局部高斯继续解释。

GP 后验协方差可生成条件于统计模型的解剖和技能候选区间；它不是 $\delta,\eta,L$ 的确定性证书。给概率覆盖率时必须报告模型、校准方法及验证数据。此处的鲁棒恢复管则条件于明示界，两类不确定性可以在预测中相加，但不能混写成一个未经验证的概率安全保证。


\subsection{跨窗口桥接与恢复证书的衔接}
\label{sec:rev:window-bridge}
被动换图时，把同一个物理状态、任务误差及反馈状态运输到相同物理度量下的共同标架；相应正交表示不改变 $S$。主动改变参考或物理实现时则必须计算新的真实误差 $e^+=y-y_r^+$，匹配可执行速度并重新检查 $S^+\le R^+$。原误差小不推出新误差小。

已有参考位置和速度的局部逆初始化仅提供桥接端点。桥接路径及其整个物理状态管须保持模型、观测、接触、速度、加速度和输入约束；若需要任务恢复保证，桥接期间也应有有效的 $q$ 与扰动合同。$C(t)$ 与 $u_{\rm ff}(t)$ 的平滑接续必须一同验收。不能仅因两端各自可行就推断中间可行，也不能通过任意插值 $C$ 自动保留 $\widehat JC=I$。

在同一个任务误差定义、共同标架和固定 $\omega$ 下，若整个接续过程仍有 $\|\Delta(t)\|\le\bar q<1$ 及统一扰动界，原公共 Lyapunov 证明允许此时变响应；其证明不需要 $\dot\Delta$。但命令变化率仍受式~\eqref{eq:rev:command-acceleration}约束。任务参考重置、度量变化或状态跳变必须单独通过初值与桥接合同，不能被时变 $\Delta$ 的结论覆盖。






\section{假设来源与结果的组合范围}
\label{sec:rev:assumption-map}
本文的确定性结论依赖明确的有效窗口。表~\ref{tab:rev:assumptions}列出每个主要量的来源以及它支撑的结论；一个未获得的误差界保持为显式前提。
\begin{table}[htbp]\centering\small
\caption{主要假设的来源与使用范围。}\label{tab:rev:assumptions}
\renewcommand{\arraystretch}{1.18}
\begin{tabularx}{\linewidth}{p{.19\linewidth}Yp{.27\linewidth}}
\toprule
对象 & 可检查的来源或假设 & 使用位置\\\midrule
$H_a,N_a$ & 解析响应及已知受控子空间；固定物理白化；主动候选分别定义 & 商空间、阈值、校准方向数\\
$\delta_a$ & 声明的完整球或物理子集；只在完整球上使用必要性 & 精确满射阈值与首次下界\\
$\eta_a$ & 结构外几何、接触、测量驱动系数与线性化误差的算子界 & $q$、前馈残差和解析实例\\
$\varepsilon_{\rm data}$ & 输出残差、实际输入记录误差及探查间变化的堆叠界 & 估计与有限库条件\\
$P^{-1}$ & 完整复合映射非退化与有效单射分支；体表使用边界图 & 材料拉回、技能运输\\
$\mathcal U,\mathcal D_u$ & 固定物理输入内近似及独立执行误差集合；仅紧化一次 & 凸修复、实际输入能力\\
$S\le R$ & 可计算初值条件与整个声明域上的统一模型、噪声、漂移界 & 恢复与正推进不变域\\
$C,b_s$ 的总导数 & 观测处理、状态反馈、身体运动和坐标重建的完整导数界 & 速度变化与加速度检查\\
窗口接续 & 参考、状态、数据来源及完整物理路径的有效桥接 & 跨窗口执行与有限完成\\
\bottomrule
\end{tabularx}\end{table}

商空间定理先在固定局部模型内给出精确边界；实际数据将该边界转为一个可计算的充分响应证书。二阶恢复定理允许有界、非对称、时变的乘性响应误差。正推进结论进一步要求输入、几何和过渡具有严格余量；解析超声例给出同时满足这些条件的非空参数区间。材料技能与 IRD 提供候选和前馈，其接受依据始终是同一个任务及其对应证书。

任务兼容主动转移、完整人体内部模型和任意长扫描具有额外状态与信息约束。本文按候选索引和有限有效窗口保留这些依赖，既不由坐标变化制造信息收益，也不把一次局部校准直接当作整个扫描的全局最优策略。



\begin{thebibliography}{99}
\small\raggedright\setlength{\itemsep}{4pt}\setlength{\parskip}{0pt}
\bibitem{pavlakos2019} G. Pavlakos et al. Expressive Body Capture: 3D Hands, Face, and Body from a Single Image. CVPR, 2019. \url{https://arxiv.org/abs/1904.05866}
\bibitem{wagenmaker2021} A. Wagenmaker, M. Simchowitz, K. Jamieson. Task-Optimal Exploration in Linear Dynamical Systems. ICML, PMLR 139:10641--10652, 2021. \url{https://proceedings.mlr.press/v139/wagenmaker21a.html}
\bibitem{kang2023} S. Kang, K. You. Minimum input design for direct data-driven property identification of unknown linear systems. Automatica 156:111130, 2023. \url{https://doi.org/10.1016/j.automatica.2023.111130}
\bibitem{grauer2022} J. A. Grauer, J. Pei. Minimum-Variance Control Allocation Considering Parametric Model Uncertainty. AIAA SciTech, 2022. NASA NTRS 20210023819. \url{https://ntrs.nasa.gov/citations/20210023819}
\bibitem{venkata2024} J. Venkatasubramanian, J. K\"ohler, J. Berberich, F. Allg\"ower. Sequential learning and control: Targeted exploration for robust performance. arXiv:2301.07995, version 3, 2024. \url{https://arxiv.org/abs/2301.07995}
\bibitem{kobayashi2026} J. Kobayashi. When Identifiability Is Not Enough: Terminal Controllability in Task-Compatible Active Calibration. IEEE Robotics and Automation Letters 11(10):11039--11045, 2026. \url{https://doi.org/10.1109/LRA.2026.3719194}
\bibitem{krupa2005} A. Krupa, F. Chaumette. Control of an ultrasound probe by adaptive visual servoing. IROS, 2005. \url{https://doi.org/10.1109/IROS.2005.1545272}
\bibitem{foucart2024} S. Foucart, C. Liao. Radius of Information for Two Intersected Centered Hyperellipsoids and Implications in Optimal Recovery from Inaccurate Data. arXiv:2401.11112, 2024. \url{https://arxiv.org/abs/2401.11112}
\bibitem{aldmp2016} A. Ude, R. Vuga, B. Nemec, J. Morimoto. Trajectory Representation by Nonlinear Scaling of Dynamic Movement Primitives. IROS, 2016. \url{https://hcr.ijs.si/resources/papers/iros16.pdf}
\bibitem{gdmp2024} F. J. Abu-Dakka, M. Saveriano, V. Kyrki. A unified formulation of geometry-aware discrete dynamic movement primitives. Neurocomputing 598:128056, 2024. \url{https://doi.org/10.1016/j.neucom.2024.128056}
\bibitem{gpdmp2026} W. Si, Z. Jin, N. Wang, C. Yang. Robot Learning System Based on Target Localization and Human Demonstration for Medical Examination. IEEE/ASME Transactions on Mechatronics, accepted manuscript, 2026. \url{https://repository.essex.ac.uk/42960/1/Robot_Learning_System_Based_on_Target_Localization_and_Human_Demonstration_for_Medical_Examination.pdf}
\bibitem{gpmanifold2021} M. Hutchinson, A. Terenin, V. Borovitskiy, S. Takao, Y. W. Teh, M. P. Deisenroth. Vector-valued Gaussian Processes on Riemannian Manifolds via Gauge Independent Projected Kernels. NeurIPS, 2021. \url{https://arxiv.org/abs/2110.14423}

\bibitem{demoorgolub1991} B. De Moor, G. H. Golub. The Restricted Singular Value Decomposition: Properties and Applications. SIAM Journal on Matrix Analysis and Applications 12(3):401--425, 1991. \url{https://doi.org/10.1137/0612029}
\bibitem{elghaoui1997} L. El Ghaoui, H. Lebret. Robust Solutions to Least-Squares Problems with Uncertain Data. SIAM Journal on Matrix Analysis and Applications 18(4):1035--1064, 1997. \url{https://doi.org/10.1137/S0895479896298130}
\bibitem{bombois2006} X. Bombois, G. Scorletti, M. Gevers, P. M. J. Van den Hof, R. Hildebrand. Least costly identification experiment for control. Automatica 42:1651--1662, 2006. \url{https://doi.org/10.1016/j.automatica.2006.05.016}
\bibitem{tanaskovic2014} M. Tanaskovi\'c, L. Fagiano, M. Morari. On the optimal worst-case experiment design for constrained linear systems. Automatica 50:3291--3298, 2014. \url{https://doi.org/10.1016/j.automatica.2014.10.058}
\bibitem{shakouri2026} A. Shakouri, H. J. van Waarde, M. K. Camlibel. Experiment Design for Set-membership Identification: From Prior Knowledge to Universal Inputs. arXiv:2607.00844, version 1, 2026. \url{https://arxiv.org/abs/2607.00844}
\bibitem{peregudin2026scalar} A. Peregudin, N. T. Dinh. Exact Noise Limits for Bounded Scalar Stabilization Experiments. arXiv:2609.39895, 2026. \url{https://arxiv.org/abs/2609.39895}
\bibitem{peregudin2026minimal} A. Peregudin, N. T. Dinh. Minimal Experiments for Robust Stabilization: Information, Spectral Geometry, and Duration. arXiv:2610.02187, version 1, 2026. \url{https://arxiv.org/abs/2610.02187}
\bibitem{plaskota2023} L. Plaskota, P. Siedlecki. Worst case tractability of linear problems in the presence of noise: linear information. arXiv:2303.16328, version 1, 2023. \url{https://arxiv.org/abs/2303.16328}
\bibitem{venkata2025robust} J. Venkatasubramanian, J. K\"ohler, M. Cannon, F. Allg\"ower. Robust targeted exploration for systems with non-stochastic disturbances. arXiv:2412.20426, version 2, 2025. \url{https://arxiv.org/abs/2412.20426}
\bibitem{tohidi2020} S. S. Tohidi, Y. Yildiz, I. Kolmanovsky. Adaptive control allocation for constrained systems. Automatica 121:109161, 2020. \url{https://doi.org/10.1016/j.automatica.2020.109161}
\bibitem{barenthin2005} M. Barenthin, H. Hjalmarsson. Identification and Control: Joint Input Design and $H_\infty$ State Feedback with Ellipsoidal Parametric Uncertainty. Proceedings of the 44th IEEE Conference on Decision and Control and the European Control Conference, pp. 6454--6459, 2005. \url{https://skoge.folk.ntnu.no/prost/proceedings/cdc-ecc05/pdffiles/papers/1777.pdf}

\end{thebibliography}

\end{document}
